\documentclass[12pt]{article}
\usepackage[T1]{fontenc}
\usepackage[utf8]{inputenc}
\usepackage{amsmath,amssymb,array,booktabs,longtable,geometry,microtype,enumitem}
\geometry{letterpaper,margin=0.90in}
\setlength{\parindent}{1.25em}
\setlength{\parskip}{0.34em}
\setlength{\emergencystretch}{3em}
\allowdisplaybreaks
\newcommand{\art}[1]{\subsection*{Art. #1}}
\newcommand{\Sec}{\S}
\newcommand{\norm}[1]{\operatorname{N}(#1)}
\newcommand{\gbrack}[1]{\left[#1\right]}
\newcolumntype{L}{>{$}l<{$}}
\newcolumntype{C}{>{$}c<{$}}
\sloppy
\begin{document}
\begin{center}
{\Large C. F. Gauss}\par
{\large Werke Band II — full cumulative English current}\par
{\large Actual beginning through printed p. 291}
\end{center}



\begin{center}
{\Large C. F. Gauss}\\[3pt]
{\large Werke Band II actual-beginning backfill cumulative}\\[3pt]
{\large English cumulative backfill through printed p. 101 pre-Art. 30}
\end{center}


\section*{A New Proof of an Arithmetic Theorem}

\begin{center}
\textsc{By Carl Friedrich Gauss}\\
\textsc{Presented to the Royal Society of Sciences, 15 January 1808}\\[3pt]
\emph{Commentationes Societatis Regiae Scientiarum Gottingensis}, Vol. XVI
\end{center}

\subsection*{1.}

Questions in higher arithmetic very often present a singular phenomenon, one that occurs much more rarely in analysis and contributes greatly to their charm. In analytic investigations, one usually cannot arrive at new truths unless one has first gained possession of the principles on which they rest and which must, as it were, open the way to them. In arithmetic, by contrast, elegant new truths frequently spring forth by induction, through a kind of unexpected good fortune; their demonstrations lie so deeply hidden and are wrapped in such darkness that they evade every attempt and refuse access to the keenest investigations. Moreover, so great and so remarkable is the connection among arithmetical truths that at first sight seem quite heterogeneous, that not rarely, while we are looking for something quite different, we reach at last the long-desired demonstration, sought in vain by long meditation, by a path far different from the one expected. In general, truths of this kind can be approached by several very different routes, and the shortest routes are not always those that present themselves first. It is therefore of high value when, after such a truth has long been handled in vain, and then demonstrated only by obscure detours, one finally succeeds in discovering the simplest and genuine path.

\subsection*{2.}

Among the questions mentioned in the preceding article, a distinguished place is held by the theorem that contains almost the whole theory of quadratic residues, and that in the \emph{Disquisitiones Arithmeticae}, Section IV, is distinguished by the name of the fundamental theorem. The illustrious Legendre must undoubtedly be regarded as the first discoverer of this most elegant theorem, although long before him the great geometers Euler and Lagrange had already discovered several special cases of it by induction. I shall not pause here to enumerate the efforts of these men toward a demonstration; those to whom it is agreeable may consult the work just mentioned. I add only, in confirmation of what was said in the preceding article, what concerns my own efforts. I came upon the theorem independently in the year 1795, when I was entirely ignorant of all that had already been worked out in higher arithmetic and entirely cut off from literary resources. For a full year it tormented me and escaped my most strenuous effort, until at last I obtained the demonstration given in Section IV of that work. Afterwards three other demonstrations, resting on wholly different principles, presented themselves to me; one of them I gave in Section V, and the others, not inferior to it in elegance, I shall make public on another occasion. Yet all these demonstrations, although they seem to leave nothing to be desired with respect to rigour, are derived from principles that are too heterogeneous, perhaps excepting the first, which nevertheless proceeds by more laborious reasoning and is burdened by more prolix operations. Thus I do not hesitate to say that until now a genuine demonstration has been lacking. Let experts now judge whether the demonstration that I have recently succeeded in discovering, and that the following pages present, deserves to be honoured with this name.

\subsection*{3. Theorem}

Let $p$ be a positive prime number; let $k$ be any integer not divisible by $p$; let $A$ be the collection of the numbers
\[
1,\ 2,\ 3,\ \ldots,\ \frac12(p-1),
\]
and let $B$ be the collection of the numbers
\[
\frac12(p+1),\ \frac12(p+3),\ \frac12(p+5),\ \ldots,\ p-1.
\]
Take the least positive residues modulo $p$ of the products of $k$ by the individual numbers in $A$. These residues are plainly all distinct, and some of them belong to $A$ while some belong to $B$. If, in all, $\mu$ of the residues are supposed to belong to $B$, then $k$ will be either a quadratic residue or a quadratic non-residue of $p$, according as $\mu$ is even or odd.

\emph{Proof.} Let the residues belonging to $A$ be $a,a',a'',\ldots$, and let the remaining residues belonging to $B$ be $b,b',b'',\ldots$. It is clear that the complements of the latter, $p-b,p-b',p-b'',\ldots$, are all distinct from $a,a',a'',\ldots$, and that together with them they complete the collection $A$. Hence
\[
1\cdot2\cdot3\cdots \frac12(p-1)
=
aa'a''\cdots(p-b)(p-b')(p-b'')\cdots .
\]
But the latter product is manifestly
\[
\equiv (-1)^\mu aa'a''\cdots bb'b''\cdots
\equiv (-1)^\mu k\cdot 2k\cdot 3k\cdots \frac12(p-1)k
\]
\[
\equiv (-1)^\mu k^{\frac12(p-1)}\,1\cdot2\cdot3\cdots\frac12(p-1)\pmod p.
\]
It follows that
\[
1\equiv (-1)^\mu k^{\frac12(p-1)}\pmod p,
\]
or
\[
k^{\frac12(p-1)}\equiv \pm1 \pmod p,
\]
according as $\mu$ is even or odd; and our theorem follows immediately.

\subsection*{4.}

The following reasoning may be considerably shortened by introducing some suitable notation. Let the character $(k,p)$ denote the number of products among
\[
k,\ 2k,\ 3k,\ \ldots,\ \frac12(p-1)k
\]
whose least positive residues modulo $p$ exceed one half of $p$. Moreover, if $x$ is any non-integral quantity, let $[x]$ denote the integer immediately less than $x$, so that $x-[x]$ is always a positive quantity lying between $0$ and $1$. The following relations are then obtained with little effort:
\[
\text{I.}\quad [x]+[-x]=-1.
\]
\[
\text{II.}\quad [x]+h=[x+h],\quad \text{whenever }h\text{ is an integer.}
\]
\[
\text{III.}\quad [x]+[h-x]=h-1.
\]
\[
\text{IV.}\quad
\begin{cases}
[2x]-2[x]=0, & \text{if }x-[x]<\frac12,\\
[2x]-2[x]=1, & \text{if }x-[x]>\frac12.
\end{cases}
\]
\[
\text{V.}\quad
\begin{cases}
\left[\frac{2h}{p}\right]-2\left[\frac{h}{p}\right]=0, & \text{if the least positive residue of }h\text{ lies below }\frac12p,\\
\left[\frac{2h}{p}\right]-2\left[\frac{h}{p}\right]=1, & \text{if that residue lies above }\frac12p.
\end{cases}
\]

VI. It follows at once that
\[
(k,p)=
\left[\frac{2k}{p}\right]+\left[\frac{4k}{p}\right]+\left[\frac{6k}{p}\right]+\cdots+\left[\frac{(p-1)k}{p}\right]
-2\left(
\left[\frac{k}{p}\right]+\left[\frac{2k}{p}\right]+\left[\frac{3k}{p}\right]+\cdots+\left[\frac{\frac12(p-1)k}{p}\right]
\right).
\]

VII. From VI and I one derives without difficulty
\[
(k,p)+(-k,p)=\frac12(p-1).
\]
It follows that $-k$ has either the same relation or the opposite relation to $p$, insofar as it is a quadratic residue or non-residue of $p$, as $+k$ has, according as $p$ is of the form $4n+1$ or $4n+3$. In the former case $-1$ is manifestly a residue, and in the latter a non-residue, of $p$.

VIII. We shall transform the formula given in VI as follows. By III,
\[
\left[\frac{(p-1)k}{p}\right]=k-1-\left[\frac{k}{p}\right],\quad
\left[\frac{(p-3)k}{p}\right]=k-1-\left[\frac{3k}{p}\right],
\]
\[
\left[\frac{(p-5)k}{p}\right]=k-1-\left[\frac{5k}{p}\right],\quad\ldots
\]
Applying these substitutions to the last $\frac{p\mp1}{4}$ terms of the upper series in that expression, we obtain, first, when $p$ is of the form $4n+1$,
\[
\begin{aligned}
(k,p)=&\ \frac14(k-1)(p-1)\\
&-2\left\{\left[\frac{k}{p}\right]+\left[\frac{3k}{p}\right]+\left[\frac{5k}{p}\right]+\cdots+\left[\frac{\frac12(p-3)k}{p}\right]\right\}\\
&-\left\{\left[\frac{k}{p}\right]+\left[\frac{2k}{p}\right]+\left[\frac{3k}{p}\right]+\cdots+\left[\frac{\frac12(p-1)k}{p}\right]\right\};
\end{aligned}
\]
second, when $p$ is of the form $4n+3$,
\[
\begin{aligned}
(k,p)=&\ \frac14(k-1)(p+1)\\
&-2\left\{\left[\frac{k}{p}\right]+\left[\frac{3k}{p}\right]+\left[\frac{5k}{p}\right]+\cdots+\left[\frac{\frac12(p-1)k}{p}\right]\right\}\\
&-\left\{\left[\frac{k}{p}\right]+\left[\frac{2k}{p}\right]+\left[\frac{3k}{p}\right]+\cdots+\left[\frac{\frac12(p-1)k}{p}\right]\right\}.
\end{aligned}
\]

IX. For the special case $k=+2$, it follows from the formulas just given that
\[
(2,p)=\frac14(p\mp1),
\]
taking the upper or lower sign according as $p$ is of the form $4n+1$ or $4n+3$. Thus $(2,p)$ is even, and hence $2$ is a quadratic residue of $p$, whenever $p$ is of the form $8n+1$ or $8n+7$; on the other hand, $(2,p)$ is odd, and $2$ is a quadratic non-residue of $p$, whenever $p$ is of the form $8n+3$ or $8n+5$.

\subsection*{5. Theorem}

Let $x$ be a positive non-integral quantity such that none of its multiples $x,2x,3x,\ldots$ up to $nx$ is an integer. Put $[nx]=h$. From this it is easily concluded that no integer is found among the multiples of the reciprocal quantity, $\frac1x,\frac2x,\frac3x,\ldots$ up to $\frac hx$. I claim that
\[
\left.
\begin{array}{l}
[x]+[2x]+[3x]+\cdots+[nx]\\
+\left[\frac1x\right]+\left[\frac2x\right]+\left[\frac3x\right]+\cdots+\left[\frac hx\right]
\end{array}
\right\}=nh.
\]

\emph{Proof.} Let the series $[x]+[2x]+[3x]+\cdots+[nx]$ be denoted by $\Omega$. Its first terms up to and including $\left[\frac1x\right]$ are manifestly all $0$; the following terms up to $\left[\frac2x\right]$ are all $1$; the following terms up to $\left[\frac3x\right]$ are all $2$, and so on. Hence
\[
\begin{aligned}
\Omega={}&0\left[\frac1x\right]
+1\left\{\left[\frac2x\right]-\left[\frac1x\right]\right\}
+2\left\{\left[\frac3x\right]-\left[\frac2x\right]\right\}\\
&+3\left\{\left[\frac4x\right]-\left[\frac3x\right]\right\}+\cdots\\
&+(h-1)\left\{\left[\frac hx\right]-\left[\frac{h-1}{x}\right]\right\}
+h\left\{n-\left[\frac hx\right]\right\}\\
={}&hn-\left[\frac1x\right]-\left[\frac2x\right]-\left[\frac3x\right]-\cdots-\left[\frac hx\right].
\end{aligned}
\]
Q.E.D.

\subsection*{6. Theorem}

Let $k,p$ denote any positive odd numbers prime to one another. Then
\[
\left.
\begin{array}{l}
\left[\frac{k}{p}\right]+\left[\frac{2k}{p}\right]+\left[\frac{3k}{p}\right]+\cdots+\left[\frac{\frac12(p-1)k}{p}\right]\\
+\left[\frac{p}{k}\right]+\left[\frac{2p}{k}\right]+\left[\frac{3p}{k}\right]+\cdots+\left[\frac{\frac12(k-1)p}{k}\right]
\end{array}
\right\}
=\frac14(k-1)(p-1).
\]

\emph{Proof.} Assuming, as we may, that $k<p$, the quantity $\frac{\frac12(p-1)k}{p}$ is smaller than $\frac12k$ but greater than $\frac12(k-1)$; hence
\[
\left[\frac{\frac12(p-1)k}{p}\right]=\frac12(k-1).
\]
It is therefore clear that the present theorem follows directly from the preceding one by putting $\frac{k}{p}=x$, $\frac12(p-1)=n$, and hence $\frac12(k-1)=h$ there.

Moreover, it can be demonstrated in a similar way that if $k$ is an even number prime to $p$, then
\[
\left.
\begin{array}{l}
\left[\frac{k}{p}\right]+\left[\frac{2k}{p}\right]+\left[\frac{3k}{p}\right]+\cdots+\left[\frac{\frac12(p-1)k}{p}\right]\\
+\left[\frac{p}{k}\right]+\left[\frac{2p}{k}\right]+\left[\frac{3p}{k}\right]+\cdots+\left[\frac{\frac12 kp}{k}\right]
\end{array}
\right\}
=\frac14k(p-1).
\]
But we shall not dwell on this proposition, since it is not necessary for our purpose.

\subsection*{7.}

Now, by combining the preceding theorem with proposition VIII of article 4, the fundamental theorem follows immediately. Namely, let $k,p$ denote any two distinct positive prime numbers, and put
\[
(k,p)+\left[\frac{k}{p}\right]+\left[\frac{2k}{p}\right]+\left[\frac{3k}{p}\right]+\cdots+\left[\frac{\frac12(p-1)k}{p}\right]=L,
\]
\[
(p,k)+\left[\frac{p}{k}\right]+\left[\frac{2p}{k}\right]+\left[\frac{3p}{k}\right]+\cdots+\left[\frac{\frac12(k-1)p}{k}\right]=M.
\]
By VIII of article 4 it is clear that $L$ and $M$ are always even numbers. But by the theorem of article 6,
\[
L+M=(k,p)+(p,k)+\frac14(k-1)(p-1).
\]
Therefore, whenever $\frac14(k-1)(p-1)$ is even, which happens if either both $k,p$, or at least one of them, is of the form $4n+1$, the two numbers $(k,p)$ and $(p,k)$ must necessarily be either both even or both odd. But whenever $\frac14(k-1)(p-1)$ is odd, which happens if both $k,p$ are of the form $4n+3$, one of the numbers $(k,p),(p,k)$ must necessarily be even and the other odd. In the former case the relation of $k$ to $p$ and the relation of $p$ to $k$, insofar as each is a residue or non-residue of the other, will be identical; in the latter case they will be opposite. Q.E.D.

\clearpage
\section*{Source title pages}
Printed pp. 9--10 are the title and blank interleaf for \emph{Summatio quarumdam serierum singularium}; the title-page scan is included in the source-scan sidecar.


\section*{Summation of Certain Singular Series}

\begin{center}
\textsc{By Carl Friedrich Gauss}\\
\textsc{Presented to the Society on 24 August 1808}\\[3pt]
\emph{Commentationes Societatis Regiae Scientiarum Gottingensis Recentiores}, Vol. I
\end{center}

\subsection*{1.}

Among the more notable truths to which the theory of the division of the circle has opened the way, no insignificant place is held by the summation proposed in the \emph{Disquisitiones Arithmeticae}, art. 356. This is so not only because of its peculiar elegance and its remarkable fruitfulness, which a later investigation will give occasion to explain more fully, but also because its rigorous and complete proof is burdened by difficulties of no ordinary kind. These difficulties would have been all the less expected, since they do not concern the theorem itself so much as a certain limitation of the theorem; if that limitation is neglected, the proof is immediately at hand and is very easily derived from the theory explained in that work.

The theorem stated there is in the following form. Suppose that $n$ is a prime number. Let all quadratic residues of $n$ lying between $1$ and $n-1$, inclusive, be denoted indefinitely by $a$, and all non-residues lying between the same limits by $b$. Let $\omega$ denote the arc $\frac{2\pi}{n}$, and let $k$ be any fixed integer not divisible by $n$. Then:

I. for a value of $n$ of the form $4m+1$,
\[
\sum \cos ak\omega=\frac{-1\pm\sqrt n}{2},\qquad
\sum \cos bk\omega=\frac{-1\mp\sqrt n}{2},
\]
and therefore
\[
\sum \cos ak\omega-\sum \cos bk\omega=\pm\sqrt n,\qquad
\sum \sin ak\omega=0,\qquad \sum \sin bk\omega=0.
\]

II. for a value of $n$ of the form $4m+3$,
\[
\sum \cos ak\omega=-\frac12,\qquad
\sum \cos bk\omega=-\frac12,
\]
\[
\sum \sin ak\omega=\pm\frac12\sqrt n,\qquad
\sum \sin bk\omega=\mp\frac12\sqrt n,
\]
and therefore
\[
\sum \sin ak\omega-\sum \sin bk\omega=\pm\sqrt n.
\]

These summations have been proved with full rigour in the cited place, and no difficulty remains here except the determination of the sign to be prefixed to the radical quantity. It can indeed be shown without trouble that this sign depends on the number $k$ only in this respect: the same sign must always hold for all values of $k$ that are quadratic residues of $n$, and the opposite sign for all values of $k$ that are quadratic non-residues of $n$. Thus the whole matter will turn on the value $k=1$, and it is clear that, as soon as the sign valid for this value is known, the signs for all other values of $k$ will immediately be available. Yet in this very question, which at first sight seems to belong among the easier ones, we encounter unexpected difficulties; the method under whose guidance we had so far advanced without impediment now entirely refuses further help.

\subsection*{2.}

Before going further, it will not be out of place to work out some examples of our summation by numerical calculation. First, however, it is useful to place before them some general observations.

I. Consider the case in which $n$ is a prime of the form $4m+1$. Let all the quadratic residues of $n$ lying between $1$ and $\frac12(n-1)$, inclusive, be denoted indefinitely by $a'$, and all the non-residues between the same limits by $b'$. Then all the numbers $n-a'$ are comprised among the $a$, and all the numbers $n-b'$ among the $b$. Since $a',b',n-a',n-b'$ plainly fill out the whole set of numbers $1,2,3,\ldots,n-1$, the $a'$ together with the $n-a'$ comprise all the $a$, and likewise the $b'$ together with the $n-b'$ comprise all the $b$. Hence
\[
\sum\cos ak\omega=\sum\cos a'k\omega+\sum\cos(n-a')k\omega,
\]
\[
\sum\cos bk\omega=\sum\cos b'k\omega+\sum\cos(n-b')k\omega,
\]
\[
\sum\sin ak\omega=\sum\sin a'k\omega+\sum\sin(n-a')k\omega,
\]
\[
\sum\sin bk\omega=\sum\sin b'k\omega+\sum\sin(n-b')k\omega.
\]
Since
\[
\cos(n-a')k\omega=\cos a'k\omega,\quad
\cos(n-b')k\omega=\cos b'k\omega,
\]
\[
\sin(n-a')k\omega=-\sin a'k\omega,\quad
\sin(n-b')k\omega=-\sin b'k\omega,
\]
it follows at once that
\[
\sum\sin ak\omega=0,\qquad \sum\sin bk\omega=0.
\]
The summation of the cosines, however, assumes the form
\[
\sum\cos ak\omega=2\sum\cos a'k\omega,\qquad
\sum\cos bk\omega=2\sum\cos b'k\omega,
\]
and therefore one must have
\[
1+4\sum\cos a'k\omega=\pm\sqrt n,\qquad
1+4\sum\cos b'k\omega=\mp\sqrt n,
\]
\[
2\sum\cos a'k\omega-2\sum\cos b'k\omega=\pm\sqrt n.
\]

II. In the case where $n$ is of the form $4m+3$, the complement of any residue $a$ to $n$ is a non-residue, and the complement of any $b$ is a residue. Hence all the $n-a$ agree with the $b$, and all the $n-b$ with the $a$. Thus
\[
\sum\cos ak\omega=\sum\cos(n-b)k\omega=\sum\cos bk\omega.
\]
Since the $a$ and $b$ together fill out all the numbers $1,2,3,\ldots,n-1$, one has
\[
\sum\cos ak\omega+\sum\cos bk\omega
=\cos k\omega+\cos2k\omega+\cdots+\cos(n-1)k\omega=-1,
\]
and therefore the summations
\[
\sum\cos ak\omega=-\frac12,\qquad \sum\cos bk\omega=-\frac12
\]
are immediate. Likewise
\[
\sum\sin ak\omega=\sum\sin(n-b)k\omega=-\sum\sin bk\omega,
\]
which shows how the two summations
\[
2\sum\sin ak\omega=\pm\sqrt n,\qquad
2\sum\sin bk\omega=\mp\sqrt n
\]
depend on one another.

\subsection*{3.}

Here is the numerical calculation for several examples.

I. For $n=5$ there is one value of $a'$, namely $a'=1$, and one value of $b'$, namely $b'=2$; moreover
\[
\cos\omega=+0.3090169944,\qquad \cos2\omega=-0.8090169944.
\]
Thus
\[
1+4\cos\omega=+\sqrt5,\qquad 1+4\cos2\omega=-\sqrt5.
\]

II. For $n=13$ there are three values of $a'$, namely $1,3,4$, and as many values of $b'$, namely $2,5,6$. We compute
\[
\begin{array}{rl@{\qquad}rl}
\cos\omega &= +0.8854560257,& \cos2\omega &= +0.5680647467,\\
\cos3\omega &= +0.1205366803,& \cos5\omega &= -0.7485107482,\\
\cos4\omega &= -0.3546048870,& \cos6\omega &= -0.9709418174.
\end{array}
\]
\[
\text{Sum}=+0.6513878190,\qquad \text{Sum}=-1.1513878189.
\]
Hence
\[
1+4\sum\cos a'\omega=+\sqrt{13},\qquad
1+4\sum\cos b'\omega=-\sqrt{13}.
\]

III. For $n=17$ we have four values of $a'$, namely $1,2,4,8$, and as many values of $b'$, namely $3,5,6,7$. The cosines are
\[
\begin{array}{rl@{\qquad}rl}
\cos\omega &= +0.9324722294,& \cos3\omega &= +0.4457383558,\\
\cos2\omega &= +0.7390089172,& \cos5\omega &= -0.2736629901,\\
\cos4\omega &= +0.0922683595,& \cos6\omega &= -0.6026346364,\\
\cos8\omega &= -0.9829730997,& \cos7\omega &= -0.8502171357.
\end{array}
\]
\[
\text{Sum}=+0.7807764064,\qquad \text{Sum}=-1.2807764065.
\]
Hence
\[
1+4\sum\cos a'\omega=+\sqrt{17},\qquad
1+4\sum\cos b'\omega=-\sqrt{17}.
\]

IV. For $n=3$ there is a single value of $a$, namely $a=1$, to which corresponds
\[
\sin\omega=+0.8660254038.
\]
Hence $2\sin\omega=+\sqrt3$.

V. For $n=7$ there are three values of $a$, namely $1,2,4$; hence the sines are
\[
\sin\omega=+0.7818314825,\quad
\sin2\omega=+0.9749279122,\quad
\sin4\omega=-0.4338837391.
\]
\[
\text{Sum}=+1.3228756556,
\]
and therefore $2\sum\sin a\omega=+\sqrt7$.

VI. For $n=11$ the values of $a$ are $1,3,4,5,9$, to which correspond the sines
\[
\begin{array}{rl}
\sin\omega &= +0.5406408175,\\
\sin3\omega &= +0.9898214419,\\
\sin4\omega &= +0.7557495744,\\
\sin5\omega &= +0.2817325568,\\
\sin9\omega &= -0.9096319954.
\end{array}
\]
\[
\text{Sum}=+1.6583123952,
\]
and so $2\sum\sin a\omega=+\sqrt{11}$.

VII. For $n=19$ the values of $a$ are $1,4,5,6,7,9,11,16,17$, to which correspond the sines
\[
\begin{array}{rl@{\qquad}rl}
\sin\omega &= +0.3246994692,& \sin9\omega &= +0.1645945903,\\
\sin4\omega &= +0.9694002659,& \sin11\omega &= -0.4759473930,\\
\sin5\omega &= +0.9965844930,& \sin16\omega &= -0.8371664783,\\
\sin6\omega &= +0.9157733267,& \sin17\omega &= -0.6142127127,\\
\sin7\omega &= +0.7357239107.&
\end{array}
\]
\[
\text{Sum}=+2.1794494718,
\]
and therefore $2\sum\sin a\omega=+\sqrt{19}$.

\subsection*{4.}

In all these examples the radical quantity obtains the positive sign, and the same is easily confirmed for larger values, $n=23$, $n=29$, and so forth. From this there already arises a strong probability that the same generally holds. But the proof of this phenomenon cannot be drawn from the principles set forth in the cited place, and it must with full right be regarded as a matter of deeper investigation. The aim of this memoir is therefore to bring forward a rigorous proof of this most elegant theorem, after it had formerly been attempted in vain for several years in various ways, and was finally completed successfully by rather singular and subtle considerations. At the same time we shall raise the theorem itself, with its elegance preserved or rather increased, to a much greater generality. Finally, as a crown, we shall show the remarkably close connection between this summation and another very important arithmetical theorem. We hope that these investigations will not only be pleasing to geometers in themselves, but that the methods by which all this could be accomplished, and which may be useful on other occasions as well, will seem worthy of their attention.

\subsection*{5.}

Our proof is drawn from the consideration of a singular kind of progressions whose terms depend on expressions such as
\[
\frac{(1-x^m)(1-x^{m-1})(1-x^{m-2})\cdots(1-x^{m-\mu+1})}
{(1-x)(1-x^2)(1-x^3)\cdots(1-x^\mu)}.
\]
For brevity we shall denote such a fraction by $(m,\mu)$, and first make some general observations about functions of this kind.

I. Whenever $m$ is a positive integer less than $\mu$, the function $(m,\mu)$ plainly vanishes, since the numerator contains the factor $1-x^0$. For $m=\mu$, the factors in the numerator are identical with the factors in the denominator in reverse order, whence $(\mu,\mu)=1$. Finally, in the case in which $m$ is a positive integer greater than $\mu$, we have
\[
(\mu+1,\mu)=\frac{1-x^{\mu+1}}{1-x}=(\mu+1,1),
\]
\[
(\mu+2,\mu)=\frac{(1-x^{\mu+2})(1-x^{\mu+1})}{(1-x)(1-x^2)}=(\mu+2,2),
\]
\[
(\mu+3,\mu)=\frac{(1-x^{\mu+3})(1-x^{\mu+2})(1-x^{\mu+1})}{(1-x)(1-x^2)(1-x^3)}=(\mu+3,3),\quad\text{etc.}
\]
or, in general,
\[
(m,\mu)=(m,m-\mu).
\]

II. It is also easy to confirm that, in general,
\[
(m,\mu+1)=(m-1,\mu+1)+x^{m-\mu-1}(m-1,\mu).
\]
Since similarly
\[
\begin{aligned}
(m-1,\mu+1)&=(m-2,\mu+1)+x^{m-\mu-2}(m-2,\mu),\\
(m-2,\mu+1)&=(m-3,\mu+1)+x^{m-\mu-3}(m-3,\mu),\\
(m-3,\mu+1)&=(m-4,\mu+1)+x^{m-\mu-4}(m-4,\mu),\quad\text{etc.},
\end{aligned}
\]
and this series can be continued until
\[
(\mu+2,\mu+1)=(\mu+1,\mu+1)+x(\mu+1,\mu)
=(\mu,\mu)+x(\mu+1,\mu),
\]
provided $m$ is a positive integer greater than $\mu+1$, we obtain
\[
\begin{aligned}
(m,\mu+1)=&(\mu,\mu)+x(\mu+1,\mu)+x^2(\mu+2,\mu)+x^3(\mu+3,\mu)+\cdots\\
&+x^{m-\mu-1}(m-1,\mu).
\end{aligned}
\]
It follows that if, for some fixed value of $\mu$, every function $(m,\mu)$ is integral, with $m$ a positive integer, then every function $(m,\mu+1)$ must also become integral. Since this supposition holds for $\mu=1$, it also holds for $\mu=2$, and then for $\mu=3$, and so on. Thus, in general, for every positive integral value of $m$, $(m,\mu)$ is an integral function; equivalently, the product
\[
(1-x^m)(1-x^{m-1})(1-x^{m-2})\cdots(1-x^{m-\mu+1})
\]
is divisible by
\[
(1-x)(1-x^2)(1-x^3)\cdots(1-x^\mu).
\]

\subsection*{6.}

We shall now consider two progressions, either of which can lead to our goal. The first progression is
\[
1-\frac{1-x^m}{1-x}
+\frac{(1-x^m)(1-x^{m-1})}{(1-x)(1-x^2)}
-\frac{(1-x^m)(1-x^{m-1})(1-x^{m-2})}{(1-x)(1-x^2)(1-x^3)}
+\cdots,
\]
or
\[
1-(m,1)+(m,2)-(m,3)+(m,4)-\cdots.
\]
For brevity we shall denote it by $f(x,m)$. It is immediately clear that, whenever $m$ is a positive integer, this series breaks off after its $(m+1)$-st term, which becomes $\pm1$, and that in this case the sum must be a finite integral function of $x$. Moreover, by art. 5.II, for any value of $m$ one has
\[
\begin{aligned}
1&=1,\\
-(m,1)&=-(m-1,1)-x^{m-1},\\
+(m,2)&=+(m-1,2)+x^{m-2}(m-1,1),\\
-(m,3)&=-(m-1,3)-x^{m-3}(m-1,2),\quad\text{etc.}
\end{aligned}
\]
and therefore
\[
\begin{aligned}
f(x,m)=&1-x^{m-1}-(1-x^{m-2})(m-1,1)\\
&+(1-x^{m-3})(m-1,2)-(1-x^{m-4})(m-1,3)+\cdots.
\end{aligned}
\]
But plainly
\[
\begin{aligned}
(1-x^{m-2})(m-1,1)&=(1-x^{m-1})(m-2,1),\\
(1-x^{m-3})(m-1,2)&=(1-x^{m-1})(m-2,2),\\
(1-x^{m-4})(m-1,3)&=(1-x^{m-1})(m-2,3),\quad\text{etc.},
\end{aligned}
\]
from which we deduce the equation
\[
f(x,m)=(1-x^{m-1})f(x,m-2). \tag{1}
\]

\subsection*{7.}

Since for $m=0$ one has $f(x,m)=1$, the formula just found gives
\[
f(x,2)=1-x,
\]
\[
f(x,4)=(1-x)(1-x^3),
\]
\[
f(x,6)=(1-x)(1-x^3)(1-x^5),
\]
\[
f(x,8)=(1-x)(1-x^3)(1-x^5)(1-x^7),\quad\text{etc.}
\]
or, in general, for any even value of $m$,
\[
f(x,m)=(1-x)(1-x^3)(1-x^5)\cdots(1-x^{m-1}). \tag{2}
\]

On the other hand, since for $m=1$ one has $f(x,m)=0$, one also has
\[
f(x,3)=0,\qquad f(x,5)=0,\qquad f(x,7)=0,\quad\text{etc.},
\]
or, in general, for any odd value of $m$,
\[
f(x,m)=0.
\]
The latter summation could already have been derived from the fact that, in the progression
\[
1-(m,1)+(m,2)-(m,3)+\cdots+(m,m-1)-(m,m),
\]
the last term destroys the first, the penultimate the second, and so on.

\subsection*{8.}

For our purpose it is enough to consider the case in which $m$ is a positive odd integer; but because of the singularity of the matter it will not be unwelcome to add a few words about the cases in which $m$ is fractional or negative. Then our series no longer breaks off, but runs out to infinity; furthermore it is easily seen to become divergent whenever a value smaller than $1$ is assigned to $x$. Consequently its summation must be restricted to values of $x$ greater than $1$.

From formula [1] of art. 6 we have
\[
f(x,-2)=\frac{1}{1-\frac1x},
\]
\[
f(x,-4)=\frac{1}{1-\frac1x}\cdot\frac{1}{1-\frac1{x^3}},
\]
\[
f(x,-6)=\frac{1}{1-\frac1x}\cdot\frac{1}{1-\frac1{x^3}}\cdot\frac{1}{1-\frac1{x^5}},\quad\text{etc.}
\]
Thus the value of the function $f(x,m)$ can be assigned in finite terms even for negative integral even values of $m$. For the remaining values of $m$, however, we shall convert the function $f(x,m)$ into an infinite product as follows.

As $m$ increases toward a negative infinite value, the function $f(x,m)$ passes into
\[
1+\frac{1}{x-1}
+\frac{1}{x-1}\cdot\frac{1}{x^2-1}
+\frac{1}{x-1}\cdot\frac{1}{x^2-1}\cdot\frac{1}{x^3-1}
+\cdots.
\]
This series is therefore equal to the infinite product
\[
\frac{1}{1-\frac1x}\cdot\frac{1}{1-\frac1{x^3}}\cdot
\frac{1}{1-\frac1{x^5}}\cdot\frac{1}{1-\frac1{x^7}}\cdots.
\]
Moreover, since in general
\[
f(x,m)=f(x,m-2\lambda)(1-x^{m-1})(1-x^{m-3})(1-x^{m-5})\cdots(1-x^{m-2\lambda+1}),
\]
one obtains
\[
f(x,m)=f(x,-\infty)(1-x^{m-1})(1-x^{m-3})(1-x^{m-5})\cdots
\]
and hence
\[
=\frac{1-x^{m-1}}{1-x^{-1}}\cdot
\frac{1-x^{m-3}}{1-x^{-3}}\cdot
\frac{1-x^{m-5}}{1-x^{-5}}\cdot
\frac{1-x^{m-7}}{1-x^{-7}}\cdots.
\]
It is clear that these factors eventually converge more and more closely to unity.

The case $m=-1$ deserves special attention; here
\[
f(x,-1)=1+x^{-1}+x^{-3}+x^{-6}+x^{-10}+\cdots.
\]
This series is therefore equal to the infinite product
\[
\frac{1-x^{-2}}{1-x^{-1}}\cdot
\frac{1-x^{-4}}{1-x^{-3}}\cdot
\frac{1-x^{-6}}{1-x^{-5}}\cdot\cdots.
\]
Writing $x$ for $x^{-1}$, this becomes
\[
1+x+x^3+x^6+\cdots
=
\frac{1-x^2}{1-x}\cdot
\frac{1-x^4}{1-x^3}\cdot
\frac{1-x^6}{1-x^5}\cdot
\frac{1-x^8}{1-x^7}\cdots.
\]
This equality between two rather abstruse expressions, to which we shall return on another occasion, is certainly very memorable.

\subsection*{9.}

In the second place we shall consider the following progression:
\[
1+x^{\frac12}\frac{1-x^m}{1-x}
+x\frac{(1-x^m)(1-x^{m-1})}{(1-x)(1-xx)}
+x^{\frac32}\frac{(1-x^m)(1-x^{m-1})(1-x^{m-2})}{(1-x)(1-xx)(1-x^3)}
+\cdots,
\]
or
\[
1+x^{\frac12}(m,1)+x(m,2)+x^{\frac32}(m,3)+xx(m,4)+\cdots,
\]
which we shall denote by $F(x,m)$. We restrict this investigation to the case in which $m$ is a positive integer, so that this series too always terminates with the $(m+1)$-st term, which is $x^{\frac12m}(m,m)$. Since
\[
(m,m)=1,\qquad (m,m-1)=(m,1),\qquad (m,m-2)=(m,2),\quad\text{etc.},
\]
the progression can also be displayed as
\[
F(x,m)=x^{\frac12m}+x^{\frac12(m-1)}(m,1)+x^{\frac12(m-2)}(m,2)+x^{\frac12(m-3)}(m,3)+\cdots.
\]
Hence
\[
\begin{aligned}
(1+x^{\frac12m+\frac12})F(x,m)
={}&1+x^{\frac12}(m,1)+x(m,2)+x^{\frac32}(m,3)+\cdots\\
&+x^{\frac12}x^m+x x^{m-1}(m,1)+x^{\frac32}x^{m-2}(m,2)+\cdots.
\end{aligned}
\]
Therefore, since by art. 5.II one has
\[
\begin{aligned}
(m,1)+x^m&=(m+1,1),\\
(m,2)+x^{m-1}(m,1)&=(m+1,2),\\
(m,3)+x^{m-2}(m,2)&=(m+1,3),\quad\text{etc.},
\end{aligned}
\]
there results
\[
(1+x^{\frac12m+\frac12})F(x,m)=F(x,m+1). \tag{3}
\]
But $F(x,0)=1$; hence
\[
F(x,1)=1+x^{\frac12},
\]
\[
F(x,2)=(1+x^{\frac12})(1+x),
\]
\[
F(x,3)=(1+x^{\frac12})(1+x)(1+x^{\frac32}),\quad\text{etc.},
\]
or, generally,
\[
F(x,m)=(1+x^{\frac12})(1+x)(1+x^{\frac32})\cdots(1+x^{\frac12m}). \tag{4}
\]

\subsection*{10.}

With these preliminary investigations in place, let us now come closer to our object. When $n$ is prime, the squares
\[
1,4,9,\ldots,\bigl(\tfrac12(n-1)\bigr)^2
\]
are all mutually incongruent modulo $n$. It is therefore clear that their least residues modulo $n$ must be identical with the numbers $a$, and hence that
\[
\sum\cos ak\omega=\cos k\omega+\cos4k\omega+\cos9k\omega+\cdots+
\cos\bigl(\tfrac12(n-1)\bigr)^2k\omega,
\]
\[
\sum\sin ak\omega=\sin k\omega+\sin4k\omega+\sin9k\omega+\cdots+
\sin\bigl(\tfrac12(n-1)\bigr)^2k\omega.
\]
Likewise, since the same squares, in reverse order, are congruent to
\[
\bigl(\tfrac12(n+1)\bigr)^2,\quad
\bigl(\tfrac12(n+3)\bigr)^2,\quad
\bigl(\tfrac12(n+5)\bigr)^2,\ldots,(n-1)^2,
\]
one also has
\[
\sum\cos ak\omega=\cos\bigl(\tfrac12(n+1)\bigr)^2k\omega+
\cos\bigl(\tfrac12(n+3)\bigr)^2k\omega+\cdots+\cos(n-1)^2k\omega,
\]
\[
\sum\sin ak\omega=\sin\bigl(\tfrac12(n+1)\bigr)^2k\omega+
\sin\bigl(\tfrac12(n+3)\bigr)^2k\omega+\cdots+\sin(n-1)^2k\omega.
\]
Putting therefore
\[
T=1+\cos k\omega+\cos4k\omega+\cos9k\omega+\cdots+\cos(n-1)^2k\omega,
\]
\[
U=\sin k\omega+\sin4k\omega+\sin9k\omega+\cdots+\sin(n-1)^2k\omega,
\]
we have
\[
1+2\sum\cos ak\omega=T,\qquad 2\sum\sin ak\omega=U.
\]
It follows that the summations proposed in art. 1 depend on summing the series $T$ and $U$. We therefore put those former sums aside and adapt our investigation to these series, doing so with enough generality to include not only prime values of $n$ but arbitrary composite values as well. We shall suppose, however, that $k$ is prime to $n$; for the case in which $k$ and $n$ had a common divisor can be reduced to this one without difficulty.

\subsection*{11.}

Denote the imaginary quantity $\sqrt{-1}$ by $i$, and put
\[
\cos k\omega+i\sin k\omega=r.
\]
Then $r^n=1$, so that $r$ is a root of the equation $x^n-1=0$. It is easy to see that the numbers
\[
k,2k,3k,\ldots,(n-1)k
\]
are not divisible by $n$ and are mutually incongruent modulo $n$. Hence the powers of $r$
\[
1,r,rr,r^3,\ldots,r^{n-1}
\]
are all unequal, while each also satisfies the equation $x^n-1=0$. For this reason these powers represent all the roots of the equation $x^n-1=0$.

These conclusions would not hold if $k$ had a common divisor with $n$. For if $\nu$ were such a common divisor, then $k\cdot\frac n\nu$ would be divisible by $n$, and therefore a power lower than $r^n$, namely $r^{\frac n\nu}$, would be equal to unity. In that case the powers of $r$ would display at most $\frac n\nu$ roots of the equation $x^n-1=0$; and they would in fact display that many distinct roots if $\nu$ is the greatest common divisor of $k,n$. In our case, where $k$ and $n$ are assumed relatively prime, $r$ may conveniently be called a proper root of the equation $x^n-1=0$. In the other case, where $k$ and $n$ had the greatest common divisor $\nu$, $r$ would be called an improper root of that equation, although then the same $r$ would plainly be a proper root of the equation $x^{\frac n\nu}-1=0$. The simplest improper root is unity, and when $n$ is prime there are no other improper roots.

\subsection*{12.}

If now we put
\[
W=1+r+r^4+r^9+\cdots+\cdots+r^{(n-1)^2},
\]
then plainly $W=T+iU$, so that $T$ is the real part of $W$, while $U$ arises from the imaginary part of $W$ after suppressing the factor $i$. The whole matter is therefore reduced to finding the sum $W$. For this purpose either the series considered in art. 6 or the one that we learned to sum in art. 9 can be used; the former, however, is less suitable when $n$ is even. Nevertheless, we hope it will be welcome to the reader if we treat the case in which $n$ is odd by both methods.

First suppose that $n$ is odd, that $r$ denotes any proper root of the equation $x^n-1=0$, and that in the function $f(x,m)$ one puts $x=r$ and $m=n-1$. It is then clear that
\[
\frac{1-x^m}{1-x}=\frac{1-r^{n-1}}{1-r}=-r^{-1},
\]
\[
\frac{1-x^{m-1}}{1-xx}=\frac{1-r^{n-2}}{1-rr}=-r^{-2},
\]
\[
\frac{1-x^{m-2}}{1-x^3}=\frac{1-r^{n-3}}{1-r^3}=-r^{-3},\quad\text{etc.},
\]
down to
\[
\frac{1-x}{1-x^m}=\frac{1-r}{1-r^m}=-r^{-m}.
\]
It is not superfluous to observe that these equations hold only insofar as $r$ is assumed to be a proper root. If $r$ were an improper root, in some of those fractions numerator and denominator would vanish at the same time, and the fractions would become indeterminate.

We therefore obtain
\[
f(r,n-1)=1+r^{-1}+r^{-3}+r^{-6}+\cdots+r^{-\frac12(n-1)n}
=(1-r)(1-r^3)(1-r^5)\cdots(1-r^{n-2}).
\]
The same equation will still hold if $r^\lambda$ is substituted for $r$, where $\lambda$ denotes any integer prime to $n$; for then $r^\lambda$ is also a proper root of $x^n-1=0$. Write therefore $r^{n-2}$, or what is the same thing $r^{-2}$, in place of $r$. Then
\[
1+r^2+r^6+r^{12}+\cdots+r^{(n-1)n}
=(1-r^{-2})(1-r^{-6})(1-r^{-10})\cdots(1-r^{-2(n-2)}).
\]
Multiply each side of this equation by
\[
r\cdot r^3\cdot r^5\cdots r^{n-2}=r^{\frac14(n-1)^2}.
\]
Using
\[
r^{2+\frac14(n-1)^2}=r^{\frac14(n-3)^2},\qquad
r^{(n-1)n+\frac14(n-1)^2}=r^{\frac14(n+1)^2},
\]
\[
r^{6+\frac14(n-1)^2}=r^{\frac14(n-5)^2},\qquad
r^{(n-2)(n-1)+\frac14(n-1)^2}=r^{\frac14(n+3)^2},
\]
\[
r^{12+\frac14(n-1)^2}=r^{\frac14(n-7)^2},\qquad
r^{(n-3)(n-2)+\frac14(n-1)^2}=r^{\frac14(n+5)^2},\quad\text{etc.},
\]
there results
\[
r^{\frac14(n-1)^2}+r^{\frac14(n-3)^2}+r^{\frac14(n-5)^2}+\cdots+r+1
+r^{\frac14(n+1)^2}+r^{\frac14(n+3)^2}+r^{\frac14(n+5)^2}+\cdots+r^{\frac14(2n-2)^2}
\]
\[
=(r-r^{-1})(r^3-r^{-3})(r^5-r^{-5})\cdots(r^{n-2}-r^{-n+2}),
\]
or, arranging the terms of the first member differently,
\[
1+r+r^4+\cdots+r^{(n-1)^2}=(r-r^{-1})(r^3-r^{-3})\cdots(r^{n-2}-r^{-n+2}). \tag{5}
\]

\subsection*{13.}

The factors of the second member of equation [5] can also be displayed as
\[
r-r^{-1}=-(r^{n-1}-r^{-n+1}),
\]
\[
r^3-r^{-3}=-(r^{n-3}-r^{-n+3}),
\]
\[
r^5-r^{-5}=-(r^{n-5}-r^{-n+5}),\quad\text{etc.},
\]
down to
\[
r^{n-2}-r^{-n+2}=-(r^2-r^{-2}).
\]
Thus the equation assumes the form
\[
W=(-1)^{\frac12(n-1)}(r^2-r^{-2})(r^4-r^{-4})(r^6-r^{-6})\cdots(r^{n-1}-r^{-n+1}).
\]
Multiplying this equation by [5] in its primitive form gives
\[
W^2=(-1)^{\frac12(n-1)}(r-r^{-1})(r^2-r^{-2})(r^3-r^{-3})\cdots(r^{n-1}-r^{-n+1}).
\]
Here $(-1)^{\frac12(n-1)}$ is either $+1$ or $-1$ according as $n$ is of the form $4\mu+1$ or $4\mu+3$. Hence
\[
W^2=\pm r^{\frac12 n(n-1)}(1-r^{-2})(1-r^{-4})(1-r^{-6})\cdots(1-r^{-2(n-1)}).
\]
But it is immediately clear that $r^{-2},r^{-4},r^{-6},\ldots,r^{-2n+2}$ exhibit all the roots of $x^n-1=0$ except the root $x=1$. Therefore the indefinite identity
\[
(x-r^{-2})(x-r^{-4})(x-r^{-6})\cdots(x-r^{-2n+2})
=x^{n-1}+x^{n-2}+x^{n-3}+\cdots+x+1
\]
must hold. Setting $x=1$ gives
\[
(1-r^{-2})(1-r^{-4})(1-r^{-6})\cdots(1-r^{-2n+2})=n.
\]
Since also $r^{\frac12 n(n-1)}=1$, our equation becomes
\[
W^2=\pm n. \tag{6}
\]
Thus in the case where $n$ is of the form $4\mu+1$,
\[
W=\pm\sqrt n,\qquad\text{and hence}\quad T=\pm\sqrt n,\quad U=0.
\]
In the other case, where $n$ is of the form $4\mu+3$,
\[
W=\pm i\sqrt n,\qquad\text{and therefore}\quad T=0,\quad U=\pm\sqrt n.
\]

\subsection*{14.}

The method of the preceding article assigns only the absolute value of the aggregates $T,U$, and leaves it ambiguous whether one should put $T$ in the former case and $U$ in the latter case equal to $+\sqrt n$ or to $-\sqrt n$. This, however, at least when $k=1$, can be decided from equation [5] as follows. Since, for $k=1$,
\[
r-r^{-1}=2i\sin\omega,\qquad
r^3-r^{-3}=2i\sin3\omega,\qquad
r^5-r^{-5}=2i\sin5\omega,\quad\text{etc.},
\]
that equation is transformed into
\[
W=(2i)^{\frac12(n-1)}\sin\omega\,\sin3\omega\,\sin5\omega\cdots\sin(n-2)\omega.
\]
Now, in the case where $n$ is of the form $4\mu+1$, the series of odd numbers
\[
1,3,5,7,\ldots,\tfrac12(n-3),\quad \tfrac12(n+1),\ldots,n-2
\]
contains $\tfrac14(n-1)$ numbers less than $\tfrac12 n$, and to these positive sines correspond. The remaining $\tfrac14(n-1)$ are greater than $\tfrac12 n$, and to these negative sines correspond. The product of all the sines is therefore to be taken as the product of a positive quantity with the multiplier $(-1)^{\frac14(n-1)}$. Hence $W$ will be the product of a positive real quantity with $i^{n-1}$, that is, with $1$, since $i^4=1$ and $n-1$ is divisible by $4$. Thus $W$ is a positive real quantity, and necessarily
\[
W=+\sqrt n,\qquad T=+\sqrt n.
\]

In the other case, where $n$ is of the form $4\mu+3$, the series of odd numbers
\[
1,3,5,7,\ldots,\tfrac12(n-1),\quad \tfrac12(n+3),\ldots,n-2
\]
has the first $\tfrac14(n+1)$ numbers less than $\tfrac12 n$ and the remaining $\tfrac14(n-3)$ greater. Hence among the sines of the arcs $\omega,3\omega,5\omega,\ldots,(n-2)\omega$, the negative ones number $\tfrac14(n-3)$, and so $W$ is the product of $i^{\frac12(n-1)}$ into a positive real quantity into $(-1)^{\frac14(n-3)}$. The third factor is $i^{\frac12(n-3)}$, which combined with the first produces $i^{n-2}=i$, since $i^{n-3}=1$. Therefore necessarily
\[
W=+i\sqrt n,\qquad\text{and}\qquad U=+\sqrt n.
\]

\subsection*{15.}

We shall now show how the same conclusions can be drawn from the progression considered in art. 9. In equation [4], write $-y^{-1}$ in place of $x^{\frac12}$. Then
\[
1-y^{-1}\frac{1-y^{-2m}}{1-y^{-2}}
+y^{-2}\frac{(1-y^{-2m})(1-y^{-2m+2})}{(1-y^{-2})(1-y^{-4})}
-y^{-3}\frac{(1-y^{-2m})(1-y^{-2m+2})(1-y^{-2m+4})}{(1-y^{-2})(1-y^{-4})(1-y^{-6})}
+\cdots
\]
down to the $(m+1)$-st term
\[
=(1-y^{-1})(1+y^{-2})(1-y^{-3})(1+y^{-4})\cdots(1\pm y^{-m}). \tag{7}
\]
If here $y$ is taken to be a proper root of $y^n-1=0$, say $r$, and if at the same time $m=n-1$, then
\[
\frac{1-y^{-2m}}{1-y^{-2}}=\frac{1-r^2}{1-r^{-2}}=-r^2,
\]
\[
\frac{1-y^{-2m+2}}{1-y^{-4}}=\frac{1-r^4}{1-r^{-4}}=-r^4,
\]
\[
\frac{1-y^{-2m+4}}{1-y^{-6}}=\frac{1-r^6}{1-r^{-6}}=-r^6,
\]
down to
\[
\frac{1-y^{-2}}{1-y^{-2m}}=\frac{1-r^{2n-2}}{1-r^{-2n+2}}=-r^{2n-2}.
\]
It should be noted that none of the denominators $1-r^{-2},1-r^{-4}$ etc. becomes zero. Hence equation [7] assumes the form
\[
1+r+r^4+r^9+\cdots+r^{(n-1)^2}
=(1-r^{-1})(1+r^{-2})(1-r^{-3})\cdots(1+r^{-n+1}).
\]
Multiplying, in the second member, the first term by the last, the second by the next-to-last, and so on, we have
\[
(1-r^{-1})(1+r^{-n+1})=r-r^{-1},
\]
\[
(1+r^{-2})(1-r^{-n+2})=r^{n-2}-r^{-n+2},
\]
\[
(1-r^{-3})(1+r^{-n+3})=r^3-r^{-3},
\]
\[
(1+r^{-4})(1-r^{-n+4})=r^{n-4}-r^{-n+4},\quad\text{etc.}
\]
From these partial products it is easy to see that the product formed is
\[
(r-r^{-1})(r^3-r^{-3})(r^5-r^{-5})\cdots(r^{n-4}-r^{-n+4})(r^{n-2}-r^{-n+2}),
\]
and therefore that it is
\[
=1+r+r^4+r^9+\cdots+r^{(n-1)^2}=W.
\]
This equation is identical with equation [5], derived in art. 12 from the first progression, and the remaining reasoning is then built in the same way as in arts. 13 and 14.

\subsection*{16.}

We pass to the other case, where $n$ is even. First let $n$ be of the form $4\mu+2$, that is, evenly divisible by $2$ but not by $4$. It is clear that the numbers
\[
\tfrac14 nn,\quad \bigl(\tfrac12 n+1\bigr)^2-1,\quad \bigl(\tfrac12 n+2\bigr)^2-4,\quad\text{etc.},
\]
or generally $\bigl(\tfrac12 n+\lambda\bigr)^2-\lambda^2$, when divided by $\tfrac12 n$, produce odd quotients, and hence become congruent modulo $n$ to $\tfrac12 n$. It follows that, if $r$ is a proper root of $x^n-1=0$, so that $r^{\frac12 n}=-1$, then
\[
r^{(\frac12 n)^2}=-1,
\]
\[
r^{(\frac12 n+1)^2}=-r,
\]
\[
r^{(\frac12 n+2)^2}=-r^4,
\]
\[
r^{(\frac12 n+3)^2}=-r^9,\quad\text{etc.}
\]
Hence, in the progression
\[
1+r+r^4+r^9+\cdots+r^{(n-1)^2},
\]
the term $r^{(\frac12 n)^2}$ cancels the first term, the following one cancels the second, and so on; therefore
\[
W=0,\qquad T=0,\qquad U=0.
\]

\subsection*{17.}

There remains the case where $n$ is of the form $4\mu$, that is, divisible by $4$. Here generally $\bigl(\tfrac12 n+\lambda\bigr)^2-\lambda^2$ is divisible by $n$, and hence
\[
r^{(\frac12 n+\lambda)^2}=r^{\lambda^2}.
\]
Thus, in the series
\[
1+r+r^4+r^9+\cdots+r^{(n-1)^2},
\]
the term $r^{(\frac12 n)^2}$ is equal to the first, the following one to the second, and so on, so that
\[
W=2\left(1+r+r^4+r^9+\cdots+r^{(\frac12 n-1)^2}\right).
\]
Now suppose that in equation [7] of art. 15 we put $m=\tfrac12 n-1$, and that for $y$ we take a proper root of $y^n-1=0$, namely $r$. Then, just as in art. 15, the equation takes the form
\[
1+r+r^4+\cdots+r^{(\frac12 n-1)^2}
=(1-r^{-1})(1+r^{-2})(1-r^{-3})\cdots(1-r^{-\frac12n+1}),
\]
or
\[
W=2(1-r^{-1})(1+r^{-2})(1-r^{-3})(1+r^{-4})\cdots(1-r^{-\frac12 n+1}). \tag{8}
\]
Furthermore, since $r^{\frac12 n}=-1$,
\[
1+r^{-2}=-r^{\frac12 n-2}(1-r^{-\frac12 n+2}),
\]
\[
1+r^{-4}=-r^{\frac12 n-4}(1-r^{-\frac12 n+4}),
\]
\[
1+r^{-6}=-r^{\frac12 n-6}(1-r^{-\frac12 n+6}),\quad\text{etc.}
\]
The product of the factors $-r^{\frac12n-2},-r^{\frac12n-4},-r^{\frac12n-6}$, and so on down to $-r^2$, is
\[
=(-1)^{\frac14n-1}r^{\frac1{16}nn-\frac14n}.
\]
The preceding equation can therefore also be written
\[
W=2(-1)^{\frac14 n-1}r^{\frac1{16}nn-\frac14n}
(1-r^{-1})(1-r^{-2})(1-r^{-3})(1-r^{-4})\cdots(1-r^{-\frac12n+1}).
\]
Since
\[
1-r^{-1}=-r^{-1}(1-r^{-n+1}),\quad
1-r^{-2}=-r^{-2}(1-r^{-n+2}),\quad
1-r^{-3}=-r^{-3}(1-r^{-n+3}),\quad\text{etc.},
\]
we have
\[
(1-r^{-1})(1-r^{-2})(1-r^{-3})\cdots(1-r^{-\frac12n+1})
\]
\[
=(-1)^{\frac12 n-1}r^{-\frac18 nn+\frac14 n}
(1-r^{-\frac12n-1})(1-r^{-\frac12n-2})(1-r^{-\frac12n-3})\cdots(1-r^{-n+1}),
\]
and consequently
\[
W=2(-1)^{\frac34 n-2}r^{-\frac1{16}nn}
(1-r^{-\frac12n-1})(1-r^{-\frac12n-2})(1-r^{-\frac12n-3})\cdots(1-r^{-n+1}).
\]
Multiplying this value of $W$ by the value found earlier, and adjoining the factor $1-r^{-\frac12n}$ on both sides, gives
\[
(1-r^{-\frac12 n})W^2
=4(-1)^{n-3}r^{-\frac14n}(1-r^{-1})(1-r^{-2})(1-r^{-3})\cdots(1-r^{-n+1}).
\]
But
\[
1-r^{-\frac12n}=2,\qquad (-1)^{n-3}=-1,\qquad r^{-\frac14n}=-r^{\frac14n},
\]
and
\[
(1-r^{-1})(1-r^{-2})(1-r^{-3})\cdots(1-r^{-n+1})=n.
\]
Thus at last one concludes that
\[
W^2=2r^{\frac14 n}n. \tag{9}
\]
It is now easy to see that $r^{\frac14 n}$ is either $+i$ or $-i$, according as $k$ is of the form $4\mu+1$ or $4\mu+3$. And since
\[
2i=(1+i)^2,\qquad -2i=(1-i)^2,
\]
in the case where $k$ is of the form $4\mu+1$ one has
\[
W=\pm(1+i)\sqrt n,\qquad\text{and hence}\quad T=U=\pm\sqrt n,
\]
while in the other case, where $k$ is of the form $4\mu+3$,
\[
W=\pm(1-i)\sqrt n,\qquad\text{and hence}\quad T=-U=\pm\sqrt n.
\]

\subsection*{18.}

The method of the preceding article has supplied the absolute values of the functions $T,U$ and has assigned the conditions under which equal or opposite signs are to be attributed to them; but it has not yet determined the signs themselves. We shall complete this for the case $k=1$ in the following manner.

Put
\[
\rho=\cos\tfrac12\omega+i\sin\tfrac12\omega,\qquad\text{so that}\quad r=\rho\rho.
\]
It is clear, since $\rho^n=-1$, that equation [8] can be displayed as
\[
W=2(1+\rho^{n-2})(1+\rho^{-4})(1+\rho^{n-6})(1+\rho^{-8})\cdots(1+\rho^{-n+4})(1+\rho^2),
\]
or, with the factors arranged in another order,
\[
W=2(1+\rho^2)(1+\rho^{-4})(1+\rho^6)(1+\rho^{-8})\cdots(1+\rho^{-n+4})(1+\rho^{n-2}).
\]
Now
\[
1+\rho^2=2\rho\cos\tfrac12\omega,
\]
\[
1+\rho^{-4}=2\rho^{-2}\cos\omega,
\]
\[
1+\rho^6=2\rho^3\cos\tfrac32\omega,
\]
\[
1+\rho^{-8}=2\rho^{-4}\cos2\omega,\quad\text{etc.},
\]
down to
\[
1+\rho^{-n+4}=2\rho^{-\frac12n+2}\cos\bigl(\tfrac14n-1\bigr)\omega,
\]
\[
1+\rho^{n-2}=2\rho^{\frac12n-1}\cos\bigl(\tfrac14n-\tfrac12\bigr)\omega.
\]
Consequently
\[
W=2^{\frac12n}\rho^{\frac14n}\cos\tfrac12\omega\cos\omega\cos\tfrac32\omega\cdots\cos\bigl(\tfrac14n-\tfrac12\bigr)\omega.
\]
The cosines entering this product are manifestly all positive, while the factor is
\[
\rho^{\frac14n}=\cos45^\circ+i\sin45^\circ=(1+i)\sqrt{\tfrac12}.
\]
Hence we conclude that $W$ is the product of $1+i$ and a positive real quantity. Therefore it must necessarily be
\[
W=(1+i)\sqrt n,\qquad T=+\sqrt n,\quad U=+\sqrt n.
\]

\subsection*{19.}

It will be worth the effort to collect here in one view all the summations so far developed. In general one has
\[
\begin{array}{c|c|c}
T= & U= & \text{according as } n\text{ is of the form}\\ \hline
\pm\sqrt n & \pm\sqrt n & 4\mu\\
\pm\sqrt n & 0 & 4\mu+1\\
0 & 0 & 4\mu+2\\
0 & \pm\sqrt n & 4\mu+3
\end{array}
\]
and in the case where $k$ is supposed equal to $1$, the positive sign must be attributed to the radical quantity. Thus the statements that we had observed by induction in art. 3 for prime values of $n$ have now been demonstrated with full rigour, and nothing remains except to determine the signs for arbitrary values of $k$ in all cases. Before we can undertake this in complete generality, however, we must first consider more closely the cases in which $n$ is a prime number or a power of a prime number.

\subsection*{20.}

First let $n$ be an odd prime number. It is clear from what we explained in art. 10 that
\[
W=1+2\sum r^a=1+2\sum R^{ak},
\]
if one puts
\[
R=\cos\omega+i\sin\omega,
\]
where $a$, as there, denotes indefinitely all quadratic residues of $n$ contained between $1$ and $n-1$. If, likewise, by $b$ we denote indefinitely all quadratic non-residues between the same limits, it is easily seen that all the numbers $ak$ become congruent modulo $n$ either to all the $a$ or to all the $b$, without regard to order, according as $k$ is a residue or a non-residue. Therefore in the former case
\[
W=1+2\sum R^a=1+R+R^4+R^9+\cdots+R^{(n-1)^2},
\]
and hence
\[
W=+\sqrt n,\quad\text{if }n\text{ is of the form }4\mu+1,\qquad
W=+i\sqrt n,\quad\text{if }n\text{ is of the form }4\mu+3.
\]
In the other case, where $k$ is a non-residue of $n$, one has
\[
W=1+2\sum R^b.
\]
Since all the $a,b$ plainly fill out the complete set of numbers $1,2,3,\ldots$, and hence
\[
\sum R^a+\sum R^b=R+R^2+R^3+\cdots+R^{n-1}=-1,
\]
it follows that
\[
W=-1-2\sum R^a=-(1+R+R^4+R^9+\cdots+R^{(n-1)^2}),
\]
and therefore
\[
W=-\sqrt n,\quad\text{if }n\text{ is of the form }4\mu+1,\qquad
W=-i\sqrt n,\quad\text{if }n\text{ is of the form }4\mu+3.
\]
Thus we conclude:

\emph{first}, if $n$ is of the form $4\mu+1$ and $k$ is a quadratic residue of $n$,
\[
T=+\sqrt n,\qquad U=0;
\]
\emph{second}, if $n$ is of the form $4\mu+1$ and $k$ is a quadratic non-residue of $n$,
\[
T=-\sqrt n,\qquad U=0;
\]
\emph{third}, if $n$ is of the form $4\mu+3$ and $k$ is a quadratic residue of $n$,
\[
T=0,\qquad U=+\sqrt n;
\]
\emph{fourth}, if $n$ is of the form $4\mu+3$ and $k$ is a quadratic non-residue of $n$,
\[
T=0,\qquad U=-\sqrt n.
\]



\subsection*{21.}

In the second place, let $n$ be the square, or a higher power, of an odd prime number $p$, and put
\[
n=p^{2\chi}q,
\]
so that $q$ is either $1$ or $p$. First it is useful to observe that, if $\lambda$ is any integer not divisible by $p^\chi$, then
\[
\begin{aligned}
&r^{\lambda\lambda}+r^{(\lambda+p^\chi q)^2}+r^{(\lambda+2p^\chi q)^2}
+r^{(\lambda+3p^\chi q)^2}+\cdots+r^{(\lambda+n-p^\chi q)^2}\\
&\quad
=r^{\lambda\lambda}\{1+r^{2\lambda p^\chi q}+r^{4\lambda p^\chi q}
+r^{6\lambda p^\chi q}+\cdots+r^{2\lambda(n-p^\chi q)}\}
=\frac{r^{\lambda\lambda}(1-r^{2\lambda n})}{1-r^{2\lambda p^\chi q}}=0.
\end{aligned}
\]
It is therefore easy to see that
\[
W=1+r^{p^{2\chi}}+r^{4p^{2\chi}}+r^{9p^{2\chi}}+\cdots+r^{(n-p^\chi)^2}.
\]
For the remaining terms of the progression
\[
1+r+r^4+r^9+\cdots+r^{(n-1)^2}
\]
can be distributed into $(p^\chi-1)q$ partial progressions, each of $p^\chi$ terms, and by the transformation just given they produce vanishing sums.

It follows that, in the case $q=1$, that is, when $n$ is a power of a prime with even exponent,
\[
W=p^\chi=+\sqrt n,\qquad\text{and hence}\qquad T=+\sqrt n,\quad U=0.
\]
On the other hand, in the case $q=p$, that is, when $n$ is a power of a prime with odd exponent, put $r^{p^{2\chi}}=\rho$. Then $\rho$ is a proper root of the equation $x^p-1=0$, namely
\[
\rho=\cos\frac{k}{p}360^\circ+i\sin\frac{k}{p}360^\circ,
\]
and then
\[
\begin{aligned}
W&=1+\rho+\rho^4+\rho^9+\cdots+\rho^{(p^{\chi+1}-1)^2}\\
&=p^\chi(1+\rho+\rho^4+\rho^9+\cdots+\rho^{(p-1)^2}).
\end{aligned}
\]
The sum of the series $1+\rho+\rho^4+\rho^9+\cdots+\rho^{(p-1)^2}$ is determined by the preceding article; hence it follows directly that
\[
W=\pm\sqrt n=T,\quad\text{if }p\text{ is of the form }4\mu+1,
\]
\[
W=\pm i\sqrt n=iU,\quad\text{if }p\text{ is of the form }4\mu+3,
\]
where the positive or negative sign holds according as $k$ is a residue or a non-residue of $p$.

\subsection*{22.}

From the material set out in Articles 20 and 21 one also easily derives the following proposition, which will be of notable use below. Put
\[
W'=1+r^h+r^{4h}+r^{9h}+\cdots+r^{h(n-1)^2},
\]
where $h$ denotes any integer not divisible by $p$. In the case where $n=p$, or where $n$ is a power of $p$ with odd exponent, one has
\[
W'=W,\quad\text{if }h\text{ is a quadratic residue of }p,
\]
\[
W'=-W,\quad\text{if }h\text{ is a quadratic non-residue of }p.
\]
For $W'$ plainly arises from $W$ if $kh$ is substituted for $k$; in the former case $k$ and $kh$ will be alike, and in the latter case unlike, in so far as they are residues or non-residues of $p$.

In the case where $n$ is a power of $p$ with even exponent, however, one plainly has $W=+\sqrt n$, and therefore always $W'=W$.

\subsection*{23.}

In Articles 20, 21, and 22 we considered odd primes and their powers. It remains to consider the case where $n$ is a power of two.

For $n=2$ one plainly has $W=1+r=0$.

For $n=4$ there results
\[
W=1+r+r^4+r^9=2+2r;
\]
hence
\[
W=2+2i,\quad\text{whenever }k\text{ is of the form }4\mu+1,
\]
and
\[
W=2-2i,\quad\text{whenever }k\text{ is of the form }4\mu+3.
\]

For $n=8$ we have
\[
W=1+r+r^4+r^9+r^{16}+r^{25}+r^{36}+r^{49}=2+4r+2r^4=4r.
\]
Thus
\[
\begin{array}{ll}
W=(1+i)\sqrt8, & \text{whenever }k\text{ is of the form }8\mu+1,\\
W=(-1+i)\sqrt8, & \text{whenever }k\text{ is of the form }8\mu+3,\\
W=(-1-i)\sqrt8, & \text{whenever }k\text{ is of the form }8\mu+5,\\
W=(1-i)\sqrt8, & \text{whenever }k\text{ is of the form }8\mu+7.
\end{array}
\]

If $n$ is a higher power of two, put
\[
n=2^{2\chi}q,
\]
where $q$ is either $1$ or $2$, and $\chi$ is greater than $1$. First it should be observed that, if $\lambda$ is any integer not divisible by $2^{\chi-1}$, then
\[
\begin{aligned}
&r^{\lambda\lambda}+r^{(\lambda+2^\chi q)^2}+r^{(\lambda+2\cdot2^\chi q)^2}
+r^{(\lambda+3\cdot2^\chi q)^2}+\cdots+r^{(\lambda+n-2^\chi q)^2}\\
&\quad
=r^{\lambda\lambda}\{1+r^{2^{\chi+1}\lambda q}
+r^{2\cdot2^{\chi+1}\lambda q}
+r^{3\cdot2^{\chi+1}\lambda q}
+\cdots+r^{(2n-2^{\chi+1}q)\lambda}\}\\
&\quad
=\frac{r^{\lambda\lambda}(1-r^{2\lambda n})}{1-r^{2^{\chi+1}\lambda q}}=0.
\end{aligned}
\]
It is therefore easy to see that
\[
W=1+r^{2^{2\chi-2}}+r^{4\cdot 2^{2\chi-2}}+r^{9\cdot 2^{2\chi-2}}+\cdots+r^{(n-2^{\chi-1})^2}.
\]
Put
\[
r^{2^{2\chi-2}}=\rho.
\]
Then $\rho$ is a root of the equation $x^{4q}-1=0$, namely
\[
\rho=\cos\frac{k}{4q}360^\circ+i\sin\frac{k}{4q}360^\circ;
\]
next
\[
\begin{aligned}
W&=1+\rho+\rho^4+\rho^9+\cdots+\rho^{(2^{\chi+1}q-1)^2}\\
&=2^{\chi-1}(1+\rho+\rho^4+\rho^9+\cdots+\rho^{(4q-1)^2}).
\end{aligned}
\]
The sum of the series $1+\rho+\rho^4+\rho^9+\cdots+\rho^{(4q-1)^2}$ is determined by what we have explained for the cases $n=4$ and $n=8$. Thus, in the case $q=1$, that is, when $n$ is a power of $4$, we obtain
\[
W=(1+i)2^\chi=(1+i)\sqrt n,\quad\text{if }k\text{ is of the form }4\mu+1,
\]
\[
W=(1-i)2^\chi=(1-i)\sqrt n,\quad\text{if }k\text{ is of the form }4\mu+3.
\]
These are exactly the formulas already given for $n=4$. In the case $q=2$, that is, when $n$ is a power of two with odd exponent greater than $3$, we obtain
\[
\begin{array}{ll}
W=(1+i)2^\chi\sqrt2=(1+i)\sqrt n, & \text{if }k\text{ is of the form }8\mu+1,\\
W=(-1+i)2^\chi\sqrt2=(-1+i)\sqrt n, & \text{if }k\text{ is of the form }8\mu+3,\\
W=(-1-i)2^\chi\sqrt2=(-1-i)\sqrt n, & \text{if }k\text{ is of the form }8\mu+5,\\
W=(1-i)2^\chi\sqrt2=(1-i)\sqrt n, & \text{if }k\text{ is of the form }8\mu+7.
\end{array}
\]
These too agree completely with the formulas given for $n=8$.

\subsection*{24.}

It will also be worth determining the ratio of the sum of the progression
\[
W'=1+r^h+r^{4h}+r^{9h}+\cdots+r^{h(n-1)^2}
\]
to $W$, where $h$ denotes any odd integer. Since $W'$ arises from $W$ by changing $k$ into $kh$, the value of $W'$ will depend on the form of the number $kh$ just as $W$ depends on the form of $k$. Put
\[
\frac{W'}{W}=l.
\]
Then it is clear that:

\emph{I.} in the case where $n=4$, or a higher power of two with even exponent,
\[
\begin{array}{ll}
l=1, & \text{if }h\text{ is of the form }4\mu+1,\\
l=-i, & \text{if }h\text{ is of the form }4\mu+3\text{ and }k\text{ is of the form }4\mu+1,\\
l=+i, & \text{if }h\text{ is of the form }4\mu+3\text{ and }k\text{ is of the same form.}
\end{array}
\]

\emph{II.} in the case where $n=8$, or a higher power of two with odd exponent,
\[
\begin{array}{ll}
l=1, & \text{if }h\text{ is of the form }8\mu+1,\\
l=-1, & \text{if }h\text{ is of the form }8\mu+5,\\
l=+i, & \text{if either }h\text{ is of the form }8\mu+3\text{ and }k\text{ of the form }4\mu+1,\\
&\text{or }h\text{ is of the form }8\mu+7\text{ and }k\text{ of the form }4\mu+3,\\
l=-i, & \text{if either }h\text{ is of the form }8\mu+3\text{ and }k\text{ of the form }4\mu+3,\\
&\text{or }h\text{ is of the form }8\mu+7\text{ and }k\text{ of the form }4\mu+1.
\end{array}
\]

By the preceding articles the determination of the sum $W$ is complete for the cases where $n$ is a prime number or a power of a prime number. It remains to settle the cases where $n$ is composed of several prime numbers; the following theorem will pave the way.

\subsection*{25.}

\textsc{Theorem.} Let $n$ be the product of two positive relatively prime integers $a,b$, and put
\[
P=1+r^{aa}+r^{4aa}+r^{9aa}+\cdots+r^{(b-1)^2aa},
\]
\[
Q=1+r^{bb}+r^{4bb}+r^{9bb}+\cdots+r^{(a-1)^2bb}.
\]
I say that
\[
W=PQ.
\]

\emph{Proof.} Let $\alpha$ range over the numbers $0,1,2,3,\ldots,a-1$, let $\beta$ range over $0,1,2,3,\ldots,b-1$, and let $\nu$ range over $0,1,2,3,\ldots,n-1$. It is then clear that
\[
P=\sum r^{aa\beta\beta},\qquad Q=\sum r^{bb\alpha\alpha},\qquad W=\sum r^{\nu\nu}.
\]
Hence
\[
PQ=\sum r^{aa\beta\beta+bb\alpha\alpha},
\]
where all values of $\alpha$ and $\beta$ are substituted, combined with one another in all ways. Since $2ab\alpha\beta=2\alpha\beta n$, this gives
\[
PQ=\sum r^{(a\beta+b\alpha)^2}.
\]
But it is immediate that the individual values of $a\beta+b\alpha$ are all distinct from one another and equal to some value of $\nu$. Therefore
\[
PQ=\sum r^{\nu\nu}=W.
\]
It should also be noted that $r^{aa}$ is a proper root of the equation $x^b-1=0$, and $r^{bb}$ a proper root of $x^a-1=0$.

\subsection*{26.}

Next let $n$ be the product of three mutually prime numbers $a,b,c$. If $bc=b'$, then $a$ and $b'$ are also relatively prime; hence $W$ is the product of the two factors
\[
1+r^{aa}+r^{4aa}+r^{9aa}+\cdots+r^{(b'-1)^2aa},
\]
\[
1+r^{b'b'}+r^{4b'b'}+r^{9b'b'}+\cdots+r^{(a-1)^2b'b'}.
\]
Since $r^{aa}$ is a proper root of the equation $x^{bc}-1=0$, the first factor itself is the product of
\[
1+\rho^{bb}+\rho^{4bb}+\rho^{9bb}+\cdots+\rho^{(c-1)^2bb},
\]
\[
1+\rho^{cc}+\rho^{4cc}+\rho^{9cc}+\cdots+\rho^{(b-1)^2cc},
\]
if $r^{aa}=\rho$. Hence $W$ is the product of the three factors
\[
1+r^{bbcc}+r^{4bbcc}+r^{9bbcc}+\cdots+r^{(a-1)^2bbcc},
\]
\[
1+r^{aacc}+r^{4aacc}+r^{9aacc}+\cdots+r^{(b-1)^2aacc},
\]
\[
1+r^{aabb}+r^{4aabb}+r^{9aabb}+\cdots+r^{(c-1)^2aabb},
\]
where $r^{bbcc}$, $r^{aacc}$, $r^{aabb}$ are respectively proper roots of
\[
x^a-1=0,\qquad x^b-1=0,\qquad x^c-1=0.
\]

\subsection*{27.}

It follows easily in general that, if $n$ is the product of any number of mutually prime factors $a,b,c$ etc., then $W$ becomes the product of the same number of factors, namely
\[
1+r^{\frac{nn}{aa}}+r^{\frac{4nn}{aa}}+r^{\frac{9nn}{aa}}+\cdots+r^{\frac{(a-1)^2nn}{aa}},
\]
\[
1+r^{\frac{nn}{bb}}+r^{\frac{4nn}{bb}}+r^{\frac{9nn}{bb}}+\cdots+r^{\frac{(b-1)^2nn}{bb}},
\]
\[
1+r^{\frac{nn}{cc}}+r^{\frac{4nn}{cc}}+r^{\frac{9nn}{cc}}+\cdots+r^{\frac{(c-1)^2nn}{cc}},
\]
and so on, where $r^{\frac{nn}{aa}},r^{\frac{nn}{bb}},r^{\frac{nn}{cc}}$ etc. are proper roots of
\[
x^a-1=0,\qquad x^b-1=0,\qquad x^c-1=0,\quad\text{etc.}
\]

\subsection*{28.}

From these principles the passage to the complete determination of $W$, for any value of $n$, is now obvious. Decompose $n$ into factors $a,b,c$ etc. which are either distinct prime numbers or powers of distinct prime numbers, and put
\[
r^{\frac{nn}{aa}}=A,\qquad r^{\frac{nn}{bb}}=B,\qquad r^{\frac{nn}{cc}}=C,\quad\text{etc.}
\]
Then $A,B,C$ etc. will be proper roots of
\[
x^a-1=0,\qquad x^b-1=0,\qquad x^c-1=0,\quad\text{etc.},
\]
and $W$ will be the product of the factors
\[
1+A+A^4+A^9+\cdots+A^{(a-1)^2},
\]
\[
1+B+B^4+B^9+\cdots+B^{(b-1)^2},
\]
\[
1+C+C^4+C^9+\cdots+C^{(c-1)^2}.
\]
These individual factors can be determined by what we have shown in Articles 20, 21, and 23, whence the value of the product also becomes known. It will not be useless to collect here, in one view, the rules for determining those factors. Since the root
\[
A=\cos\frac{kn}{a}\cdot\frac{360^\circ}{a}+i\sin\frac{kn}{a}\cdot\frac{360^\circ}{a},
\]
the aggregate
\[
1+A+A^4+A^9+\cdots+A^{(a-1)^2},
\]
which we shall denote by $L$, will be determined by the number $\frac{kn}{a}$ just as, in our general investigation, $W$ is determined by $k$. Twelve cases are now to be distinguished.

\begin{enumerate}[label=\Roman*.,leftmargin=2.5em]
\item If $a$ is a prime number of the form $4\mu+1$, say $p$, or a power of such a prime with odd exponent, and at the same time $\frac{kn}{a}$ is a quadratic residue of $p$, then $L=+\sqrt a$.
\item If, the remaining assumptions being as above, $\frac{kn}{a}$ is a quadratic non-residue of $p$, then $L=-\sqrt a$.
\item If $a$ is a prime number of the form $4\mu+3$, say $p$, or a power of such a prime with odd exponent, and at the same time $\frac{kn}{a}$ is a quadratic residue of $p$, then $L=+i\sqrt a$.
\item If, the remaining assumptions being as in III, $\frac{kn}{a}$ is a quadratic non-residue of $p$, then $L=-i\sqrt a$.
\item If $a$ is a square, or a higher power of an odd prime with even exponent, then $L=+\sqrt a$.
\item If $a=2$, then $L=0$.
\item If $a=4$, or a higher power of two with even exponent, and at the same time $\frac{kn}{a}$ is of the form $4\mu+1$, then $L=(1+i)\sqrt a$.
\item If, the remaining assumptions being as in VII, $\frac{kn}{a}$ is of the form $4\mu+3$, then $L=(1-i)\sqrt a$.
\item If $a=8$, or a higher power of two with odd exponent, and at the same time $\frac{kn}{a}$ is of the form $8\mu+1$, then $L=(1+i)\sqrt a$.
\item If, the remaining assumptions being as in IX, $\frac{kn}{a}$ is of the form $8\mu+3$, then $L=(-1+i)\sqrt a$.
\item If, the remaining assumptions being as above, $\frac{kn}{a}$ is of the form $8\mu+5$, then $L=(-1-i)\sqrt a$.
\item If, the remaining assumptions being as above, $\frac{kn}{a}$ is of the form $8\mu+7$, then $L=(1-i)\sqrt a$.
\end{enumerate}

\subsection*{29.}

For example, let
\[
n=2520=8\cdot9\cdot5\cdot7,\qquad k=13.
\]
Then
\[
\begin{array}{ll}
\text{for }a=8,\text{ by case XII,} & L=(1-i)\sqrt8,\\
\text{for the factor }9,\text{ by case V, the corresponding sum is} & =\sqrt9,\\
\text{for the factor }5,\text{ by case II, the corresponding sum is} & =-\sqrt5,\\
\text{for the factor }7,\text{ by case III, the corresponding sum is} & =+i\sqrt7.
\end{array}
\]
Hence
\[
W=(1-i)\cdot(-i)\cdot\sqrt{2520}=(-1-i)\sqrt{2520}.
\]
For the same value of $n$, let $k=1$. Then there corresponds
\[
\begin{array}{ll}
\text{to the factor }8 & \text{the sum }(-1+i)\sqrt8,\\
\text{to the factor }9 & \text{the sum }\sqrt9,\\
\text{to the factor }5 & \text{the sum }\sqrt5,\\
\text{to the factor }7 & \text{the sum }-i\sqrt7.
\end{array}
\]
Thus the product formed is
\[
W=(1+i)\sqrt{2520}.
\]

\subsection*{30.}

Another method of determining the sum $W$ in general is derived from what was set out in Articles 22 and 24. Put
\[
\cos\omega+i\sin\omega=\rho,
\]
and
\[
\rho^{\frac{nn}{aa}}=\alpha,\qquad
\rho^{\frac{nn}{bb}}=\beta,\qquad
\rho^{\frac{nn}{cc}}=\gamma,\quad\text{etc.}
\]
so that
\[
r=\rho^k,\qquad A=\alpha^k,\qquad B=\beta^k,\qquad C=\gamma^k,\quad\text{etc.}
\]
Then
\[
1+\rho+\rho^4+\rho^9+\cdots+\rho^{(n-1)^2}
\]
is the product of the factors
\[
1+\alpha+\alpha^4+\alpha^9+\cdots+\alpha^{(a-1)^2},
\]
\[
1+\beta+\beta^4+\beta^9+\cdots+\beta^{(b-1)^2},
\]
\[
1+\gamma+\gamma^4+\gamma^9+\cdots+\gamma^{(c-1)^2},\quad\text{etc.}
\]
Hence $W$ is the product of the factors
\[
w=1+\rho+\rho^4+\rho^9+\cdots+\rho^{(n-1)^2},
\]
\[
\mathfrak A=
\frac{1+A+A^4+A^9+\cdots+A^{(a-1)^2}}
{1+\alpha+\alpha^4+\alpha^9+\cdots+\alpha^{(a-1)^2}},
\]
\[
\mathfrak B=
\frac{1+B+B^4+B^9+\cdots+B^{(b-1)^2}}
{1+\beta+\beta^4+\beta^9+\cdots+\beta^{(b-1)^2}},
\]
\[
\mathfrak C=
\frac{1+C+C^4+C^9+\cdots+C^{(c-1)^2}}
{1+\gamma+\gamma^4+\gamma^9+\cdots+\gamma^{(c-1)^2}}.
\]
The first factor $w$ has already been determined by the investigations above (Article 19). The remaining factors $\mathfrak A,\mathfrak B,\mathfrak C$ etc. are obtained from the formulas of Articles 22 and 24, which we collect here again so that all may stand together. Twelve cases are to be distinguished:

\begin{enumerate}[label=\Roman*.,leftmargin=2.5em]
\item If $a$ is an odd prime number $p$, or a power of such a prime with odd exponent, and $k$ is a quadratic residue of $p$, then the corresponding factor is $\mathfrak A=+1$.
\item If, the remaining assumptions being as above, $k$ is a quadratic non-residue of $p$, then $\mathfrak A=-1$.
\item If $a$ is the square of an odd prime, or a higher power of it with even exponent, then $\mathfrak A=+1$.
\item If $a=4$, or a higher power of two with even exponent, and at the same time $k$ is of the form $4\mu+1$, then $\mathfrak A=+1$.
\item If, the remaining assumptions being as in IV, $k$ is of the form $4\mu+3$, and $\frac{n}{a}$ of the form $4\mu+1$, then $\mathfrak A=-i$.
\item If, the remaining assumptions being as in IV, $k$ is of the form $4\mu+3$, and $\frac{n}{a}$ of the form $4\mu+3$, then $\mathfrak A=+i$.
\item If $a=8$, or a higher power of two with odd exponent, and $k$ is of the form $8\mu+1$, then $\mathfrak A=+1$.
\item If, the remaining assumptions being as in VII, $k$ is of the form $8\mu+5$, then $\mathfrak A=-1$.
\item If, the remaining assumptions being as in VII, $k$ is of the form $8\mu+3$, and $\frac{n}{a}$ of the form $4\mu+1$, then $\mathfrak A=+i$.
\item If, the remaining assumptions being as in VII, $k$ is of the form $8\mu+3$, and $\frac{n}{a}$ of the form $4\mu+3$, then $\mathfrak A=-i$.
\item If, the remaining assumptions being as in VII, $k$ is of the form $8\mu+7$, and $\frac{n}{a}$ of the form $4\mu+1$, then $\mathfrak A=-i$.
\item If, the remaining assumptions being as in VII, $k$ is of the form $8\mu+7$, and $\frac{n}{a}$ of the form $4\mu+3$, then $\mathfrak A=+i$.
\end{enumerate}

We pass over the case $a=2$; here $\mathfrak A$ would be $0/0$, that is, indeterminate, but then always $W=0$.

The remaining factors $\mathfrak B,\mathfrak C$ etc. depend on $b,c$ etc. in the same way as $\mathfrak A$ depends on $a$, insofar as those quantities enter into their determination.

\subsection*{31.}

By this second method, the first example of Article 29 is as follows:
\[
\begin{array}{ll}
\text{the factor }w\text{ becomes} & =(1+i)\sqrt{2520},\\
\text{for }a=8\text{ the corresponding factor }\mathfrak A\text{ is, by case VIII,} & =-1,\\
\text{to the second factor of }n,\text{ namely }9,\text{ there corresponds the factor} & +1\text{ (by case III)},\\
\text{to the factor }5\text{ there corresponds the factor} & -1\text{ (by case II)},\\
\text{to the factor }7\text{ there corresponds the factor} & -1\text{ (by case II)}.
\end{array}
\]
Thus the product formed is
\[
W=(-1-i)\sqrt{2520},
\]
as in Article 29.

\subsection*{32.}

Since the value of $W$ can be determined by two methods, one of which rests on the relations of the numbers
\[
\frac{nk}{a},\quad \frac{nk}{b},\quad \frac{nk}{c},\quad\text{etc.}
\]
to the numbers $a,b,c$ etc., while the other depends on the relations of $k$ itself to the numbers $a,b,c$ etc., there must be a certain conditional connection among all these relations, so that any one can be determined from the others. Suppose that all the numbers $a,b,c$ etc. are odd primes and take $k=1$. Distribute the factors $a,b,c$ etc. into two classes: one containing those of the form $4\mu+1$, denoted by $p,p',p''$ etc.; the other consisting of those of the form $4\mu+3$, denoted by $q,q',q''$ etc. Let the number of the latter be denoted by $m$. With this done, we first observe that $n$ is of the form $4\mu+1$ if $m$ is even, including also the case in which the second class has no factors at all, that is, $m=0$; on the other hand $n$ is of the form $4\mu+3$ if $m$ is odd.

Now the determination of $W$ by the first method is carried out as follows. Let the numbers $P,P',P''$ etc., $Q,Q',Q''$ etc. depend on the relations of the numbers
\[
\frac{n}{p},\frac{n}{p'},\frac{n}{p''},\ldots,\qquad
\frac{n}{q},\frac{n}{q'},\frac{n}{q''},\ldots
\]
to the numbers $p,p',p''$ etc., $q,q',q''$ etc., respectively, so that
\[
P=+1,\quad\text{if }\frac{n}{p}\text{ is a quadratic residue of }p,
\]
\[
P=-1,\quad\text{if }\frac{n}{p}\text{ is a quadratic non-residue of }p,
\]
and similarly for the rest. Then $W$ is the product of the factors
\[
P\sqrt p,\quad P'\sqrt{p'},\quad P''\sqrt{p''},\ldots,\qquad
iQ\sqrt q,\quad iQ'\sqrt{q'},\quad iQ''\sqrt{q''},\ldots,
\]
and hence
\[
W=PP'P''\cdots QQ'Q''\cdots\, i^m\sqrt n.
\]
By the second method, or rather directly by the rules of Article 19,
\[
W=+\sqrt n,\quad\text{if }n\text{ is of the form }4\mu+1,\text{ or equivalently if }m\text{ is even},
\]
\[
W=+i\sqrt n,\quad\text{if }n\text{ is of the form }4\mu+3,\text{ or if }m\text{ is odd}.
\]
Both cases may be comprised in the single formula
\[
W=i^{mm}\sqrt n.
\]
It follows that
\[
PP'P''\cdots QQ'Q''\cdots=i^{mm-m}.
\]
But $i^{mm-m}$ is $1$ whenever $m$ is of the form $4\mu$ or $4\mu+1$, and $-1$ whenever $m$ is of the form $4\mu+2$ or $4\mu+3$. Thus we derive the following very elegant theorem.

\textsc{Theorem.} Let $a,b,c$ etc. denote distinct positive odd primes whose product is put equal to $n$, and suppose that among them $m$ are of the form $4\mu+3$, the rest being of the form $4\mu+1$. Then the number of those among $a,b,c$ etc. for which $\frac{n}{a},\frac{n}{b},\frac{n}{c}$ etc. are respectively non-residues is even when $m$ is of the form $4\mu$ or $4\mu+1$, and odd when $m$ is of the form $4\mu+2$ or $4\mu+3$.

Thus, for example, putting $a=3,b=5,c=7,d=11$, we have three numbers of the form $4\mu+3$, namely $3,7,11$. But
\[
5\cdot7\cdot11\;R3,\quad
3\cdot7\cdot11\;R5,\quad
3\cdot5\cdot11\;R7,\quad
3\cdot5\cdot7\;N11,
\]
so that the single $\frac{n}{d}$ is a non-residue of $d$.

\subsection*{33.}

The celebrated fundamental theorem concerning quadratic residues is nothing other than a special case of the theorem just developed. Namely, by limiting the number of numbers $a,b,c$ etc. to two, it is clear that if only one of them, or neither, is of the form $4\mu+3$, either both $aRb,bRa$ must hold or both $aNb,bNa$ must hold. On the other hand, if both are of the form $4\mu+3$, one of them must be a non-residue of the other, while the latter is a residue of the former. Thus we have a fourth demonstration of this most important theorem: the first and second demonstrations were given in the \emph{Disquisitiones Arithmeticae}, the third recently in a special memoir (\emph{Comment. T. XVI}); two further demonstrations, again relying on entirely different principles, we shall present later. It is indeed highly remarkable that this beautiful theorem, which at first so obstinately eluded every attempt, could afterwards be approached by so many routes, all widely distant from one another.

\subsection*{34.}

The other theorems, which form as it were supplements to the fundamental theorem and by which one recognizes the prime numbers whose residues or non-residues are $-1,+2$, and $-2$, can also be derived from the same principles. We begin with the residue $+2$.

Put $n=8a$, where $a$ is a prime number, and take $k=1$. By the method of Article 28, $W$ is the product of two factors. One of them is $+\sqrt a$, or $+i\sqrt a$, if $8$, or equivalently $2$, is a quadratic residue of $a$; on the other hand it is $-\sqrt a$ or $-i\sqrt a$ if $2$ is a non-residue of $a$. The second factor is
\[
\begin{array}{ll}
(1+i)\sqrt8, & \text{if }a\text{ is of the form }8\mu+1,\\
(-1+i)\sqrt8, & \text{if }a\text{ is of the form }8\mu+3,\\
(-1-i)\sqrt8, & \text{if }a\text{ is of the form }8\mu+5,\\
(1-i)\sqrt8, & \text{if }a\text{ is of the form }8\mu+7.
\end{array}
\]
But by Article 18 one always has
\[
W=(1+i)\sqrt n.
\]
Dividing this value by the four values of the second factor, it is clear that the first factor must be
\[
\begin{array}{ll}
+\sqrt a, & \text{if }a\text{ is of the form }8\mu+1,\\
-i\sqrt a, & \text{if }a\text{ is of the form }8\mu+3,\\
-\sqrt a, & \text{if }a\text{ is of the form }8\mu+5,\\
+i\sqrt a, & \text{if }a\text{ is of the form }8\mu+7.
\end{array}
\]
It follows at once that in the first and fourth cases $2$ must be a residue of $a$, while in the second and third cases it is a non-residue.

\subsection*{35.}

The prime numbers for which $-1$ is a residue or a non-residue are easily recognized with the aid of the following theorem, which is also quite noteworthy in itself.

\textsc{Theorem.} The product of the two factors
\[
W'=1+r^{-1}+r^{-4}+\cdots+r^{-(n-1)^2},
\]
\[
W=1+r+r^4+\cdots+r^{(n-1)^2}
\]
is equal to $n$ if $n$ is odd; to $0$ if $n$ is singly even; and to $2n$ if $n$ is doubly even.

\emph{Proof.} Since plainly
\[
W=r+r^4+r^9+\cdots+r^{nn}
=r^4+r^9+\cdots+r^{(n+1)^2}
=r^9+\cdots+r^{(n+2)^2}\quad\text{etc.},
\]
the product $WW'$ can also be exhibited as
\[
\begin{aligned}
&1+r+r^4+r^9+\cdots+r^{(n-1)^2}\\
&+r^{-1}(r+r^4+r^9+r^{16}+\cdots+r^{nn})\\
&+r^{-4}(r^4+r^9+r^{16}+r^{25}+\cdots+r^{(n+1)^2})\\
&+r^{-9}(r^9+r^{16}+r^{25}+r^{36}+\cdots+r^{(n+2)^2})\\
&+\text{etc.}\\
&+r^{-(n-1)^2}(r^{(n-1)^2}+r^{nn}+r^{(n+1)^2}+r^{(n+2)^2}+\cdots+r^{(2n-2)^2}).
\end{aligned}
\]
Summed vertically, this aggregate produces
\[
\begin{aligned}
&n\\
&+r(1+r^n+r^{2n}+r^{3n}+\cdots+r^{nn-n})\\
&+r^4(1+r^{2n}+r^{4n}+r^{6n}+\cdots+r^{2(n-1)^2})\\
&+r^9(1+r^{3n}+r^{6n}+r^{9n}+\cdots+r^{3nn-3n})\\
&+\text{etc.}\\
&+r^{(n-1)^2}(1+r^{2n-2}+r^{4n-4}+r^{6n-6}+\cdots+r^{2(n-1)^2}).
\end{aligned}
\]
If $n$ is odd, all parts of this aggregate except the first, $n$, are zero: the second, for example, is plainly
\[
\frac{r(1-r^{2n})}{1-rr},
\]
the third
\[
\frac{r^4(1-r^{4n})}{1-r^4},
\]
and so on. But when $n$ is even one must also except the part
\[
r^{\frac14 nn}(1+r^n+r^{2n}+r^{3n}+\cdots+r^{nn-n}),
\]
which becomes $nr^{\frac14 nn}$. In the former case, therefore, $WW'=n$; in the latter it is
\[
n+nr^{\frac14 nn}.
\]
But $r^{\frac14 nn}=+1$ if $n$ is doubly even, and then $WW'=2n$; while $r^{\frac14 nn}=-1$ if $n$ is singly even, and then $WW'=0$. Q.E.D.

\subsection*{36.}

By Article 22 we know that, if $n$ is an odd prime number,
\[
\frac{W'}{W}
\]
becomes $+1$ or $-1$ according as $-1$ is a residue or a non-residue of $n$. Hence in the former case one must have
\[
W^2=+n,
\]
and in the latter
\[
W^2=-n.
\]
Therefore, by Article 13, we conclude that the former case can occur only when $n$ is of the form $4\mu+1$, and the latter when $n$ is of the form $4\mu+3$.

Finally, from combining the conditions found for the residues $+2$ and $-1$, it follows directly that $-2$ is a residue of every prime number of the form $8\mu+1$ or $8\mu+3$, and a non-residue of every prime number of the form $8\mu+5$ or $8\mu+7$.



\section*{New Demonstrations and Extensions of the Fundamental Theorem in the Theory of Quadratic Residues}

\begin{center}
\textsc{By Carl Friedrich Gauss}\\
\textsc{Presented to the Royal Society of Sciences, 10 February 1817}\\[2pt]
\emph{Commentationes Societatis Regiae Scientiarum Gottingensis Recentiores}, Vol. IV, Goettingen, 1818
\end{center}

The fundamental theorem on quadratic residues, which is counted among the most beautiful truths of higher arithmetic, was indeed easily discovered by induction, but was much more difficult to prove. In this kind of subject it often happens that very simple truths, which seem almost to offer themselves spontaneously to the investigator by induction, have proofs hidden very deeply; only after many fruitless attempts can they finally be brought to light, and often by a path far different from the one by which they were sought. Then, once a first path has been found, it is not uncommon for several others to open toward the same goal: some shorter and more direct, others beginning as if from the side and from principles so different that one would scarcely have suspected any connection between them and the proposed question. Such a remarkable connection among more recondite truths not only gives these investigations a special charm, but also deserves careful investigation and analysis, since new aids or advances of the science often flow from it.

Thus, although the arithmetical theorem treated here may seem to have been completely settled by earlier work, which supplied four proofs entirely different from one another,\footnote{Two were set out in sections four and five of the \emph{Disquisitiones Arithmeticae}; the third in a separate memoir (\emph{Commentt. Soc. Gotting.}, Vol. XVI); the fourth was inserted in the memoir \emph{Summatio quarundam serierum singularium} (\emph{Commentt. Recentiores}, Vol. I).} I nevertheless return to the same argument and add two other proofs, which will certainly shed new light on the matter. The first is in a certain way akin to the third, since it starts from the same lemma; afterwards, however, it follows a different course, so that it may rightly be regarded as a new proof, not inferior to the third in elegance, even if not superior to it. By contrast, the sixth proof rests on a wholly different and subtler principle, and gives a new example of the surprising connection among arithmetical truths that, at first sight, appear very far removed from each other. To these two proofs is added a very simple new algorithm for deciding whether a given integer is a quadratic residue or non-residue of a given prime.

There was also another reason why I now especially wished to publish the new proofs that had already been promised nine years ago. When, from the year 1805, I began to investigate the theory of cubic and biquadratic residues, a much more difficult subject, I experienced almost the same fortune as I had formerly experienced in the theory of quadratic residues. The theorems that completely exhaust those questions, and in which a remarkable analogy with the theorems concerning quadratic residues stands out, were at once discovered by induction as soon as a suitable way had been sought; but all attempts to obtain fully perfected proofs of them remained fruitless for a long time. This itself was the stimulus that made me strive so eagerly to add other proofs to the proofs already known concerning quadratic residues, in the hope that, among many diverse methods, one or another might contribute something toward illuminating the allied argument. That hope was by no means vain, and tireless labour was at last followed by favorable success. It will soon be possible to publish the fruits of these vigils; but before undertaking that arduous work, I decided to return once again to the theory of quadratic residues, to finish whatever still remains to be done there, and thus, as it were, to take leave of this part of higher arithmetic.

\section*{Fifth proof of the fundamental theorem in the theory of quadratic residues}

\subsection*{1. Lemma}

In the introduction we already stated that the fifth and the third proof start from the same lemma. For convenience, and with the notation adapted to the present investigation, it seemed useful to repeat it here.

Let $m$ be an odd positive prime number, and let $M$ be an integer not divisible by $m$. Take the least positive residues of the numbers
\[
M,\,2M,\,3M,\,4M,\ldots,\frac12(m-1)M
\]
modulo $m$. Some of them are less than $\frac12m$ and some greater; let the number of the latter be $n$. Then $M$ is a quadratic residue of $m$, or a non-residue, according as $n$ is even or odd.

\emph{Proof.} Let the residues less than $\frac12m$ be $a,b,c,d$ and so on, and let the remaining residues, greater than $\frac12m$, be $a',b',c',d'$ and so on. The complements of the latter to $m$,
\[
m-a',\quad m-b',\quad m-c',\quad m-d',\ldots
\]
are all plainly less than $\frac12m$ and are distinct both from one another and from the residues $a,b,c,d,\ldots$. Together with them, in a changed order, they are therefore identical with all the numbers
\[
1,2,3,4,\ldots,\frac12(m-1).
\]
Thus, putting
\[
1\cdot2\cdot3\cdot4\cdots \frac12(m-1)=P,
\]
one has
\[
P=abcd\cdots(m-a')(m-b')(m-c')(m-d')\cdots
\]
and therefore
\[
(-1)^nP=abcd\cdots(a'-m)(b'-m)(c'-m)(d'-m)\cdots .
\]
Moreover, modulo $m$,
\[
PM^{\frac12(m-1)}
 \equiv abcd\cdots a'b'c'd'\cdots
 \equiv abcd\cdots(a'-m)(b'-m)(c'-m)(d'-m)\cdots,
\]
hence
\[
PM^{\frac12(m-1)}\equiv P(-1)^n.
\]
It follows that
\[
M^{\frac12(m-1)}\equiv \pm1,
\]
where the upper or lower sign is taken according as $n$ is even or odd. With the aid of the theorem proved in article 106 of the \emph{Disquisitiones Arithmeticae}, the truth of the lemma follows at once.

\subsection*{2. Theorem}

Let $m,M$ be positive odd integers prime to each other. Let $n$ be the number among the least positive residues of
\[
M,2M,3M,\ldots,\frac12(m-1)M
\]
modulo $m$ that are greater than $\frac12m$; and similarly let $N$ be the number among the least positive residues of
\[
m,2m,3m,\ldots,\frac12(M-1)m
\]
modulo $M$ that are greater than $\frac12M$. Then the three numbers
\[
n,\qquad N,\qquad \frac14(m-1)(M-1)
\]
are either all even, or one is even and the other two are odd.

\emph{Proof.} Let
\[
\begin{array}{ll}
f  & \text{denote the set }1,2,3,\ldots,\frac12(m-1),\\
f' & \text{denote the set }m-1,m-2,m-3,\ldots,\frac12(m+1),\\
F  & \text{denote the set }1,2,3,\ldots,\frac12(M-1),\\
F' & \text{denote the set }M-1,M-2,M-3,\ldots,\frac12(M+1).
\end{array}
\]
Thus $n$ indicates how many of the numbers $Mf$ have their least positive residues modulo $m$ in the set $f'$, and similarly $N$ indicates how many of the numbers $mF$ have their least positive residues modulo $M$ in the set $F'$.

Finally, let
\[
\begin{array}{ll}
\varphi  & \text{denote the set }1,2,3,\ldots,\frac12(mM-1),\\
\varphi' & \text{denote the set }mM-1,mM-2,mM-3,\ldots,\frac12(mM+1).
\end{array}
\]
Every integer not divisible by $m$ is, modulo $m$, congruent either to a residue from $f$ or to one from $f'$, and similarly every integer not divisible by $M$ is, modulo $M$, congruent either to a residue from $F$ or to one from $F'$. Hence all the numbers $\varphi$, among which none is divisible by both $m$ and $M$, may be distributed into the following eight classes.

\begin{enumerate}[label=\Roman*.]
\item Numbers congruent modulo $m$ to some number of $f$, and modulo $M$ to some number of $F$. Let their number be $\alpha$.
\item Numbers congruent modulo $m,M$ respectively to numbers from $f,F'$. Let their number be $\beta$.
\item Numbers congruent modulo $m,M$ respectively to numbers from $f',F$. Let their number be $\gamma$.
\item Numbers congruent modulo $m,M$ respectively to numbers from $f',F'$. Let their number be $\delta$.
\item Numbers divisible by $m$ and congruent modulo $M$ to residues from $F$.
\item Numbers divisible by $m$ and congruent modulo $M$ to residues from $F'$.
\item Numbers divisible by $M$ and congruent modulo $m$ to residues from $f$.
\item Numbers divisible by $M$ and congruent modulo $m$ to residues from $f'$.
\end{enumerate}

Classes V and VI together plainly contain all the numbers $mF$. The number of numbers contained in VI is $N$, and therefore the number contained in V is $\frac12(M-1)-N$. Similarly, classes VII and VIII together contain all the numbers $Mf$; in class VIII there are $n$ numbers, and in class VII there are $\frac12(m-1)-n$.

In exactly the same way all the numbers $\varphi'$ are distributed into the eight classes IX--XVI. If the same order is kept, it is easy to see that the numbers contained in classes
\[
\mathrm{IX},\mathrm{X},\mathrm{XI},\mathrm{XII},\mathrm{XIII},\mathrm{XIV},\mathrm{XV},\mathrm{XVI}
\]
are respectively the complements to $mM$ of the numbers contained in classes
\[
\mathrm{IV},\mathrm{III},\mathrm{II},\mathrm{I},\mathrm{VI},\mathrm{V},\mathrm{VIII},\mathrm{VII}.
\]
Thus there are $\delta$ numbers in class IX, $\gamma$ in class X, and so on.

If all numbers of the first class are joined with all numbers of the ninth class, one obtains all numbers below $mM$ that are congruent modulo $m$ to some number of $f$ and modulo $M$ to some number of $F$. Their number is plainly equal to the number of all combinations of each element of $f$ with each element of $F$. Thus
\[
\alpha+\delta=\frac14(m-1)(M-1),
\]
and similarly
\[
\beta+\gamma=\frac14(m-1)(M-1).
\]
Joining all numbers of classes II, IV, and VI, we clearly obtain all numbers below $\frac12mM$ that are congruent modulo $M$ to a residue from $F'$. These same numbers can also be displayed as
\[
F',\quad M+F',\quad 2M+F',\quad 3M+F',\ldots,\frac12(m-3)M+F',
\]
so their total number is $\frac14(m-1)(M-1)$, or
\[
\beta+\delta+N=\frac14(m-1)(M-1).
\]
Similarly, from the union of all numbers of classes III, IV, and VIII one obtains
\[
\gamma+\delta+n=\frac14(m-1)(M-1).
\]
From these four equations follow
\[
\begin{aligned}
2\alpha&=\frac14(m-1)(M-1)+n+N,\\
2\beta&=\frac14(m-1)(M-1)+n-N,\\
2\gamma&=\frac14(m-1)(M-1)-n+N,\\
2\delta&=\frac14(m-1)(M-1)-n-N.
\end{aligned}
\]
Any one of these equations proves the theorem.

\subsection*{3.}

If now we suppose that $m$ and $M$ are prime numbers, the fundamental theorem follows immediately from the preceding theorem together with the lemma of article 1. For:
\begin{enumerate}[label=\Roman*.]
\item whenever both $m,M$, or even only one of them, is of the form $4k+1$, the number $\frac14(m-1)(M-1)$ is even, and hence $n$ and $N$ are either both even or both odd. Therefore either $m$ and $M$ are each quadratic residues of the other, or each is a non-residue of the other;
\item whenever both $m,M$ are of the form $4k+3$, $\frac14(m-1)(M-1)$ is odd, so one of $n,N$ is even and the other odd. Therefore one of the numbers $m,M$ is a quadratic residue of the other, and the other is a non-residue of it.
\end{enumerate}
Q.E.D.

\section*{Sixth proof of the fundamental theorem in the theory of quadratic residues}

\subsection*{1. Theorem}

Let $p$ denote an odd positive prime, $n$ a positive integer not divisible by $p$, and $x$ an indeterminate quantity. Then the function
\[
1+x^n+x^{2n}+x^{3n}+\cdots+x^{np-n}
\]
is divisible by
\[
1+x+xx+x^3+\cdots+x^{p-1}.
\]

\emph{Proof.} Choose a positive integer $g$ such that $gn\equiv1\pmod p$, and put $gn=1+hp$. Then
\[
\begin{aligned}
\frac{1+x^n+x^{2n}+x^{3n}+\cdots+x^{np-n}}{1+x+xx+x^3+\cdots+x^{p-1}}
&=\frac{(1-x^{np})(1-x)}{(1-x^n)(1-x^p)}\\
&=\frac{(1-x^{np})(1-x^{gn}-x+x^{hp+1})}{(1-x^n)(1-x^p)}\\
&=\frac{1-x^{np}}{1-x^p}\cdot\frac{1-x^{gn}}{1-x^n}
-\frac{x(1-x^{np})}{1-x^n}\cdot\frac{1-x^{hp}}{1-x^p}.
\end{aligned}
\]
This is plainly an integral function. Q.E.D. Therefore any integral function of $x$ divisible by $\frac{1-x^{np}}{1-x^n}$ is also divisible by $\frac{1-x^p}{1-x}$.

\subsection*{2.}

Let $\alpha$ denote a positive primitive root modulo $p$, that is, a positive integer such that the least positive residues of
\[
1,\alpha,\alpha\alpha,\alpha^3,\ldots,\alpha^{p-1}
\]
modulo $p$ are, without regard to order, identical with $1,2,3,4,\ldots,p-1$. If $fx$ denotes the function
\[
x+x^\alpha+x^{\alpha\alpha}+x^{\alpha^3}+\cdots+x^{\alpha^{p-2}}+1,
\]
then $fx-1-x-xx-x^3-\cdots-x^{p-1}$ is divisible by $1-x^p$, and hence a fortiori by
\[
\frac{1-x^p}{1-x}=1+x+xx+x^3+\cdots+x^{p-1}.
\]
Thus $fx$ itself is divisible by this function. Since $x$ is an indeterminate, it follows that $f(x^n)$ is divisible by $\frac{1-x^{np}}{1-x^n}$ and therefore, by the preceding article, also by $\frac{1-x^p}{1-x}$ whenever $n$ is not divisible by $p$. If, on the other hand, $n$ is divisible by $p$, each term of $f(x^n)$ less unity is divisible by $1-x^p$; hence in this case $f(x^n)-p$ is divisible by $1-x^p$, and therefore also by $\frac{1-x^p}{1-x}$.

\subsection*{3. Theorem}

Put
\[
x-x^\alpha+x^{\alpha\alpha}-x^{\alpha^3}+x^{\alpha^4}-\cdots-x^{\alpha^{p-2}}=\xi.
\]
Then $\xi\xi\mp p$ is divisible by $\frac{1-x^p}{1-x}$, the upper sign being used when $p$ is of the form $4k+1$ and the lower sign when $p$ is of the form $4k+3$.

\emph{Proof.} From the $p-1$ functions displayed below,
\[
\begin{aligned}
&+x\xi-xx+x^{\alpha+1}-x^{\alpha\alpha+1}+\cdots+x^{\alpha^{p-2}+1},\\
&-x^\alpha\xi-x^{2\alpha}+x^{\alpha\alpha+\alpha}-x^{\alpha^3+\alpha}+\cdots+x^{\alpha^{p-1}+\alpha},\\
&+x^{\alpha\alpha}\xi-x^{2\alpha\alpha}+x^{\alpha^3+\alpha\alpha}-x^{\alpha^4+\alpha\alpha}+\cdots+x^{\alpha^p+\alpha\alpha},\\
&-x^{\alpha^3}\xi-x^{2\alpha^3}+x^{\alpha^4+\alpha^3}-x^{\alpha^5+\alpha^3}+\cdots+x^{\alpha^{p+1}+\alpha^3},\\
&\qquad\vdots\\
&-x^{\alpha^{p-2}}\xi-x^{2\alpha^{p-2}}+x^{\alpha^{p-1}+\alpha^{p-2}}-x^{\alpha^p+\alpha^{p-2}}+\cdots+x^{\alpha^{2p-4}+\alpha^{p-2}},
\end{aligned}
\]
it is easy to see that the first is $0$ and all the rest are divisible by $1-x^p$. Hence their sum is also divisible by $1-x^p$, and this sum is
\[
\begin{aligned}
\Omega=\xi\xi
&-(f(xx)-1)+(f(x^{\alpha+1})-1)-(f(x^{\alpha\alpha+1})-1)\\
&+(f(x^{\alpha^3+1})-1)-\cdots+f(x^{\alpha^{p-2}+1})-1\\
=\xi\xi
&-f(xx)+f(x^{\alpha+1})-f(x^{\alpha\alpha+1})
 +f(x^{\alpha^3+1})-\cdots+f(x^{\alpha^{p-2}+1}).
\end{aligned}
\]
Thus the expression $\Omega$ is also divisible by $\frac{1-x^p}{1-x}$.

Among the exponents
\[
2,\quad \alpha+1,\quad \alpha\alpha+1,\quad \alpha^3+1,\ldots,\alpha^{p-2}+1
\]
only one is divisible by $p$, namely $\alpha^{\frac12(p-1)}+1$. Therefore, by the preceding article, all the terms of $\Omega$ of the form
\[
f(xx),\quad f(x^{\alpha+1}),\quad f(x^{\alpha\alpha+1}),\quad f(x^{\alpha^3+1}),\ldots
\]
except the single term $f(x^{\alpha^{\frac12(p-1)}+1})$ are divisible by $\frac{1-x^p}{1-x}$. Those terms may therefore be erased, leaving still divisible by $\frac{1-x^p}{1-x}$ the function
\[
\xi\xi\mp f(x^{\alpha^{\frac12(p-1)}+1}),
\]
where the upper or lower sign is used according as $p$ is of the form $4k+1$ or $4k+3$. Since $f(x^{\alpha^{\frac12(p-1)}+1})-p$ is also divisible by $\frac{1-x^p}{1-x}$, it follows that $\xi\xi\mp p$ is divisible by the same divisor. Q.E.D.

To remove any ambiguity from the double sign, let $\varepsilon$ denote $+1$ or $-1$ according as $p$ is of the form $4k+1$ or $4k+3$. Thus
\[
Z=\frac{(1-x)(\xi\xi-\varepsilon p)}{1-x^p}
\]
is an integral function of $x$.

\subsection*{4.}

Let $q$ be a positive odd number, so that $\frac12(q-1)$ is an integer. Then
\[
(\xi\xi)^{\frac12(q-1)}-(\varepsilon p)^{\frac12(q-1)}
\]
is divisible by $\xi\xi-\varepsilon p$, and therefore also by $\frac{1-x^p}{1-x}$. Put
\[
\varepsilon^{\frac12(q-1)}=\delta,\qquad
\xi^{q-1}-\delta p^{\frac12(q-1)}=\frac{1-x^p}{1-x}\,Y.
\]
Then $Y$ is an integral function of $x$, and $\delta=+1$ whenever one of $p,q$, or both, is of the form $4k+1$; while $\delta=-1$ whenever both $p,q$ are of the form $4k+3$.

\subsection*{5.}

Now suppose that $q$ too is a prime number distinct from $p$. By the theorem proved in article 51 of the \emph{Disquisitiones Arithmeticae}, the expression
\[
\xi^q-\left(x^q-x^{q\alpha}+x^{q\alpha\alpha}-x^{q\alpha^3}+\cdots-x^{q\alpha^{p-2}}\right)
\]
is divisible by $q$, say equal to $qX$, where $X$ is an integral function of $x$ also with respect to its numerical coefficients. The same is to be understood of the remaining integral functions $Z,Y,W$ occurring here. Let the index of $q$ modulo $p$ with respect to the primitive root $\alpha$ be denoted by $\mu$, so that $q\equiv\alpha^\mu\pmod p$. Then the numbers
\[
q,q\alpha,q\alpha\alpha,q\alpha^3,\ldots,q\alpha^{p-2}
\]
are congruent modulo $p$ respectively to
\[
\alpha^\mu,\alpha^{\mu+1},\alpha^{\mu+2},\ldots,\alpha^{p-2},1,\alpha,\alpha\alpha,\ldots,\alpha^{\mu-1}.
\]
Hence the differences
\[
\begin{aligned}
&x^q-x^{\alpha^\mu},\quad x^{q\alpha}-x^{\alpha^{\mu+1}},\quad
x^{q\alpha\alpha}-x^{\alpha^{\mu+2}},\ldots,\\
&x^{q\alpha^{p-\mu-2}}-x^{\alpha^{p-2}},\quad
x^{q\alpha^{p-\mu-1}}-x,\quad
x^{q\alpha^{p-\mu}}-x^\alpha,\ldots,\quad
x^{q\alpha^{p-2}}-x^{\alpha^{\mu-1}}
\end{aligned}
\]
are divisible by $1-x^p$. Taking these quantities alternately with positive and negative signs and adding, one sees that
\[
x^q-x^{q\alpha}+x^{q\alpha\alpha}-x^{q\alpha^3}+\cdots-x^{q\alpha^{p-2}}\mp \xi
\]
is divisible by $1-x^p$, the upper or lower sign holding according as $\mu$ is even or odd, that is, according as $q$ is a quadratic residue or non-residue of $p$. We therefore put
\[
x^q-x^{q\alpha}+x^{q\alpha\alpha}-x^{q\alpha^3}+\cdots-x^{q\alpha^{p-2}}-\gamma\xi=(1-x^p)W,
\]
where $\gamma=+1$ or $\gamma=-1$ according as $q$ is a quadratic residue or non-residue of $p$. Then $W$ is an integral function.

\subsection*{6.}

After these preparations, a combination of the preceding equations gives
\[
q\xi X=\varepsilon p\left(\delta p^{\frac12(q-1)}-\gamma\right)
+\frac{1-x^p}{1-x}\left(Z\left(\delta p^{\frac12(q-1)}-\gamma\right)+Y\xi\xi-W\xi(1-x)\right).
\]
Suppose that division of $\xi X$ by
\[
x^{p-1}+x^{p-2}+x^{p-3}+\cdots+x+1
\]
gives a quotient $U$ and a remainder $T$, so that
\[
\xi X=\frac{1-x^p}{1-x}U+T,
\]
where $U$ and $T$ are integral functions, also with respect to their numerical coefficients, and $T$ is certainly of lower order than the divisor. Thus
\[
qT-\varepsilon p\left(\delta p^{\frac12(q-1)}-\gamma\right)
=\frac{1-x^p}{1-x}\left(Z\left(\delta p^{\frac12(q-1)}-\gamma\right)+Y\xi\xi-W\xi(1-x)-qU\right).
\]
This equation plainly cannot hold unless both the left and right members vanish separately. Therefore $\varepsilon p(\delta p^{\frac12(q-1)}-\gamma)$ is divisible by $q$. If $\beta$ denotes $+1$ or $-1$ according as $p$ is a quadratic residue or non-residue of $q$, then
\[
p^{\frac12(q-1)}-\beta
\]
is divisible by $q$, and so also is $\beta-\gamma\delta$. This cannot occur unless $\beta=\gamma\delta$. From this the fundamental theorem follows at once:
\begin{enumerate}[label=\Roman*.]
\item If both $p,q$, or one of them, is of the form $4k+1$, then $\delta=+1$, so $\beta=\gamma$. Hence either $q$ is a quadratic residue of $p$ and $p$ is a quadratic residue of $q$, or $q$ is a non-residue of $p$ and $p$ is a non-residue of $q$.
\item If both $p,q$ are of the form $4k+3$, then $\delta=-1$, so $\beta=-\gamma$. Hence either $q$ is a quadratic residue of $p$ and $p$ is a non-residue of $q$, or $q$ is a non-residue of $p$ and $p$ is a residue of $q$.
\end{enumerate}
Q.E.D.

\section*{A new algorithm for deciding whether a given positive integer is a quadratic residue or non-residue of a given positive prime}

\subsection*{1.}

Before we set out the new solution of this problem, we briefly recall the solution given in the \emph{Disquisitiones Arithmeticae}. It is carried out quite expeditiously with the help of the fundamental theorem and the following known theorems:
\begin{enumerate}[label=\Roman*.]
\item The relation of the number $a$ to the number $b$, in so far as $a$ is a quadratic residue or non-residue of $b$, is the same as the relation of $c$ to $b$, if $a\equiv c\pmod b$.
\item If $a$ is a product of factors $\alpha,\beta,\gamma,\delta$ etc., and $b$ is prime, the relation of $a$ to $b$ depends on the relations of these factors to $b$ in such a way that $a$ is a quadratic residue or non-residue of $b$ according as the number of those factors that are non-residues of $b$ is even or odd. Thus, whenever a factor is a square, it need not be considered in this examination; if a factor is a power of an integer with odd exponent, that integer itself may take its place.
\item The number $2$ is a quadratic residue of every prime of the form $8m+1$ or $8m+7$, and a non-residue of every prime of the form $8m+3$ or $8m+5$.
\end{enumerate}

Given a number $a$ whose relation to the prime $b$ is sought, if $a>b$ one first substitutes its least positive residue modulo $b$. Resolving that residue into prime factors, the question is reduced by theorem II to finding the relation of each of these factors to $b$. The relation of the factor $2$, if it occurs once, three times, five times, and so on, is known from theorem III. The relation of the remaining factors depends, by the fundamental theorem, on the relation of $b$ to the individual factors. Thus, instead of one relation of a given number to the prime $b$, one has to investigate several relations of $b$ to odd primes smaller than $b$; those problems are lowered in the same way to smaller moduli, and the successive reductions plainly come to an end.

\subsection*{2. Example}

As an illustration, seek the relation of $103$ to $379$. Since $103$ is already less than $379$ and is itself prime, the fundamental theorem is applied immediately. It says that the desired relation is the opposite of the relation of $379$ to $103$. This is equal to the relation of $70$ to $103$, which depends on the relations of $2,5,7$ to $103$. The first of these is known from theorem III. The second depends, by the fundamental theorem, on the relation of $103$ to $5$, which by theorem I is equal to the relation of $3$ to $5$; this again, by the fundamental theorem, depends on the relation of $5$ to $3$, which by theorem I is equal to the relation of $2$ to $3$, known by theorem III. Similarly, the relation of $7$ to $103$ depends, by the fundamental theorem, on the relation of $103$ to $7$, which by theorem I is equal to the relation of $5$ to $7$; this again depends, by the fundamental theorem, on the relation of $7$ to $5$, which is equal by theorem I to the relation of $2$ to $5$, known by theorem III. If this analysis is now transformed into a synthesis, the decision of the question is referred to fourteen steps, which we write out in full so that the greater elegance of the new solution may appear more clearly.

\begin{enumerate}
\item $2$ is a quadratic residue of $103$ (theorem III).
\item $2$ is a quadratic non-residue of $3$ (theorem III).
\item $5$ is a quadratic non-residue of $3$ (from I and 2).
\item $3$ is a quadratic non-residue of $5$ (fundamental theorem and 3).
\item $103$ is a quadratic non-residue of $5$ (I and 4).
\item $5$ is a quadratic non-residue of $103$ (fundamental theorem and 5).
\item $2$ is a quadratic non-residue of $5$ (theorem III).
\item $7$ is a quadratic non-residue of $5$ (I and 7).
\item $5$ is a quadratic non-residue of $7$ (fundamental theorem and 8).
\item $103$ is a quadratic non-residue of $7$ (I and 9).
\item $7$ is a quadratic residue of $103$ (fundamental theorem and 10).
\item $70$ is a quadratic non-residue of $103$ (II, 1, 6, 11).
\item $379$ is a quadratic non-residue of $103$ (I and 12).
\item $103$ is a quadratic residue of $379$ (fundamental theorem and 13).
\end{enumerate}

In what follows, for brevity, we use the notation introduced in \emph{Comment. Gotting.} Vol. XVI. Namely, $[x]$ denotes the quantity $x$ itself when $x$ is an integer, and otherwise the nearest integer less than $x$, so that $x-[x]$ is always non-negative and less than unity.

\subsection*{3. Problem}

Let $a,b$ be positive integers prime to one another, and put $\left[\frac12a\right]=a'$. Find the sum
\[
\left[\frac{b}{a}\right]+\left[\frac{2b}{a}\right]+\left[\frac{3b}{a}\right]+\left[\frac{4b}{a}\right]+\cdots+\left[\frac{a'b}{a}\right].
\]

\emph{Solution.} For brevity denote such a sum by $\varphi(a,b)$, so that also
\[
\varphi(b,a)=\left[\frac{a}{b}\right]+\left[\frac{2a}{b}\right]+\left[\frac{3a}{b}\right]+\cdots+\left[\frac{b'a}{b}\right]
\]
if $\left[\frac12b\right]=b'$. In the third proof of the fundamental theorem it was shown that, in the case where $a$ and $b$ are odd,
\[
\varphi(a,b)+\varphi(b,a)=a'b',
\]
and by following the same method this proposition is easily extended to the case where one of $a,b$ is odd, as was already indicated there. Divide $a$ by $b$ as in the algorithm for the greatest common divisor; let $\beta$ be the quotient and $c$ the remainder. Next divide $b$ by $c$, and so on, so that
\[
a=\beta b+c,\qquad b=\gamma c+d,\qquad c=\delta d+e,\qquad d=\varepsilon e+f,\ \ldots
\]
In this way, since $a$ and $b$ are prime to one another, the decreasing sequence $b,c,d,e,f,\ldots$ eventually reaches unity, and the last equation is of the form
\[
k=\lambda l+1.
\]
Since plainly
\[
\left[\frac{a}{b}\right]=\left[\beta+\frac{c}{b}\right]=\beta+\left[\frac{c}{b}\right],
\]
\[
\left[\frac{2a}{b}\right]=2\beta+\left[\frac{2c}{b}\right],\qquad
\left[\frac{3a}{b}\right]=3\beta+\left[\frac{3c}{b}\right],\quad\ldots,
\]
we have
\[
\varphi(b,a)=\varphi(b,c)+\frac12\beta(b'b'+b')
\]
and therefore
\[
\varphi(a,b)=a'b'-\frac12\beta(b'b'+b')-\varphi(b,c).
\]
Similar reasoning gives, with $\left[\frac12c\right]=c'$, $\left[\frac12d\right]=d'$, $\left[\frac12e\right]=e'$, and so on,
\[
\begin{aligned}
\varphi(b,c)&=b'c'-\frac12\gamma(c'c'+c')-\varphi(c,d),\\
\varphi(c,d)&=c'd'-\frac12\delta(d'd'+d')-\varphi(d,e),\\
\varphi(d,e)&=d'e'-\frac12\varepsilon(e'e'+e')-\varphi(e,f),
\end{aligned}
\]
and so on, down to
\[
\varphi(k,l)=k'l'-\frac12\lambda(l'l'+l')-\varphi(l,1).
\]
Since $\varphi(l,1)=0$, we obtain
\[
\begin{aligned}
\varphi(a,b)= {}& a'b'-b'c'+c'd'-d'e'+\cdots\pm k'l'\\
&-\frac12\beta(b'b'+b')+\frac12\gamma(c'c'+c')-\frac12\delta(d'd'+d')\\
&+\frac12\varepsilon(e'e'+e')-\cdots\mp\frac12\lambda(l'l'+l').
\end{aligned}
\]

\subsection*{4.}

From what was set out in the third proof it is now easy to infer that the relation of $b$ to $a$, whenever $a$ is prime, is obtained directly from the value of the sum $\varphi(a,2b)$. According as this sum is even or odd, $b$ is a quadratic residue or non-residue of $a$. The sum $\varphi(a,b)$ itself can also be used for the same purpose, but with the following distinction between the case where $b$ is odd and the case where $b$ is even:
\begin{enumerate}[label=\Roman*.]
\item If $b$ is odd, then $b$ is a quadratic residue or non-residue of $a$ according as $\varphi(a,b)$ is even or odd.
\item If $b$ is even, the same rule holds provided that $a$ is of the form $8n+1$ or $8n+7$. If, for an even value of $b$, the modulus $a$ is of the form $8n+3$ or $8n+5$, the opposite rule must be applied: $b$ is a quadratic residue of $a$ when $\varphi(a,b)$ is odd, and a non-residue when $\varphi(a,b)$ is even.
\end{enumerate}
All this is derived very easily from article 4 of the third proof.

\subsection*{5. Example}

If the relation of $103$ to the prime $379$ is sought, then, to compute $\varphi(379,103)$, we have
\[
\begin{array}{r|r|r}
a=379&a'=189&\\
b=103&b'=51&\beta=3\\
c=70&c'=35&\gamma=1\\
d=33&d'=16&\delta=2\\
e=4&e'=2&\varepsilon=8
\end{array}
\]
and hence
\[
\varphi(379,103)=9639-1785+560-32-3978+630-272+24=4786.
\]
Therefore $103$ is a quadratic residue of $379$. If instead one prefers to use the sum $\varphi(379,206)$ for the same purpose, the scheme is
\[
\begin{array}{r|r|r}
379&189&\\
206&103&1\\
173&86&1\\
33&16&5\\
8&4&4
\end{array}
\]
from which we derive
\[
\varphi(379,206)=19467-8858+1376-64-5356+3741-680+40=9666.
\]
Thus again $103$ is a quadratic residue of $379$.

\subsection*{6.}

Since, in order to decide the relation of $b$ to $a$, it is not necessary to compute the individual parts of the sum $\varphi(a,b)$, but only to know how many of them are odd, our rule may also be presented as follows. As above, put
\[
a=\beta b+c,\qquad b=\gamma c+d,\qquad c=\delta d+e,\quad\ldots,
\]
until the sequence $a,b,c,d,e,\ldots$ has reached unity. Put
\[
\left[\frac12a\right]=a',\qquad \left[\frac12b\right]=b',\qquad \left[\frac12c\right]=c',\quad\ldots.
\]
Let $\mu$ be the number of odd numbers in the sequence $a',b',c',\ldots$ which are immediately followed by an odd number. Let $\nu$ be the number of odd numbers in the sequence $\beta,\gamma,\delta,\ldots$ for which the corresponding number in the sequence $b',c',d',\ldots$ is of the form $4n+1$ or $4n+2$. Then $b$ is a quadratic residue or non-residue of $a$ according as $\mu+\nu$ is even or odd, with the single exception of the case in which $b$ is even and $a$ is of the form $8n+3$ or $8n+5$; in that case the opposite rule holds.

In our example the sequence $a',b',c',d',e'$ gives two successions of odd numbers, hence $\mu=2$. In the sequence $\beta,\gamma,\delta,\varepsilon$ there are indeed two odd terms, but the corresponding numbers in the sequence $b',c',d',e'$ are of the form $4n+3$, and hence $\nu=0$. Therefore $\mu+\nu$ is even, and so $103$ is a quadratic residue of $379$.



\section*{Theory of Biquadratic Residues. First Memoir}

\begin{center}
\textsc{By Carl Friedrich Gauss}\\
\textsc{Presented to the Royal Society, 5 April 1825}\\[3pt]
\emph{Recent Commentaries of the Royal Society of Sciences at Göttingen}, Vol. VI\\
Göttingen, 1828
\end{center}

\subsection*{1.}

The theory of quadratic residues is reduced to a few fundamental theorems, to be counted among the most beautiful treasures of higher arithmetic. These were first easily discovered by induction and then proved in many different ways so fully that nothing further has been left to desire.

The theory of cubic and biquadratic residues is of far deeper investigation. When, beginning in 1805, we set out to examine it, besides the matters lying almost on the threshold, several special theorems presented themselves, very remarkable both for their simplicity and for the difficulty of their proofs. Soon, however, we found that the principles of arithmetic hitherto in use were by no means sufficient for establishing the general theory; rather, the theory necessarily requires that the field of higher arithmetic be, as it were, promoted infinitely many times. How this is to be understood will become quite clear in the continuation of these investigations.

As soon as we entered this new field, the path to knowing the simplest theorems, which exhaust the whole theory, immediately lay open by induction; but their proofs were hidden so deeply that only after many fruitless attempts could they at last be brought to light.

Now that we are preparing to publish these studies, we shall begin with the theory of biquadratic residues. In this first memoir we shall explain those investigations that could already be completed this side of the enlarged field of arithmetic; they in a sense pave the way to it, and at the same time add some new increments to the theory of the division of the circle.

\subsection*{2.}

We introduced the notion of a biquadratic residue in \emph{Disquisitiones Arithmeticae}, art. 115: namely, an integer $a$, positive or negative, is called a biquadratic residue of the integer $p$ if $a$ can be made congruent to a biquadrate modulo $p$; it is likewise called a biquadratic non-residue if no such congruence exists.

In all the following investigations, unless the contrary is stated explicitly, we shall suppose the modulus $p$ to be a prime number, positive and odd, and $a$ not divisible by $p$, since all remaining cases can very easily be reduced to this one.

\subsection*{3.}

It is clear that every biquadratic residue of the number $p$ is also a quadratic residue of the same number, and hence that every quadratic non-residue is also a biquadratic non-residue. This proposition may also be converted whenever $p$ is a prime number of the form $4n+3$.

For if, in this case, $a$ is a quadratic residue of $p$, put
\[
a\equiv bb \pmod p.
\]
Then $b$ is either a quadratic residue of $p$ or a non-residue. In the first case put $b\equiv cc$, whence $a\equiv c^4$, that is, $a$ is a biquadratic residue of $p$. In the second case $-b$ becomes a quadratic residue of $p$, since $-1$ is a non-residue of every prime of the form $4n+3$; putting $-b\equiv cc$, we have as before $a\equiv c^4$, and $a$ is a biquadratic residue of $p$.

It is also easily seen that in this case the congruence
\[
x^4\equiv a\pmod p
\]
has no solutions besides these two, $x\equiv c$ and $x\equiv -c$. Since these obvious propositions exhaust the entire theory of biquadratic residues for prime moduli of the form $4n+3$, we shall exclude such moduli entirely from our investigation, or rather restrict it to prime moduli of the form $4n+1$.

\subsection*{4.}

Thus, when $p$ is a prime of the form $4n+1$, the proposition of the preceding article cannot be converted. There can exist quadratic residues that are not at the same time biquadratic residues; this happens whenever a quadratic residue is congruent to the square of a quadratic non-residue.

Indeed, put $a\equiv bb$, where $b$ is a quadratic non-residue of $p$. If the congruence $x^4\equiv a$ could be satisfied by $x\equiv c$, then $c^4\equiv bb$, so the product
\[
(cc-b)(cc+b)
\]
would be divisible by $p$. Hence $p$ would divide either the factor $cc-b$ or the other factor $cc+b$, that is, either $+b$ or $-b$ would be a quadratic residue of $p$, and consequently both would be, since $-1$ is a quadratic residue; this contradicts the hypothesis.

Therefore all integers not divisible by $p$ could be distributed into three classes: the first containing the biquadratic residues, the second the biquadratic non-residues that are nevertheless quadratic residues, and the third the quadratic non-residues. Plainly it is enough to subject only the numbers $1,2,3,\ldots,p-1$ to such a classification. One half of them is reduced to the third class, while the other half is distributed between the first and second classes.

\subsection*{5.}

It will be preferable, however, to establish four classes, whose nature is as follows. Let $A$ be the complex of all biquadratic residues of $p$ lying between $1$ and $p-1$, inclusive, and let $e$ be a quadratic non-residue of $p$, chosen arbitrarily. Let $B$ be the complex of the least positive residues arising from the products $eA$ modulo $p$, and likewise let $C,D$ be the complexes of least positive residues arising from the products $eeA,e^3A$ modulo $p$, respectively.

With this done, it is easy to see that the individual numbers in $B$ are mutually distinct, and similarly those in $C$ and in $D$; moreover zero cannot occur among any of these numbers. It is also clear that all numbers contained in $A$ and $C$ are quadratic residues of $p$, while all numbers in $B$ and $D$ are quadratic non-residues. Thus the complexes $A,C$ certainly cannot have any number in common with $B$ or $D$. But neither can $A$ have any number in common with $C$, nor $B$ with $D$.

Suppose first that some number from $A$, say $\alpha$, is also found in $C$, where it has arisen from a product $eea$ congruent to it, with $a$ a number from $A$. Put $\alpha\equiv \alpha_0^4$, $a\equiv c^4$, and choose an integer $\theta$ such that $\theta\alpha_0\equiv1$. Then
\[
ee c^4\equiv \alpha_0^4,
\]
and hence, multiplying by $\theta^4$,
\[
ee\equiv c^4\theta^4.
\]
Thus $ee$ is a biquadratic residue, and therefore $e$ a quadratic residue, contrary to the hypothesis.

Similarly, suppose that some number is common to the complexes $B,D$ and has arisen from products $e\alpha,e^3a$, where $\alpha,a$ are numbers from $A$. From the congruence $e\alpha\equiv e^3a$ it follows that $\alpha\equiv eea$, so that there would be a number arising from a product $eea$ that belongs simultaneously to $C$ and to $A$, which has just been shown to be impossible.

It is also easy to prove that all quadratic residues of $p$ lying between $1$ and $p-1$ must occur either in $A$ or in $C$, and that all quadratic non-residues of $p$ within the same limits must occur either in $B$ or in $D$. For every such quadratic residue that is also a biquadratic residue is, by hypothesis, found in $A$. Let $h<p$ be a quadratic residue that is nevertheless a biquadratic non-residue, and put $h=gg$, where $g$ is a quadratic non-residue. Choose an integer $\gamma$ such that $e\gamma\equiv g$; then $\gamma$ is a quadratic residue of $p$, and we put $\gamma=kk$. Hence
\[
h\equiv gg\equiv ee\gamma\gamma\equiv ee k^4.
\]
Since the least residue of $k^4$ is found in $A$, the number $h$, which arises from its product by $ee$, must necessarily be contained in $C$. Finally, let $h$ denote a quadratic non-residue of $p$ between $1$ and $p-1$, and find within the same limits an integer $g$ such that $eg\equiv h$. Then $g$ is a quadratic residue, and hence is contained either in $A$ or in $C$. In the first case $h$ is plainly found among the numbers $B$, and in the second among the numbers $D$.

From all this it follows that all the numbers $1,2,3,\ldots,p-1$ are distributed among the four series $A,B,C,D$ in such a way that each of them is found in exactly one of these series; hence each series must contain $\frac14(p-1)$ numbers. In this classification, the classes $A$ and $C$ possess their numbers essentially, but the distinction between the classes $B$ and $D$ is arbitrary to the extent that it depends on the choice of the number $e$, which itself is always to be referred to $B$. Therefore, if instead of it another number from class $D$ is adopted, the classes $B,D$ are interchanged.

\subsection*{6.}

Since $-1$ is a quadratic residue of $p$, put
\[
-1\equiv ff\pmod p.
\]
Then the four roots of the congruence $x^4\equiv1$ are $1,f,-1,-f$. Thus if $a$ is a biquadratic residue of $p$, say $a\equiv\alpha^4$, the four roots of the congruence $x^4\equiv a$ are
\[
\alpha,\quad f\alpha,\quad -\alpha,\quad -f\alpha,
\]
and it is easy to see that they are incongruent to one another. Hence, if the least positive residues of the biquadrates
\[
1,\;16,\;81,\;256,\ldots,(p-1)^4
\]
are collected, they will always occur in fours equal to one another, so that $\frac14(p-1)$ distinct biquadratic residues form the complex $A$. If the least residues of the biquadrates are collected only up to $\left(\frac12(p-1)\right)^4$, each will occur twice.

\subsection*{7.}

The product of two biquadratic residues is plainly a biquadratic residue. Thus, from multiplying two numbers of class $A$, there always arises a product whose least positive residue belongs to the same class. Likewise, products of a number from $B$ with a number from $D$, or of a number from $C$ with a number from $C$, have their least residues in $A$.

The residues of the products $A\cdot B$ and $C\cdot D$ fall in $B$; those of $A\cdot C$, $B\cdot B$, and $D\cdot D$ fall in $C$; finally those of $A\cdot D$ and $B\cdot C$ fall in $D$. The proofs are so obvious that it is enough to indicate one. Let, for instance, $c$ and $d$ be numbers from $C$ and $D$, with $c\equiv eea$ and $d\equiv e^3a'$, where $a,a'$ are numbers from $A$. Then $e^4aa'$ is a biquadratic residue, that is, its least residue is referred to $A$. Since the product $cd$ is congruent to $e\cdot e^4aa'$, its least residue is contained in $B$.

It can now also be decided easily to which class a product of several factors is to be referred. Assigning to the classes $A,B,C,D$ the characters $0,1,2,3$, respectively, the character of the product is equal to the sum of the characters of the individual factors, or to its least residue modulo $4$.

\subsection*{8.}

It seemed worth the effort to develop these elementary propositions without the aid of the theory of residues of powers. With that theory called in, everything so far can be proved much more easily.

Let $g$ be a primitive root modulo $p$, that is, a number such that in the series of powers $g,gg,g^3,\ldots$ none before $g^{p-1}$ becomes congruent to unity modulo $p$. Then the least positive residues of the numbers
\[
1,\;g,\;gg,\;g^3,\ldots,g^{p-2}
\]
coincide, in another order, with $1,2,3,\ldots,p-1$, and are distributed into the four classes as follows:
\[
\begin{array}{c|l}
A&1,\;g^4,\;g^8,\ldots,g^{p-5}\\
B&g,\;g^5,\;g^9,\ldots,g^{p-4}\\
C&g^2,\;g^6,\;g^{10},\ldots,g^{p-3}\\
D&g^3,\;g^7,\;g^{11},\ldots,g^{p-2}.
\end{array}
\]
From this all the preceding propositions flow of themselves. Moreover, just as the numbers $1,2,3,\ldots,p-1$ are here distributed into four classes, whose complexes we denote by $A,B,C,D$, so any integer not divisible by $p$ may be assigned to one of these classes according to its least residue modulo $p$.

\subsection*{9.}

We shall denote by $f$ the least residue of the power $g^{\frac14(p-1)}$ modulo $p$. Since
\[
ff\equiv g^{\frac12(p-1)}\equiv -1
\]
(\emph{Disq. Arithm.}, art. 62), it is clear that the character $f$ here signifies the same as in art. 6. Therefore the power
\[
g^{\lambda\frac14(p-1)},
\]
where $\lambda$ is a positive integer, is congruent modulo $p$ to $1,f,-1,-f$, respectively, according as $\lambda$ is of the form $4m,4m+1,4m+2,4m+3$; equivalently, according as the least residue of $g^\lambda$ is found in $A,B,C,D$, respectively.

Hence we obtain a very simple criterion for deciding to which class a given number $h$, not divisible by $p$, is to be referred: namely, $h$ belongs to $A,B,C$, or $D$ according as the power
\[
h^{\frac14(p-1)}
\]
is congruent modulo $p$ to $1,f,-1$, or $-f$.

As a corollary it follows that $-1$ is always to be referred to class $A$ when $p$ is of the form $8n+1$, and to class $C$ when $p$ is of the form $8n+5$. A proof of this theorem independent of the theory of residues of powers can easily be arranged from what we taught in \emph{Disquisitiones Arithmeticae}, art. 115, III.

\subsection*{10.}

Since all primitive roots modulo $p$ arise from the residues of powers $g^\lambda$, taking for $\lambda$ all numbers relatively prime to $p-1$, it is easily seen that they will be divided equally between the complexes $B$ and $D$, with the base $g$ always contained in $B$.

If, in place of the number $g$, another primitive root from the complex $B$ is taken as the base, the classification remains the same; but if a primitive root from the complex $D$ is adopted as the base, the classes $B$ and $D$ are interchanged. If the classification is based on the criterion given in the preceding article, the distinction between classes $B$ and $D$ will depend on which root of the congruence
\[
xx\equiv -1\pmod p
\]
we adopt as the characteristic number $f$.

\subsection*{11.}

So that the subtler investigations on which we are now entering may be illustrated by examples, we append here the construction of the classes for all moduli below $100$. For each, we have adopted the least primitive root.

\begin{longtable}{>{\raggedright\arraybackslash}p{0.16\textwidth}p{0.78\textwidth}}
\toprule
Modulus & Classes\\
\midrule
$p=5$, $g=2$, $f=2$ &
$A$: 1;\quad $B$: 2;\quad $C$: 4;\quad $D$: 3.\\[2pt]
$p=13$, $g=2$, $f=8$ &
$A$: 1, 3, 9;\quad $B$: 2, 5, 6;\quad $C$: 4, 10, 12;\quad $D$: 7, 8, 11.\\[2pt]
$p=17$, $g=3$, $f=13$ &
$A$: 1, 4, 13, 16;\quad $B$: 3, 5, 12, 14;\quad $C$: 2, 8, 9, 15;\quad $D$: 6, 7, 10, 11.\\[2pt]
$p=29$, $g=2$, $f=12$ &
$A$: 1, 7, 16, 20, 23, 24, 25;\quad $B$: 2, 3, 11, 14, 17, 19, 21;\quad $C$: 4, 5, 6, 9, 13, 22, 28;\quad $D$: 8, 10, 12, 15, 18, 26, 27.\\[2pt]
$p=37$, $g=2$, $f=31$ &
$A$: 1, 7, 9, 10, 12, 16, 26, 33, 34;\quad $B$: 2, 14, 15, 18, 20, 24, 29, 31, 32;\quad $C$: 3, 4, 11, 21, 25, 27, 28, 30, 36;\quad $D$: 5, 6, 8, 13, 17, 19, 22, 23, 35.\\[2pt]
$p=41$, $g=6$, $f=32$ &
$A$: 1, 4, 10, 16, 18, 23, 25, 31, 37, 40;\quad $B$: 6, 14, 15, 17, 19, 22, 24, 26, 27, 35;\quad $C$: 2, 5, 8, 9, 20, 21, 32, 33, 36, 39;\quad $D$: 3, 7, 11, 12, 13, 28, 29, 30, 34, 38.\\[2pt]
$p=53$, $g=2$, $f=30$ &
$A$: 1, 10, 13, 15, 16, 24, 28, 36, 42, 44, 46, 47, 49;\quad $B$: 2, 3, 19, 20, 26, 30, 31, 32, 35, 39, 41, 45, 48;\quad $C$: 4, 6, 7, 9, 11, 17, 25, 29, 37, 38, 40, 43, 52;\quad $D$: 5, 8, 12, 14, 18, 21, 22, 23, 27, 33, 34, 50, 51.\\[2pt]
$p=61$, $g=2$, $f=11$ &
$A$: 1, 9, 12, 13, 15, 16, 20, 22, 25, 34, 42, 47, 56, 57, 58;\quad $B$: 2, 7, 18, 23, 24, 26, 30, 32, 33, 40, 44, 50, 51, 53, 55;\quad $C$: 3, 4, 5, 14, 19, 27, 36, 39, 41, 45, 46, 48, 49, 52, 60;\quad $D$: 6, 8, 10, 11, 17, 21, 28, 29, 31, 35, 37, 38, 43, 54, 59.\\[2pt]
$p=73$, $g=5$, $f=27$ &
$A$: 1, 2, 4, 8, 9, 16, 18, 32, 36, 37, 41, 55, 57, 64, 65, 69, 71, 72;\quad $B$: 5, 7, 10, 14, 17, 20, 28, 33, 34, 39, 40, 45, 53, 56, 59, 63, 66, 68;\quad $C$: 3, 6, 12, 19, 23, 24, 25, 27, 35, 38, 46, 48, 49, 50, 54, 61, 67, 70;\quad $D$: 11, 13, 15, 21, 22, 26, 29, 30, 31, 42, 43, 44, 47, 51, 52, 58, 60, 62.\\[2pt]
$p=89$, $g=3$, $f=34$ &
$A$: 1, 2, 4, 8, 11, 16, 22, 25, 32, 39, 44, 45, 50, 57, 64, 67, 73, 78, 81, 85, 87, 88;\quad $B$: 3, 6, 7, 12, 14, 23, 24, 28, 33, 41, 43, 46, 48, 56, 61, 65, 66, 75, 77, 82, 83, 86;\quad $C$: 5, 9, 10, 17, 18, 20, 21, 34, 36, 40, 42, 47, 49, 53, 55, 68, 69, 71, 72, 79, 80, 84;\quad $D$: 13, 15, 19, 26, 27, 29, 30, 31, 35, 37, 38, 51, 52, 54, 58, 59, 60, 62, 63, 70, 74, 76.\\[2pt]
$p=97$, $g=5$, $f=22$ &
$A$: 1, 4, 6, 9, 16, 22, 24, 33, 35, 36, 43, 47, 50, 54, 61, 62, 64, 73, 75, 81, 88, 91, 93, 96;\quad $B$: 5, 13, 14, 17, 19, 20, 21, 23, 29, 30, 41, 45, 52, 56, 67, 68, 74, 76, 77, 78, 80, 83, 84, 92;\quad $C$: 2, 3, 8, 11, 12, 18, 25, 27, 31, 32, 44, 48, 49, 53, 65, 66, 70, 72, 79, 85, 86, 89, 94, 95;\quad $D$: 7, 10, 15, 26, 28, 34, 37, 38, 39, 40, 42, 46, 51, 55, 57, 58, 59, 60, 63, 69, 71, 82, 87, 90.\\
\bottomrule
\end{longtable}

\subsection*{12.}

Since the number $2$ is a quadratic residue of all prime numbers of the form $8n+1$, and a non-residue of all those of the form $8n+5$, for prime moduli of the former form $2$ is found in class $A$ or $C$, and for moduli of the latter form in class $B$ or $D$. Since the distinction between classes $B$ and $D$ is not essential, depending only on the choice of the number $f$, we shall set the moduli of the form $8n+5$ aside for a time.

Subjecting the moduli of the form $8n+1$ to induction, we find that $2$ belongs to $A$ for
\[
p=73,89,113,233,257,281,337,353,\ldots;
\]
whereas $2$ belongs to $C$ for
\[
p=17,41,97,137,193,241,313,401,409,433,449,457,\ldots.
\]
Moreover, since for a prime modulus of the form $8n+1$ the number $-1$ is a biquadratic residue, it is clear that $-2$ is always to be referred to the same class as $+2$.

\subsection*{13.}

If the examples of the preceding article are compared with one another, at first sight no simple criterion seems to offer itself by which the former moduli may be distinguished from the latter. Nevertheless two such criteria exist, distinguished by elegance and simplicity. The following considerations will prepare the way for one of them.

The modulus $p$, as a prime number of the form $8n+1$, can be reduced, and in only one way, to the form
\[
p=a^2+2b^2
\]
(\emph{Disquiss. Arithm.}, art. 182, II). We suppose the roots $a,b$ to be taken positive. Manifestly $a$ is odd and $b$ is even. Put
\[
b=2^\lambda c,
\]
where $c$ is odd.

We first observe that since $p\equiv a^2\pmod c$, the number $p$ is a quadratic residue of $c$, and hence of each prime factor into which $c$ is resolved. Conversely, by the fundamental theorem, each of these prime factors is a quadratic residue of $p$, and therefore their product $c$ is also a quadratic residue of $p$. Since this also holds for the number $2$, it is clear that $b$ is a quadratic residue of $p$, and consequently that $b^2$, and also $-b^2$, is a biquadratic residue.

It follows that $-2b^2$ must be referred to the same class in which the number $2$ is found. Since $a^2\equiv -2b^2$, it is clear that $2$ is found either in class $A$ or in class $C$ according as $a$ is a quadratic residue or a quadratic non-residue of $p$.

Now suppose that $a$ has been resolved into its prime factors. Let those that are of the forms $8m+1$ or $8m+7$ be denoted by $\alpha,\alpha',\alpha''$ etc., and those of the forms $8m+3$ or $8m+5$ by $\beta,\beta',\beta''$ etc.; let the number of the latter be $\mu$. Since $p\equiv 2b^2\pmod \alpha$, $p$ is a quadratic residue of those prime factors of $a$ for which $2$ is a quadratic residue, namely the factors $\alpha,\alpha',\alpha''$ etc., and is a quadratic non-residue of those factors for which $2$ is a quadratic non-residue, namely $\beta,\beta',\beta''$ etc. Therefore, conversely, by the fundamental theorem, each of $\alpha,\alpha',\alpha''$ etc. is a quadratic residue of $p$, while each of $\beta,\beta',\beta''$ etc. is a quadratic non-residue. Hence the product $a$ is a quadratic residue or non-residue of $p$ according as $\mu$ is even or odd.

But it is easily confirmed that the product of all $\alpha,\alpha',\alpha''$ etc. is of the form $8m+1$ or $8m+7$, and the same holds for the product of all $\beta,\beta',\beta''$ etc. if the number of these is even; so in this case the product $a$ must also be of the form $8m+1$ or $8m+7$. On the other hand, when the number of the factors $\beta,\beta',\beta''$ etc. is odd, their product is of the form $8m+3$ or $8m+5$, and the same therefore holds in this case for the product $a$.

From all this the following elegant theorem is obtained: whenever $a$ is of the form $8m+1$ or $8m+7$, the number $2$ is contained in the complex $A$; whenever $a$ is of the form $8m+3$ or $8m+5$, the number $2$ is found in the complex $C$.

This is confirmed by the examples enumerated in the preceding article. The former moduli are split as follows:
\[
\begin{gathered}
73=1+2\cdot36,\quad 89=81+2\cdot4,\quad 113=81+2\cdot16,\\
233=225+2\cdot4,\quad 257=225+2\cdot16,\quad 281=81+2\cdot100,\\
337=49+2\cdot144,\quad 353=225+2\cdot64;
\end{gathered}
\]
the latter as follows:
\[
\begin{gathered}
17=9+2\cdot4,\quad 41=9+2\cdot16,\quad 97=25+2\cdot36,\\
137=9+2\cdot64,\quad 193=121+2\cdot36,\quad 241=169+2\cdot36,\\
313=25+2\cdot144,\quad 401=9+2\cdot196,\quad 409=121+2\cdot144,\\
433=361+2\cdot36,\quad 449=441+2\cdot4,\quad 457=169+2\cdot144.
\end{gathered}
\]

\subsection*{14.}

Since the decomposition of $p$ into a simple square and a double square has revealed such a striking connection with the classification of the number $2$, it seems worth trying whether the decomposition into two squares, to which the number $p$ is also subject, may perhaps provide a similar success. Here, then, are the decompositions of the numbers $p$ for which $2$ belongs to the class:
\[
\begin{array}{c|c}
A & C\\
\hline
9+64&1+16\\
25+64&25+16\\
49+64&81+16\\
169+64&121+16\\
1+256&49+144\\
25+256&225+16\\
81+256&169+144\\
289+64&1+400\\
&9+400\\
&289+144\\
&49+400\\
&441+16
\end{array}
\]
First of all we observe that, of the two squares into which $p$ is divided, one must be odd, which we shall put equal to $a^2$, and the other even, which we shall put equal to $b^2$. Since $a^2$ is of the form $8n+1$, it is clear that values of $b$ that are odd-even correspond to values of $p$ of the form $8n+5$, excluded here from our induction because they would have the number $2$ in class $B$ or $D$. For values of $p$ that are of the form $8n+1$, however, $b$ must be evenly even; and if the induction suggested by the displayed scheme is to be trusted, the number $2$ is to be referred to class $A$ for all moduli for which $b$ is of the form $8n$, and to class $C$ for all moduli for which $b$ is of the form $8n+4$.

But this theorem is of far deeper investigation than the one found in the preceding article. Several preliminary investigations must be put before its proof, concerning the order in which the numbers of the complexes $A,B,C,D$ follow one another.



\section*{Theory of Biquadratic Residues. First Memoir\\
Articles 15--23}

\subsection*{15.}

Let the number of numbers from the complex $A$ which are immediately followed by a number from $A,B,C,D$, respectively, be denoted by $(00),(01),(02),(03)$. Likewise let the number of numbers from the complex $B$ followed by a number from $A,B,C,D$ be denoted by $(10),(11),(12),(13)$; similarly let $(20),(21),(22),(23)$ be the corresponding numbers in $C$, and $(30),(31),(32),(33)$ the corresponding numbers in $D$. Our object is to determine these sixteen numbers a priori. To make it easier to compare the general reasoning with the examples, we record here the numerical values of the terms of the scheme $S$
\[
\begin{array}{llll}
(00),&(01),&(02),&(03)\\
(10),&(11),&(12),&(13)\\
(20),&(21),&(22),&(23)\\
(30),&(31),&(32),&(33)
\end{array}
\]
for the individual moduli for which the classifications were given in Article 11.

\[
\begin{array}{c|rrrr|c|rrrr|c|rrrr|c|rrrr}
&\multicolumn{4}{c|}{p=5}&&\multicolumn{4}{c|}{p=13}&&\multicolumn{4}{c|}{p=17}&&\multicolumn{4}{c}{p=29}\\
&0&1&0&0&&0&1&2&0&&0&2&1&0&&2&3&0&2\\
&0&0&0&1&&1&1&0&1&&2&0&1&1&&1&1&2&3\\
&0&0&0&0&&0&1&0&1&&1&1&1&1&&2&1&2&1\\
&0&0&1&0&&1&0&1&1&&0&1&1&2&&1&2&3&1
\end{array}
\]
\[
\begin{array}{c|rrrr|c|rrrr|c|rrrr|c|rrrr}
&\multicolumn{4}{c|}{p=37}&&\multicolumn{4}{c|}{p=41}&&\multicolumn{4}{c|}{p=53}&&\multicolumn{4}{c}{p=61}\\
&2&1&2&4&&0&4&3&2&&2&3&6&2&&4&3&2&6\\
&2&2&4&1&&4&2&2&2&&4&4&2&3&&3&3&6&3\\
&2&2&2&2&&3&2&3&2&&2&4&2&4&&4&3&4&3\\
&2&4&1&2&&2&2&2&4&&4&2&3&4&&3&6&3&3
\end{array}
\]
\[
\begin{array}{c|rrrr|c|rrrr|c|rrrr}
&\multicolumn{4}{c|}{p=73}&&\multicolumn{4}{c|}{p=89}&&\multicolumn{4}{c}{p=97}\\
&5&6&4&2&&3&8&6&4&&2&6&7&8\\
&6&2&5&5&&8&4&5&5&&6&8&5&5\\
&4&5&4&5&&6&5&6&5&&7&5&7&5\\
&2&5&5&6&&4&5&5&8&&8&5&5&6
\end{array}
\]

Since moduli of the forms $8n+1$ and $8n+5$ behave differently, the two cases must be treated separately. We begin with the former.

\subsection*{16.}

The character $(00)$ indicates in how many different ways the equation $\alpha+1=\alpha'$ can be satisfied, where $\alpha,\alpha'$ range over numbers of the complex $A$. Since for a modulus of the form $8n+1$, as is here understood, $\alpha'$ and $p-\alpha'$ belong to the same complex, we may say more neatly that $(00)$ expresses the number of different ways of satisfying
\[
1+\alpha+\alpha'=p;
\]
in place of this equation we may plainly use the congruence
\[
1+\alpha+\alpha'\equiv0\pmod p.
\]
Similarly, $(01)$ denotes the number of solutions of $1+\alpha+\beta\equiv0$, $(02)$ of $1+\alpha+\gamma\equiv0$, $(03)$ of $1+\alpha+\delta\equiv0$, $(11)$ of $1+\beta+\beta'\equiv0$, and so on, where $\beta,\beta'$ are numbers from $B$, $\gamma$ from $C$, and $\delta$ from $D$. From this one immediately obtains
\[
(01)=(10),\quad (02)=(20),\quad (03)=(30),\quad
(12)=(21),\quad (13)=(31),\quad (23)=(32).
\]

From any solution of $1+\alpha+\beta\equiv0$ there follows a solution of $\delta+\delta'+1\equiv0$ by choosing $\delta$ so that $\beta\delta\equiv1$ and $\delta'$ as the least positive residue of $\alpha\delta$. The inverse passage is equally clear; hence
\[
(01)=(33).
\]
The same argument gives
\[
(02)=(22),\qquad (03)=(11),\qquad (12)=(13)=(23).
\]
Thus eleven equations have been found among the sixteen unknowns, and the scheme $S$ may be written
\[
S=\begin{array}{rrrr}
h&i&k&l\\
i&l&m&m\\
k&m&k&m\\
l&m&m&i
\end{array}.
\]
The three independent row-sum conditions are
\[
h+i+k+l=2n-1,\qquad i+l+2m=2n,\qquad k+m=n.
\]
These reduce the sixteen unknowns to two.

\subsection*{17.}

To complete the determination, consider the number of solutions of
\[
1+\alpha+\beta+\gamma\equiv0\pmod p,
\]
where $\alpha,\beta,\gamma$ range through the complexes $A,B,C$. The value $\alpha=p-1$ is inadmissible, since then one could not have $\beta+\gamma\equiv0$. As $\alpha$ runs through the remaining values, the sums $1+\alpha$ fall in $A,B,C,D$ in the numbers $h,i,k,l$. Counting the corresponding solutions gives
\[
hm+ii+kl+lm.
\]
Counting the same congruence by first varying $\beta$ gives
\[
ik+lm+km+mm.
\]
The same number is obtained if the analysis is based on the possible values of $1+\gamma$.

\subsection*{18.}

Equating the two expressions for this number gives
\[
0=hm+ii+kl-ik-km-mm.
\]
Eliminating $h$ by means of $h=2m-k-1$ gives
\[
0=(k-m)^2+ii+kl-ik-kk-m.
\]
Since the last two equations of Article 16 give $k=\frac12(l+i)$, the part $ii+kl-ik-kk$ becomes $\frac14(l-i)^2$. Hence
\[
0=4(k-m)^2+(l-i)^2-4m.
\]
Using $4m=2(k+m)-2(k-m)=2n-2(k-m)$, one obtains
\[
2n=4(k-m)^2+2(k-m)+(l-i)^2,
\]
or
\[
8n+1=(4(k-m)+1)^2+4(l-i)^2.
\]
Putting
\[
4(k-m)+1=a,\qquad 2l-2i=b,
\]
we have
\[
p=a^2+b^2.
\]
The prime $p$ can be decomposed into two squares in only one way, with one square odd and the other even; thus $a^2$ and $b^2$ are completely determined. The sign of $a$ is also fixed: the root is taken positive or negative according as the positive root is of the form $4M+1$ or $4M+3$. The sign of $b$ will be determined next.

Combining these equations with those of Article 16 gives
\[
\begin{aligned}
8h&=4n-3a-5,&\quad 8i&=4n+a-2b-1,\\
8k&=4n+a-1,& 8l&=4n+a+2b-1,\\
8m&=4n-a+1.
\end{aligned}
\]
Equivalently, if the modulus $p$ is used in place of $n$, then the terms of $S$, multiplied by $16$, are
\[
16S=\begin{array}{rrrr}
p-6a-11&p+2a-4b-3&p+2a-3&p+2a+4b-3\\
p+2a-4b-3&p+2a+4b-3&p-2a+1&p-2a+1\\
p+2a-3&p-2a+1&p+2a-3&p-2a+1\\
p+2a+4b-3&p-2a+1&p-2a+1&p+2a-4b-3
\end{array}.
\]

\subsection*{19.}

It remains to assign the sign of $b$. The distinction between $B$ and $D$ depends on the choice of the number $f$, one of the two roots of $xx\equiv-1$. If the other root is chosen instead, $B$ and $D$ are interchanged. Inspection of the preceding scheme shows that this interchange corresponds to changing the sign of $b$, so there must be a relation between the sign of $b$ and $f$.

We first observe that, if $\mu$ is a non-negative integer and $z$ runs over $1,2,\ldots,p-1$, then modulo $p$ one has either $\sum z^\mu\equiv0$ or $\sum z^\mu\equiv-1$, according as $\mu$ is not divisible by $p-1$ or is divisible by $p-1$. The latter case is immediate. For the former, let $g$ be a primitive root; then the $z$ are the least residues of $g^y$, $y=0,1,\ldots,p-2$, and
\[
\sum g^{\mu y}=\frac{g^{\mu(p-1)}-1}{g^\mu-1},\qquad
(g^\mu-1)\sum z^\mu\equiv0.
\]
Since $g^\mu\not\equiv1$ when $\mu$ is not divisible by $p-1$, the sum is $0$. This proves the lemma.

Expanding $(x^4+1)^{\frac12(p-1)}$ and applying the lemma gives
\[
\sum (z^4+1)^{\frac12(p-1)}\equiv -2 \pmod p.
\]
The least residues of all $z^4$ reproduce all numbers in $A$, each four times; hence, by the criteria for the complexes,
\[
\sum (z^4+1)^{\frac12(p-1)}
 \equiv 4(00)+4f(01)-4(02)-4f(03),
\]
so that
\[
-2\equiv4(00)+4f(01)-4(02)-4f(03).
\]
Substituting the values found in the preceding article yields
\[
a+bf\equiv0\pmod p,
\]
or, after multiplication by $f$,
\[
b\equiv af\pmod p.
\]
This congruence determines the sign of $b$ once $f$ has been selected, or conversely determines $f$ if the sign of $b$ is prescribed.

\subsection*{20.}

The case $p\equiv5\pmod8$ is dealt with in the same manner. Here $-1$ belongs to $C$, and complements to $p$ of the numbers in $A,B,C,D$ belong respectively to $C,D,A,B$. This gives
\[
(00)=(22),\quad (01)=(32),\quad (03)=(12),\quad
(10)=(23),\quad (11)=(33),\quad (21)=(30).
\]
A further multiplication argument gives
\[
(00)=(20),\quad (01)=(13),\quad (03)=(31),\quad (10)=(11)=(21).
\]
Thus
\[
S=\begin{array}{rrrr}
h&i&k&l\\
m&m&l&i\\
h&m&h&m\\
m&l&i&m
\end{array}.
\]
The independent row-sum equations are
\[
h+i+k+l=2n+1,\qquad 2m+i+l=2n+1,\qquad h+m=n.
\]
Counting the solutions of $1+\alpha+\beta+\gamma\equiv0$ in two ways gives
\[
hm+il+ik+lm=hm+mm+hl+im,
\]
or
\[
0=mm+hl+im-il-ik-lm.
\]
Using $k=2m-h$ and $l+i=1+2h$, this becomes
\[
0=4mm-4m-1-8hm+4hh+(i-l)^2.
\]
Substituting $4m=2(h+m)-2(h-m)=2n-2(h-m)$ yields
\[
8n+5=(4(h-m)+1)^2+4(i-l)^2.
\]
Thus, putting
\[
4(h-m)+1=a,\qquad 2i-2l=b,
\]
one again has $p=a^2+b^2$. In this case the sign of $a$ is fixed by requiring $a\equiv1\pmod4$, and the sign of $b$ by $b\equiv af\pmod p$.

The resulting expressions are
\[
\begin{aligned}
8h&=4n+a-1,&\quad 8i&=4n+a+2b+3,\\
8k&=4n-3a+3,& 8l&=4n+a-2b+3,\\
8m&=4n-a+1.
\end{aligned}
\]
In terms of $p$, the entries of $S$, multiplied by $16$, are
\[
16S=\begin{array}{rrrr}
p+2a-7&p+2a+4b+1&p-2a+1&p+2a-4b+1\\
p-2a-3&p-2a-3&p+2a-4b+1&p+2a+4b+1\\
p+2a-7&p-2a-3&p+2a-7&p-2a-3\\
p-2a-3&p+2a-4b+1&p+2a+4b+1&p-2a-3
\end{array}.
\]

\subsection*{21.}

We now return to the principal question, namely the determination of the complex to which the number $2$ belongs.

First suppose $p$ is of the form $8n+1$. The number $2$ belongs either to $A$ or to $C$. In the former case $\frac12(p-1)$ and $\frac12(p+1)$ also belong to $A$; in the latter case they belong to $C$. If $\alpha$ and $\alpha+1$ are consecutive numbers of the complex $A$, then $p-\alpha-1$ and $p-\alpha$ have the same property; hence such numbers occur in associated pairs, except when one is associated with itself. Thus $(00)$ is odd precisely when $2$ belongs to $A$, and even when $2$ belongs to $C$. But
\[
16(00)=a^2+b^2-6a-11.
\]
Writing $a=4q+1$, $b=4r$, one obtains
\[
(00)=q^2-q+r^2-r-1.
\]
Since $q^2-q$ is always even, $(00)$ is odd or even according as $r$ is even or odd. Therefore $2$ belongs to $A$ or to $C$ according as $b$ is of the form $8m$ or $8m+4$, exactly as was found by induction in Article 14.

Second, when $p$ is of the form $8n+5$, the number $2$ belongs either to $B$ or to $D$. The same pairing argument shows that $(13)$ is even when $2$ belongs to $D$ and odd when $2$ belongs to $B$. Since
\[
16(13)=a^2+b^2+2a+4b+1,
\]
and, writing $a=4q+1$, $b=4r+2$,
\[
(13)=q^2+q+r^2+2r+1,
\]
one obtains that $2$ belongs to $B$ when $b$ has the form $8m+2$, and to $D$ when $b$ has the form $8m+6$.

The result may be stated compactly as follows: the number $2$ belongs to the complex $A,B,C,$ or $D$ according as $\frac12 b$ is of the form $4m,4m+1,4m+2,$ or $4m+3$.

\subsection*{22.}

In the \emph{Disquisitiones Arithmeticae} we explained the general theory of the division of the circle and the solution of $x^p-1=0$. In particular, if $\mu$ divides $p-1$, the function
\[
\frac{x^p-1}{x-1}
\]
can be resolved into $\mu$ factors of degree $(p-1)/\mu$ by means of an auxiliary equation of degree $\mu$. The cases $\mu=2$ and $\mu=3$ were also treated separately there, and the auxiliary equation was assigned a priori, without first developing the scheme of least residues of powers of a primitive root modulo $p$. Attentive readers will see the close connection of the next case, $\mu=4$, with the investigations of Articles 15--20. With their aid that case could also be completely solved. We reserve this treatment for another occasion and, in the present memoir, have preferred to keep the investigation in a purely arithmetical form, without bringing in the theory of the equation $x^p-1=0$. As a closing addition, however, we add some new purely arithmetical theorems closely connected with what has just been treated.

\subsection*{23.}

When $(x^4+1)^{\frac12(p-1)}$ is expanded by the binomial theorem, three terms occur whose exponents of $x$ are divisible by $p-1$: the first, middle, and last. If $P$ denotes the middle coefficient
\[
P=\frac{\frac12(p-1)\cdot\frac12(p-3)\cdot\frac12(p-5)\cdots\frac14(p+3)}
{1\cdot2\cdot3\cdots\frac14(p-1)},
\]
then, substituting for $x$ the numbers $1,2,3,\ldots,p-1$, Article 19 gives
\[
\sum (x^4+1)^{\frac12(p-1)}\equiv -2-P \pmod p.
\]
On the other hand, since the numbers in $A,B,C,D$ raised to the power $\frac12(p-1)$ are respectively congruent to $+1,-1,+1,-1$, we get
\[
\sum (x^4+1)^{\frac12(p-1)}
 \equiv 4(00)-4(01)+4(02)-4(03).
\]
Using the schemes at the ends of Articles 18 and 20, this is
\[
\sum (x^4+1)^{\frac12(p-1)}\equiv -2a-2.
\]
Comparison gives the elegant congruence
\[
P\equiv 2a\pmod p.
\]

Let the four products
\[
\begin{gathered}
1\cdot2\cdot3\cdots\frac14(p-1),\\
\frac14(p+3)\cdot\frac14(p+7)\cdot\frac14(p+11)\cdots\frac12(p-1),\\
\frac12(p+1)\cdot\frac12(p+3)\cdot\frac12(p+5)\cdots\frac34(p-1),\\
\frac14(3p+1)\cdot\frac14(3p+5)\cdot\frac14(3p+9)\cdots(p-1)
\end{gathered}
\]
be denoted respectively by $q,r,s,t$. The preceding theorem is then
\[
2a\equiv \frac{r}{q}\pmod p.
\]
Each factor of $q$ has its complement to $p$ in $t$; hence $q\equiv t$ for $p\equiv1\pmod8$ and $q\equiv-t$ for $p\equiv5\pmod8$. Similarly $r\equiv s$ in the former case and $r\equiv -s$ in the latter. In both cases $qr\equiv st$; and since $qrst\equiv-1$, it follows that $q^2r^2\equiv-1$, so that
\[
qr\equiv\pm f\pmod p.
\]
Combining this with the preceding congruence gives
\[
r^2\equiv\pm 2af,
\]
and therefore, by Articles 19 and 20,
\[
2b\equiv\pm r^2\pmod p.
\]
It is noteworthy that the decomposition of $p$ into two squares can be found by entirely direct operations: the root of the odd square is the absolutely least residue of $\frac{r}{2q}$, and the root of the even square is the absolutely least residue of $\frac14r^2$ modulo $p$. The expression $\frac{r}{2q}$, whose value is $1$ for $p=5$, may for larger $p$ be written
\[
\frac{6\cdot10\cdot14\cdot18\cdots(p-3)}
{2\cdot3\cdot4\cdot5\cdots\frac12(p-1)}.
\]
Since we know the sign with which the odd root arises from this formula, namely so that it is always of the form $4m+1$, it is noteworthy that no analogous general criterion has yet been found for the sign of the even root. Anyone who finds and communicates such a criterion would put us greatly in his debt.

We add here the values of $a,b,f$ for $p<200$, as obtained from the least residues of $\frac{r}{2q}$, $\frac14r^2$, and $qr$:
\[
\begin{array}{r|r|r|r}
p&a&b&f\\\hline
5&+1&+2&2\\
13&-3&-2&5\\
17&+1&-4&13\\
29&+5&+2&12\\
37&+1&-6&31\\
41&+5&+4&9\\
53&-7&-2&23\\
61&+5&-6&11\\
73&-3&-8&27\\
89&+5&-8&34\\
97&+9&+4&22\\
101&+1&-10&91\\
109&-3&+10&33\\
113&-7&+8&15\\
137&-11&+4&37\\
149&-7&-10&44\\
157&-11&-6&129\\
173&+13&+2&80\\
181&+9&+10&162\\
193&-7&+12&81\\
197&+1&-14&183
\end{array}
\]
and
\[
\{(a\mp b)q\}^2\equiv a\equiv\left(\frac{r-qr^2}{2}\right)^2.
\]



\section*{Theory of Biquadratic Residues. Second Memoir\\
Title page and Articles 24--29}

\begin{center}
\textsc{Theory of Biquadratic Residues}\\
\textsc{Second Memoir}\\[2pt]
\textsc{By Carl Friedrich Gauss}\\
\textsc{Presented to the Royal Society, 15 April 1831}\\
\emph{Commentationes Societatis Regiae Scientiarum Gottingensis recentiores}, Vol. VII.\\
Göttingen, 1832.
\end{center}

\subsection*{24.}

In the first memoir, all that is required for the biquadratic classification of the number $+2$ was completely settled. We divide all numbers not divisible by a modulus $p$, where $p$ is a prime of the form $4n+1$, among four complexes $A,B,C,D$ according as they become congruent modulo $p$ to $+1,+f,-1,-f$ when raised to the exponent $\frac14(p-1)$, where $f$ denotes either root of $ff\equiv-1\pmod p$. We found that the decision as to the complex to which $+2$ belongs depends on the decomposition of $p$ into two squares. Thus, if $p=a^2+b^2$, where $a^2$ is the odd square and $b^2$ the even square, and if the signs of $a$ and $b$ are chosen so that $a\equiv1\pmod4$ and $b\equiv af\pmod p$, then $+2$ belongs to $A,B,C,D$ according as $\frac12 b$ has the form $4n,4n+1,4n+2,4n+3$, respectively.

The rule for classifying $-2$ follows at once. Since $-1$ belongs to $A$ when $\frac12 b$ is even, and to $C$ when it is odd, Article 7 gives that $-2$ belongs to $A,B,C,D$ according as $\frac12 b$ has the form $4n,4n+3,4n+2,4n+1$, respectively. Equivalently:
\[
\begin{array}{c|cc}
\text{Complex} & +2 & -2\\
\hline
A & b\equiv 0 & b\equiv 0\\
B & b\equiv 2a & b\equiv 6a\\
C & b\equiv 4a & b\equiv 4a\\
D & b\equiv 6a & b\equiv 2a
\end{array}
\qquad(\bmod\,8).
\]
It is easy to see that, stated in this way, the theorems no longer depend on the condition $a\equiv1\pmod4$; they remain valid also when $a\equiv3\pmod4$, provided the other condition $af\equiv b\pmod p$ is retained. The whole result can also be compressed into the single formula
\[
b^{\frac12ab}\equiv a^{\frac12ab}\,2^{\frac14(p-1)}\pmod p,
\]
when $a$ and $b$ are taken positive.

\subsection*{25.}

We next ask how far induction indicates the classification of the number $3$. Continuing the table of Article 11, always taking the least primitive root, shows that $+3$ belongs to
\[
\begin{array}{c|rr@{\quad}c|rr@{\quad}c|rr@{\quad}c|rr}
\multicolumn{3}{c}{A\ \text{for}}&
\multicolumn{3}{c}{B\ \text{for}}&
\multicolumn{3}{c}{C\ \text{for}}&
\multicolumn{3}{c}{D\ \text{for}}\\
p&a&b&p&a&b&p&a&b&p&a&b\\\hline
13&-3&+2&17&+1&-4&37&+1&-6&5&+1&+2\\
109&-3&+10&29&+5&+2&61&+5&-6&41&+5&-4\\
181&+9&+10&53&-7&+2&73&-3&-8&149&-7&+10\\
193&-7&-12&89&+5&-8&97&+9&+4&173&+13&+2\\
229&-15&+2&101&+1&+10&157&-11&-6&&&\\
277&+9&+14&113&-7&-8&241&-15&-4&&&\\
&&&137&-11&-4&&&&&&\\
&&&197&+1&-14&&&&&&\\
&&&233&+13&+8&&&&&&\\
&&&257&+1&-16&&&&&&\\
&&&269&+13&+10&&&&&&\\
&&&281&+5&+16&&&&&&\\
&&&293&+17&+2&&&&&&
\end{array}
\]
At first sight no simple connection appears among the values $a,b$ that correspond to the same complex. But if we recall that, in the theory of quadratic residues, the analogous decision is simpler for $-3$ than for $+3$, there is hope for an equally successful outcome here. We find that $-3$ belongs to
\[
\begin{array}{c|rr@{\quad}c|rr@{\quad}c|rr@{\quad}c|rr}
\multicolumn{3}{c}{A\ \text{for}}&
\multicolumn{3}{c}{B\ \text{for}}&
\multicolumn{3}{c}{C\ \text{for}}&
\multicolumn{3}{c}{D\ \text{for}}\\
p&a&b&p&a&b&p&a&b&p&a&b\\\hline
37&+1&-6&5&+1&+2&13&-3&+2&29&+5&+2\\
61&+5&-6&17&+1&-4&73&-3&-8&41&+5&-4\\
157&-11&-6&89&+5&-8&97&+9&+4&53&-7&+2\\
193&-7&-12&113&-7&-8&109&-3&+10&101&+1&+10\\
&&&137&-11&-4&181&+9&+10&197&+1&-14\\
&&&149&-7&+10&229&-15&+2&269&+13&+10\\
&&&173&+13&+2&241&-15&-4&293&+17&+2\\
&&&233&+13&+8&277&+9&+14&&&\\
&&&257&+1&-16&&&&&&\\
&&&281&+5&+16&&&&&&
\end{array}
\]
and here the inductive law presents itself immediately:
\[
\begin{array}{l|l}
A& b\equiv0\pmod3\\
B& a+b\equiv0\pmod3,\ \text{that is } b\equiv2a\pmod3\\
C& a\equiv0\pmod3\\
D& a-b\equiv0\pmod3,\ \text{that is } b\equiv a\pmod3.
\end{array}
\]

\subsection*{26.}

The number $+5$ is found to belong to
\[
\begin{array}{l|l}
A& p=101,109,149,181,269\\
B& p=13,17,73,97,157,193,197,233,277,293\\
C& p=29,41,61,89,229,241,281\\
D& p=37,53,113,137,173,257.
\end{array}
\]
Taking into account the values of $a,b$ belonging to each $p$, the law is just as easy to grasp as for the classification of $-3$:
\[
\begin{array}{l|l}
A& b\equiv0\pmod5\\
B& b\equiv a\pmod5\\
C& a\equiv0\pmod5\\
D& b\equiv4a\pmod5.
\end{array}
\]
These rules plainly cover all cases; for if $b\equiv2a$ or $b\equiv3a\pmod5$, then $a^2+b^2\equiv0$, contrary to the hypothesis that $p$ is a prime different from $5$.

\subsection*{27.}

Applying induction in the same way to the numbers $-7,-11,+13,+17,-19,-23$, and carrying it far enough, gives the following rules:
\[
\begin{array}{c|l}
\multicolumn{2}{c}{\text{For }-7}\\\hline
A& a\equiv0,\ \text{or }b\equiv0\pmod7\\
B& b\equiv4a,\ \text{or }5a\\
C& b\equiv a,\ \text{or }6a\\
D& b\equiv2a,\ \text{or }3a
\end{array}
\]
\[
\begin{array}{c|l}
\multicolumn{2}{c}{\text{For }-11}\\\hline
A& b\equiv0,5a,\ \text{or }6a\pmod{11}\\
B& b\equiv a,3a,\ \text{or }4a\\
C& a\equiv0,\ \text{or }b\equiv2a,\ \text{or }9a\\
D& b\equiv7a,8a,\ \text{or }10a
\end{array}
\]
\[
\begin{array}{c|l}
\multicolumn{2}{c}{\text{For }+13}\\\hline
A& b\equiv0,4a,9a\pmod{13}\\
B& b\equiv6a,11a,12a\\
C& a\equiv0;\ b\equiv3a,10a\\
D& b\equiv a,2a,7a
\end{array}
\]
\[
\begin{array}{c|l}
\multicolumn{2}{c}{\text{For }+17}\\\hline
A& a\equiv0;\ b\equiv0,a,16a\pmod{17}\\
B& b\equiv2a,6a,8a,14a\\
C& b\equiv5a,7a,10a,12a\\
D& b\equiv3a,9a,11a,15a
\end{array}
\]
\[
\begin{array}{c|l}
\multicolumn{2}{c}{\text{For }-19}\\\hline
A& b\equiv0,2a,5a,14a,17a\pmod{19}\\
B& b\equiv3a,7a,11a,13a,18a\\
C& a\equiv0;\ b\equiv4a,9a,10a,15a\\
D& b\equiv a,6a,8a,12a,16a
\end{array}
\]
\[
\begin{array}{c|l}
\multicolumn{2}{c}{\text{For }-23}\\\hline
A& a\equiv0;\ b\equiv0,7a,10a,13a,16a\pmod{23}\\
B& b\equiv2a,3a,4a,11a,15a,17a\\
C& b\equiv a,5a,9a,14a,18a,22a\\
D& b\equiv6a,8a,12a,19a,20a,21a.
\end{array}
\]

\subsection*{28.}

The special theorems obtained by induction in this way are confirmed as far as the induction is continued, and they exhibit a very beautiful form of criteria. When they are compared in order to draw general conclusions from them, the following observations occur at once.

The criteria for deciding to which complex the prime number $\pm q$ is to be referred, taking the upper or lower sign according as $q$ has the form $4n+1$ or $4n+3$, depend on the forms of $a,b$ compared with one another modulo $q$.

First, when $a\equiv0\pmod q$, $\pm q$ belongs to a fixed complex: namely to $A$ for $q=7,17,23$, and to $C$ for $q=3,11,13,19$. This suggests that the former case holds in general when $q$ is of the form $8n\pm1$, and the latter when $q$ is of the form $8n\pm3$. Moreover, $B$ and $D$ are excluded without induction when $a$ is divisible by $q$, since then $p\equiv b^2\pmod q$, so that $p$ is a quadratic residue of $q$, and therefore by the fundamental theorem $\pm q$ must be a quadratic residue of $p$.

Second, when $a$ is not divisible by $q$, the criterion depends on the value of $\frac{b}{a}\pmod q$. This expression admits $q$ different values, $0,1,2,\ldots,q-1$; but when $q$ is of the form $4n+1$, the two values of $\sqrt{-1}\pmod q$ must be excluded, since they plainly cannot be values of $\frac{b}{a}$ under the standing assumption that $p=a^2+b^2$ is a prime distinct from $q$. Thus the number of admissible values of $\frac{b}{a}\pmod q$ is $q-2$ for $q\equiv1\pmod4$, while it remains $q$ for $q\equiv3\pmod4$.

These values are distributed into four classes. Let certain values, denoted indefinitely by $\alpha$, correspond to $A$; others, denoted by $\beta$, to $B$; others, $\gamma$, to $C$; and the remaining values, $\delta$, to $D$. Thus $\pm q$ is assigned to $A,B,C,D$ according as $b\equiv\alpha a$, $b\equiv\beta a$, $b\equiv\gamma a$, $b\equiv\delta a\pmod q$. The law of this distribution appears more hidden, although certain general facts can be noticed at once. Three of the classes have the same number of values, $\frac14(q-1)$ or $\frac14(q+1)$, while in one class, namely the class corresponding to the complex with criterion $a\equiv0$, the number is smaller by one. Hence the total number of distinct criteria for each complex is the same, namely $\frac14(q-1)$ or $\frac14(q+1)$.

We also observe that $0$ is always in the first class, among the $\alpha$, and that the complements to $q$ of the numbers $\alpha,\beta,\gamma,\delta$ lie respectively in the first, fourth, third, and second classes. Finally, the values of $1/\alpha,1/\beta,1/\gamma,1/\delta\pmod q$ are seen to belong to the first, fourth, third, and second classes when the criterion $a\equiv0$ corresponds to $A$; but to the third, second, first, and fourth classes, respectively, when that criterion is referred to $C$. This is about the limit of what induction can provide, unless we boldly anticipate results that will below be drawn from their genuine sources.

\subsection*{29.}

Before proceeding further, it is useful to observe that the criteria for prime numbers, taken positive when they are of the form $4n+1$ and negative when they are of the form $4n+3$, are enough to decide the case of all remaining numbers, provided Article 7 and the criteria for $-1$ and $\pm2$ are also used. For example, if criteria are wanted for $+3$, the criteria given in Article 25 for $-3$ remain valid for $+3$ when $\frac12b$ is even; but when $\frac12b$ is odd, the complexes $A,B,C,D$ must be permuted with $C,D,A,B$. Hence:
\[
\begin{array}{c|l}
\multicolumn{2}{c}{+3\ \text{belongs}}\\\hline
A& b\equiv0\pmod{12};\ \text{or also }a\equiv0\pmod3,\ b\equiv2\pmod4\\
B& b\equiv8a,\ \text{or }10a\pmod{12}\\
C& b\equiv6a\pmod{12};\ \text{or also }a\equiv0\pmod3,\ b\equiv0\pmod4\\
D& b\equiv2a,\ \text{or }4a\pmod{12}.
\end{array}
\]
Likewise the criteria for $\pm6$ are obtained by combining the criteria for $\mp2$ and $-3$:
\[
\begin{array}{c|l}
\multicolumn{2}{c}{+6\ \text{belongs}}\\\hline
A& b\equiv0,2a,22a\pmod{24};\ \text{or also }a\equiv0\pmod3,\ b\equiv4a\pmod8\\
B& b\equiv4a,6a,8a\pmod{24};\ \text{or also }a\equiv0\pmod3,\ b\equiv2a\pmod8\\
C& b\equiv10a,12a,14a\pmod{24};\ \text{or also }a\equiv0\pmod3,\ b\equiv0\pmod8\\
D& b\equiv16a,18a,20a\pmod{24};\ \text{or also }a\equiv0\pmod3,\ b\equiv6a\pmod8,
\end{array}
\]
whereas $-6$ belongs as follows:
\[
\begin{array}{c|l}
\multicolumn{2}{c}{-6\ \text{belongs}}\\\hline
A& b\equiv0,10a,14a\pmod{24};\ \text{or also }a\equiv0\pmod3,\ b\equiv4a\pmod8\\
B& b\equiv4a,8a,18a\pmod{24};\ \text{or also }a\equiv0\pmod3,\ b\equiv6a\pmod8\\
C& b\equiv2a,12a,22a\pmod{24};\ \text{or also }a\equiv0\pmod3,\ b\equiv0\pmod8\\
D& b\equiv6a,16a,20a\pmod{24};\ \text{or also }a\equiv0\pmod3,\ b\equiv2a\pmod8.
\end{array}
\]
In the same way, criteria for $+21$ will be assembled from the criteria for $-3$ and $-7$; criteria for $-105$ from those for $-1,-3,+5,-7$, and so on.


\clearpage
\art{30}

Induction therefore opens an exceedingly rich harvest of special theorems akin to the theorem for the number $2$; but the common bond is lacking, and rigorous demonstrations are lacking, since the method by which, in the first memoir, we settled the number $2$ does not admit further application. There are indeed other methods by which one could obtain demonstrations for particular cases, especially those concerning the distribution of quadratic residues among the complexes $A$ and $C$; but we shall not dwell on them, since what is wanted is a general theory embracing all cases.

When, from the year 1805 onward, we began to devote our reflections to this matter, we soon became certain that the genuine source of the general theory must be sought in a promoted, or extended, field of arithmetic, as we already indicated in Article~1.

Just as higher arithmetic, in the questions treated up to this point, ranges only over real integers, so the theorems concerning biquadratic residues shine forth in their full simplicity and genuine beauty only when the field of arithmetic is extended to imaginary quantities, so that, without restriction, its objects are the numbers of the form $a+bi$, where $i$ denotes, as usual, the imaginary quantity $\sqrt{-1}$, and $a$ and $b$ range indefinitely over all real integers from $-\infty$ to $+\infty$. We shall call such numbers \textit{complex integral numbers}, in such a way that real numbers are not opposed to complex numbers, but are instead regarded as a species included among them. The present memoir will set forth both the elementary doctrine of complex numbers and the first beginnings of the theory of biquadratic residues; we shall undertake in a following continuation to make that theory complete on every side.\footnote{It remains useful to observe here, at least in passing, that the field thus defined is especially adapted to the theory of biquadratic residues. The theory of cubic residues is similarly to be built upon the consideration of numbers of the form $a+bh$, where $h$ is an imaginary root of the equation $h^3-1=0$, for instance $h=-\frac{1}{2}+\sqrt{\frac{3}{4}}\,i$; and likewise the theory of residues of higher powers will require the introduction of other imaginary quantities.}

\art{31}

First of all we put forward certain names, whose introduction will serve brevity and clarity. The field of complex numbers $a+bi$ contains

\begin{enumerate}[label=\Roman*.]
\item real numbers, where $b=0$, and among these, according to the nature of $a$,
    \begin{enumerate}[label=\arabic*)]
    \item zero,
    \item positive numbers,
    \item negative numbers;
    \end{enumerate}
\item imaginary numbers, where $b$ is unequal to zero. These are again divided into
    \begin{enumerate}[label=\arabic*)]
    \item imaginary numbers without real part, that is, where $a=0$,
    \item imaginary numbers with real part, where neither $b$ nor $a$ is zero.
    \end{enumerate}
\end{enumerate}

The former may, if desired, be called \textit{pure imaginary numbers}; the latter, \textit{mixed imaginary numbers}.

In this doctrine we use four units, $+1$, $-1$, $+i$, $-i$, which will be called, simply, the positive, negative, positive imaginary, and negative imaginary units. The three products of any complex number by $-1$, $+i$, $-i$ we shall call its \textit{associates}, or the numbers \textit{associated} with it. Except for zero, which is associated with itself, four unequal associated numbers therefore always occur.

Conversely, we call \textit{conjugate} to a complex number the number that arises from it by replacing $i$ by $-i$. Thus among imaginary numbers there are always two unequal conjugates, whereas real numbers are conjugate to themselves, if one chooses to extend the name to them.

The product of a complex number by its conjugate we call the \textit{norm} of both. Thus the norm of a real number is to be taken as its square. In general we have eight connected numbers, namely
\[
\begin{array}{r|r}
 a+bi & a-bi\\
 -b+ai & -b-ai\\
 -a-bi & -a+bi\\
 b-ai & b+ai
\end{array}
\]
in which we see two quaternions of associated numbers and four pairs of conjugates; the common norm of all is $aa+bb$. But these eight numbers reduce to four unequal ones whenever either $a=\pm b$ or one of the numbers $a$, $b$ is zero.

From these definitions the following facts follow immediately. The conjugate of a product of two complex numbers is the product of the numbers conjugate to them. The same is true of a product of several factors, and also of quotients. The norm of a product of two complex numbers is equal to the product of their norms. This theorem likewise extends to products of arbitrarily many factors and to quotients. The norm of any complex number, except zero, an exception that we shall henceforth usually understand tacitly, is a positive number.

There is, moreover, nothing to prevent our definitions from being extended to fractional or even irrational values of $a$ and $b$; but $a+bi$ will be called a \textit{complex integral number} only when both $a$ and $b$ are integers, and it will be called \textit{rational} only when both $a$ and $b$ are rational.

\art{32}

The algorithm of arithmetic operations on complex numbers is commonly known. Division is reduced to multiplication by the introduction of the norm, since
\[
\frac{a+bi}{c+di}=(a+bi)\cdot\frac{c-di}{cc+dd}=\frac{ac+bd}{cc+dd}+\frac{bc-ad}{cc+dd}\,i.
\]
Extraction of a square root is effected with the aid of the formula
\[
\sqrt{a+bi}=\pm\left(\sqrt{\frac{\sqrt{aa+bb}+a}{2}}+i\sqrt{\frac{\sqrt{aa+bb}-a}{2}}\right)
\]
if $b$ is positive, or of the formula
\[
\sqrt{a+bi}=\pm\left(\sqrt{\frac{\sqrt{aa+bb}+a}{2}}-i\sqrt{\frac{\sqrt{aa+bb}-a}{2}}\right)
\]
if $b$ is negative. There is no need here to dwell on the use of transforming the complex quantity $a+bi$ into $r(\cos\varphi+i\sin\varphi)$ in order to make calculations easier.

\art{33}

A complex integer that can be resolved into two factors different from units\footnote{Or, what is the same thing, factors whose norms are greater than unity.} we call a \textit{composite complex number}; on the other hand, a number that admits no such resolution into factors will be called a \textit{complex prime}. It follows at once that every composite real number is also composite as a complex number.

But a real prime number may be composite as a complex number. This is true of the number $2$, and of all positive real primes of the form $4n+1$ (apart from the number $1$), since they are known to be decomposable into two positive squares; for example,
\[
2=(1+i)(1-i),\quad 5=(1+2i)(1-2i),\quad 13=(3+2i)(3-2i),\quad 17=(1+4i)(1-4i),\quad \text{etc.}
\]
On the other hand, positive real primes of the form $4n+3$ are always complex primes. For if such a number $q$ were equal to $(a+bi)(\alpha+\beta i)$, then also $q=(a-bi)(\alpha-\beta i)$, and hence
\[
qq=(aa+bb)(\alpha\alpha+\beta\beta).
\]
But $qq$ can be resolved in only one way into positive factors greater than unity, namely as $q\times q$; hence one would have to have $q=aa+bb=\alpha\alpha+\beta\beta$, which is absurd, since a sum of two squares cannot be of the form $4n+3$.

Negative real numbers plainly retain the same denominations as the corresponding positive ones, and the same is true of pure imaginary numbers.

It remains to teach how, among mixed imaginary numbers, composites are to be distinguished from primes. This is done by the following theorem.

\medskip
\noindent\textsc{Theorem.} Every mixed imaginary complex integer $a+bi$ is either a complex prime or a composite number, according as its norm is a real prime number or a composite number.

\medskip
\noindent\textsc{Proof.} I. Since the norm of a composite complex number is always composite, it is clear that a complex number whose norm is a real prime number must necessarily be a complex prime. This proves the first assertion.

II. If the norm $aa+bb$ is composite, let $p$ be a positive real prime measuring it. Two cases must now be distinguished.

1) If $p$ is of the form $4n+3$, it is known that $aa+bb$ cannot be divisible by $p$ unless $p$ measures $a$ and $b$ themselves; hence $a+bi$ will be composite.

2) If $p$ is not of the form $4n+3$, it can certainly be decomposed into two squares. We therefore put $p=\alpha\alpha+\beta\beta$. Since
\[
(a\alpha+b\beta)(a\alpha-b\beta)=aa(\alpha\alpha+\beta\beta)-\beta\beta(aa+bb)
\]
is therefore divisible by $p$, $p$ must measure one of the two factors $a\alpha+b\beta$ and $a\alpha-b\beta$. Moreover, since
\[
(a\alpha+b\beta)^2+(b\alpha-a\beta)^2=(a\alpha-b\beta)^2+(b\alpha+a\beta)^2=(aa+bb)(\alpha\alpha+\beta\beta)
\]
is divisible by $pp$, it is clear that in the first case $b\alpha-a\beta$, and in the second case $b\alpha+a\beta$, must also be divisible by $p$. Therefore in the first case
\[
\frac{a+bi}{\alpha+\beta i}=\frac{a\alpha+b\beta}{p}+\frac{b\alpha-a\beta}{p}\,i
\]
will be a complex integer, while in the second case
\[
\frac{a+bi}{\alpha-\beta i}=\frac{a\alpha-b\beta}{p}+\frac{b\alpha+a\beta}{p}\,i
\]
will be an integer. Since the proposed number is thus divisible either by $\alpha+\beta i$ or by $\alpha-\beta i$, and since by hypothesis the norm of the quotient, $\frac{aa+bb}{p}$, is different from unity, it is clear in both cases that $a+bi$ is a composite complex number. This proves the second assertion.

\art{34}

The whole range of complex prime numbers is therefore exhausted by the following four species:
\begin{enumerate}[label=\arabic*)]
\item the four units $1$, $+i$, $-1$, $-i$, which, however, when we deal with prime numbers, we shall usually understand to be tacitly excluded;
\item the number $1+i$, together with its three associates $-1+i$, $-1-i$, $1-i$;
\item positive real primes of the form $4n+3$, together with their three associates;
\item complex numbers whose norms are real primes of the form $4n+1$, greater than unity; moreover, to each such given norm there always correspond exactly eight complex primes, since such a norm can be decomposed into two squares in one way only.
\end{enumerate}

\art{35}

Just as real integers are divided into even and odd numbers, and the even numbers again into evenly even and unevenly even, so among complex numbers an equally essential distinction presents itself. Namely, they are either

\textit{not} divisible by $1+i$, as are the numbers $a+bi$ where one of $a$, $b$ is odd and the other even;

\textit{or} divisible by $1+i$ but not by $2$, whenever both $a$ and $b$ are odd;

\textit{or} divisible by $2$, whenever both $a$ and $b$ are even.

Numbers of the first class may conveniently be called \textit{odd complex numbers}; those of the second, \textit{semieven}; those of the third, \textit{even}.

A product of several complex factors will always be odd whenever all the factors are odd; semieven whenever one factor is semieven and the rest are odd; and even whenever among the factors there occur either at least two semieven factors or at least one even factor.

The norm of every odd complex number is of the form $4n+1$; the norm of a semieven number is of the form $8n+2$; finally, the norm of an even number is the product of a number of the form $4n+1$ by $4$ or by a higher power of two.

\art{36}

Since the connection among four associated complex numbers is analogous to the connection between two opposite real numbers, that is, numbers absolutely equal and affected with opposite signs, and since among such real numbers the positive one is commonly and rightly considered primary, the question arises whether a similar distinction can be established among four associated complex numbers, and whether it should be considered useful.

In order to decide this, one must observe that the principle of distinction ought to be so constituted that the product of two numbers that are primary among their associates is always primary among its associates. But we soon become certain that no such principle exists at all unless the distinction is restricted to integral numbers; indeed, the useful distinction must be limited to the odd numbers. For these, however, the proposed aim can be attained in two ways:

I. The product of two numbers $a+bi$ and $a'+b'i$ so arranged that $a$ and $a'$ are of the form $4n+1$ and $b$ and $b'$ are even will enjoy the same property, so that the real part is $\equiv 1 \pmod{4}$ and the imaginary part is even. It is easy to see that among the four odd associates exactly one is contained under that form.

II. If the number $a+bi$ is so arranged that $a-1$ and $b$ are either both evenly even or both unevenly even, then its product by a complex number of the same form will enjoy the same form; and again it is easy to see that among the four associated odd numbers exactly one is contained under this form.

Of these two nearly equally suitable principles we shall adopt the latter: among the four associated odd complex numbers, we shall regard as primary that one which, modulo $2+2i$, becomes congruent to the positive unit. In this way several noteworthy theorems can be stated with greater neatness. Thus, for example, the primary complex primes are $-1+2i$, $-1-2i$, $+3+2i$, $+3-2i$, $+1+4i$, $+1-4i$, etc., and also the real primes $-3$, $-7$, $-11$, $-19$, etc., which must plainly always be taken with the negative sign. The conjugate of an odd primary complex number will also be primary.

For semieven and even numbers in general, a similar distinction would be too arbitrary and of little use. From the associated prime numbers $1+i$, $1-i$, $-1+i$, $-1-i$, we can indeed choose one as primary before the others; but we shall not extend such a distinction to composites.

\art{37}

If among the factors of a composite complex number there are some that are themselves composite, and if these are again resolved into their own factors, then one plainly comes at last to prime factors; that is, every composite number is resoluble into prime factors. If among these there are any that are not primary, substitute in place of each the product of the associated primary number by $i$, $-1$, or $-i$. In this way it is clear that every composite complex number $M$ can be reduced to the form
\[
M=i^{\mu}A^{\alpha}B^{\beta}C^{\gamma}\ldots,
\]
where $A$, $B$, $C$, etc. are unequal primary complex primes, and $\mu=0$, $1$, $2$, or $3$. Concerning this resolution the theorem presents itself that it can be made in one way only. Considered casually, this theorem might seem evident in itself, but it certainly needs a demonstration. The following prepares the way for it.

\medskip
\noindent\textsc{Theorem.} The product $M=A^{\alpha}B^{\beta}C^{\gamma}\ldots$, where $A$, $B$, $C$, etc. denote distinct primary complex primes, cannot be divisible by any primary complex prime that is not found among $A$, $B$, $C$, etc.

\medskip
\noindent\textsc{Proof.} Let $P$ be a primary complex prime not contained among $A$, $B$, $C$, etc., and let $p$, $a$, $b$, $c$, etc. be the norms of $P$, $A$, $B$, $C$, etc. It follows easily that the norm of $M$ will be $a^{\alpha}b^{\beta}c^{\gamma}$, etc.; hence, if $M$ were divisible by $P$, this number would have to be divisible by $p$.

Since the individual norms are either real prime numbers, from the series $2$, $5$, $13$, $17$, etc., or squares of real prime numbers, from the series $9$, $49$, $121$, etc., it is immediately clear that this cannot occur unless $p$ becomes identical with one of the norms $a$, $b$, $c$, etc. We therefore suppose $p=a$. But since $P$ and $A$ are, by hypothesis, non-identical primary complex primes, it is easily seen that these conditions can coexist only if $P$ and $A$ are conjugate imaginary complex numbers, and consequently $p=a$ is an odd real prime, not the square of a prime. We therefore suppose $A=k+li$ and $P=k-li$.

It follows, extending the notion and sign of congruence to complex integers, that $A\equiv 2k \pmod{P}$, whence one easily obtains
\[
M\equiv 2^{\alpha}k^{\alpha}B^{\beta}C^{\gamma}\ldots \pmod{P}.
\]
Thus, while $M$ is assumed to be divisible by $P$, the number $2^{\alpha}k^{\alpha}B^{\beta}C^{\gamma}\ldots$ will also be divisible by $P$, and therefore the norm of this number, namely $2^{2\alpha}k^{2\alpha}b^{\beta}c^{\gamma}\ldots$, will be divisible by $p$. But since neither $2$ nor $k$ is divisible by $p$, it follows that $p$ must be identical with one of the numbers $b$, $c$, etc.; let it be, for example, $p=b$. From this, however, we conclude that either $B=k+li$ or $B=k-li$, that is, either $B=A$ or $B=P$, both contrary to the hypothesis.

From this theorem the other, namely that the resolution into prime factors can be carried out in one way only, is derived very easily, by arguments entirely analogous to those used for real numbers in the \textit{Disquisitiones Arithmeticae}, Article~16. It would therefore be superfluous to dwell on them here.

\art{38}

We now proceed to the congruence of numbers with respect to complex moduli. But at the threshold of this investigation it is useful to indicate how the domain of complex quantities can be made subject to intuition.

Just as every real quantity can be expressed by a portion of a straight line, infinite in both directions, taken from an arbitrary origin and measured according to an arbitrary segment accepted as unit, and therefore can be represented by a point, with points on one side of the origin representing positive quantities and those on the other side negative quantities, so every complex quantity can be represented by some point in an infinite plane, in which a fixed straight line is referred to the real quantities. Specifically, the complex quantity $x+iy$ is represented by the point whose abscissa is $x$ and whose ordinate, taken positively on one side of the line of abscissas and negatively on the other, is $y$.

In this way one may say that any complex quantity measures the difference between the position of the point to which it is referred and the position of the initial point; the positive unit denotes an arbitrary fixed displacement in an arbitrary fixed direction, the negative unit an equal displacement in the opposite direction, and finally the imaginary units equal displacements in the two normal lateral directions. In this way the metaphysics of the quantities that we call imaginary is strikingly illuminated.

If the initial point is denoted by $(0)$, and two complex quantities $m$ and $m'$ are referred to points $M$ and $M'$, whose positions relative to $(0)$ they express, then the difference $m-m'$ will be nothing other than the position of $M$ relative to the point $M'$. Conversely, if the product $mm'$ represents the position of the point $N$ relative to $(0)$, one easily sees that this position is determined by the position of $M$ with respect to $(0)$ in the same way as the position of $M'$ is determined by the position of the point corresponding to the positive unit; thus one may quite aptly say that the positions of the points corresponding to the complex quantities $mm'$, $m$, $m'$, and $1$ form a proportion.

But we reserve a fuller treatment of this matter for another occasion. The difficulties in which the theory of the quantities we call imaginary is thought to be involved have for the most part arisen from ill-suited names, since some have even used the absurd name ``impossible quantities.'' If, taking our start from the concepts offered by varieties of two dimensions, which are seen in their greatest purity in the intuitions of space, we had called positive quantities direct, negative quantities inverse, and imaginary quantities lateral, then in place of tangles there would have been simplicity, and in place of darkness, clarity.

\art{39}

What was set forth in the preceding article refers to continuous complex quantities. In arithmetic, which is concerned only with integers, the diagram of complex numbers will be a system of equidistant points arranged on equidistant straight lines in such a way that the infinite plane is divided into infinitely many equal squares.

All numbers divisible by a given complex number $a+bi=m$ will likewise form infinitely many squares, whose sides are $=\sqrt{aa+bb}$ and whose areas are $=aa+bb$; these latter squares will be inclined to the former ones whenever neither $a$ nor $b$ is zero.

To every number not divisible by the modulus $m$ there will correspond a point either situated within such a square or on a side common to two contiguous squares. The latter case, however, cannot occur unless $a$ and $b$ have a common divisor. Moreover, it is clear that numbers congruent modulo $m$ occupy congruent positions in their squares.

From this it is easy to conclude that, if all numbers lying within a fixed square are collected, together with any that may lie on two of its non-opposite sides, and if to these one appends the number divisible by $m$, then one has a complete system of incongruent residues modulo $m$; that is, every integer must be congruent to one and only one of them. It would not be difficult to show that the number of these residues is equal to the norm of the modulus, namely $aa+bb$. But it seems advisable to prove this very important theorem in another, purely arithmetic way.

\art{40}

\noindent\textsc{Theorem.} With respect to a given complex modulus $m=a+bi$, whose norm is $aa+bb=p$, and for which $a$ and $b$ are relatively prime, every complex integer is congruent to one of the residues in the series $0$, $1$, $2$, $3\ldots p-1$, and to no more than one.

\medskip
\noindent\textsc{Proof.} I. Let $\alpha$ and $\beta$ be integers such that $\alpha a+\beta b=1$. Then
\[
i=\alpha b-\beta a+m(\beta+\alpha i).
\]
Thus for a proposed complex integer $A+Bi$ we have
\[
A+Bi=A+(\alpha b-\beta a)B+m(\beta B+\alpha B i).
\]
Therefore, denoting by $h$ the least nonnegative residue of the number $A+(\alpha b-\beta a)B$ modulo $p$, and putting
\[
A+(\alpha b-\beta a)B=h+kp=h+m(ak-bki),
\]
we obtain
\[
A+Bi=h+m(\beta B+ak+(\alpha B-bk)i),
\]
or $A+Bi\equiv h \pmod{m}$. This proves the first assertion.

II. Whenever two real numbers $h$ and $h'$ are congruent to the same complex number modulo $m$, they are also congruent to one another. Put $h-h'=m(c+di)$; then
\[
(h-h')(a-bi)=p(c+di),
\]
and hence
\[
(h-h')a=pc,\quad (h-h')b=-pd.
\]
Moreover, since $a\alpha+b\beta=1$,
\[
h-h'=p(c\alpha-d\beta),\quad \text{that is,}\quad h\equiv h' \pmod{p}.
\]
Consequently, if $h$ and $h'$ are unequal, they cannot both be contained in the set of numbers $0$, $1$, $2$, $3\ldots p-1$. This proves the second assertion.

\art{41}

\noindent\textsc{Theorem.} With respect to the complex modulus $m=a+bi$, whose norm is $aa+bb=p$, and for which $a$ and $b$ are not relatively prime but have greatest common divisor $\lambda$, taken positive, every complex number is congruent to a residue $x+yi$ such that $x$ is one of the numbers
\[
0,1,2,3\ldots \frac{p}{\lambda}-1,
\]
and $y$ is one of $0$, $1$, $2$, $3\ldots \lambda-1$; moreover, it is congruent to exactly one among all the $p$ residues that have this form.

\medskip
\noindent\textsc{Proof.} I. Choose integers $\alpha,\beta$ such that $\alpha a+\beta b=\lambda$. Then
\[
\lambda i=\alpha b-\beta a+m(\beta+\alpha i).
\]
Now let $A+Bi$ be the proposed complex number, let $y$ be the least nonnegative residue of $B$ modulo $\lambda$, and let $x$ be the least nonnegative residue of
\[
A+(\alpha b-\beta a)\cdot\frac{B-y}{\lambda}
\]
modulo $\frac{p}{\lambda}$; and put
\[
A+(\alpha b-\beta a)\cdot\frac{B-y}{\lambda}=x+\frac{p}{\lambda}\cdot k.
\]
It follows that
\[
\begin{aligned}
A+Bi-(x+yi)
&=\frac{p}{\lambda}\cdot k+(B-y)i-(\alpha b-\beta a)\frac{B-y}{\lambda}\\
&=\frac{p}{\lambda}\cdot k+\frac{B-y}{\lambda}\cdot m(\beta+\alpha i)\\
&=\left(\frac{a}{\lambda}-\frac{b}{\lambda}\,i\right)km+\frac{B-y}{\lambda}(\beta+\alpha i)m,
\end{aligned}
\]
that is, the difference is divisible by $m$, or $A+Bi\equiv x+yi \pmod{m}$. This proves the first assertion.

II. Suppose that two numbers $x+yi$ and $x'+y'i$ are congruent modulo $m$ to the same complex number. Then they are also congruent to one another modulo $m$. A fortiori they are congruent modulo $\lambda$, and therefore $y\equiv y' \pmod{\lambda}$. If each of $y$ and $y'$ is supposed to be contained among the numbers $0$, $1$, $2$, $3\ldots \lambda-1$, one must have $y=y'$.

But then also $x\equiv x' \pmod{m}$; that is, $x-x'$ is divisible by $m$, and consequently the integer $\frac{x-x'}{\lambda}$ is divisible by $\frac{a}{\lambda}+\frac{b}{\lambda}\,i$, or
\[
\frac{x-x'}{\lambda}\equiv 0\pmod{\frac{a}{\lambda}+\frac{b}{\lambda}\,i}.
\]
Since $\frac{a}{\lambda}$ and $\frac{b}{\lambda}$ are relatively prime, it follows by the second part of the theorem of the preceding article that $\frac{x-x'}{\lambda}$ is also divisible by the norm of $\frac{a}{\lambda}+\frac{b}{\lambda}\,i$, that is, by $\frac{p}{\lambda\lambda}$, and hence that $x-x'$ is divisible by $\frac{p}{\lambda}$. Therefore, if each of $x$ and $x'$ is also supposed to be contained in the set of numbers $0$, $1$, $2$, $3\ldots \frac{p}{\lambda}-1$, it is necessary that $x=x'$; that is, the residues $x+yi$ and $x'+y'i$ are identical. This proves the second assertion.

It is, moreover, clear at once that here too one must include the case in which the modulus is a real number, say $b=0$ and hence $\lambda=\pm a$, as well as the case in which the modulus is purely imaginary, say $a=0$ and hence $\lambda=\pm b$. In both cases one has $\frac{p}{\lambda}=\lambda$.

\art{42}
Thus, if all complex integers that are congruent to one another with respect to a given modulus are referred to the same class, and incongruent integers to different classes, there will be altogether $p$ classes exhausting the whole range of integers, where $p$ denotes the norm of the modulus. A collection of just as many numbers, taken one from each class, will give a complete system of incongruent residues, as we assigned in Articles 40 and 41.

In that system the choice of the residues that represent their classes was based on this principle: in each class one adopted a residue $x+yi$ for which $y$ had the least value; and, among all residues having this same least value of $y$, one chose the one for which the value of $x$ was least, negative values being excluded both for $x$ and for $y$.

For other purposes it will be better to use other principles. Especially to be noted is the method in which those residues are adopted which, when divided by the modulus, give the simplest quotients. Plainly, if
\[
\alpha+\beta i,\quad \alpha'+\beta'i,\quad \alpha''+\beta''i,\ldots
\]
are quotients arising from division of congruent numbers by the modulus, then the differences among the quantities $\alpha,\alpha',\alpha'',\ldots$ are integers, and so also are the differences among $\beta,\beta',\beta'',\ldots$. Hence there is always a single residue for which $\alpha$ and $\beta$ lie between $0$ and $1$, the former limit included and the latter excluded. Such a residue we call simply a \emph{minimum residue}.

If preferred, instead of those limits one may adopt $-\frac12$ and $+\frac12$, admitting one and excluding the other. The residue corresponding to this restriction we shall call \emph{absolutely minimum}. The following problems present themselves concerning these minimum residues.

\art{43}
The minimum residue of a given complex number $A+Bi$ with respect to the modulus $a+bi$, whose norm is $p$, is found as follows. If $x+yi$ is the required minimum residue, then $(x+yi)(a-bi)$ will be the minimum residue of the product $(A+Bi)(a-bi)$ with respect to the modulus $(a+bi)(a-bi)$, that is, with respect to the modulus $p$.

Therefore put
\[
aA+bB=Fp+f,\qquad aB-bA=Gp+g,
\]
so that $f$ and $g$ are minimum residues of $aA+bB$ and $aB-bA$ modulo $p$. Then
\[
x+yi=\frac{f+gi}{a-bi},
\]
or
\[
 x=\frac{af-bg}{p}=A-aF+bG,
\qquad
 y=\frac{ag+bf}{p}=B-aG-bF.
\]
Obviously the minimum residues $f,g$ must be taken either between $0$ and $p-1$, or between $-\frac12p$ and $+\frac12p$, according as a simply minimum or an absolutely minimum complex residue is desired.

\art{44}
A complete system of minimum residues for a given modulus can be constructed in several ways. The first method proceeds by first determining the limits within which the real terms must lie, and then, for each value lying within those limits, assigning the limits for the imaginary parts.

The general criterion for a minimum residue $x+yi$ with respect to the modulus $a+bi$ consists in this: both $ax+by=\xi$ and $ay-bx=\eta$ must lie between the limits $0$ and $aa+bb$ when simply minimum residues are in question, or between the limits $-\frac12(aa+bb)$ and $+\frac12(aa+bb)$ when absolutely minimum residues are desired, one of the two endpoints being excluded.

Special rules would require a distinction of the cases introduced by the various possible signs of $a$ and $b$. Since that development presents no difficulty, we shall not dwell on it here; it will be enough to display the nature of the method by a single example.

For the modulus $5+2i$, the simply minimum residues $x+yi$ must be so chosen that both $5x+2y=\xi$ and $5y-2x=\eta$ are equal to one of the numbers $0,1,2,3,\ldots,28$.

The equation $29x=5\xi-2\eta$ shows that positive values of $x$ cannot be greater than $\frac{5\cdot28}{29}$, and that negative values, disregarding the sign, cannot be greater than $\frac{2\cdot28}{29}$. Thus all admissible values of $x$ are $-1,0,1,2,3,4$.

For $x=-1$, $2y$ must be equal to one of the numbers $5,6,7,\ldots,33$, while $5y$ must be equal to one of the numbers $-2,-1,0,1,\ldots,26$; hence the least value of $y$ is $+3$, and the greatest is $+5$.

Treating the remaining values of $x$ in the same way, one obtains the following scheme of all minimum residues:
\[
\begin{array}{r|l}
x&y\\ \hline
-1&3,4,5\\
0&0,1,2,3,4,5\\
+1&1,2,3,4,5,6\\
+2&1,2,3,4,5,6\\
+3&2,3,4,5,6\\
+4&2,3,4
\end{array}
\]

Similarly, for absolutely minimum residues, $\xi$ and $\eta$ must be equal to one of the numbers $-14,-13,-12,\ldots,+14$. Hence $29x$ cannot fall outside the limits $-7\cdot14$ and $+7\cdot14$, and therefore $x$ must be equal to one of the numbers $-3,-2,-1,0,1,2,3$.

For $x=-3$, $2y=\xi-5x=\xi+15$ will be equal to one of the numbers $1,2,3,\ldots,29$, while $5y=\eta+2x=\eta-6$ will be equal to one of the numbers $-20,-19,-18,\ldots,+8$. Hence there results for $y$ the single value $+1$.

Treating the remaining values of $x$ in the same way, we obtain the scheme of all absolutely minimum residues:
\[
\begin{array}{r|l}
x&y\\ \hline
-3&+1\\
-2&-2,-1,0,+1,+2\\
-1&-3,-2,-1,0,+1,+2\\
0&-2,-1,0,+1,+2\\
+1&-2,-1,0,+1,+2,+3\\
+2&-2,-1,0,+1,+2\\
+3&-1
\end{array}
\]

\art{45}
In applying the second method it will be convenient to distinguish two cases.

In the first case, where $a$ and $b$ have no common divisor, put
\[
\alpha a+\beta b=1,
\]
and let $k$ be the least positive residue of $\beta a-\alpha b$ modulo $p$. From the identities
\[
a(\beta a-\alpha b)=\beta p-b(\alpha a+\beta b),\qquad
b(\beta a-\alpha b)=-\alpha p+a(\alpha a+\beta b)
\]
it follows that
\[
ak\equiv -b,\qquad bk\equiv a\pmod p.
\]
Putting, as above,
\[
ax+by=\xi,\qquad ay-bx=\eta,
\]
we therefore have
\[
\eta\equiv k\xi,\qquad \xi\equiv -k\eta\pmod p.
\]

Accordingly, all numbers $\xi+\eta i$ to which simply minimum residues $x+yi$ correspond will be obtained either by taking for $\xi$ successively the values $0,1,2,3,\ldots,p-1$ and for $\eta$ the least positive residues of the products $k\xi$ modulo $p$, or, in the other order, by taking those values for $\eta$ and for $\xi$ the least residues of the products $-k\eta$.

From each $\xi+\eta i$ the corresponding $x+yi$ is then found by the formula
\[
x+yi=\frac{\xi+\eta i}{a-bi}
=\frac{a\xi-b\eta}{p}+\frac{a\eta+b\xi}{p}i.
\]
Moreover, it is clear that when $\xi$ is increased by one, $\eta$ undergoes either an increase by $k$ or a decrease by $p-k$, and therefore $x+yi$ undergoes either the change
\[
\frac{a-kb}{p}+\frac{ak+b}{p}i
\]
or this one:
\[
\frac{a-kb}{p}+b+\left(\frac{ak+b}{p}-a\right)i.
\]
This observation is useful for making the construction easier.

Finally, if absolutely minimum residues $x+yi$ are wanted, the preceding rules need be changed only to this extent: the values successively assigned to $\xi$ should lie between $-\frac12p$ and $+\frac12p$, while $\eta$ should be taken as the absolutely minimum residue of the products $k\xi$.

Here is the display of the minimum residues for the modulus $5+2i$ arranged by this method.

\begin{center}
\small
\textbf{Simply minimum residues.}
\begin{longtable}{L@{\quad}L|L@{\quad}L|L@{\quad}L}
\toprule
\xi+\eta i & x+yi & \xi+\eta i & x+yi & \xi+\eta i & x+yi\\
\midrule
0&0&10+25i&+5i&20+21i&+2+5i\\
1+17i&-1+3i&11+13i&+1+3i&21+9i&+3+3i\\
2+5i&+i&12+i&+2+i&22+26i&+2+6i\\
3+22i&+1+4i&13+18i&+1+4i&23+14i&+3+4i\\
4+10i&+2i&14+6i&+2+2i&24+2i&+4+2i\\
5+27i&-1+5i&15+23i&+1+5i&25+19i&+3+5i\\
6+15i&+3i&16+11i&+2+3i&26+7i&+4+3i\\
7+3i&+1+i&17+28i&+1+6i&27+24i&+3+6i\\
8+20i&+4i&18+16i&+2+4i&28+12i&+4+4i\\
9+8i&+1+2i&19+4i&+3+2i&&\\
\bottomrule
\end{longtable}
\end{center}

\begin{center}
\small
\textbf{Absolutely minimum residues.}
\begin{longtable}{L@{\quad}L|L@{\quad}L|L@{\quad}L}
\toprule
\xi+\eta i & x+yi & \xi+\eta i & x+yi & \xi+\eta i & x+yi\\
\midrule
-14-6i&-2-2i&-4-10i&-2i&+5-2i&+1\\
-13+11i&-3+i&-3+7i&-1+i&+6-14i&+2-2i\\
-12-i&-2-i&-2-5&-i&+7+3i&+1+i\\
-11-13i&-1-3i&-1+12i&-1+2i&+8-9i&+2-i\\
-10+4i&-2&0&0&+9+8i&+1+2i\\
-9-8i&-1-2i&+1-12i&+1-2i&+10-4i&+2\\
-8+9i&-2+i&+2+5i&+i&+11+13i&+1+3i\\
-7-3i&-1-i&+3-7i&+1-i&+12+i&+2+i\\
-6+14i&-2+2i&+4+10i&+2i&+13-11i&+3-i\\
-5+2i&-1&&&+14+6i&+2+2i\\
\bottomrule
\end{longtable}
\end{center}

The second case, where $a,b$ are not relatively prime, is easily reduced to the preceding case. Let $\lambda$ be the greatest common divisor of $a,b$, and let $a=\lambda a'$, $b=\lambda b'$. Let $F$ denote, indefinitely, a minimum residue for the modulus $\lambda$, considered as a complex number; that is, let it denote a number $x+yi$ such that $x,y$ lie either between $0$ and $\lambda$, or between $-\frac12\lambda$ and $+\frac12\lambda$, according as one is dealing with simply or absolutely minimum residues. Let $F'$ likewise denote, indefinitely, a minimum residue for the modulus $a'+b'i$. Then
\[
(a'+b'i)F+F'
\]
will be, indefinitely, a minimum residue for the modulus $a+bi$; and a complete system of such residues will be produced by combining every $F$ with every $F'$.

\art{46}
Two complex numbers are said to be relatively prime if, apart from the units, they admit no other common divisors. Whenever such common divisors are present, those common divisors whose norm is greatest are called greatest common divisors.

If the decomposition of the two proposed numbers into prime factors is at hand, the determination of the greatest common divisor is carried out exactly as for real integers (Disquisitiones Arithmeticae, Art. 18). At the same time it follows that every common divisor of the two given numbers must divide the greatest common divisor found in this way.

Since it is already obvious that the three associates of such a number are also common divisors, exactly four numbers, and no more, will always be called greatest common divisors; their norm will be a multiple of the norm of any other common divisor.

If the decomposition of the two proposed numbers into simple factors is not available, the greatest common divisor is extracted by means of an algorithm similar to the one used for real integers. Let $m,m'$ be the two proposed numbers, and form by repeated division the sequence $m'',m''',\ldots$, so that $m''$ is the absolutely minimum residue of $m$ modulo $m'$, then $m'''$ the absolutely minimum residue of $m'$ modulo $m''$, and so on.

Denoting the norms of the numbers $m,m',m'',m''',\ldots$ by $p,p',p'',p''',\ldots$ respectively, $p''/p'$ is the norm of the quotient $m''/m'$ and therefore, by the definition of the absolutely minimum residue, is certainly not greater than $\frac12$. The same holds for $p'''/p''$, and so forth. Consequently the positive real integers $p',p'',p''',\ldots$ form a continually decreasing sequence, and therefore one must finally reach the term $0$.

Let the preceding final divisor be $m^{(n+1)}$, and put
\[
\begin{aligned}
 m&=km'+m'',\\
 m'&=k'm''+m''',\\
 m''&=k''m'''+m'''',
\end{aligned}
\]
and so on, up to
\[
m^{(n)}=k^{(n)}m^{(n+1)}.
\]
Running through these equations in the reverse order, it is clear that $m^{(n+1)}$ divides all preceding terms $m^{(n)},\ldots,m'',m',m$; running through them in the direct order, it is clear that any common divisor of $m,m'$ also divides all following terms. The first conclusion says that $m^{(n+1)}$ is a common divisor of $m,m'$; the second says that this divisor is greatest.

Finally, whenever the last residue $m^{(n+1)}$ turns out to be equal to one of the four units $1,-1,i,-i$, this is a sign that $m$ and $m'$ are relatively prime.

\art{47}
If the equations of the preceding article, with the last omitted, are combined so that $m'',m''',m'''',\ldots,m^{(n)}$ are eliminated, an equation of the form
\[
m^{(n+1)}=hm+h'm'
\]
will arise, where $h,h'$ are integers. In the notation introduced in Disq. Arith. Art. 27,
\[
\begin{aligned}
h&=\pm[k',k'',k''',\ldots,k^{(n-1)}]
=\pm[k^{(n-1)},k^{(n-2)},\ldots,k'',k'],\\
h'&=\mp[k,k',k'',k''',\ldots,k^{(n-1)}]
=\mp[k^{(n-1)},k^{(n-2)},\ldots,k'',k',k],
\end{aligned}
\]
the upper or lower signs holding according as $n$ is even or odd.

We state this theorem as follows: the greatest common divisor of two complex numbers $m,m'$ can be reduced to the form $hm+h'm'$, where $h,h'$ are integers.

For this holds not only for the greatest common divisor to which the algorithm of the preceding article has led, but also for its three associates. For those, in place of the coefficients $h,h'$ one must take either $hi,h'i$, or $-h,-h'$, or $-hi,-h'i$.

Therefore whenever $m,m'$ are relatively prime, it will be possible to satisfy the equation
\[
1=hm+h'm'.
\]
For example, suppose the proposed numbers are $31+6i=m$ and $11-20i=m'$. We find
\[
\begin{array}{cl@{\qquad}cl}
k&=i,& m''&=+11-5i,\\
k'&=+1-i,& m'''&=+5-4i,\\
k''&=+2,& m''''&=+1+3i,\\
k'''&=-1-2i,& m'''''&=+i,\\
k''''&=+3-i.&&
\end{array}
\]
Hence
\[
[k',k'',k''']=-6-5i,
\qquad
[k,k',k'',k''']=+4-10i,
\]
and therefore
\[
m'''''=i=(6+5i)m+(4-10i)m',
\]
as well as
\[
1=(5-6i)m+(-10-4i)m',
\]
which is confirmed by carrying out the calculation.

\art{48}
All that is required for the theory of first-degree congruences in the arithmetic of complex integers has now been prepared. Since that theory does not essentially differ from the corresponding theory for real integers, and since it has been fully explained in the Disquisitiones Arithmeticae, it will be enough to record the principal points here.

I. The congruence
\[
mt\equiv 1 \pmod {m'}
\]
is equivalent to the indeterminate equation
\[
mt+m'u=1.
\]
If this equation is satisfied by $t=h,u=h'$, the solution of the congruence is expressed generally by
\[
t\equiv h\pmod {m'}.
\]
The condition of solvability is that the modulus $m'$ have no common divisor with the coefficient $m$.

II. The solution of the congruence
\[
ax+b\equiv c\pmod M
\]
in the case where $a,M$ are relatively prime depends on the solution of
\[
at\equiv1\pmod M.
\]
If this is satisfied by $t=h$, the general solution of the proposed congruence is contained in
\[
x\equiv(c-b)h\pmod M.
\]

III. The congruence $ax+b\equiv c\pmod M$ in the case where $a,M$ have the common divisor $\lambda$ is equivalent to
\[
\frac a\lambda x\equiv \frac{c-b}{\lambda}\pmod{M/\lambda}.
\]
Therefore, when $\lambda$ is taken to be the greatest common divisor of $a,M$, the solution of the proposed congruence is reduced to the preceding case; and it is clear that, for solvability, it is necessary and sufficient that $\lambda$ also divide the difference $c-b$.

\art{49}
Up to this point we have touched only elementary matters, which nevertheless could not be omitted because of the connection. In higher investigations the arithmetic of complex integers is like the arithmetic of real integers in this respect: more elegant and simpler theorems appear when only such moduli as are prime numbers are admitted. In truth, the extension of these theorems to composite moduli is usually longer rather than more difficult, more a matter of labor than of art. Therefore, in what follows, prime moduli will chiefly be considered.

\art{50}
Let $X$ denote a function of the indeterminate $x$ of the form
\[
Ax^n+Bx^{n-1}+Cx^{n-2}+\text{etc.}+Mx+N,
\]
where $n$ is a positive real integer, $A,B,C,\ldots$ are real or imaginary integers, and $m$ is a complex integer. We shall here also call a root of the congruence
\[
X\equiv0\pmod m
\]
any integer which, when substituted for $x$, gives to $X$ a value divisible by the modulus $m$. Solutions by roots congruent modulo the modulus will not be regarded as different.

Whenever the modulus is prime, such a congruence of order $n$ can here too admit no more than $n$ different solutions.

Let $\alpha$ be any fixed complex integer. By division by $x-\alpha$, the function $X$ can be reduced indefinitely to the form
\[
X=(x-\alpha)X'+h,
\]
where $h$ is a fixed integer and $X'$ is a function of order $n-1$ with integral coefficients. Whenever $\alpha$ is a root of the congruence $X\equiv0\pmod m$, $h$ will plainly be divisible by $m$, or
\[
X\equiv(x-\alpha)X'\pmod m.
\]

Similarly, if for a fixed integer $\beta$ the function $X'$ is reduced to $(x-\beta)X''+h'$, then $X''$ is of order $n-2$ with integral coefficients. If $\beta$ is also assumed to be a root of the congruence $X\equiv0$, it must satisfy $(\beta-\alpha)X'\equiv0$, and also $X'\equiv0$, provided the roots $\alpha,\beta$ are incongruent. Therefore $h'$ too must be divisible by $m$, or
\[
X\equiv(x-\alpha)(x-\beta)X''\pmod m.
\]

In the same way, with a third root $\gamma$ incongruent to the preceding ones, we obtain indefinitely
\[
X\equiv(x-\alpha)(x-\beta)(x-\gamma)X''',
\]
where $X'''$ has order $n-3$ and integral coefficients. Continuing in the same manner, and noting that the coefficient of the highest term in each function is $A$ - which may be assumed not divisible by $m$, since otherwise the congruence $X\equiv0$ would essentially belong to a lower order - we get, whenever there are $n$ incongruent roots $\alpha,\beta,\gamma,\ldots,\nu$,
\[
X\equiv A(x-\alpha)(x-\beta)(x-\gamma)\cdots(x-\nu)\pmod m.
\]
Thus substitution of a new value incongruent to all of $\alpha,\beta,\gamma,\ldots,\nu$ would certainly give $X$ a value not divisible by $m$. The theorem follows at once. This proof agrees essentially with that given in Disq. Arith. Art. 43, and each of its steps is valid for complex integers just as for real ones.

\art{51}
What was set forth in the third section of the Disquisitiones Arithmeticae concerning power residues remains valid, with only slight changes, for the most part also in the arithmetic of complex integers. Indeed, the demonstrations of the theorems could usually be retained. So that nothing be lacking, however, we shall give the principal theorems supported by concise proofs, always understanding that the modulus is a prime number.

\emph{Theorem.} Let $k$ be an integer not divisible by the modulus $m$, whose norm is $p$. Then
\[
k^{p-1}\equiv1\pmod m.
\]
\emph{Proof.} Let $a,b,c,\ldots$ constitute a complete system of incongruent residues modulo $m$, with the residue divisible by $m$ omitted. Thus the number of these residues, which we denote collectively by $C$, is $p-1$. Let $C'$ be the collection of products $ka,kb,kc,\ldots$.

None of these products is divisible by $m$ by hypothesis. Therefore each will have a congruent residue in $C$, say
\[
ak\equiv a',\quad bk\equiv b',\quad ck\equiv c',\ldots\pmod m,
\]
where $a',b',c',\ldots$ are themselves found in $C$. Denote this collection by $C''$.

Let $P,P',P''$ be the products of the numbers in the collections $C,C',C''$ respectively, that is,
\[
P=abc\cdots,
\qquad
P'=k^{p-1}abc\cdots=k^{p-1}P,
\qquad
P''=a'b'c'\cdots.
\]
Since the numbers of $C''$ are successively congruent to the numbers of $C'$, we have $P''\equiv P'$, or $P''\equiv k^{p-1}P$. But it is easy to see that any two numbers in $C''$ are incongruent, and hence all distinct; necessarily the collection $C''$ agrees exactly with $C$, only in changed order. Thus $P''=P$.

Consequently $(k^{p-1}-1)P$ is divisible by $m$. Since $m$ is prime and does not divide any single factor of $P$, it necessarily follows that $k^{p-1}-1$ must be divisible by $m$. Q.E.D.

\art{52}
\emph{Theorem.} Let $k$, as in the preceding article, be an integer not divisible by the modulus $m$, and let $t$ be the least positive exponent for which
\[
k^t\equiv1\pmod m.
\]
Then $t$ divides any other exponent $u$ for which $k^u\equiv1\pmod m$.

\emph{Proof.} If $t$ did not divide $u$, let $gt$ be the multiple of $t$ immediately greater than $u$; then $gt-u$ is a positive integer smaller than $t$. From $k^t\equiv1$ and $k^u\equiv1$ it follows that
\[
0\equiv k^{gt}-k^u\equiv k^u(k^{gt-u}-1),
\]
so that $k^{gt-u}\equiv1$. This gives a power of $k$ congruent to unity with exponent smaller than $t$, contrary to the hypothesis.

As a corollary, $t$ certainly divides $p-1$. Those numbers $k$ for which $t=p-1$ we shall here also call primitive roots modulo $m$; that such numbers actually exist we now show.

\art{53}
Resolve the number $p-1$ into its prime factors, so that
\[
p-1=a^\alpha b^\beta c^\gamma\cdots,
\]
where $a,b,c,\ldots$ denote distinct positive real primes. Let $A,B,C,\ldots$ be complex integers not divisible by $m$, and respectively not satisfying the congruences
\[
x^{(p-1)/a}\equiv1,
\quad
x^{(p-1)/b}\equiv1,
\quad
x^{(p-1)/c}\equiv1,
\quad\text{etc.}
\]
modulo $m$; that such choices exist is clear from the theorem of Article 50. Finally let $h$ be congruent modulo $m$ to the product
\[
A^{(p-1)/a^\alpha}B^{(p-1)/b^\beta}C^{(p-1)/c^\gamma}\cdots.
\]
I say that $h$ will be a primitive root.

\emph{Proof.} Denote by $t$ the exponent of the least power $h^t$ congruent to unity. If $h$ were not a primitive root, $t$ would be a proper divisor of $p-1$, so $(p-1)/t$ would be an integer greater than one. Plainly this integer has among its prime factors one of $a,b,c,\ldots$. We may therefore suppose that $(p-1)/t$ is divisible by $a$, and put $p-1=atu$.

Since $h^t\equiv1$, also $h^{tu}\equiv1$, that is,
\[
A^{\frac{p-1}{a^\alpha}\cdot\frac{p-1}{a}}
B^{\frac{p-1}{b^\beta}\cdot\frac{p-1}{a}}
C^{\frac{p-1}{c^\gamma}\cdot\frac{p-1}{a}}
\cdots\equiv1.
\]
But $(p-1)/(ab^\beta)$ is plainly an integer, and hence
\[
B^{\frac{p-1}{b^\beta}\cdot\frac{p-1}{a}}
=(B^{p-1})^{(p-1)/(ab^\beta)}\equiv1.
\]
The same is true of the remaining factors except the power of $A$. Therefore we must have
\[
A^{\frac{p-1}{a^\alpha}\cdot\frac{p-1}{a}}\equiv1.
\]
Now choose a positive integer $\lambda$ such that
\[
\lambda b^\beta c^\gamma\cdots\equiv1\pmod a,
\]
which can be done because the prime $a$ does not divide $b^\beta c^\gamma\cdots$; write
\[
\lambda b^\beta c^\gamma\cdots=1+a\mu.
\]
Then
\[
A^{\lambda\cdot\frac{p-1}{a^\alpha}\cdot\frac{p-1}{a}}\equiv1.
\]
But
\[
\lambda\cdot\frac{p-1}{a^\alpha}\cdot\frac{p-1}{a}
=(1+a\mu)\frac{p-1}{a}=(p-1)\mu+\frac{p-1}{a}.
\]
Thus
\[
A^{(p-1)\mu}A^{(p-1)/a}\equiv1.
\]
Since $A^{(p-1)\mu}\equiv1$, it follows that
\[
A^{(p-1)/a}\equiv1,
\]
contrary to the hypothesis. Therefore the supposition that $t$ is a proper divisor of $p-1$ cannot stand; necessarily $h$ is a primitive root.


\bigskip
\begin{center}{\large Continuation: Articles 54--65}\end{center}

\art{54}
Let $h$ denote a primitive root modulo $m$, whose norm is $p$. The terms of the progression
\[
1,\ h,\ hh,\ h^3,\ldots,h^{p-2}
\]
will be mutually incongruent. Hence it is easily inferred that every integer not divisible by the modulus must be congruent to one of these; in other words, this series gives a complete system of incongruent residues, zero excluded. The exponent of the power to which a given number is congruent may be called its \textit{index}, while $h$ is regarded as the \textit{base}. Here are some examples, in which the absolutely least residue has been put next to each index.

\begin{center}\textit{First example.}\par
$m=5+4i,\quad p=41,\quad h=1+2i$\end{center}
\begin{center}\small
\begin{longtable}{L@{\quad}L|L@{\quad}L|L@{\quad}L|L@{\quad}L|L@{\quad}L}
\toprule
\text{Ind.} & \text{Residuum} & \text{Ind.} & \text{Residuum} & \text{Ind.} & \text{Residuum} & \text{Ind.} & \text{Residuum} & \text{Ind.} & \text{Residuum}\\
\midrule
0 & +1 & 8 & -4 & 16 & -2+2i & 24 & +2i & 32 & +1+i\\
1 & +1+2i & 9 & -3+i & 17 & -1+2i & 25 & -3i & 33 & +1+3i\\
2 & +1-i & 10 & -i & 18 & +4i & 26 & +2+2i & 34 & +2\\
3 & +3+i & 11 & +2-i & 19 & +1+3i & 27 & +2+i & 35 & -3\\
4 & -2i & 12 & -1-i & 20 & -1 & 28 & +4 & 36 & +2-2i\\
5 & +3i & 13 & +1-3i & 21 & -1-2i & 29 & +3-i & 37 & +1-2i\\
6 & -2-2i & 14 & -2 & 22 & -1+i & 30 & +i & 38 & -4i\\
7 & -2-i & 15 & +3 & 23 & -3-i & 31 & -2+i & 39 & -1-3i\\
\bottomrule
\end{longtable}
\end{center}

\begin{center}\textit{Second example.}\par
$m=7,\quad p=49,\quad h=1+2i$\end{center}
\begin{center}\small
\begin{longtable}{L@{\quad}L|L@{\quad}L|L@{\quad}L|L@{\quad}L|L@{\quad}L}
\toprule
\text{Ind.} & \text{Residuum} & \text{Ind.} & \text{Residuum} & \text{Ind.} & \text{Residuum} & \text{Ind.} & \text{Residuum} & \text{Ind.} & \text{Residuum}\\
\midrule
0 & +1 & 10 & -1-i & 20 & +2i & 30 & +2-2i & 40 & +3\\
1 & +1+2i & 11 & +1-3i & 21 & +3+2i & 31 & -1+2i & 41 & +3-i\\
2 & -3-3i & 12 & -i & 22 & -1+i & 32 & +2 & 42 & -2-2i\\
3 & +3-2i & 13 & +2-i & 23 & -3-i & 33 & +2-3i & 43 & +2+i\\
4 & -3i & 14 & -3+3i & 24 & -1 & 34 & +1+i & 44 & -2i\\
5 & -1-3i & 15 & -2-3i & 25 & -1-2i & 35 & -1+3i & 45 & -3-2i\\
6 & -2+2i & 16 & -3 & 26 & +3+3i & 36 & +i & 46 & +1-i\\
7 & +1-2i & 17 & -3+i & 27 & -3+2i & 37 & -2+i & 47 & +3+i\\
8 & -2 & 18 & +2+2i & 28 & +3i & 38 & +3-3i &  & \\
9 & -2+3i & 19 & -2-i & 29 & +1+3i & 39 & +2+3i &  & \\
\bottomrule
\end{longtable}
\end{center}

\art{55}
We add a few observations about primitive roots and the algorithm of indices, omitting the proofs because of their ease.

\begin{enumerate}[label=\Roman*.]
\item Indices congruent modulo $p-1$ correspond, in the given system, to residues congruent modulo $m$, and conversely.
\item The residues corresponding to indices prime to $p-1$ are themselves primitive roots, and conversely.
\item If, with the primitive root $h$ taken as base, the index of another primitive root $h'$ is $t$, and conversely $t'$ is the index of $h$ when $h'$ is taken as base, then $tt'\equiv1\pmod{p-1}$; and, under the same assumptions, if the indices of some other number in these two systems are respectively $u$ and $u'$, then $tu'\equiv u$ and $t'u\equiv u'\pmod{p-1}$.
\item When the numbers $1$, $1+i$, and their three associates are excluded from the moduli under consideration, as being too meagre, there remain those primes that we placed third and fourth in Article~34. The norms of the latter will be real primes of the form $4n+1$; the norms of the former are squares of odd real primes. Thus in both cases $p-1$ is divisible by $4$.
\item If the index of the number $-1$ is denoted by $u$, then $2u\equiv0\pmod{p-1}$, and hence either $u\equiv0$ or $u\equiv\frac12(p-1)$. But since index $0$ corresponds to the residue $+1$, the index of $-1$ must necessarily be $\frac12(p-1)$.
\item Similarly, if $u$ denotes the index of $i$, then $2u\equiv\frac12(p-1)\pmod{p-1}$, and hence either $u\equiv\frac14(p-1)$ or $u\equiv\frac34(p-1)$. This ambiguity depends on the choice of primitive root. Namely, if, with primitive root $h$ taken as base, the index of $i$ is $\frac14(p-1)$, the index becomes $\frac34(p-1)$ when $h^\mu$ is taken as base, where $\mu$ is a positive integer of the form $4n+3$ and prime to $p-1$, for instance the number $p-2$ itself, and conversely. Thus one half of the primitive roots gives $i$ the index $\frac14(p-1)$, the other half the index $\frac34(p-1)$; plainly, for the former bases $-i$ will have index $\frac34(p-1)$, and for the latter it will have index $\frac14(p-1)$.
\item Whenever the modulus is a positive real prime of the form $4n+3$, say $q$, so that $p=qq$, the indices of all real numbers are divisible by $q+1$. For if $t$ denotes the index of the real number $k$, then from $k^{q-1}\equiv1\pmod q$ we get $(q-1)t\equiv0\pmod{qq-1}$, and hence $t/(q+1)$ is an integer. Similarly the indices of purely imaginary numbers such as $ki$ will be divisible by $\frac12(q+1)$. Thus primitive roots for such moduli must be sought only among mixed numbers.
\item By contrast, for a modulus $m$ that is a mixed complex prime, whose norm $p$ is therefore a real prime of the form $4n+1$, primitive roots can even be chosen among the real numbers, among which a complete system of incongruent residues may be displayed (Article~40). Plainly every real number that is a primitive root modulo the complex modulus $m$ is at the same time, in the arithmetic of real numbers, a primitive root modulo $p$, and conversely.
\end{enumerate}

\art{56}
Although the theory of quadratic residues and non-residues in the arithmetic of complex numbers is contained in the very theory of biquadratic residues, nevertheless, before passing to the latter, we shall present separately here the principal theorems of the former. For brevity we shall speak here only of the principal case, in which the modulus is an odd complex prime.

Let $m$ be such a modulus, and let $p$ be its norm. Plainly every integer not divisible by $m$ - which will always be understood here - either can or cannot be made congruent to a square modulo $m$, according as its index, with some primitive root taken as base, is even or odd. In the former case that integer will be called a quadratic residue of $m$; in the latter, a non-residue. It follows that among the $p-1$ numbers forming a complete system of incongruent residues not divisible by $m$, one half are quadratic residues and the other half quadratic non-residues. Every number outside that system is to be given, in this respect, the same character as the number of the system to which it is congruent.

It also follows from the same source that the product of two quadratic residues, as well as the product of two non-residues, is a quadratic residue; whereas the product of a quadratic residue and a non-residue is a non-residue. More generally, a product of any number of factors is a quadratic residue or a non-residue according as the number of non-residue factors is even or odd.

For distinguishing quadratic residues from non-residues, the following general criterion at once presents itself:
\[
 k\ \text{is a quadratic residue or non-residue}\quad
 \Longleftrightarrow\quad
 k^{\frac12(p-1)}\equiv 1\ \text{or}\ -1\pmod m.
\]
The truth of this theorem follows immediately from the fact that, after any primitive root has been accepted as base, the index of the power $k^{\frac12(p-1)}$ becomes either congruent to $0$ or to $\frac12(p-1)$, according as the index of $k$ is even or odd.

\art{57}
For a given modulus it is easy enough to divide a complete system of incongruent residues into two classes, namely quadratic residues and non-residues; in this way all remaining numbers are assigned to their classes automatically. But the question of criteria distinguishing the moduli for which a given number is a quadratic residue from those for which it is a non-residue is much deeper.

As for the real units $+1$ and $-1$, these are in fact squares in the arithmetic of complex numbers, and hence are quadratic residues for every modulus. Just as easily, from the criterion of the preceding article, it follows that the number $i$, and likewise $-i$, is a quadratic residue of every modulus whose norm $p$ is of the form $8n+1$, but a non-residue of every modulus whose norm is of the form $8n+5$.

Since it plainly makes no difference whether the number $m$ or one of the numbers associated with it, $im$, $-m$, or $-im$, is adopted as modulus, we may suppose that the modulus is the primary one among its associates (Article~36, II). Thus, putting the modulus equal to $a+bi$, $a$ is odd and $b$ is even. In this arrangement, since always $aa\equiv1\pmod 8$, while $bb$ is either $\equiv0$ or $\equiv4\pmod 8$ according as $b$ is divisible by $4$ or is even but not divisible by $4$, it is clear that $+i$ and $-i$ are quadratic residues of the modulus in the former case and non-residues in the latter.

\art{58}
Since judging the character of a composite number, whether it is a quadratic residue or a non-residue, depends on the characters of its factors, it is plainly enough to limit the development of criteria distinguishing the moduli for which a given number $k$ is a quadratic residue from those for which it is a non-residue to values of $k$ that are primes and, moreover, primary among their associates.

In this investigation induction immediately supplies very elegant theorems. We begin with the number $1+i$, which is found to be a quadratic residue of the moduli
\[
-1+2i,
+3-2i,
-5-2i,
-1-6i,
+5+4i,
+5-4i,
-7,
+7+2i,
-5+6i,\ldots
\]
and a quadratic non-residue of the following:
\[
-1-2i,
-3,
+3+2i,
+1+4i,
+1-4i,
-5+2i,
-1+6i,
+7-2i,
-5-6i,
-3+8i,
-3-8i,
+5+8i,
+5-8i,
+9+4i,
+9-4i,\ldots
\]
If we examine this list carefully, always taking from each set of four associated moduli the primary one, we readily observe that the moduli $a+bi$ in the first class are all those for which $a+b\equiv+1\pmod 8$, and those in the second are those for which $a+b\equiv-3\pmod 8$. Plainly, if instead of the primary modulus $m$ we adopt its associate $-m$, this criterion must be changed so that for the moduli of the first class $a+b\equiv-1$, and for those of the second $a+b\equiv+3\pmod 8$.

Therefore, if induction has not deceived us, in general, for a prime number $a+bi$ in which $a$ is odd and $b$ is even, $1+i$ becomes its quadratic residue or non-residue according as $a+b\equiv\pm1$ or $\equiv\pm3\pmod 8$.

For the number $-1-i$ the same rule holds as for $1+i$. On the other hand, regarding $1-i$ as the product of $-i$ and $1+i$, it is clear that $1-i$ has the same character as $1+i$ whenever $b$ is divisible by $4$, and the opposite character whenever $b$ is even but not divisible by $4$. Hence one easily obtains that $1-i$ is a quadratic residue of the prime number $a+bi$ whenever $a-b\equiv\pm1$, and a non-residue whenever $a-b\equiv\pm3\pmod 8$, always supposing that $a$ is odd and $b$ even.

This second proposition may also be deduced from the first by means of the more general theorem stated as follows:

\begin{quote}
In the theory of quadratic residues, the character of the number $\alpha+\beta i$ with respect to the modulus $a+bi$ is the same as that of $\alpha-\beta i$ with respect to the modulus $a-bi$.
\end{quote}
The proof of this theorem is drawn from the fact that the two moduli have the same norm $p$, and that whenever $(\alpha+\beta i)^{\frac12(p-1)}-1$ is divisible by $a+bi$, then $(\alpha-\beta i)^{\frac12(p-1)}-1$ also becomes divisible by $a-bi$; and whenever $(\alpha+\beta i)^{\frac12(p-1)}+1$ is divisible by $a+bi$, then $(\alpha-\beta i)^{\frac12(p-1)}+1$ must also be divisible by $a-bi$.

\art{59}
Let us proceed to odd prime numbers. We find that the number $-1+2i$ is a quadratic residue of the moduli
\[
+3+2i,
+1-4i,
-5+2i,
-5-2i,
-1-6i,
+7-2i,
-3+8i,
+5+8i,
+5-8i,
+9+4i,\ldots
\]
and a non-residue of the moduli
\[
-1-2i,
-3,
+3-2i,
+1+4i,
-1+6i,
+5+4i,
+5-4i,
-7,
+7+2i,
-5+6i,
-5-6i,
-3-8i,
+9-4i,\ldots
\]
Reducing the moduli of the first class to their absolutely least residues modulo $-1+2i$, we find only $+1$ and $-1$, for instance $+3+2i\equiv-1$, $+1-4i\equiv-1$, $-5+2i\equiv+1$, $-5-2i\equiv-1$, and so on. On the other hand, all moduli of the second class are found, modulo $-1+2i$, to be congruent either to $+i$ or to $-i$.

But the numbers $+1$ and $-1$ themselves are quadratic residues of the modulus $-1+2i$, while $+i$ and $-i$ are its non-residues. Therefore, so far as one may trust induction, the following theorem emerges: the number $-1+2i$ is a quadratic residue or non-residue of the prime number $a+bi$ according as that number is a quadratic residue or non-residue of $-1+2i$, provided that $a+bi$ is the primary one among its four associates, or rather, provided that $a$ is odd and $b$ even.

From this theorem there follow at once the analogous theorems concerning the numbers $+1-2i$, $-1-2i$, and $+1+2i$.

\art{60}
Carrying out a similar induction for the number $-3$ or $+3$, we find that both are quadratic residues of the moduli
\[
+3+2i,
+3-2i,
-1+6i,
-1-6i,
-7,
-5+6i,
-5-6i,
-3+8i,
-3-8i,
+9+4i,
+9-4i,\ldots
\]
and non-residues of the following:
\[
-1+2i,
-1-2i,
+1+4i,
+1-4i,
-5+2i,
-5-2i,
+5+4i,
+5-4i,
+7+2i,
+7-2i,
+5+8i,
+5-8i,\ldots
\]
The former are, modulo $3$, congruent to one of the four numbers $+1$, $-1$, $+i$, $-i$; the latter to one of $+1+i$, $+1-i$, $-1+i$, $-1-i$. The former are precisely the quadratic residues of the number $3$, the latter the non-residues.

This induction therefore teaches that a prime number $a+bi$, with $a$ odd and $b$ even, has with respect to $-3$ and also to $+3$ the same relation that $-3$ has to it, namely as regards whether one is the quadratic residue or non-residue of the other.

Extending a similar induction to other prime numbers, we find everywhere this very elegant law of reciprocity confirmed, and we are led to the following fundamental theorem concerning quadratic residues in the arithmetic of complex numbers:

\begin{quote}
If $a+bi$ and $A+Bi$ denote prime numbers such that $a$ and $A$ are odd, while $b$ and $B$ are even, then either each is a quadratic residue of the other, or each is a non-residue of the other.
\end{quote}
Yet, despite the extreme simplicity of this theorem, its proof is beset with great difficulties. We shall not linger over them here, since the theorem itself is only a special case of a more general theorem that almost exhausts the whole theory of biquadratic residues. We now pass to that theory.

\art{61}
What was set forth in Article~2 of the preceding memoir concerning the notion of a biquadratic residue and non-residue we extend also to the arithmetic of complex numbers; and, as there, here too we restrict the inquiry to moduli that are prime numbers. At the same time it will usually be tacitly understood that the modulus is taken to be primary among its associates, for example congruent to $1$ modulo $2+2i$, and also that the numbers whose character is under consideration, as biquadratic residues or non-residues, are not divisible by the modulus.

For a given modulus, therefore, the numbers not divisible by it could be divided into three classes: the first containing the biquadratic residues, the second the biquadratic non-residues that are nevertheless quadratic residues, and the third the quadratic non-residues.

But here too it is preferable, in place of the third class, to establish two classes, so that there are four in all. With any primitive root taken as base, the biquadratic residues will have indices divisible by $4$, that is, of the form $4n$; the non-residues that are quadratic residues will have indices of the form $4n+2$; finally, the indices of the quadratic non-residues will be partly of the form $4n+1$ and partly of the form $4n+3$. In this way four classes do arise, but the distinction between the last two would not be absolute; it would depend on the choice of the primitive root taken as base. For it is easy to see that one half of the primitive roots assign to a given quadratic non-residue an index of the form $4n+1$, and the other half an index of the form $4n+3$.

To remove this ambiguity, we shall always suppose such a primitive root to be adopted that the index $\frac14(p-1)$ belongs to the number $+i$ (cf. Article~55, VI). This produces a classification that can more neatly be stated independently of primitive roots as follows. The first class contains those numbers $k$ for which $k^{\frac14(p-1)}\equiv1$; these are the biquadratic residues of the modulus. The second class contains those for which $k^{\frac14(p-1)}\equiv i$. The third contains those for which $k^{\frac14(p-1)}\equiv-1$. Finally, the fourth contains those for which $k^{\frac14(p-1)}\equiv-i$.

The third class comprises the biquadratic non-residues that are quadratic residues; the quadratic non-residues are distributed between the second and fourth. To the numbers of these classes we shall assign respectively the biquadratic characters $0$, $1$, $2$, and $3$. If we define the character $\lambda$ of the number $k$ modulo $m$ to be the exponent of that power of $i$ to which $k^{\frac14(p-1)}$ is congruent, then plainly characters congruent modulo $4$ are to be regarded as equivalent. For the present this notion is limited to prime moduli; in the continuation of these investigations we shall show how it can also be adapted to composite moduli.

\art{62}
In order that abundant induction concerning the characters of numbers may be built up more easily, we add here a compact table. With its aid the character of any proposed number with respect to a modulus whose norm does not exceed $157$ can be obtained with little work, provided the following observations are kept in mind.

Since the character of a composite number is equal, or congruent modulo $4$, to the sum of the characters of its individual factors, it suffices to be able to assign, for a given modulus, the characters of prime numbers. Moreover, since the characters of the units $-1$, $i$, and $-i$ are plainly congruent to the numbers $\frac12(p-1)$, $\frac14(p-1)$, and $\frac34(p-1)$ modulo $4$, it will also suffice to have displayed the characters of the numbers primary among their associates. Finally, since moduli congruent modulo $m$ have the same character, it is enough to put into the table the characters of numbers contained in the system of absolutely least residues.

Furthermore, by reasoning similar to that of Article~58 one proves that, if for the modulus $a+bi$ the character of $A+Bi$ is $\lambda$, while for the modulus $a-bi$ the character of $A-Bi$ is $\lambda'$, then always $\lambda\equiv-\lambda'\pmod4$, or $\lambda+\lambda'$ is divisible by $4$. Therefore it is enough to include in the table moduli for which $b$ is either $0$ or positive.

For example, if the character of $11-6i$ with respect to the modulus $-5-6i$ is sought, we substitute for these numbers $11+6i$ and $-5+6i$. We then determine (Article~43) the absolutely least residue of $11+6i$ modulo $-5+6i$, which is $-1-4i=-1\times(1+4i)$. Since for the modulus $-5+6i$ the character of $-1$ is $30$, while the character of $1+4i$ is, from the table, $2$, the character of $11+6i$ modulo $-5+6i$ is $32$, that is, $0$; and hence, by the last observation, this is also the character of $11-6i$ modulo $-5-6i$.

Similarly, if the character of $-5+6i$ with respect to the modulus $11+6i$ is sought, its absolutely least residue $1-5i$ is resolved into the factors $-i$, $1+i$, and $3-2i$, to which correspond the characters $117$, $0$, and $1$. Hence the sought character is $118$, that is, $2$. The same character is also to be assigned to $-5-6i$ with respect to the modulus $11-6i$.

\begin{center}\small

\begin{longtable}{L C L}
\toprule
\text{Modulus} & \text{Character} & \text{Numeri}\\
\midrule
-3 & 3 & 1+i\\
+3+2i & 3 & 1+i\\
+1+4i & 1 & -1+2i\\
 & 3 & 1+i\\
-5+2i & 0 & -1-2i\\
 & 1 & 1+i\\
 & 2 & -1+2i\\
-1+6i & 0 & -3\\
 & 1 & 1+i, -1+2i\\
 & 2 & -1-2i\\
+5+4i & 0 & 1+i\\
 & 1 & -3\\
 & 3 & -1+2i, -1-2i\\
-7 & 0 & -3\\
 & 1 & -1+2i, 3-2i\\
 & 2 & 1+i\\
 & 3 & -1-2i, 3+2i\\
+7+2i & 0 & 1+i, 3+2i, 3-2i, 1-4i\\
 & 1 & -3\\
 & 2 & -1-2i, 1+4i\\
 & 3 & -1+2i\\
-5+6i & 0 & 1+i, -3, 3+2i, 3-2i\\
 & 1 & 1-4i\\
 & 2 & 1+4i\\
 & 3 & -1+2i, -1-2i\\
-3+8i & 0 & -1+2i, 3-2i, 1-4i\\
 & 1 & 1+i, 3+2i\\
 & 2 & -3\\
 & 3 & -1-2i, 1+4i, -5+2i\\
+5+8i & 0 & -1-2i\\
 & 1 & -5-2i, -1+6i\\
 & 2 & -1+2i, 3-2i\\
 & 3 & 1+i, -3, 3+2i, 1+4i, 1-4i\\
+9+4i & 0 & -1+2i, 3+2i\\
 & 1 & 1+i, -1-2i, 3-2i\\
 & 2 & -3, 1+4i\\
 & 3 & 1-4i, -5+2i\\
-1+10i & 0 & 1+i, -1+2i, -1-2i, 3+2i\\
 & 1 & -3\\
 & 2 & 3-2i, -5+2i, 5-4i\\
 & 3 & 1+4i, 1-4i\\
+3+10i & 1 & 1+i, -1-2i, 1-4i\\
 & 2 & -3, 3+2i, 1+4i, -5-2i\\
 & 3 & -1+2i, 3-2i\\
-7+8i & 0 & 1+i, -7\\
 & 1 & 3+2i, 3-2i, 1-4i, -5-2i\\
 & 2 & -1-2i, 1+4i, -5+2i, -1-6i\\
 & 3 & -1+2i, -3, -1+6i\\
-11 & 0 & -3\\
 & 1 & 1+i, 3-2i, 1+4i, -5+2i, 5+4i\\
 & 2 & -1+2i, -1-2i\\
 & 3 & 3+2i, 1-4i, -5-2i, 5-4i\\
-11+4i & 0 & 1+i, -1+2i, 3+2i, 5+4i\\
 & 1 & -1-2i, -1+6i\\
 & 2 & -5+2i\\
 & 3 & -3, 3-2i, 1+4i, 1-4i, -5-2i\\
+7+10i & 0 & 1+4i, 1-4i, -1+6i, -1-6i\\
 & 1 & -1+2i, 3+2i, -5+2i\\
 & 2 & 1+i, 3-2i\\
 & 3 & -1-2i, -3, -5-2i\\
+11+6i & 0 & 1+i, -1+2i, -3, 1+4i, 1-4i, -7\\
 & 1 & -1-2i, 3+2i, 3-2i\\
 & 2 & -5-2i, -1+6i, 5-4i\\
 & 3 & -5+2i, 5+4i, 7-2i\\
\bottomrule
\end{longtable}
\end{center}

\art{63}
We shall now take pains to detect by induction the common criteria for the moduli for which a given prime number has the same character. We always suppose the moduli to be primary among their associates, that is, numbers $a+bi$ for which either $a\equiv1$, $b\equiv0$, or $a\equiv3$, $b\equiv2\pmod4$.

With respect to the number $1+i$, from which we begin, the law of induction is more easily grasped if we separate the moduli of the first kind, for which $a\equiv1$, $b\equiv0$, from those of the second kind, for which $a\equiv3$, $b\equiv2$. With the aid of the table in the preceding article we find the following correspondence.

\begin{center}\small
\begin{tabular}{C L}
\toprule
\text{character} & \text{moduli of the first kind}\\
\midrule
0 & 5+4i, -7+8i, -7-8i, -11+4i\\
1 & 1-4i, -3+8i, -3-8i, 9+4i, -11\\
2 & 5-4i, -7, -11-4i\\
3 & -3, 1+4i, 5+8i, 5-8i, 9-4i\\
\bottomrule
\end{tabular}
\end{center}

If these seventeen examples are considered attentively, in all of them we find the character
\[
\equiv \frac14(a-b-1) \pmod 4.
\]
Similarly there corresponds

\begin{center}\small
\begin{tabular}{C L}
\toprule
\text{character} & \text{moduli of the second kind}\\
\midrule
0 & 3-2i, -1-6i, 7+2i, -5+6i, -1+10i, 11+6i\\
1 & -5+2i, -1+6i, 7-2i, -1-10i, 3+10i\\
2 & -1+2i, -5-2i, 3-10i, 7+10i\\
3 & -1-2i, 3+2i, -5-6i, 7-10i, 11-6i\\
\bottomrule
\end{tabular}
\end{center}

In all these twenty examples, with slight attention, the character is found to be
\[
\equiv \frac14(a-b-5) \pmod 4.
\]
These two rules can easily be contracted into one, valid for both kinds of moduli, if we note that $\frac14bb$ is congruent to $0$ for moduli of the first kind and to $1$ for those of the second kind, modulo $4$. Thus the character of $1+i$ with respect to any prime modulus primary among its associates is
\[
\equiv \frac14(a-b-1-bb) \pmod 4.
\]
It is useful to note in passing that, since $(b+1)^2$ is always of the form $8n+1$, or $\frac14(2b+bb)$ is even, this character is always even or odd according as $\frac14(a+b-1)$ is even or odd; this agrees with the rule for the quadratic character given in Article~58.

Since $\frac14(a-b-1)$ and $\frac14(a-b+3)$ are integers one of which is even and the other odd, their product will be even; that is,
\[
\frac18(a-b-1)(a-b+3)\equiv0\pmod4.
\]
Therefore, in place of the expression just given for the biquadratic character, the following may also be adopted:
\[
\frac14(a-b-1-bb)-\frac18(a-b-1)(a-b+3)=\frac18(-aa+2ab-3bb+1).
\]
This form recommends itself also for the reason that it is not restricted to primary moduli, but only supposes $a$ to be odd and $b$ even. For under this supposition either $a+bi$ or $-a-bi$ will be primary among its associates, and the value of this formula is the same for both moduli.

\art{64}
Starting from the last rule found in the preceding article, we find that
\[
\begin{array}{c|c}
\text{number} & \text{character }\equiv\\
\hline
-1+i & \frac18(aa+2ab-bb-1)\\
-1-i & \frac18(-aa+2ab+bb+1)\\
+1-i & \frac18(aa+2ab+3bb-1)
\end{array}
\]
This follows immediately from the fact that the character of $i$ is $\frac14(aa+bb-1)$, while the character of $-1$ is $\frac12(aa+bb-1)\equiv\frac12bb$, since $aa-1$ is always of the form $8n$.

Plainly these four rules, although up to this point borrowed from induction, are so linked with one another that as soon as the proof of one has been completed, the other three will have been proved at the same time. It hardly needs to be said that in these rules also one supposes only that $a$ is odd and $b$ even.

If one does not dislike using formulas restricted to primary moduli, we may use this form:
\[
\begin{array}{c|c}
\text{number} & \text{character }\equiv\\
\hline
-1+i & \frac14(-a-b+1-bb)\\
-1-i & \frac14(a-b-1+bb)\\
+1-i & \frac14(-a-b+1+bb)
\end{array}
\]
The simplest formulas appear if, as we did at the beginning of our induction, we distinguish moduli of the first and second kind. Namely, the character is
\[
\begin{array}{c|c|c}
\text{number} & \text{for moduli of the first kind} & \text{for moduli of the second kind}\\
\hline
-1+i & \frac14(-a-b+1) & \frac14(-a-b-3)\\
-1-i & \frac14(a-b-1) & \frac14(a-b+3)\\
+1-i & \frac14(-a-b+1) & \frac14(-a-b+5)
\end{array}
\]

\art{65}
For the number $-1+2i$, to which we now proceed, we shall use the same distinction between moduli $a+bi$ for which $a\equiv1$, $b\equiv0$, and those for which $a\equiv3$, $b\equiv2$. The table of Article~62 teaches that, with respect to this number, there corresponds

\begin{center}\small
\begin{tabular}{C L}
\toprule
\text{character} & \text{moduli of the first kind}\\
\midrule
0 & -3+8i, +5-8i, +9+4i, -11+4i\\
1 & +1+4i, +5-4i, -7, -3-8i\\
2 & +1-4i, +5+8i, -7-8i, -11\\
3 & -3, +5+4i, +9-4i, -7+8i, -11-4i\\
\bottomrule
\end{tabular}
\end{center}

Reducing each of these moduli to its absolutely least residue modulo $-1+2i$, we notice that all those to which character $0$ corresponds are congruent to $1$; those to which character $1$ corresponds are congruent to $i$; those whose character is $2$ become congruent to $-1$; finally, all those whose character is $3$ become congruent to $-i$. But the characters of the numbers $1$, $i$, $-1$, and $-i$ modulo $-1+2i$ are respectively $0$, $1$, $2$, and $3$. Hence in all these seventeen examples the character of $-1+2i$ with respect to a modulus $a+bi$ of the first kind is identical with the character of that modulus with respect to the modulus $-1+2i$.

Similarly, with the aid of the table, there is found to correspond

\begin{center}\small
\begin{tabular}{C L}
\toprule
\text{character} & \text{moduli of the second kind}\\
\midrule
0 & +3+2i, -5-2i, -1+10i, -1-10i, +11+6i\\
1 & +3-2i, -1+6i, -5-6i, +7+10i, +7-10i\\
2 & -5+2i, -1-6i, +7-2i\\
3 & -1-2i, +7+2i, -5+6i, +3+10i, +3-10i, +11-6i\\
\bottomrule
\end{tabular}
\end{center}

Reducing these moduli to least residues modulo $-1+2i$, all those to which the characters $0$, $1$, $2$, and $3$ respectively correspond are found congruent to the numbers $-1$, $-i$, $+1$, and $+i$; but to these very numbers, if conversely $-1+2i$ is adopted as modulus, belong the characters $2$, $3$, $0$, and $1$, respectively. Therefore, in all these nineteen examples, the character of $-1+2i$ with respect to a modulus of the second kind differs by two units from the character of that modulus with respect to the number $-1+2i$ taken as modulus.

It is moreover seen without any difficulty that entirely similar facts will hold with respect to the number $-1-2i$.

\art{66}
For the number $-3$ we omit the distinction between moduli of the first and of the second kind, since the result shows that here it is superfluous. Thus there corresponds
\begin{center}\small
\begin{tabular}{c p{5.4in}}
\toprule
character & moduli\\
\midrule
$0$ & $-1+6i,\ -1-6i,\ -7,\ -5+6i,\ -5-6i,\ -11,\ 11+6i,\ 11-6i$\\
$1$ & $-1-2i,\ 1-4i,\ -5+2i,\ 5+4i,\ 7+2i,\ 5-8i,\ -1+10i,\ -7-8i,\ -11-4i,\ 7-10i$\\
$2$ & $3+2i,\ 3-2i,\ -3+8i,\ -3-8i,\ 9+4i,\ 3+10i,\ 3-10i$\\
$3$ & $-1+2i,\ 1+4i,\ -5-2i,\ 5-4i,\ 7-2i,\ 5+8i,\ -1-10i,\ -7+8i,\ -11+4i,\ 7+10i$\\
\bottomrule
\end{tabular}
\end{center}

Reducing these moduli to least residues modulo $3$, we see that those to which character $0$ corresponds are partly $\equiv1$ and partly $\equiv-1$; those whose character is $1$ become either $\equiv1-i$ or $\equiv-1+i$; those whose character is $2$ become either $\equiv i$ or $\equiv-i$; finally, those to which character $3$ belongs are either $\equiv1+i$ or $\equiv-1-i$. From this induction we therefore infer that the character of the number $-3$ with respect to a modulus which is a primary prime among its associates is identical with the character of this same number when $3$, or, what comes to the same thing, $-3$, is considered as modulus.

\art{67}
By a similar induction carried out about other prime numbers, we find that the numbers $3\pm2i$, $-1\pm6i$, $7\pm2i$, $-5\pm6i$, etc., furnish theorems similar to the one at which we arrived in Article~65 with respect to the number $-1+2i$; on the other hand, the numbers $1\pm4i$, $5\pm4i$, $-3\pm8i$, $5\pm8i$, $9\pm4i$, etc., behave in the same way as the number $-3$.

Induction therefore leads to an exceedingly elegant theorem, which, after the model of the theory of quadratic residues in the arithmetic of real numbers, may be called the \textit{fundamental theorem of the theory of biquadratic residues}, namely:

Let $a+bi$ and $a'+b'i$ denote distinct prime numbers, each primary among its associates, that is, congruent to unity modulo $2+2i$. The biquadratic character of the number $a+bi$ with respect to the modulus $a'+b'i$ will be identical with the character of the number $a'+b'i$ with respect to the modulus $a+bi$ if either both of the numbers $a+bi$, $a'+b'i$, or at least one of them, belongs to the first kind, that is, is congruent to unity modulo $4$. On the other hand, those characters will differ from one another by two units if neither of the numbers $a+bi$, $a'+b'i$ belongs to the first kind, that is, if each is congruent to $3+2i$ modulo $4$.

Notwithstanding the great simplicity of this theorem, however, its proof must be counted among the most deeply hidden mysteries of higher arithmetic, so that, at least as matters now stand, it can be unravelled only by very subtle investigations, which would far exceed the limits of the present memoir. We therefore reserve for the third memoir the publication of this proof, as well as the development of the connection between this theorem and the results which at the beginning of this memoir we had begun to establish by induction. As a kind of crown to the present discussion, however, we shall give here what is needed for the proof of the theorems proposed in Articles~63 and 64.

\art{68}
We begin with prime numbers $a+bi$ such that $b=0$ (the third species of Article~34). Thus, in order that the number be primary among its associates, $a$ must be a negative real prime number of the form $-(4n+3)$; for it we shall write $-q$, such as $-3$, $-7$, $-11$, $-19$, and so on.

Let $\lambda$ denote the character of the number $1+i$ when that number is taken as modulus. We must have
\[
i^\lambda \equiv (1+i)^{\frac14(qq-1)}
       \equiv 2^{\frac18(qq-1)}\cdot i^{\frac18(qq-1)}
       \pmod q.
\]
But it is known that $2$ is a quadratic residue, or a quadratic non-residue, of $q$ according as $q$ is of the form $8n+7$, or of the form $8n+3$. Hence we infer that, in general,
\[
2^{\frac12(q-1)}\equiv (-1)^{\frac14(q+1)}
      \equiv i^{\frac12(q+1)}\pmod q,
\]
and therefore, raising to the power whose exponent is $\frac14(q+1)$,
\[
2^{\frac18(qq-1)}\equiv i^{\frac18(q+1)^2}\pmod q.
\]
The preceding equation therefore takes the form
\[
i^\lambda\equiv i^{\frac18(q+1)^2+\frac18(qq-1)}
      \equiv i^{\frac14(qq+q)}\pmod q,
\]
from which it follows that
\[
\lambda\equiv \frac14(qq+q)
      \equiv \frac14(q+1)^2-\frac14(q+1)\pmod4,
\]
or, since $\frac14(q+1)^2\equiv0\pmod4$,
\[
\lambda\equiv-\frac14(q+1)\equiv\frac14(a-1)\pmod4.
\]
This is precisely the theorem of Article~63 for the case $b=0$.

\art{69}
Far more difficult are those moduli $a+bi$ for which $b$ is not $0$ (numbers of the fourth species in Article~34), and several investigations must be put in place first. We shall denote by $p$ the norm $aa+bb$, which will be a real prime number of the form $4n+1$.

Let $S$ denote the complex of all simply least residues modulo $a+bi=m$, zero excluded; thus the number of entries contained in $S$ is $p-1$. Let $x+yi$ denote, indefinitely, a number of this system, and put
\[
ax+by=\xi,\qquad ay-bx=\eta.
\]
Then $\xi$ and $\eta$ will be integers contained exclusively between the limits $0$ and $p$. In the present case, where $a$ and $b$ are relatively prime, the formulas of Article~45, namely
\[
\eta\equiv k\xi,\qquad \xi\equiv-k\eta\pmod p,
\]
show that neither of the numbers $\xi$, $\eta$ can be $0$, unless the other vanishes at the same time, and hence $x=0$, $y=0$, a combination that we have already excluded. Thus the criterion for a number $x+yi$ contained in $S$ consists in the positivity of the four numbers $\xi$, $\eta$, $p-\xi$, $p-\eta$.

We further observe that for no such number can $\xi=\eta$ hold; for then it would follow that
\[
p(x+y)=a(\xi+\eta)+b(\xi-\eta)=2a\xi,
\]
which is absurd, since none of the factors $2$, $a$, $\xi$ is divisible by $p$. In a similar way the equation
\[
p(x-y+a+b)=2a\xi+(a+b)(p-\xi-\eta)
\]
shows that $\xi+\eta=p$ cannot hold.

Therefore, since the numbers $\xi-\eta$ and $p-\xi-\eta$ must be either positive or negative, we derive from this a subdivision of the system $S$ into four complexes $C$, $C'$, $C''$, $C'''$, namely by putting into them as follows:
\[
\begin{array}{c|l}
\text{into the complex} & \text{the numbers for which}\\
\hline
C & \xi-\eta\ \text{is positive, } p-\xi-\eta\ \text{is positive}\\
C' & \xi-\eta\ \text{is positive, } p-\xi-\eta\ \text{is negative}\\
C'' & \xi-\eta\ \text{is negative, } p-\xi-\eta\ \text{is negative}\\
C''' & \xi-\eta\ \text{is negative, } p-\xi-\eta\ \text{is positive}
\end{array}
\]
Thus the criterion for a number of the complex $C$ is properly sixfold: the six numbers $\xi$, $\eta$, $p-\xi$, $p-\eta$, $\xi-\eta$, $p-\xi-\eta$ must be positive. But plainly conditions 2, 5, and 6 already imply the rest. Similar statements hold for the complexes $C'$, $C''$, $C'''$, so that the complete criteria are triple:
\[
\begin{array}{c|l}
\text{for the complex} & \text{the numbers that must be positive}\\
\hline
C & \eta,\ \xi-\eta,\ p-\xi-\eta\\
C' & p-\xi,\ \xi-\eta,\ \xi+\eta-p\\
C'' & p-\eta,\ \eta-\xi,\ \xi+\eta-p\\
C''' & \xi,\ \eta-\xi,\ p-\xi-\eta
\end{array}
\]
Moreover, even without our saying so, anyone will readily understand that, in the figurative representation of complex numbers (see Article~39), the numbers of the system $S$ are contained inside the square whose sides join the points representing the numbers $0$, $a+bi$, $(1+i)(a+bi)$, and $i(a+bi)$, and that the subdivision of the system $S$ corresponds to the partition of the square by its diagonal lines. But at this point we have preferred to use purely arithmetical reasonings, leaving the illustration through geometric intuition to the experienced reader for the sake of brevity.

\art{70}
If four complex numbers $r=x+yi$, $r'=x'+y'i$, $r''=x''+y''i$, $r'''=x'''+y'''i$ are connected with one another so that
\[
r'=m+ir,\qquad r''=m+ir'=(1+i)m-r,\qquad
r'''=m+ir''=im-ir,
\]
and the first, $r$, is supposed to belong to the complex $C$, then the remaining numbers $r'$, $r''$, $r'''$ will belong respectively to the complexes $C'$, $C''$, $C'''$.

Indeed, putting $\xi=ax+by$, $\eta=ay-bx$, $\xi'=ax'+by'$, $\eta'=ay'-bx'$, $\xi''=ax''+by''$, $\eta''=ay''-bx''$, $\xi'''=ax'''+by'''$, $\eta'''=ay'''-bx'''$, one finds
\[
\begin{aligned}
\eta=p-\xi'&=p-\eta''=\xi''',\\
\xi-\eta=\xi'+\eta'-p&=\eta''-\xi''=p-\xi'''-\eta''',\\
p-\xi-\eta=\xi'-\eta'&=\xi''+\eta''-p=\eta'''-\xi''' .
\end{aligned}
\]
The truth of the theorem then follows at once with the aid of the criteria.

And since again $r=m+ir'''$, it is easily seen that if $r$ is supposed to belong to $C'$, then the numbers $r'$, $r''$, $r'''$ belong respectively to $C''$, $C'''$, $C$; if $r$ belongs to $C''$, they belong to $C'''$, $C$, $C'$; and finally, if $r$ belongs to $C'''$, they belong to $C$, $C'$, $C''$. At the same time it follows that in each of the complexes $C$, $C'$, $C''$, $C'''$ there are equally many numbers, namely $\frac14(p-1)$.

\art{71}
\textit{Theorem.} Let $k$ denote an integer not divisible by $m$. Multiply each number of the complex $C$ by $k$, and distribute the simply least residues of the products modulo $m$ among the complexes $C$, $C'$, $C''$, $C'''$. If the numbers of those residues which belong to the four complexes are denoted respectively by $c$, $c'$, $c''$, $c'''$, then the character of the number $k$ with respect to the modulus $m$ will be
\[
\equiv c'+2c''+3c'''\pmod4.
\]

\textit{Proof.} Let the $c$ least residues belonging to $C$ be $\alpha$, $\beta$, $\gamma$, $\delta$, etc.; then let the $c'$ residues belonging to $C'$ be $m+i\alpha'$, $m+i\beta'$, $m+i\gamma'$, $m+i\delta'$, etc.; further, let the $c''$ residues belonging to $C''$ be $(1+i)m-\alpha''$, $(1+i)m-\beta''$, $(1+i)m-\gamma''$, $(1+i)m-\delta''$, etc.; and finally let the $c'''$ residues belonging to $C'''$ be $im-i\alpha'''$, $im-i\beta'''$, $im-i\gamma'''$, $im-i\delta'''$, etc.

Now consider four products, namely:

\begin{enumerate}[label=\arabic*.]
\item the product of all the $\frac14(p-1)$ numbers that constitute the complex $C$;
\item the product of the products obtained by multiplying each of these numbers by $k$;
\item the product of the least residues of these products, namely of the numbers $\alpha$, $\beta$, $\gamma$, $\delta$, etc., $m+i\alpha'$, $m+i\beta'$, etc. etc.;
\item the product of all the $c+c'+c''+c'''$ numbers $\alpha$, $\beta$, $\gamma$, $\delta$, etc., $\alpha'$, $\beta'$, $\gamma'$, etc., $\alpha''$, $\beta''$, $\gamma''$, etc., $\alpha'''$, $\beta'''$, $\gamma'''$, etc.
\end{enumerate}

Denote these four products, in order, by $P$, $P'$, $P''$, $P'''$. Plainly
\[
P'=k^{\frac14(p-1)}P,\qquad
P'\equiv P'',\qquad
P''\equiv P'''i^{c'+2c''+3c'''}\pmod m,
\]
and hence
\[
Pk^{\frac14(p-1)}\equiv P'''i^{c'+2c''+3c'''}\pmod m.
\]
But it is easily seen that the numbers $\alpha'$, $\beta'$, $\gamma'$, etc., $\alpha''$, $\beta''$, $\gamma''$, etc., $\alpha'''$, $\beta'''$, $\gamma'''$, etc. all belong to the complex $C$, and that they are distinct both from one another and from the numbers $\alpha$, $\beta$, $\gamma$, $\delta$, etc., just as these latter are mutually distinct. All these numbers taken together, and abstracting from their order, must therefore be completely identical with all the numbers constituting the complex $C$. Hence we infer that $P=P'''$, and so
\[
Pk^{\frac14(p-1)}\equiv Pi^{c'+2c''+3c'''}\pmod m.
\]
Finally, since none of the individual factors of the product $P$ is divisible by $m$, it follows that
\[
k^{\frac14(p-1)}\equiv i^{c'+2c''+3c'''}\pmod m,
\]
and hence $c'+2c''+3c'''$ is the character of the number $k$ with respect to the modulus $m$. Q. E. D.

\art{72}
In order that the general theorem of the preceding article may be applied to the number $1+i$, the complex $C$ must again be subdivided into two smaller complexes $G$ and $G'$. We shall refer to $G$ those numbers $x+yi$ for which $ax+by=\xi$ is less than $\frac12p$, and to $G'$ those for which $\xi$ is greater than $\frac12p$. We shall denote by $g$ and $g'$ respectively the number of entries contained in $G$ and $G'$, whence $g+g'=\frac14(p-1)$.

The complete criterion for numbers belonging to $G$ will therefore be that the three numbers $\eta$, $\xi-\eta$, $p-2\xi$ are positive; for the third condition for the complex $C$, according to which $p-\xi-\eta$ must be positive, is already implicitly contained in these, since
\[
p-\xi-\eta=(\xi-\eta)+(p-2\xi).
\]
Similarly, the complete criterion for numbers belonging to $G'$ consists in positive values of the three numbers $\eta$, $p-\xi-\eta$, and $2\xi-p$.

It follows easily that the product of any number of the complex $G$ by the number $1+i$ belongs to the complex $C'''$. For if we put
\[
\begin{aligned}
(x+yi)(1+i)&=x'+y'i, \quad
ax'+by'=\xi',\quad ay'-bx'=\eta',\\
\text{one finds}\quad
\xi'&=\xi-\eta,\qquad
\eta'-\xi'=2\eta,\qquad
p-\xi'-\eta'=p-2\xi,
\end{aligned}
\]
that is, the criterion for a number $x+yi$ placed in the complex $G$ is identical with the criterion for the number $x'+y'i$ belonging to the complex $C'''$. In exactly the same way it is shown that the product of any number of $G'$ by $1+i$ belongs to the complex $C''$.

Therefore, if in the preceding article we give $k$ the value $1+i$, then $c=0$, $c'=0$, $c''=g'$, $c'''=g$, and hence the character of the number $1+i$ becomes
\[
3g+2g'=\frac12(p-1)+g.
\]
And since the characters of the numbers $i$ and $-1$ are $\frac14(p-1)$ and $\frac12(p-1)$, respectively, the characters of the numbers $-1+i$, $-1-i$, $1-i$ are respectively
\[
\frac34(p-1)+g,\qquad g,\qquad \frac14(p-1)+g.
\]
Thus the whole crux of the matter now turns on the investigation of the number $g$.

\art{73}
What we have set out in Articles~69--72 is properly independent of the supposition that $m$ is a primary number. From now on, however, we shall at least suppose that $a$ is odd, $b$ even, and moreover that $a$, $b$, and $a-b$ are positive numbers.

Above all, the limits of the values of $x$ in the complex $G$ must be established. Putting
\[
ay-bx=\eta,\qquad (a+b)x-(a-b)y=\zeta,\qquad p-2ax-2by=\theta,
\]
the criterion for numbers $x+yi$ belonging to the complex $G$ consists in the three conditions that $\eta$, $\zeta$, and $\theta$ are positive numbers.

Since
\[
px=(a-b)\eta+a\zeta,\qquad p(a-2x)=a\theta+2b\eta,
\]
it is plain that $x$ and $2a-x$ must be positive numbers; that is, $x$ must be equal to one of the numbers $1$, $2$, $3,\ldots,\frac12(a-1)$.

Further, since
\[
(a-b)\theta=2b\zeta+p(a-b-2x),
\]
it is clear that as long as $x$ is less than $\frac12(a-b)$, the second condition (according to which $\zeta$ must be positive) already implies the third (that $\theta$ must be positive); on the other hand, whenever $x$ is greater than $\frac12(a-b)$, the second condition is already contained under the third. Therefore, for the values of $x$ equal to $1$, $2$, $3,\ldots,\frac12(a-b-1)$, one need only ensure that $\eta$ and $\zeta$ become positive; that is, that $y$ is greater than $\frac{bx}{a}$ and less than $\frac{(a+b)x}{a-b}$. Thus, for such a given value of $x$, there will be exactly
\[
\gbrack{\frac{(a+b)x}{a-b}}-\gbrack{\frac{bx}{a}}
\]
numbers $x+yi$, if we use brackets in the same meaning in which we have already used them in many other places.

On the other hand, for the values of $x$ equal to $\frac12(a-b+1)$, $\frac12(a-b+3),\ldots,\frac12(a-1)$, it will suffice to make $\eta$ and $\theta$ positive; that is, $y$ must be greater than $\frac{bx}{a}$ and less than
\[
\frac{p-2ax}{2b}=\frac12b+\frac{aa-2ax}{2b}.
\]
Hence, for such a given value of $x$, there will be exactly
\[
\gbrack{\frac12b+\frac{aa-2ax}{2b}}-\gbrack{\frac{bx}{a}}
\]
numbers $x+yi$. We therefore infer that the number of entries of the complex $G$ is
\[
g=\Sigma\gbrack{\frac{(a+b)x}{a-b}}
 +\Sigma\gbrack{\frac12b+\frac{aa-2ax}{2b}}
 -\Sigma\gbrack{\frac{bx}{a}},
\]
where in the first term the summation is to be extended over all integral values of $x$ from $1$ to $\frac12(a-b-1)$; in the second from $\frac12(a-b+1)$ to $\frac12(a-1)$; in the third from $1$ to $\frac12(a-1)$.

If we use the characteristic $\varphi$ in the same meaning as in the cited place, namely
\[
\varphi(t,u)=\gbrack{\frac{u}{t}}+\gbrack{\frac{2u}{t}}+\gbrack{\frac{3u}{t}}+\cdots+\gbrack{\frac{t'u}{t}},
\]
where $t$ and $u$ denote any positive numbers and $t'$ denotes the number $\gbrack{\frac12t}$, then the first term becomes $\varphi(a-b,a+b)$, and the third becomes $-\varphi(a,b)$. The second becomes
\[
\frac14bb+\sum\gbrack{\frac{aa-2ax}{2b}}.
\]
But, writing the terms in the reverse order,
\[
\Sigma\gbrack{\frac{aa-2ax}{2b}}
=\gbrack{\frac{a}{2b}}+\gbrack{\frac{3a}{2b}}+\gbrack{\frac{5a}{2b}}+\cdots+\gbrack{\frac{(b-1)a}{2b}}
=\varphi(2b,a)-\varphi(b,a).
\]
Thus our formula takes the form
\[
g=\varphi(a-b,a+b)+\varphi(2b,a)-\varphi(a,b)-\varphi(b,a)+\frac14bb.
\]

First consider the term $\varphi(a-b,a+b)$, which is at once transformed into
\[
\varphi(a-b,2b)+1+2+3+\cdots+\frac12(a-b-1),
\]
or into
\[
\varphi(a-b,2b)+\frac18((a-b)^2-1).
\]
Next, since by the general theorem
\[
\varphi(t,u)+\varphi(u,t)=\gbrack{\frac12t}\cdot\gbrack{\frac12u},
\]
when $t$ and $u$ are positive relatively prime integers, we have
\[
\varphi(a-b,2b)=\frac12b(a-b-1)-\varphi(2b,a-b),
\]
and therefore
\[
\varphi(a-b,a+b)=\frac18(aa+2ab-3bb-4b-1)-\varphi(2b,a-b).
\]
Let us arrange the parts of $\varphi(2b,a-b)$ as follows:
\[
\begin{aligned}
&\gbrack{\frac{a-b}{2b}}+\gbrack{\frac{3(a-b)}{2b}}+\gbrack{\frac{5(a-b)}{2b}}+\cdots+\gbrack{\frac{(b-1)(a-b)}{2b}}\\
+{}&\gbrack{\frac{a-b}{b}}+\gbrack{\frac{2(a-b)}{b}}+\gbrack{\frac{3(a-b)}{b}}+\cdots+\gbrack{\frac{\frac12b(a-b)}{b}}.
\end{aligned}
\]
The second series plainly becomes
\[
\varphi(b,a-b)=\varphi(b,a)-1-2-3-\cdots-\frac12b
=\varphi(b,a)-\frac18(bb+2b).
\]
We write the first series in the reverse order of its terms as
\[
\begin{aligned}
&\gbrack{\frac12(a+1-b)-\frac{a}{2b}}
+\gbrack{\frac12(a+3-b)-\frac{3a}{2b}}\\
&+\gbrack{\frac12(a+5-b)-\frac{5a}{2b}}
+\cdots+\gbrack{\frac12(a-1)-\frac{(b-1)a}{2b}},
\end{aligned}
\]
which expression, since in general, if $t$ denotes an integer and $u$ a fraction, $[t-u]=t-1-[u]$, is changed into
\[
\begin{aligned}
&\frac18b(2a-4-b)-\gbrack{\frac{a}{2b}}-\gbrack{\frac{3a}{2b}}
-\gbrack{\frac{5a}{2b}}-\cdots-\gbrack{\frac{(b-1)a}{2b}}\\
={}&\frac18b(2a-4-b)-\varphi(2b,a)+\varphi(b,a).
\end{aligned}
\]
Hence
\[
\varphi(2b,a-b)=2\varphi(b,a)-\varphi(2b,a)+\frac14b(a-3-b),
\]
and consequently
\[
\varphi(a-b,a+b)=\varphi(2b,a)-2\varphi(b,a)+\frac18(aa-bb+2b-1).
\]
Substituting this value in the formula for $g$ given above, and also using
\[
\varphi(a,b)+\varphi(b,a)=\frac14b(a-1),
\]
we obtain
\[
g=2\varphi(2b,a)-2\varphi(b,a)+\frac18(aa-2ab+bb+4b-1).
\]

\art{74}
By entirely similar reasonings one handles the case in which $a$ and $b$ remain positive but $a-b$ is negative, or $b-a$ positive. The equations
\[
p(a-2x)=2b\eta+a\theta,\qquad
p(b-a+2x)=2b\zeta+(b-a)\theta
\]
show that $\frac12a-x$ and $x+\frac12(b-a)$ are positive, and hence that $x$ must be equal to one of the numbers
\[
-\frac12(b-a-1),\ -\frac12(b-a-3),\ -\frac12(b-a-5),\ldots,+\frac12(a-1).
\]

Further, from the equation $px+(b-a)\eta=a\zeta$ it follows that, for negative values of $x$, the condition by which $\eta$ must be positive is already contained under the condition by which $\zeta$ must be positive; the opposite happens whenever a positive value is assigned to $x$. Hence the values of $y$, for a fixed negative value of $x$, must be contained between $\frac{(a+b)x}{a-b}$ and $\frac{p-2ax}{2b}$, while for a positive value of $x$ they must be contained between $\frac{bx}{a}$ and $\frac{p-2ax}{2b}$. Plainly for $x=0$ these limits are $0$ and $\frac{p-2ax}{2b}$, the value $y=0$ itself being excluded.

It follows that
\[
g=-\sum\gbrack{\frac{(a+b)x}{a-b}}
+\sum\gbrack{\frac12b+\frac{aa-2ax}{2b}}
-\sum\gbrack{\frac{bx}{a}},
\]
where in the first term the summation is to be extended over all negative values of $x$ from $-1$ to $-\frac12(b-a-1)$; in the second over all values of $x$ from $-\frac12(b-a-1)$ to $\frac12(a-1)$; in the third over all positive values of $x$ from $+1$ to $\frac12(a-1)$. In this way the first summation gives $-\varphi(b-a,b+a)$; the second, as in the preceding article, gives $\frac14bb+\varphi(2b,a)-\varphi(b,a)$; and the third gives $-\varphi(a,b)$. Thus
\[
g=-\varphi(b-a,b+a)+\varphi(2b,a)-\varphi(b,a)-\varphi(a,b)+\frac14bb.
\]
Now, in the same way as in the preceding article,
\[
\begin{aligned}
\varphi(b-a,b+a)
&=\varphi(b-a,2b)-\frac18((b-a)^2-1)\\
&=\frac18(3bb-2ab-aa-4b+1)-\varphi(2b,b-a),
\end{aligned}
\]
and also
\[
\varphi(2b,b-a)=\varphi(2b,a)-2\varphi(b,a)+\frac14b(b-1-a),
\]
so that
\[
\varphi(b-a,b+a)=2\varphi(b,a)-\varphi(2b,a)+\frac18(bb-aa-2b+1),
\]
and finally
\[
g=2\varphi(2b,a)-2\varphi(b,a)+\frac18(aa-2ab+bb+4b-1).
\]
It has therefore been established that the same formula for $g$ holds whether $a-b$ is positive or negative, provided $a$ and $b$ are positive.

\art{75}
To obtain the further reduction, put
\[
\begin{aligned}
L&=\gbrack{\frac{a}{2b}}+\gbrack{\frac{2a}{2b}}+\gbrack{\frac{3a}{2b}}+\cdots+\gbrack{\frac{\frac12ba}{2b}},\\
M&=\gbrack{\frac{(\frac12b+1)a}{2b}}+\gbrack{\frac{(\frac12b+2)a}{2b}}+\gbrack{\frac{(\frac12b+3)a}{2b}}+\cdots+\gbrack{\frac{ba}{2b}},\\
N&=\gbrack{\frac{a+b}{2b}}+\gbrack{\frac{2a+b}{2b}}+\gbrack{\frac{3a+b}{2b}}+\cdots+\gbrack{\frac{\frac12ba+b}{2b}}.
\end{aligned}
\]
Since it is easily seen that, in general, $[u]+\left[u+\frac12\right]=[2u]$ for any real quantity $u$, we get $L+N=\varphi(b,a)$; and since plainly $L+M=\varphi(2b,a)$, we have
\[
\varphi(2b,a)-\varphi(b,a)=M-N.
\]
Further, it is obvious that the sum of the first term of the series $N$ with the penultimate term of the series $M$, namely
\[
\gbrack{\frac{a+b}{2b}}+\gbrack{\frac{(b-1)a}{2b}},
\]
becomes $\frac12(a-1)$, and that the same sum is produced from the second term of the series $N$ with the antepenultimate term of $M$, and so on. Hence, since also the last term of the series $M$ becomes $\frac12(a-1)$, while the last term of the series $N$ is
\[
\gbrack{\frac{a+2}{4}}=\frac14(a\mp1),
\]
with the upper or lower sign holding according as $a$ is of the form $4n+1$ or $4n-1$, we get
\[
M+N=\frac14(a-1)b+\frac14(a\mp1),
\]
and hence
\[
\varphi(2b,a)-\varphi(b,a)=\frac14(a-1)b+\frac14(a\mp1)-2N.
\]
Thus the formula for $g$ found in Articles~73 and 74 passes into
\[
g=\frac18((a+b)^2-1)+2n-4N,
\]
putting $a\mp1=4n$, where $n$ is an integer. But since from this we have $1=16nn-8an+aa$, the formula may also be written in the following way:
\[
g=\frac18(-aa+2ab+bb+1)+4\left(\frac12(a+1)n-nn-N\right).
\]
Consequently, since $g$ is the character of the number $-1-i$ with respect to the modulus $a+bi$, this character becomes
\[
\equiv \frac18(-aa+2ab+bb+1)\pmod4.
\]
This is precisely the theorem found above by induction in Article~64, and from it the theorems about the characters of the numbers $1+i$, $1-i$, and $-1+i$ follow at once. Therefore these four theorems have now been rigorously proved for the case in which $a$ and $b$ are positive.

\art{76}
If, while $a$ remains positive, $b$ is negative, put $b=-b'$, so that $b'$ is positive. Since it has now been proved that, for the modulus $a+b'i$, the character of the number $-1-i$ is
\[
\equiv \frac18(-aa+2ab'+b'b'+1)\pmod4,
\]
the character of the number $-1+i$ for the modulus $a-b'i$, by the theorem stated in Article~62, will be
\[
\equiv \frac18(aa-2ab'-b'b'-1).
\]
That is, the character of the number $-1+i$ for the modulus $a+bi$ becomes
\[
\equiv \frac18(aa+2ab-bb-1).
\]
But this is exactly the theorem stated in Article~64, from which the remaining three statements about the characters of $1+i$, $1-i$, and $-1-i$ follow at once. Therefore these theorems have also been proved for the case in which $b$ is negative, that is, for all cases in which $a$ is positive.

Finally, if $a$ is negative, put $a=-a'$, $b=-b'$. Since by what has already been proved the character of the number $1+i$ with respect to the modulus $a'+b'i$ is
\[
\equiv \frac18(-a'a'+2a'b'-3b'b'+1)\pmod4,
\]
and since it makes no difference whether we take $a'+b'i$ or its opposite $-a'-b'i$ as modulus, it is plain that the character of the number $1+i$ with respect to the modulus $a+bi$ is
\[
\equiv \frac18(-aa+2ab-3bb+1),
\]
and similar statements hold about the characters of the numbers $1-i$, $-1+i$, and $-1-i$. It follows from this that the proof of the theorems about the characters of the numbers $1+i$, $1-i$, $-1+i$, and $-1-i$ (Articles~63 and 64) is no longer subject to any limitation.


\clearpage

\section*{Notices of Gauss's Own Writings}

\begin{center}
\textsc{Göttingische gelehrte Anzeigen. 12 May 1808.}
\end{center}

\subsection*{Theorematis arithmetici demonstratio nova}

A memoir submitted by Prof. Gauss to the Royal Society of Sciences on 15 January of this year,
\begin{center}\textit{Theorematis arithmetici demonstratio nova},\end{center}
for which we still have to supply a notice of its contents, concerns the celebrated fundamental theorem in the theory of quadratic residues, a theorem that plays so important a role both in the whole of higher arithmetic and in the adjacent parts of analysis. As is well known, an integer $a$ is called a quadratic residue of the integer $b$ if there are numbers of the form $xx-a$ that are divisible by $b$; in the opposite case $a$ is called a quadratic non-residue of $b$. The number $a$ may be positive or negative, whereas $b$ is always regarded as positive. Higher arithmetic teaches that all prime numbers $b$ for which a given number $a$ is a quadratic residue are comprised under certain linear forms, and that other linear forms likewise contain all prime numbers of which $a$ is a non-residue. Thus, for example, $-1$ is a quadratic residue of all primes of the form $4n+1$, and a quadratic non-residue of all primes of the form $4n+3$; similarly $+2$ is a quadratic residue of all primes of the forms $8n+1$ and $8n+7$, but a quadratic non-residue of all primes of the forms $8n+3$ and $8n+5$. There is an infinite multitude of similar special theorems, but they can all be derived by combining the two just mentioned with the following general theorem: two distinct positive odd prime numbers $p,q$ always have the same reciprocal relation to one another - that is, one is a quadratic residue or non-residue of the other according as the other is a residue or non-residue of the first - if either both are of the form $4n+1$, or at least one of them is; on the other hand their reciprocal relation is opposite, that is, one is a non-residue of the other when the other is a residue of the first, and conversely, whenever both are at once of the form $4n+3$. This is the fundamental theorem mentioned above, which can be expressed in more than one form. The form chosen here is the one in which it is newly proved in Prof. Gauss's memoir.

The most beautiful theorems of higher arithmetic, and especially those now under discussion, have the peculiar property that they are easily discovered by induction, while their proofs lie extremely hidden and can be tracked down only by very deep investigations. This is precisely what gives higher arithmetic that magical charm which has made it the favourite science of the foremost geometers, quite apart from its inexhaustible wealth, in which it so far surpasses all other parts of pure mathematics. The two special theorems mentioned above were already known to Fermat, who claimed also to possess their proofs. Whether he was mistaken in this we cannot decide, since he never made anything of them public; but we may certainly regard it as possible, since there are several examples of self-deception among other great geometers, namely Euler, Legendre, and even Fermat himself. Euler gave the first proof of the first of those theorems; but this great geometer, despite eager efforts continued over many years, did not succeed in proving the other. It was reserved for Lagrange to fill this gap. Both geometers also proved several other special propositions, but a larger number that they found by induction always resisted their efforts at proof. It is, however, a remarkable play of chance that neither geometer was led by induction to the general fundamental theorem, which admits such a simple formulation. This theorem was first presented, although in a somewhat different form, by Legendre in the \textit{Histoire de l'Académie des Sciences de Paris} for 1785. Both there and later in his \textit{Essai d'une théorie des nombres}, that excellent analyst sought to base the proof on very ingenious investigations, which nevertheless did not lead to the desired goal - a goal which, unless we are mistaken, could not be reached by that route.

The author of the memoir to which this notice is devoted entered the path of higher arithmetic at a time when all earlier work of other geometers in this science was entirely unknown to him. It is chiefly to this circumstance that one should ascribe the fact that everywhere he took a wholly distinctive course. He found the fundamental theorem itself very early by induction, but only a whole year later, after many difficulties and fruitless attempts, did he succeed in finding the first fully rigorous proof, developed in the fourth section of his \textit{Disquisitiones Arithmeticae}. That proof, however, rests on very laborious and lengthy discussions. Later he found three other proofs, free from that inconvenience, but requiring other very deep investigations that are quite heterogeneous in content. One of these proofs is also communicated in the cited work, art. 262; the other two will be made public in due course. Thus there still remained the wish that it might be possible to discover a shorter proof, independent of extraneous investigations. The author therefore hopes that friends of higher arithmetic will be pleased to see a fifth proof, which is presented in the present memoir in less than five pages and seems in every respect to leave nothing to be desired. Given the compressed brevity in which this proof is written, we can give here only an imperfect idea of its course; more would be superfluous here in any case, since the sixteenth volume of the \textit{Commentationes}, in which it is already printed, will soon appear.

The foundation of the proof is the following new theorem. Let $p$ be a positive odd prime and let $k$ be any integer not divisible by $p$. Among the residues arising on division by $p$ of the $\frac12(p-1)$ products
\[
 k,\;2k,\;3k,\ldots,\frac12(p-1)k,
\]
suppose there are $\mu$ residues greater than $\frac12p$, and hence $\frac12(p-1)-\mu$ residues smaller than $\frac12p$. Then $k$ is a quadratic residue of $p$ if $\mu$ is even, and a quadratic non-residue if $\mu$ is odd. The number $\mu$, depending only on $k$ and $p$, may be denoted by $(k,p)$. By a series of inferences that does not admit an excerpt, it is then shown that if $k$ and $p$ are two odd numbers having no common divisor, then
\[
(k,p)+(p,k)+\frac12(k-1)(p-1)
\]
is always even. Hence, whenever $k$ and $p$ are both of the form $4n+3$, one of the numbers $(k,p)$, $(p,k)$ must be even and the other odd; in all other cases, that is, whenever both $k$ and $p$, or at least one of them, have the form $4n+1$, the numbers $(k,p)$ and $(p,k)$ must be either both even or both odd. In conjunction with the above theorem, the truth of the fundamental theorem follows at once. By the same route by which these results are found, the memoir at the same time gives a new proof of the two special theorems mentioned above: it is easily shown that $(-1,p)=\frac12(p-1)$, and hence is even or odd according as $p$ has the form $4n+1$ or $4n+3$; similarly $(2,p)=\frac14(p-1)$ if $p$ has the form $4n+1$, and $(2,p)=\frac14(p+1)$ if $p$ has the form $4n+3$. Thus $(2,p)$ is even whenever $p$ has the form $8n+1$ or $8n+7$, and odd whenever $p$ has the form $8n+3$ or $8n+5$.

\begin{center}
\textsc{Göttingische gelehrte Anzeigen. 19 September 1808.}
\end{center}

\subsection*{Summatio quarumdam serierum singularium}

A lecture submitted by Prof. Gauss to the Royal Society of Sciences,
\begin{center}\textit{Summatio quarumdam serierum singularium},\end{center}
has as its aim to develop more fully and in greater generality a remarkable investigation belonging to the division of the circle, whose foundation had already been laid in the \textit{Disquisitiones Arithmeticae}; to provide it with complete proofs; and to show its unexpected connection with other important truths. Let $n$ be a prime number, let $k$ be any integer not divisible by $n$, let $\omega$ denote the arc $360^\circ/n$, and let the distinct quadratic residues of $n$ among the numbers $1,2,3,4,\ldots,n-1$ be represented by $a,a',a'',\ldots$, while the numbers remaining after these are excluded, that is, the quadratic non-residues of $n$, are represented by $b,b',b'',\ldots$. Then in the cited work, art. 356, it is proved that in the case where $n$ is of the form $4m+1$,
\begin{align*}
&\cos ak\omega+\cos a'k\omega+\cos a''k\omega+\cdots
 -\cos bk\omega-\cos b'k\omega-\cos b''k\omega-\cdots = \pm\sqrt n,\\
&\sin ak\omega+\sin a'k\omega+\sin a''k\omega+\cdots
 -\sin bk\omega-\sin b'k\omega-\sin b''k\omega-\cdots =0,
\end{align*}
whereas in the case where $n$ is of the form $4m+3$, the sum of the first series is $0$ and that of the second is $\pm\sqrt n$. The sign to be prefixed to the radical depends on the value of $k$, or rather on its relation to $n$, and it is easy to determine it for all values of $k$ for a given value of $n$ as soon as it is determined for one value. One can show that for all values of $k$ which are quadratic residues of $n$ exactly the same sign holds, and the opposite sign holds for all those which are quadratic non-residues of $n$. Since in the work just cited the investigation had already been carried this far, and only the determination of the sign for some value of $k$ still remained, one might have thought that, after disposing of the main point, this more precise determination would be easy to supply, especially since induction at once gives an extremely simple result: for $k=1$, or for all values that are quadratic residues of $n$, the radical in the above formulae must always be taken positively. Yet in seeking the proof of this observation we meet with wholly unexpected difficulties, and the procedure which led so satisfactorily to the determination of the absolute value of those series is found to be altogether insufficient when it is a matter of the complete determination of the signs. The metaphysical ground of this phenomenon, to use the expression customary among French geometers, must be sought in the fact that in the division of the circle analysis makes no distinction between the arcs $\omega,2\omega,3\omega,
\ldots,(n-1)\omega$, but includes them all in the same way. Since this gives the investigation a new interest, Prof. Gauss found in it, as it were, a summons to leave nothing untried in order to overcome the difficulty. Only after many varied and fruitless attempts did he succeed, by a route remarkable even in itself.

He begins with the summation of certain series whose terms are included under the following form:
\[
(m,\mu)=\frac{(1-x^m)(1-x^{m-1})(1-x^{m-2})\cdots(1-x^{m-\mu+1})}{(1-x)(1-x^2)(1-x^3)\cdots(1-x^\mu)}.
\]
For brevity denote such a function by $(m,\mu)$; as is shown in the memoir, it is always an integral function of $x$. Then, provided $m$ is a positive integer, the series
\begin{align*}
1-(m,1)+(m,2)-(m,3)+\cdots,\\
1+x^{1/2}(m,1)+x(m,2)+x^{3/2}(m,3)+\cdots
\end{align*}
break off after the $(m+1)$st term, and the sum of the first series, for even values of $m$, becomes
\[
(1-x)(1-x^3)(1-x^5)\cdots(1-x^{m-1}),
\]
and is $0$ for odd values of $m$; the sum of the second series is always
\[
(1+x^{1/2})(1+x)(1+x^{3/2})\cdots(1+x^{m/2}).
\]
For fractional and negative values of $m$ also the summation of these series leads to interesting results, although they are not needed for the present purpose. We content ourselves with citing just one of them. The infinite series
\[
1+x+x^3+x^6+x^{10}+\cdots,
\]
where the exponents are triangular numbers, is the product of the factors
\[
\frac{1-xx}{1-x}\cdot\frac{1-x^4}{1-x^3}\cdot\frac{1-x^6}{1-x^5}\cdot\frac{1-x^8}{1-x^7}\cdots,
\]
or, if one prefers, of
\[
\frac{(1+x)^2(1+x^2)^2(1+x^3)^2(1+x^4)^2\cdots}{(1-x)(1-xx)(1-x^3)(1-x^4)\cdots}.
\]

The development of the way in which these summations are applied to the main subject would take us too far here; we may all the more readily refer readers to the memoir itself, since it will soon appear in print. The summations cited above are only a special application of the summation of the following series:
\begin{align*}
1+\cos k\omega+\cos4k\omega+\cos9k\omega+\cdots+\cos(n-1)^2k\omega&=T,\\
\sin k\omega+\sin4k\omega+\sin9k\omega+\cdots+\sin(n-1)^2k\omega&=U,
\end{align*}
which is taught in the memoir for all values of $k$, without the restriction that $n$ be prime. It is shown that
\[
T=\pm\sqrt n,
\quad T=\pm\sqrt n,
\quad T=0,
\quad T=0
\]
and
\[
U=\pm\sqrt n,
\quad U=0,
\quad U=0,
\quad U=\pm\sqrt n
\]
according as $n$ is respectively of the form $4m$, $4m+1$, $4m+2$, $4m+3$. The sign of the radical again depends on $k$, and the determination of it, which makes it necessary to distinguish many individual cases, is developed and proved in two different ways so completely that nothing remains to be desired. Comparison of these two routes with one another leads further to the following very remarkable theorem. Let $n$ be the product of any number of distinct odd prime numbers $a,b,c,d$, and so on, among which altogether $\mu$ are of the form $4m+3$; and suppose further that among those factors altogether $\nu$ occur such that for each of them the product of the others - respectively $n/a,n/b,n/c,n/d$, and so on - is a quadratic non-residue. Then $\nu$ is even whenever $\mu$ is of the form $4m$ or $4m+1$, and odd whenever $\mu$ is of the form $4m+2$ or $4m+3$. The familiar fundamental theorem on quadratic residues is only a special case of this theorem, just as conversely the latter can easily be derived from the former. Thus through these investigations one is at the same time in possession of a fourth proof of this important theorem, which the author had first proved by two entirely different routes in the \textit{Disquisitiones Arithmeticae} and recently by a third, equally different route, in a separate memoir.

\begin{center}
\textsc{Göttingische gelehrte Anzeigen. 10 March 1817.}
\end{center}

\subsection*{Theorematis fundamentalis in doctrina de residuis quadraticis demonstrationes et ampliationes novae}

On 10 February the Royal Society received from Court Councillor Gauss a lecture entitled
\begin{center}\textit{Theorematis fundamentalis in doctrina de residuis quadraticis demonstrationes et ampliationes novae}.\end{center}
It is a peculiarity of higher arithmetic that so many of its most beautiful theorems can be discovered with the greatest ease by induction, while their proofs are anything but near at hand and are often found only after many fruitless attempts, with the aid of deeply penetrating investigations and fortunate combinations. This remarkable phenomenon springs from the often wonderful interlacing of heterogeneous doctrines in that part of mathematics; and from this also it comes that theorems for which proof had at first been sought in vain for years can later be proved by several wholly different routes. Once a new theorem has been discovered by induction, the discovery of some proof is, of course, the first requirement. But after one proof has succeeded, one should not always regard the investigation in higher arithmetic as closed, nor the search for other proofs as a superfluous luxury. For one usually does not arrive first at the most beautiful and simplest proofs; and besides, precisely the insight into the wonderful concatenation of truth in higher arithmetic is one of the chief attractions of this study and often in turn leads to the discovery of new truths. For these reasons the discovery of new proofs of already known truths is often to be regarded as at least as important as the discovery of the truth itself. Experts in higher arithmetic know these reflections already. It is known that a great part of Euler's services to it consisted in finding proofs of theorems that Fermat, it seems, had already found by induction.

The theory of quadratic residues gives a clear illustration of what has just been said. It rests chiefly on the so-called fundamental theorem, which says that the reciprocal relations of two positive odd primes to one another, insofar as one is a quadratic residue or non-residue of the other, are the same whenever one of the primes or both of them is of the form $4k+1$, and opposite whenever both primes are of the form $4k+3$. For readers less familiar with higher arithmetic, we recall that an integer is called a quadratic residue of another if the former, increased by a multiple of the latter, can give a square; otherwise it is called a non-residue. We will not repeat here the full history of this beautiful theorem, so easily found by induction, but only note that the author of the present memoir, after rather long and initially fruitless investigations, has gradually given four proofs of it that differ completely from one another. Two are contained in the \textit{Disquisitiones Arithmeticae}; the third forms the subject of a separate memoir in the sixteenth volume of the \textit{Commentationes}; and the fourth is woven into the memoir \textit{Summatio quarumdam serierum singularium} in the first volume of the \textit{Commentationes recentiores}. For those two memoirs one may consult our notices of 12 May and 19 September 1808, where fuller historical information is also found. That the author did not rest content with these four proofs, although each of them by itself leaves nothing to be desired in respect of rigor, requires no justification for friends of higher arithmetic; nevertheless, he would probably not have exerted himself so eagerly to add still others to them had not a special circumstance led him to do so, and that circumstance must be mentioned here.

Since 1805 he had begun to occupy himself with the theories of cubic and biquadratic residues, which are still richer and more interesting than the theory of quadratic residues. In those investigations the same phenomena appeared as in the latter theory, only, as it were, on a magnified scale. By induction, as soon as the right path to it had been found, a number of very simple theorems immediately appeared, exhausting those theories, having a surprising resemblance to the theorems that hold for quadratic residues, and in particular presenting the counterpart to the fundamental theorem. But the difficulties in finding fully satisfactory proofs of those theorems proved far greater here; only after many attempts continued through a considerable series of years did the author finally succeed in reaching his goal. The great analogy between the theorems themselves for quadratic and for higher residues suggested that there must also be analogous proofs for the two cases. Yet the proof methods first found for quadratic residues allowed no application at all to higher residues; and precisely this circumstance was the motive for seeking still other new proofs for the former. The author therefore wishes the present memoir, which opens some new resources for the theory of quadratic residues, to be regarded as the forerunner of the theory of cubic and biquadratic residues, which he intends to make known in the future and which belong among the most difficult subjects of higher arithmetic.

The present memoir consists of three mutually independent parts. It contains the fifth and sixth proofs of the fundamental theorem, and a new method, connected with the third proof, for deciding whether a given integer is a quadratic residue or non-residue of a given prime. Among the first four proofs, the third was unquestionably the one that combined the greatest simplicity with independence from extraneous investigations; for that reason Legendre included it in the new edition of his \textit{Essai d'une théorie des nombres}. The fifth proof seems at least to equal the third in both respects. The two proofs have a certain kinship, inasmuch as they start from one and the same lemma, but they are wholly different in their further execution. This lemma is the following. If $m$ is a positive odd prime and $M$ is an integer not divisible by $m$; and if among the residues that arise on division by $m$ of the products
\[
M,\;2M,\;3M,\;4M,\ldots,\frac12(m-1)M
\]
the number of those greater than $\frac12m$ is denoted by $n$, then $M$ is a quadratic residue or non-residue of $m$ according as $n$ is even or odd. To arrive at the proof of the fundamental theorem, one assumes that $M$ is itself an odd positive prime and that $N$, in relation to $M$ and $m$, means the same thing as $n$ expresses in relation to $m$ and $M$, so that $N$ decides, by being even or odd, whether $m$ is a quadratic residue or non-residue of $M$. By a very short series of inferences the author shows that the number of all positive integers that are at the same time less than $\frac12mM$, that give on division by $m$ a remainder less than $\frac12m$, and on division by $M$ a remainder less than $\frac12M$, is
\[
\frac18(m-1)(M-1)+\frac12n+\frac12N,
\]
and hence that
\[
\frac14(m-1)(M-1)+n+N
\]
is always an even number. Thus whenever at least one of the numbers $m,M$ is of the form $4k+1$, so that $\frac14(m-1)(M-1)$ is even, $n+N$ will also be even, and consequently either $n$ and $N$ are both even or both odd. But if both $m$ and $M$ are of the form $4k+3$, then necessarily $n+N$ is odd, and consequently one of the numbers $n,N$ is even and the other odd. In conjunction with the lemma above, the fundamental theorem follows at once.

The sixth proof is of the same brevity and concision as the fifth, but rests on somewhat more artificial combinations. The limited space of these pages allows us only to touch on its chief point, passing over the details. Let $p,q$ denote two distinct positive odd primes; let $a$ be a so-called \textit{radix primitiva} for the modulus $p$, that is, a positive integer not divisible by $p$ such that no lower power than $a^{p-1}$ is congruent to unity modulo $p$; let $x$ be an indeterminate; and let $\xi$ denote the function
\[
x-x^\zeta+x^\eta-x^\theta+x^\iota-\cdots-x^\lambda,
\]
where, for typographical convenience, $\zeta,\eta,\theta,
\ldots,\lambda$ stand for the numbers $aa,a^3,a^4,\ldots,a^{p-2}$. Let $\varepsilon$ be the unit taken positively if $p$ is of the form $4k+1$, negatively if $p$ is of the form $4k+3$; let $\delta$ be the unit taken positively if at least one of $p,q$ is of the form $4k+1$, negatively if both are of the form $4k+3$; let $\gamma$ be the unit taken positively if $q$ is a quadratic residue of $p$, negatively if $q$ is a quadratic non-residue of $p$; and let $\vartheta$ be the unit taken positively if $p$ is a quadratic residue of $q$, negatively if $p$ is a quadratic non-residue of $q$. After these preparations it follows easily from art. 51 of the \textit{Disquisitiones Arithmeticae} that, when the function
\[
\xi^q-x^q+x^{q\zeta}-x^{q\eta}+x^{q\theta}-x^{q\iota}+\cdots+x^{q\lambda}
\]
is expanded, all its coefficients become divisible by $q$; therefore, if this function is put equal to $qX$, then $X$ is an integral function also with respect to its coefficients. By arguments into which it would be too lengthy to enter here, the memoir proves that division of the function $qX\xi$ by
\[
x^{p-1}+x^{p-2}+x^{p-3}+x^{p-4}+\cdots+x+1
\]
gives the remainder
\[
\varepsilon p(\vartheta p^{\frac12(q-1)}-\gamma),
\]
so that division of $X\xi$ by the same divisor gives the remainder
\[
\frac{\varepsilon p(\vartheta p^{\frac12(q-1)}-\gamma)}{q}.
\]
This quantity must therefore necessarily be an integer, whence, since $\delta\delta=1$, it is easily concluded that
\[
p^{\frac12(q-1)}-\gamma\delta
\]
must be divisible by $q$. Since also $p^{\frac12(q-1)}-\vartheta$ is divisible by $q$ by a known theorem, it follows necessarily that $\vartheta=\gamma\delta$, from which the fundamental theorem again follows at once.

The fundamental theorem, combined with a few known lemmas, can indeed furnish a rather short solution of the problem of deciding whether a given positive integer is a quadratic residue or non-residue of a given prime, as is shown in detail in the memoir. But on further reflection about the third proof of the fundamental theorem the author found a still more flexible solution, which forms the third section of the memoir. We transcribe here only the final rule, omitting for brevity the development of its grounds. To decide whether the positive integer $b$, not divisible by the prime $a$, is a quadratic residue or non-residue of $a$, form, in exactly the same way as if one were seeking the greatest common divisor of $a$ and $b$, the equations
\[
\begin{aligned}
a&=\beta b+c,\\
b&=\gamma c+d,\\
c&=\delta d+e,\\
d&=\varepsilon e+f\quad\text{and so on,}
\end{aligned}
\]
until in the series $a,b,c,d,e,f,
\ldots$ one arrives at unity. Denote the numbers $\frac12a,\frac12b,\frac12c,\frac12d,
\ldots$, omitting the attached fraction $\frac12$ wherever some of the numbers $a,b,c,d,
\ldots$ are odd, by $a',b',c',d',
\ldots$; let $\mu$ be the number of occurrences in the series $a',b',c',d',
\ldots$ of two odd numbers immediately following one another; and finally let $\nu$ be the number of those odd numbers in the series $\beta,\gamma,\delta,\varepsilon,
\ldots$ to which, in order, there corresponds in the series $b',c',d',e',
\ldots$ a number of the form $4k+1$ or $4k+2$. With this understood, $b$ is a quadratic residue or non-residue of $a$ according as $\mu+\nu$ is even or odd, with the single exception of the case in which $b$ is even and $a$ is of the form $8k+3$ or $8k+5$; in that case the opposite of the rule holds, so that an even $\mu+\nu$ indicates that $b$ is a quadratic non-residue of $a$, and an odd $\mu+\nu$ indicates that $b$ is a quadratic residue of $a$.

\begin{center}
\textsc{Göttingische gelehrte Anzeigen. 11 April 1825.}
\end{center}

\subsection*{Theoria Residuorum Biquadraticorum, Commentatio prima}

On 5 April Court Councillor Gauss submitted to the Royal Society a lecture entitled
\begin{center}\textit{Theoria Residuorum Biquadraticorum, Commentatio prima}.\end{center}
The theory of quadratic residues is known to form one of the most interesting parts of higher arithmetic, which after repeated investigations may now be regarded as complete and closed. Historical notices concerning it are found in these pages for 12 May and 19 September 1808, and for 10 March 1817. In the latter place some preliminary information was also given about the researches which the author of the present memoir had undertaken since 1805 on the related, equally fruitful and interesting, but much more difficult theory of cubic and biquadratic residues. Although already at that time in possession of the essential elements of this theory, he was prevented by other work from making anything of it public; only now has it become possible for him to occupy himself with the elaboration of part of these investigations. A beginning has now been made with the theory of biquadratic residues, which is more closely related to the theory of quadratic residues than is the theory of cubic residues. Meanwhile the present memoir is by no means intended to exhaust this exceedingly rich subject. The development of the general theory, which requires a very peculiar enlargement of the field of higher arithmetic, is rather reserved for the future continuation, while this first memoir includes those investigations which could be completely presented without such an enlargement. Of the results, only part can be singled out in this notice.

An integer $a$ is called a biquadratic residue of the integer $p$ if there are numbers of the form $x^4-a$ divisible by $p$; a biquadratic non-residue if no numbers of that form can be divisible by $p$. Clearly all biquadratic residues of $p$ are at the same time quadratic residues of that number, and hence all quadratic non-residues are also biquadratic non-residues. But not all quadratic residues are also biquadratic residues. It is sufficient to restrict the investigations to the case where $p$ is a prime of the form $4n+1$ and $a$ is not divisible by $p$, since all other cases can easily be reduced to this one.

The investigations on this subject divide into two parts according as $p$ or $a$ is regarded as given. The first is of much less difficulty than the second and, compared with it, is to be regarded as entirely elementary. The memoir contains completely all that is essential to say about it. From the second part, however, only a few special cases have so far been treated, cases that could be disposed of without excessive preparations and that can serve as preparations for the general theory to be given in the future. These are the cases where $a=-1$ and $a=+2$. The first case has no difficulty at all. It had already been shown in the \textit{Disquisitiones Arithmeticae} that $-1$ is a biquadratic residue of $p$ whenever $p$ has the form $8n+1$, but is merely a quadratic residue and a biquadratic non-residue of $p$ when $p$ has the form $8n+5$. The case $a=+2$ behaves quite differently. It has long been known that $+2$ and $-2$ are quadratic and therefore also biquadratic non-residues of $p$ when $p$ has the form $8n+5$, and at least quadratic residues when $p$ has the form $8n+1$; also that for this form of $p$, either $+2$ and $-2$ are both biquadratic residues, or both are biquadratic non-residues. But distinguishing which of these two cases occurs is an investigation of a much higher kind, and two different criteria are developed for it in the memoir.

The first criterion is connected with the decomposition of $p$ into a single and a double square, which, as already noted, is always possible and possible in only one way, assuming $p$ is prime. If one writes
\[
p=gg+2hh,
\]
then $+2$ is a biquadratic residue of $p$ when $g$ is of the form $8n+1$ or $8n+7$, and a biquadratic non-residue when $g$ is of the form $8n+3$ or $8n+5$. The second criterion is connected with the decomposition of $p$ into two squares, which likewise is always possible and possible in only one way. If one writes $p=ee+ff$, and assumes that $ee$ denotes the odd square and $ff$ the even square, then the assumed form $p=8n+1$ already entails that $\frac12f$ is also even, so that $f$ is either of the form $8m$ or of the form $8m+4$. In the first case $+2$ is a biquadratic residue, in the other a biquadratic non-residue of $p$.

We merely point here to the observation, for which higher arithmetic so often gives occasion, that it is not so much the beauty and simplicity of the theorems themselves as the difficulty of their foundations that makes them especially remarkable. Once one is led to suspect the existence of a connection between the behaviour of the number $+2$ and the two decompositions of $p$ just mentioned, it is extremely easy actually to discover this connection by induction. But even for the first criterion the proof is not altogether easy to conduct; for the second it lies much more deeply hidden, where it is intimately interwoven with other subtle auxiliary investigations that in turn lead to a remarkable extension of the theory of the division of the circle. This wonderful concatenation of truths is especially what, as has often been remarked, gives higher arithmetic such a distinctive charm. These foundations themselves naturally admit no excerpt here and must be consulted in the memoir itself. But a couple of other new arithmetical theorems, likewise intimately connected with the foundation of the second criterion, deserve to be especially singled out here for their simplicity.

If $p$ is a prime of the form $4k+1$ and $p=ee+ff$, with $ee$ denoting the odd square and $ff$ the even square; and if one further puts
\[
1\cdot2\cdot3\cdots k=q,
\qquad
(k+1)(k+2)(k+3)\cdots2k=r,
\]
then $\pm e$ is always the least residue obtained by dividing $r/(2q)$ by $p$, and $\pm f$ the least residue obtained from division of $\frac12rr$ by $p$, ``least residue'' always understood as one taken between the limits $-\frac12p$ and $+\frac12p$. The number $r/(2q)$, which for $p=5$ has the value $1$, can for larger values of $p$ also be written in the form
\[
\frac{6\cdot10\cdot14\cdot18\cdots(p-3)}{2\cdot3\cdot4\cdot5\cdots k}.
\]
It is very remarkable that in this way the decomposition of $p$ into two squares can be obtained by a completely direct route. But an accompanying circumstance is almost still more remarkable. By this procedure one always finds the root of the odd square, $e$, with positive sign when $e$, taken positively, is of the form $4m+1$, and with negative sign when $e$, taken positively, is of the form $4m+3$. On the other hand, for the sign with which the root of the even square, $f$, emerges from that operation, no general rule at all has yet been found, either a priori or by induction. At the close of the memoir the author therefore recommends this subject to friends of higher arithmetic for further investigation, convinced that with its successful solution a rich source of new extensions of this beautiful part of mathematics will at the same time be opened.

\begin{center}
\textsc{Göttingische gelehrte Anzeigen. 23 April 1831.}
\end{center}

\subsection*{Theoria residuorum biquadraticorum, commentatio secunda}

A lecture submitted on 15 April by Court Councillor Gauss to the Royal Society,
\begin{center}\textit{Theoria residuorum biquadraticorum, commentatio secunda},\end{center}
is the continuation of the memoir already printed in the sixth volume of the \textit{Commentationes novae}, for which a notice was also given in these pages at the time, on 11 April 1825. Even this continuation, although more than twice as extensive as the first memoir, does not yet exhaust the exceedingly rich subject, and the completion of the whole will be reserved for a future third memoir.

Although the fundamental concepts of these doctrines and the content of the first memoir may be assumed known to all who have made a study of higher arithmetic, we will nevertheless recall them briefly here for the convenience of those friends of this part of mathematics who do not have the first memoir immediately at hand. Relative to any integer $p$, another integer $k$ is called a biquadratic residue if there are numbers of the form $x^4-k$ divisible by $p$; in the opposite case it is called a biquadratic non-residue of $p$. It is sufficient to restrict oneself here to the case in which $p$ is a prime of the form $4n+1$ and $k$ is not divisible by it, since all other cases are either clear in themselves or reducible to this one.

For such a given value of $p$, all numbers not divisible by $p$ split into four classes. One contains the biquadratic residues; a second contains those biquadratic non-residues that are quadratic residues of $p$; and the two remaining classes contain the biquadratic non-residues that are at the same time quadratic non-residues. The principle of this division consists in the fact that exactly one of $k^n-1$, $k^n-f$, $k^n+1$, $k^n+f$ will be divisible by $p$, where $f$ denotes an integer that makes $ff+1$ divisible by $p$. Anyone acquainted with the elementary terminology will immediately see how these explanations can be clothed in it.

The theory of this classification, not only for the case $k=-1$ lying on the surface, but also for the cases $k=\pm2$, which require subtle auxiliary investigations, is fully completed in the first memoir. At the beginning of the present memoir Gauss proceeds to larger values of $k$; but at first it is necessary to consider only those that are themselves prime numbers, and the result shows that the simplest form of the results is obtained when these values are taken positive or negative according as, in absolute value, they are of the form $4m+1$ or $4m+3$. Induction immediately yields here, with great ease, a rich harvest of new theorems, of which we mention only a few. The numbering of the classes as $1,2,3,4$ is referred to the cases where $k^n$ is congruent to the numbers $1,f,-1,-f$; at the same time, for the number $f$ the value is always chosen that makes $a+bf$ divisible by $p$, when $aa+bb$ represents the decomposition of $p$ into an odd and an even square. Thus induction finds that the number $-3$ always belongs to class $1,2,3,4$ according as $b,a+b,a,a-b$ is divisible by $3$; that the number $+5$ belongs successively to those classes according as $b,a-b,a,a+b$ is divisible by $5$; and that $-7$ falls into class $1$ when $a$ or $b$, into class $2$ when $a-2b$ or $a-3b$, into class $3$ when $a-b$ or $a+b$, and into class $4$ when $a+2b$ or $a+3b$ is divisible by $7$. Similar theorems arise for the numbers $-11,+13,+17,-19,-23$, and so on. But however easy it is to discover all such special theorems by induction, it seems hard by this route to find a general law for these forms, even if much in common soon catches the eye; and it is harder still to find proofs for these theorems. The methods used in the first memoir for the numbers $+2$ and $-2$ no longer admit application here, and although other methods could also serve to settle what concerns the first and third classes, they nevertheless prove unsuitable for establishing complete proofs.

One therefore soon recognizes that one can penetrate this rich domain of higher arithmetic only by wholly new routes. Already in the first memoir the author had indicated that a peculiar enlargement of the entire field of higher arithmetic is essentially required for this. To explain precisely the nature of this enlargement is one of the purposes of the present memoir. Whereas higher arithmetic hitherto had to do only with real integers, here the consideration must be extended to imaginary integers of the form $a+bi$, where $a$ and $b$ are arbitrary real integers and $i$ denotes the usual imaginary unit $\sqrt{-1}$. The real integers thus form only a special case of the complex integers, namely the case $b=0$. In this domain there are four units, $+1,-1,+i,-i$. A complex integer is called composite if it is the product of two integer factors different from a unit; whereas a complex number that admits no such decomposition into factors is called a complex prime. Thus, for example, the real number $3$, even when considered as a complex number, is a prime, while $5$ is composite as a complex number, $5=(1+2i)(1-2i)$. Just as in the higher arithmetic of real numbers, prime numbers play the leading role in the enlarged field of these complex numbers.

Let $a+bi$ be a complex prime whose norm $aa+bb$ is expressed by a real odd prime $p$, and let $k$ be a complex integer not divisible by $a+bi$. Then $k^{\frac14(p-1)}$ will, modulo $a+bi$, be congruent to one of the numbers $+1,+i,-1,-i$; thereby all numbers not divisible by $a+bi$ are divided into four classes, to which the biquadratic characters $0,1,2,3$ are respectively assigned. Clearly character $0$ refers to the biquadratic residues, the others to the biquadratic non-residues, and in such a way that character $2$ corresponds also to quadratic residues, while characters $1$ and $3$ correspond to quadratic non-residues.

One sees easily that the main thing is to be able to determine this character only for such values of $k$ as are themselves complex primes; and here induction at once leads to extremely simple results. First putting $k=1+i$, it appears that the character of this number is always
\[
\equiv \frac18(-aa+2ab-3bb+1)\pmod 4,
\]
and similar expressions are found for the cases $k=1-i$, $k=-1+i$, $k=-1-i$. If, however, $k=\alpha+\beta i$ is such a prime, with $\alpha$ odd and $\beta$ even, induction very readily gives a reciprocity law entirely analogous to the fundamental theorem for quadratic residues. It can be expressed most simply as follows. If both $\alpha+\beta-1$ and $a+b-1$ are divisible by $4$ - to which case all others can easily be reduced - and if the character of the number $\alpha+\beta i$ relative to the modulus $a+bi$ is denoted by $\lambda$, while the character of $a+bi$ relative to the modulus $\alpha+\beta i$ is denoted by $l$, then $\lambda=l$ if one of the numbers $\beta,b$ (or both) is divisible by $4$; but $\lambda=l+2$ if neither $\beta$ nor $b$ is divisible by $4$.

These theorems contain, in essence, all that is essential in the theory of biquadratic residues. But as easy as it was to discover them by induction, it is difficult to give rigorous proofs of them, especially of the second, the fundamental theorem of biquadratic residues. Because the present memoir had already grown to a large extent, the author found it necessary to reserve the exposition of the proof of the latter theorem, which he has possessed for twenty years, for a future third memoir. In the present memoir, however, the complete proof is given for the former theorem concerning the number $1+i$, on which the other theorems for $1-i$, $-1+i$, and $-1-i$ depend; this proof can already give some idea of the complexity of the subject.

We must now add a few general remarks. The transfer of the doctrine of biquadratic residues into the domain of complex numbers might perhaps seem offensive and unnatural to someone less familiar with the nature of imaginary quantities and caught in false notions about them, and might lead to the opinion that the investigation is thereby, as it were, suspended in the air, acquires an unstable posture, and departs entirely from intuition. Nothing would be less well founded than such an opinion. On the contrary, the arithmetic of complex numbers admits the most vivid visualization; and although in his present exposition the author has followed a purely arithmetical treatment, he has nevertheless given the necessary indications for this visualization, which makes insight more lively and is therefore much to be recommended, and which will suffice for readers who think for themselves. Just as the absolute integers are represented by a series of points arranged at equal distances on a straight line, in which the initial point represents the number $0$, the next point the number $1$, and so on; and just as the representation of negative numbers then requires only an unlimited prolongation of this series on the opposite side of the initial point; so, for the representation of the complex integers, one needs only add that this line is regarded as lying in a fixed unbounded plane, and that parallel to it on both sides an unlimited number of similar lines at equal distances from one another is assumed. Thus, instead of one series of points, we have before us a system of points that can be ordered in a double way into rows of rows and can serve to divide the whole plane into equal squares. The nearest point to $0$ in the first neighbouring row on one side of the row representing the real numbers then corresponds to the number $i$, just as the nearest point to $0$ in the first neighbouring row on the other side corresponds to $-i$, and so on. With this representation the performance of arithmetical operations with complex quantities, congruence, the formation of a complete system of incongruent numbers for a given modulus, and so forth, become capable of a visualization that leaves nothing to be desired.

From another side this casts the true metaphysics of imaginary quantities in a new and bright light. Our general arithmetic, whose scope so far surpasses the geometry of the ancients, is entirely the creation of modern times. Starting originally from the concept of absolute integers, it has gradually enlarged its field: fractions were added to integers, irrational numbers to rational ones, negative numbers to positive ones, imaginary numbers to real ones. But this progress always at first proceeded with fearful and hesitant steps. The first algebraists still called the negative roots of equations false roots, and they are indeed false when the problem to which they refer is stated in such a form that the nature of the sought quantity admits no opposite. But just as little as one hesitates in general arithmetic to admit fractions, although there are many countable things for which a fraction has no meaning, so little should equal rights with positive numbers be denied to negative numbers because innumerable things admit no opposite. The reality of negative numbers is sufficiently justified by the fact that in innumerable other cases they find an adequate substrate. This has long been clear. But the imaginary quantities opposed to real quantities - formerly, and here and there still now, though improperly, called impossible - are still less naturalized than merely tolerated. They therefore appear more as an intrinsically empty play of symbols, to which one absolutely denies any thinkable substrate, while nevertheless not wishing to disdain the rich tribute which this play of symbols ultimately contributes to the treasury of relations among real quantities.

For many years the author has considered this highly important part of mathematics from a different point of view, according to which an object can be put beneath imaginary quantities just as well as beneath negative ones. Until now there had been no occasion to state this view publicly in a definite way, although attentive readers will easily find traces of it in his 1799 work on equations and in his prize essay on the transformation of surfaces. In the present memoir its basic features are briefly indicated; they are as follows. Positive and negative numbers can be applied only where the thing counted has an opposite which, thought as united with it, is to be regarded as annihilation. More precisely, this presupposition occurs only where what is counted is not substances, objects thinkable in themselves, but relations between pairs of objects. It is postulated that these objects are ordered in a determinate way into a series, for example $A,B,C,D,\ldots$, and that the relation of $A$ to $B$ can be regarded as equal to the relation of $B$ to $C$, and so on. Here, the concept of opposition requires nothing more than the interchange of the members of the relation, so that if the relation, or transition, from $A$ to $B$ counts as $+1$, then the relation from $B$ to $A$ must be represented by $-1$. Thus, insofar as such a series is unlimited on both sides, every real integer represents the relation of a member chosen arbitrarily as initial to a determinate member of the series.

But if the objects are such that they cannot be ordered in one series, even an unbounded one, but only in rows of rows, or, what is the same, if they form a manifold of two dimensions; and if the relations of one row to another, or the transitions from one row into another, behave in a way analogous to the transitions considered above from one member of a row to another member of the same row, then it is plainly necessary, in order to measure the transition from one member of the system to another, to have besides the previous units $+1$ and $-1$ two other units, also mutually opposite, $+i$ and $-i$. It must moreover be postulated that the unit $i$ always denotes the transition from a given member of a row to a determinate member of the immediately adjacent row. In this way, then, the system can be ordered in a double way into rows of rows.

The mathematician abstracts completely from the nature of the objects and the content of their relations. He has to do only with numbering and comparing the relations among themselves. In this respect, just as he attributes homogeneity in themselves to the relations denoted by $+1$ and $-1$, he is entitled to extend such homogeneity to all four elements $+1,-1,+i,-i$.

These relations can be brought to intuition only by a representation in space, and the simplest case is the one in which there is no reason to arrange the symbols of the objects otherwise than quadratically: that is, one divides an unbounded plane into squares by two systems of parallel lines crossing one another at right angles, and chooses the points of intersection as symbols. Each such point $A$ has four neighbours. If the relation of $A$ to one neighbouring point is denoted by $+1$, then the relation to be denoted by $-1$ is determined of itself, while either of the two others may be chosen for $+i$, or the point referring to $+i$ may be taken to the right or to the left at will. This distinction between right and left, as soon as forward and backward in the plane, and above and below with respect to the two sides of the plane, have once been fixed at will, is completely determined in itself, although we can communicate our intuition of this distinction to others only by indicating actual material things.\footnote{Both observations were already made by Kant; but one cannot understand how that acute philosopher could believe that he found in the former a proof of his opinion that space is only the form of our outer intuition, since the latter so clearly proves the contrary, namely that space must have a real meaning independent of our mode of intuition.} But even after deciding this last point, one sees that it still depended on our choice which of the two rows crossing at one point we wished to regard as the main row, and which direction in it we wished to regard as referring to positive numbers. One further sees that if one wishes to take the relation previously treated as $+i$ for $+1$, then one must necessarily take the relation previously denoted by $-1$ for $+i$. In the language of mathematicians this means that $+i$ is a mean proportional between $+1$ and $-1$, or corresponds to the sign $\sqrt{-1}$. We intentionally do not say \emph{the} mean proportional, since $-i$ plainly has an equal claim. Here, then, the demonstrability of an intuitive meaning of $\sqrt{-1}$ is fully justified, and nothing more is needed in order to admit this quantity into the domain of the objects of arithmetic.

We have believed that we would serve the friends of mathematics by this brief presentation of the chief points of a new theory of the so-called imaginary quantities. If this subject has until now been considered from a false point of view and a mysterious darkness found in it, this is largely due to the ill-suited names. Had $+1,-1,\sqrt{-1}$ not been called positive, negative, imaginary - or even impossible - units, but perhaps direct, inverse, and lateral units, there would hardly have been any question of such darkness. The author has reserved for himself the fuller future treatment of the subject, which in the present memoir is touched only incidentally. There too the question will find its answer why the relations among things that present a manifold of more than two dimensions cannot furnish still other kinds of quantities admissible in general arithmetic.


\clearpage

\begin{center}
{\Large C. F. Gauss}\par
{\large Notices of Writings Not by Gauss}\par
{\normalsize Gauss, Werke, Volume II, printed pp. 179--196}\par
\end{center}

\section*{Notices of Writings Not by Gauss}

\begin{center}
\textsc{Göttingische gelehrte Anzeigen. 11 March 1809.}
\end{center}

\subsection*{\textit{Researches on the Arithmetic and Geometric Irreducibility of Numbers and Their Powers}}

\begin{center}
\textit{Recherches sur l'irréductibilité Arithmétique et Géométrique des nombres et de leurs puissances.} 1808. (No place of printing. 25 pages, large quarto.)
\end{center}

A work whose purpose is to represent irrational root quantities in the form of rational quantities. We must content ourselves with drawing the attention of friends of mathematics to this little work, since the limits of these pages do not allow us here to enter more extensively into the presentation and examination of the author's peculiar point of view and of his treatment of root quantities, which departs entirely from the customary one.

\begin{center}
\textsc{Göttingische Gelehrte Anzeigen. 23 March 1812.}
\end{center}

\subsection*{\textit{Cribrum Arithmeticum}}

\begin{center}
\textit{Cribrum Arithmeticum, or a table containing the prime numbers separated from the composite numbers occurring in the series of numbers proceeding from unity up to one million, and beyond this to twenty thousand (1,020,000). To composite numbers not divisible by 2, 3, or 5 are appended the simple divisors, not only the least, but all of them. Prepared by Ladislaus Chernac, Pannonian, A. L. M., Doctor of Philosophy and Medicine, Professor of Philosophy in the Lyceum of Deventer.} Deventer 1811. (At the author's expense, printed by J. H. Lange. xxii and 1022 pages, large quarto.)
\end{center}

The full title of this important and very meritorious work already indicates its contents sufficiently. It is a table, computed by an equally careful and laborious effort over several years, of all simple factors of all numbers from 1 to 1,020,000 that are not divisible by 2, 3, or 5; it is neatly printed and, as far as we have found by occasional checking, very correctly printed. Everyone who has much to do with calculations involving larger numbers will easily judge how valuable such a gift to arithmetic is. The author deserves double thanks: first for the highly laborious work itself, by which he has joined his name to the unforgettable names of Rhaeticus, Pitiscus, Brigg, Vlacq, Wolfram, Taylor, and others; and secondly for the undoubtedly very considerable expense incurred for printing, for which otherwise a publisher could scarcely have been found. Tables of this kind have often been calculated before, though mostly on a smaller scale, but they have either remained entirely in manuscript or have not been completed in print. Lambert, as is well known, formerly encouraged as strongly as he could the continuation of Pell's table, which went to 100,000 and was often printed; and according to information given by Bernoulli in Lambert's correspondence, Oberreit had carried it to 500,000, a copy of which had come into Schulze's hands. Anton Felkel, as reported in the \textit{Monatl. Correspondenz}, vol. 2, p. 223, had completed it in manuscript to two million, and later wished to carry it to 2,460,000; but what had already been printed in Vienna at public expense was, because no buyers could be found, used as cartridge paper in the Turkish war. Thus a meritorious work of many years was lost to the public. For that reason we considered it all the more a duty to announce the appearance of the present work here. The first million is now available for everyone's use; whoever has opportunity and zeal for this subject may therefore direct his labor to the further continuation.

\begin{center}
\textsc{Göttingische gelehrte Anzeigen. 3 November 1814.}
\end{center}

\subsection*{Burckhardt, \textit{Tables of Divisors}, second million}

\begin{center}
\textit{Tables of divisors for all numbers of the second million, or more exactly from 1,020,000 to 2,028,000, with the prime numbers contained therein.} By J. Ch. Burckhardt, member of the Imperial Institute, of the Bureau des Longitudes of France, and of several other learned societies. Paris, 1814. Mme Ve Courcier. (viii and 112 folio pages.)
\end{center}

Earlier than we had dared to hope when announcing Chernac's factor table for the first million, we can already report the completion and appearance of a similar table for the second million. The deserving author, whose name already guarantees the greatest care and accuracy, has placed all friends of arithmetic under great obligation by this laborious work. Chernac's table for the first million gives all simple factors; Burckhardt's table for the second gives only the smallest divisor each time. The complete factorization of a number of the second million therefore requires division by the smallest divisor and looking up the quotient in Chernac's table. But this small effort is of no significance in comparison with the great advantage of possessing the table in so much smaller a space, while the prospect remains that in time the table may be extended even to ten million. The author has made the compression into the small volume possible partly by restricting the entries to the smallest divisor, and partly by the most economical printing practicable. If $a$ denotes indeterminately each of the eighty numbers below 300 that are not divisible by 2, 3, or 5, then every number not divisible by 2, 3, or 5 is contained in the form $300n+a$. All eighty numbers for which $n$ has the same value are found in one vertical column, and each page contains thirty such columns. Thus each page includes, among nine thousand numbers progressing in natural order, all those not divisible by 2, 3, or 5.

The method by which Mr. Burckhardt constructed his table deserves special mention here. He had a grid engraved on copper, in which 81 horizontal and 78 vertical lines formed a rectangle divided into $80\times77$, that is, 6160 small squares, and had the required number of impressions made from it. At the side the eighty values of $a$ could be engraved at once; the values of $300n$ in consecutive order were written with the pen above the 77 vertical columns. Thus each sheet represents all numbers not divisible by 2, 3, or 5 among every 23,100 consecutive numbers in natural order, and 44 sheets suffice to cover an entire million. It is easy to see that the numbers whose smallest divisor is 7 or 11 recur on each following sheet in the same order, so that these divisors could be engraved at once on the copper plate and therefore appeared automatically in the proper places on every sheet. To enter the next divisors, for example 13, Mr. B. took from an extra sheet only 13 columns across its width and, treating this as the beginning of his table, cut out all the squares that had to contain the divisor 13. He therefore only had to lay this lattice on the first thirteen columns of the first sheet, then on the next thirteen, and so on, in order immediately to see all places which, insofar as they did not already contain 7 or 11, had to be filled with 13. The same procedure was then followed with divisor 17, and so forth. Up to divisor 73 the extra sheets sufficed in this way; for the larger divisors 79, 83, and so on, Mr. B. seems to have assembled the frame from two or more parts. For divisors above 500, however, Mr. B. preferred to find the multiples by addition, and for the other factor he needed to take only the prime numbers. We find this entire procedure highly appropriate, and would recommend it for imitation to all who might be inclined to continue the table further. For the third and fourth millions the author himself has meanwhile already carried out a large part of the calculations, so that we have well-founded hope that these too will soon be made known in print.

\begin{center}
\textsc{Göttingische gelehrte Anzeigen. 7 November 1816.}
\end{center}

\subsection*{Burckhardt, \textit{Tables of Divisors}, third million}

\begin{center}
\textit{Tables of divisors for all numbers of the third million, or more exactly from 2,028,000 to 3,036,000, with the prime numbers contained therein;} by J. Chr. Burckhardt, member of the Royal Academy of Sciences, of the Bureau des Longitudes of France, and of several other learned societies. Paris 1816. Mme Ve Courcier. (112 folio pages.)
\end{center}

Since, in the notice of the table for the factors of the second million, we already described in detail the method of computation used by the deserving author and the arrangement of the table itself, we may here content ourselves with the bare announcement of the appearance of the table for the third million. We are also shortly to expect from the author the table for the first million, presented in the same manner, so that the whole table to beyond three million will then form only a moderate volume. For this the author is owed the thanks of all friends of arithmetic, who through this laborious work see a need satisfied on a scale that far surpasses everything one could still have dared to hope for only a few years ago.

\begin{center}
\textsc{Göttingische gelehrte Anzeigen. 9 August 1817.}
\end{center}

\subsection*{Burckhardt, \textit{Tables of Divisors}, first million}

\begin{center}
\textit{Tables of divisors for all numbers of the first million, or more exactly from 1 to 1,020,000, with the prime numbers contained therein;} by J. Chr. Burckhardt, member of the Academy of Sciences in the Royal Institute, of the Bureau des Longitudes of France, and of several other learned societies. Paris 1817. Mme Ve Courcier. (114 folio pages.)
\end{center}

Referring here to the notices of the tables for the second and third millions, we now merely announce the actual appearance of these factor tables for the first million. We therefore now possess a connected whole for the first three millions. For the present first million the author used partly Chernac's \textit{Cribrum Arithmeticum}, partly a manuscript table by Schenmark owned by the library of the Royal Institute. The latter, however, had not been constructed with all the care one could wish, and the decision, in cases where the two differed, which of the two was correct was often rather laborious. In Chernac's table only a very small number of errors appeared, which Mr. Burckhardt has communicated here.

For the fourth, fifth, and sixth millions also the author already has the materials for the most part ready, and he offers to supply this continuation if the publisher is encouraged by sufficient sales of the first three millions. It would indeed be much to be regretted if the fruits of so laborious and useful a work were withheld from the world.

\begin{center}
\textsc{Göttingische gelehrte Anzeigen. 19 December 1825.}
\end{center}

\subsection*{Erchinger, Geometric Construction of the Regular Seventeen-gon}

A short treatise has been submitted to the Royal Society by Mr. Erchinger of Thuningen in the Kingdom of Württemberg, concerning the
\begin{center}
\textit{Geometric Construction of the Regular Seventeen-gon.}
\end{center}
It is well known that the general theory of regular polygons has taken on a new form and enlargement through the intimate connection into which it has been brought with higher arithmetic. One comparatively small part of this theory is the theory of those polygons that can be described geometrically. Since the age of the Greeks it had been known that the triangle, pentagon, pentadecagon, and all polygons arising from these by doubling or repeated doubling of the number of sides have that property; and it was believed, and sometimes expressly asserted, that these were the only ones. Higher arithmetic has taught that this was an error. By laying open the true sources of the completely general theory, it followed at once that, besides the polygons named, there are infinitely many others that can be constructed geometrically, of which the seventeen-gon is the simplest. The superiority of analysis, which grasps the most general and the particular with equal ease, over geometry, which must always remain with the particular and, in passing from simpler cases to more composite ones, is held back by ever-increasing complication and would scarcely ever have reached the familiar next case without outside help, appears here in the clearest light. Meanwhile it is always important, interesting, and desirable that purely geometric treatments also continue to be cultivated, and that geometry appropriate at least part of the new fields conquered by analysis. The reviewer is not aware that anyone has publicly treated the construction of the seventeen-gon before, except Mr. Pauker in the writings of the Courland Society and in his geometry. Different from that, and carried out more in a purely geometric spirit, is the construction of Mr. Erchinger, which is as follows. (The accompanying figure, a straight line on which the points $D B G A I F C E$ lie in succession, can be drawn by anyone for himself.) Let an arbitrarily assumed straight line $AB$ be prolonged backward and forward to $C$ and $D$ so that $AC\cdot BC=AD\cdot BD=4AB\cdot AB$; further, determine the points $E$, $G$ on the two sides of the prolonged line $CA$ so that $AE\cdot EC=AG\cdot CG=AB\cdot AB$, and the point $F$ on the side $A$ of the prolonged line $BA$ so that $AF\cdot DF=AB\cdot AB$; finally divide $AE$ at $I$ so that $AI\cdot EI=AB\cdot AF$, where $AI$ is the smaller and $EI$ the larger segment of $AE$. Then make a triangle in which two sides are each $=AB$, and the third $=AI$. If a circle is described about this triangle, then $AI$ will be the side of the regular seventeen-gon inscribed in the circle.

If one checks the correctness of this construction by comparison with the theory of the seventeen-gon set out as an example in Article 354 of the \textit{Disquisitiones Arithmeticae}, one easily notices that it is nothing other than the geometric translation of those equations to which the application of the general theory leads. Indeed, the distances of the points $C,D,E,F,G,I$ from $A$ are precisely the quantities there denoted by $(8.1)$, $(8.3)$, $(4.1)$, $(4.3)$, $(4.9)$, $(2.1)$, if the positive and negative signs are expressed by the position and the distance of point $B$ from $A$, taken in the same sense, is set equal to $-1$. But the truly meritorious part of Mr. Erchinger's treatise consists not so much in setting up the construction itself, since analysis had already indicated the simplest path, as in the purely geometric grounding of its correctness. This is carried out with such exemplary, painstaking care to avoid everything not purely elementary that it does the author honor and gives rise to the wish that his genuinely uncommon mathematical talent may find every encouragement.

\begin{center}
\textsc{Göttingische gelehrte Anzeigen. 9 July 1831.}
\end{center}

\subsection*{Seeber, Investigations on Positive Ternary Quadratic Forms}

\begin{center}
\textit{Investigations on the Properties of Positive Ternary Quadratic Forms} by Ludwig August Seeber, Doctor of Philosophy, ordinary Professor of Physics at the University of Freiburg. Freiburg im Breisgau 1831. (248 pages in quarto.)
\end{center}

Functions of two indeterminate quantities $x$ and $y$ of the form $axx+2bxy+cyy$, where $a,b,c$ denote fixed integers, form, as is well known, under the name of quadratic forms, or, where a further distinction is needed, binary quadratic forms, one of the most interesting and richest subjects of higher arithmetic. The first problems that arise here are: to decide whether such a given form contains another $a'x'x'+2b'x'y'+c'y'y'$ under it, that is, whether it can be transformed into it by a substitution $x=\alpha x'+\beta y'$, $y=\gamma x'+\delta y'$, in which $\alpha,\beta,\gamma,\delta$ are integers; whether such a relation of two forms is reciprocal, in which case the forms are called equivalent; further, in both cases, to give all possible transformations of the one into the other; and finally to find all possible representations of a given integer by a given form through integer values of the indeterminate quantities. These problems are completely solved in the \textit{Disquisitiones Arithmeticae}, but they constitute by far the smaller part of the section of that work devoted to quadratic forms. The more refined investigations that follow required in part a preliminary treatment of a field one step higher and offering far greater difficulties, namely the doctrine of similar functions of three indeterminate quantities $x,y,z$, which therefore have the form $axx+byy+czz+2a'yz+2b'xz+2c'xy$ and are called ternary quadratic forms. The solution of the main problems concerning these ternary forms is developed in the work mentioned, but only as far as was necessary for the indicated purpose. After an interval of thirty years, the author of the present work has now been the first to take up these investigations again, and with respect to one principal class of ternary forms, namely the positive ones, has brought to completion what had been left unfinished in the \textit{Disquisitiones Arithmeticae}. For those who have made a deeper study of higher arithmetic, we could indicate in a few words what is newly accomplished in the present work. But in order to be intelligible also to others, we must allow ourselves somewhat greater detail; and we do this all the more readily because these investigations also have a distinctive interest outside the domain of higher arithmetic.

The properties of a binary form $axx+2bxy+cyy$ depend principally on the number $bb-ac$, which is therefore called the determinant of that form. Two equivalent forms always have equal determinants. But not all forms with a given determinant are for that reason already equivalent. Rather, such forms split into a smaller or larger, but always finite, number of classes, so that those belonging to the same class are equivalent among themselves, while those belonging to different classes are not equivalent. Forms whose determinant is positive can represent positive and negative numbers without distinction; by contrast, forms with negative determinant can represent only such numbers as have the same sign as $a$ and $c$, and here positive and negative forms are therefore distinguished. The simplest forms in each class satisfy definite criteria, are called reduced forms, and may be regarded as representatives of the whole class.

Analogous relations concerning ternary forms are demonstrated in the \textit{Disquisitiones Arithmeticae}. The determinant of the ternary form
\[
axx+byy+czz+2a'yz+2b'xz+2c'xy
\]
is the number
\[
aa'a'+bb'b'+cc'c'-abc-2a'b'c'.
\]
Here too equality of determinants is necessary for the equivalence of two forms, but not sufficient. All forms with a fixed determinant split into a finite number of classes, in each of which the simplest forms may be called reduced and may represent all the others. But the distinction between positive and negative forms is different here from the case of binary forms. For every given determinant, whether positive or negative, there are in part forms by which positive and negative numbers can be represented without distinction (indifferent forms), and in part forms by which either only positive or only negative numbers can be represented (positive or negative forms). Yet positive forms exist only for negative determinants, and negative ones only for positive determinants. Moreover, it is clear of itself that the qualification of a form as indifferent, positive, or negative belongs at the same time to the whole class to which it belongs. The present work is restricted to the positive forms, whose determinants must therefore be negative. Evidently everything that holds for these transfers at once to the negative forms, whereas the indifferent forms, entirely excluded from the work, require a completely different treatment.

In the \textit{Disquisitiones Arithmeticae}, as already mentioned, the theory of ternary forms had been developed only as far as necessary for the purpose there, and therefore the problem of deciding the equivalence of two given ternary forms had not yet been solved in full generality. It had indeed been shown there how to find for any prescribed form an equivalent one of the simplest kind, and that for every given determinant there can be only a finite number of such reduced forms. But since in each class there are several such reduced forms, which do not in every case immediately appear to be equivalent, there was still lacking a criterion by which the equivalence or non-equivalence of such forms could be recognized with certainty. The author of the present work has now removed this need completely and with exemplary thoroughness as regards positive forms. His procedure is clothed somewhat differently from the way we have just expressed the matter, and from the way it would have to stand if, in the concept of reduced positive forms, one admitted only the most essential conditions of greatest simplicity. In the case of positive forms these conditions are that the numbers $a,b,c$, positive by their nature, must not be smaller than, respectively, $b'$ or $c'$, $a'$ or $c'$, and $a'$ or $b'$, without regard to signs. Mr. Seeber has added to the concept of reduced forms further modifications such that in every class there can be only one of that kind, and there must be one. Because of a fine theorem found by Mr. Seeber by induction, to be mentioned below, we state here the principal conditions which Mr. S. has included in the concept of reduced forms. They are: 1) among the numbers $a',b',c'$, no two may have opposite signs; 2) without regard to sign, $2b'$ and $2c'$ may not be greater than $a$, furthermore $a'$ and $2a'$ may not be greater than $b$, and $b'$ not greater than $c$; 3) in the case where $a',b',c'$ are all negative, twice the sum of these numbers may not be greater than $a+b$. We may pass over the other modifications that are added for some special cases.

The main content of the work first consists in the solution of the problem of finding, for every given positive form, an equivalent form which has the character of a reduced form according to the established definition; and then in the strict proof of the theorem that two non-identical reduced forms cannot be equivalent, or, what is the same, that there is only one reduced form in each class. We must do full justice to the spirit of thoroughness with which these subjects are carried out; and if we must at the same time regret that a very great, and perhaps to some discouraging, expansiveness is connected with it, since the solution of the problem takes 41 pages and the proof of the theorem 91 pages, we do not wish this to be regarded in any way as a reproach. When a difficult problem or theorem is to be solved or proved, the first step, always to be recognized with due thanks, is that a solution or proof be found at all; and the question whether this could not have been done in an easier and simpler way remains idle so long as the possibility is not at the same time decided by the deed. We therefore consider it untimely to dwell here on this question. The remaining part of the work chiefly contains the solutions, carried out with the same thoroughness, of the problems of deciding whether a given form contains under it another given non-equivalent form; of finding all possible transformations of a given form into a given equivalent form, or into one merely contained under it; and finally of giving, for a given determinant, all possible classes of positive ternary forms.

We must also note that Mr. Seeber has taken the form of ternary forms somewhat differently from the way it was done in the \textit{Disquisitiones Arithmeticae}, where, deliberately, the coefficients of the products $yz,xz,xy$ were assumed to be even numbers, whereas Mr. S. also allows odd ones, and therefore denotes by $a',b',c'$ what above was denoted by $2a',2b',2c'$. Evidently the greater generality thereby attained is only apparent, or at least superfluous, since everything that can be said of such forms with odd coefficients also follows of itself if one considers their doubles instead. We therefore cannot approve this alteration, which moreover entails some loss of simplicity. One consequence has been that what Mr. Seeber calls determinant is always four times the number that bears this name in the \textit{Disquisitiones Arithmeticae}. In the present notice we have retained the terminology of the \textit{Disquisitiones Arithmeticae}.

In the last-mentioned problem, that of giving all possible reduced forms for each given determinant, Mr. Seeber used, in order to have bounds for the first three coefficients, a theorem according to which their product $abc$ cannot be greater than three times the determinant. This theorem is rigorously proved by Mr. Seeber. But in the preface he remarks that among more than 600 cases investigated by him he had not found a single one in which that product exceeded twice the determinant, and he therefore regards it as highly probable that this narrower bound is generally valid; however, he had not succeeded in finding a strict proof. Since this theorem, found by Mr. Seeber by induction, is noteworthy in itself and important for shortening the solution of the problem mentioned, we wish here, in order to contribute on our side in this notice to the completion of this theory, to add a very simple proof. Two cases must be distinguished.

I. If none of the numbers $a',b',c'$ is negative, set
\[
\begin{array}{lll}
b-2a'=d, & c-2b'=e, & a-2c'=f,\\
c-2a'=g, & a-2b'=h, & b-2c'=i.
\end{array}
\]
From the definition of reduced positive forms it follows immediately that, if
\[
axx+byy+czz+2a'yz+2b'xz+2c'xy
\]
is such a form, none of these six numbers is negative; it is also self-evident that $a,b,c$ are positive. If the negative determinant of the form is denoted by $-D$, then, as one is easily convinced by expansion, one has the identity
\[
2D-abc=aa'd+bb'e+cc'f+a'hi+b'gi+c'gh+ghi,
\]
in which none of the seven terms on the right can be negative, and consequently $abc$ is not greater than $2D$. The same follows in the same way from the identity
\[
2D-abc=aa'g+bb'h+cc'i+a'ef+b'df+c'de+def.
\]

II. If none of the numbers $a',b',c'$ is positive, set
\[
\begin{array}{lll}
b+2a'=d, & c+2b'=e, & a+2c'=f,\\
c+2a'=g, & a+2b'=h, & b+2c'=i,\\[2pt]
b+c+2a'+2b'+2c'=k,\\
a+c+2a'+2b'+2c'=l,\\
a+b+2a'+2b'+2c'=m.
\end{array}
\]
Again let the determinant of the form be $-D$. By the definition of reduced positive forms none of the nine numbers $d,e,f,g,h,i,k,l,m$ can be negative. Thus the identity
\[
6D-3abc=-aa'(d+2k)-bb'(e+2l)-cc'(f+2m)-a'hi-b'gi-c'gh+def+2ghi
\]
gives, since $a',b',c'$ are not positive but negative or zero, that all terms on the right are positive or zero; hence $3abc$ cannot be greater than $6D$, or $abc$ cannot be greater than $2D$. The same follows likewise from the identity
\[
6D-3abc=-aa'(g+2k)-bb'(h+2l)-cc'(i+2m)-a'ef-b'df-c'de+2def+ghi.
\]
Both equations are symmetric. If one gives up complete symmetry, the proof may be carried out with an even smaller number of terms, for example by the identity
\[
8D-4abc=-2aa'(g+k)-2bb'(e+l)-4cc'm+(c+e)df+(c+g)hi.
\]

We now add something about the meaning of positive binary and ternary quadratic forms outside the domain of higher arithmetic. It is unnecessary to treat the negative ones separately, and the indifferent ones withdraw entirely from this treatment.

The positive binary form $axx+2bxy+cyy$ generally represents the square of the distance between two indeterminate points in a plane, whose coordinates, with respect to two axes inclined to one another at an angle whose cosine is $\frac{b}{\sqrt{ac}}$, differ by $x\sqrt a$ and $y\sqrt c$. Insofar as $x$ and $y$ are to denote only integers, the form therefore refers to a system of parallelogrammatically arranged points lying at the intersections of two systems of parallel lines. The lines of each system are at equal distances from one another; those of one system, when measured parallel to the lines of the second, are $=\sqrt a$, and the distances of the other, measured parallel to the lines of the first, are $=\sqrt c$; the inclination of the two systems to one another is the one given above. In this way the plane appears divided into equal parallelograms, whose vertices form the point system, without any one of the points falling inside a parallelogram. The determinant taken with positive sign, thus $ac-bb$, denotes the square of the area of an elementary parallelogram. One and the same system of such points can be divided parallelogrammatically in infinitely many different ways, and thus referred back to just as many different forms. But all these different forms are what the technical language calls equivalent, and the area of an elementary parallelogram remains always the same. Two forms that are not equivalent, but one of which contains the other under it, refer to the same system of points, but the first form refers to the whole system and the second to a part. Two forms that, in the technical language, are called improperly equivalent refer to two equal but oppositely situated systems of points, as if the plane were thought of as turned over, and so on.

In the same way the positive ternary form
\[
axx+byy+czz+2a'yz+2b'xz+2c'xy
\]
generally denotes the square of the distance between two indeterminate points in space, whose coordinates with respect to three axes $(1),(2),(3)$ have the differences $x\sqrt a,y\sqrt b,z\sqrt c$. The cosines of the angles between axes $(2)$ and $(3)$, $(1)$ and $(3)$, $(1)$ and $(2)$ are respectively $\frac{a'}{\sqrt{bc}}$, $\frac{b'}{\sqrt{ac}}$, $\frac{c'}{\sqrt{ab}}$. Insofar as $x,y,z$ are to denote only integers, the form refers to a system of parallelepipedically arranged points, that is, points arising from the intersections of three systems of parallel equidistant planes. The whole of space thus appears divided into equal parallelepipeds whose vertices form that system of points, and the square of the volume of an elementary parallelepiped is equal to the determinant of the ternary form, taken with positive sign. Equivalent forms represent one and the same system of points, merely referred to different axes or fundamental planes. In the same way all the other main features of the theory of ternary forms find their geometric meaning here: the containment of one form under another, the representation of a definite number or of an indeterminate binary form by a ternary one, the doctrine of the adjoint ternary forms (\textit{formae adiunctae}), the disappearance of the distinction between proper and improper equivalence, the nature of reduced forms, and so on. But we must restrict ourselves to the indications above, especially since the present work, which considers ternary forms solely from the purely arithmetical point of view, has given only indirect occasion for them. One will at least recognize from this what a rich field is opened here for investigations which not only have high theoretical interest in themselves but can also be used for an equally convenient and general treatment of all relations among crystal forms. This is not the place to go into the details of that use. Yet we must not omit the remark that, although it is originally assumed that $a,b,c,a',b',c'$ denote integers, the greater part of the theory of ternary forms, and in particular what is required for that use, remains valid independently of that assumption. Indeed, Haüy's data for most crystal classes lead to very simple integer values of the coefficients in the ternary forms corresponding to the associated arrangement of the point system. But the more exact later measurements of Wollaston, Malus, Biot, Kupffer, and others are in contradiction with this, and make it doubtful whether rational ratios for those coefficients are everywhere natural. In any case, if one does not wish to omit from the theory the restriction to integer values of the coefficients, since what matters here is not absolute values but only their ratios to one another, integers can always be found that come as close to the measurement results as one wishes.

Finally, we give the geometric meaning of Seeber's theorem cited above. If a parallelepiped is so constituted that none of its twelve edges, among which four at a time are equal to one another, is greater either than any of the twelve diagonals of faces, which are pairwise equal, or than any of the four diagonals of the parallelepiped, then the volume of it multiplied by $\sqrt2$ is not smaller than the volume of a rectangular parallelepiped formed from the same edges.



\clearpage

\section*{Manuscript Nachlass}

\begin{center}
{\large Solution of the congruence $x^m-1 \equiv 0$}\par
{\large Analysis of residues. Chapter six. First part.}
\end{center}

\subsection*{237.}
In Chapter III we showed that the congruence $x^n\equiv 1$, when the modulus is a prime number $p$, has $\mu$ roots when $\mu$ is the greatest common divisor of $n$ and $p-1$, and that these roots coincide exactly with the roots of the congruence $x^\mu\equiv 1$. It is therefore enough to consider the case in which $n$ is a divisor of $p-1$. Moreover, that the solution not only of the congruence $x^n-1\equiv 0$ but of any other congruence, for arbitrary moduli, can be derived from the solution for prime moduli has already been shown in various places and will be taught more fully below, in Chapter VIII.

\subsection*{238.}
Nor is it necessary to stop here. In the same Chapter III we showed that the solution of $x^n\equiv 1$ depends on the resolution of similar congruences $x^a\equiv 1$, $x^b\equiv 1$, etc., where $a$, $b$, etc. are primes or powers of primes and $n$ is their product. Namely, if $A$, $B$, etc. are arbitrary respective roots of $x^a\equiv 1$, $x^b\equiv 1$, etc., then the product $AB\ldots$ will be one of the roots of $x^n\equiv 1$. Our investigations will therefore be restricted to the solution of $x^n\equiv 1\pmod p$, where $p$ is prime, $n$ is a prime or a power of a prime, and at the same time $n$ is a divisor of $p-1$.

\subsection*{239.}
It follows further from Chapter III that among the roots of $x^n\equiv 1$ there are always some through whose powers all the others can be exhibited. Thus, if $r$ denotes such a root (we called it \textit{primitive} above when $n=p-1$, and we shall retain this expression here, although in a broader sense), all roots of the proposed congruence will be
\[
1,\quad r,\quad rr,\quad r^3,\quad \ldots,\quad r^{n-1}.
\]
Such roots must therefore be sought with particular care, since once they have been found the remaining roots become evident of themselves. For brevity we denote any power of $r$ by enclosing its exponent in parentheses, so that $(0)$ denotes unity, $(1)$ denotes any primitive root of $x^n\equiv 1$, $(2)$ the square of $(1)$, and so on. Thus the series $(0)$, $(1)$, $(2)$, $(3)$, $\ldots$, $(n-1)$ contains all the roots. It is known, moreover, that $(k)$ will always be such a primitive root whenever $k$ is prime to $n$; that is, in our case, where $n$ is a power $t^\nu$ of a prime $t$, whenever $t$ does not divide $k$. The signs $(1)$, $(2)$, etc. are plainly indeterminate in themselves; but once a determinate value has been assigned to $(1)$, all the others become determined.

\subsection*{240.}
Since the aim is to investigate primitive roots before the others, they must first be separated from the rest. This is done by removing from the series $(0)$, $(1)$, $(2)\ldots(n-1)$ all terms $(k)$ for which $k$ is divisible by $t$; but if $n$ is prime, or $\nu=1$, only $(0)$ is removed. Before proceeding to the investigation of the remaining roots, we urge the reader to make several examples, so that everything which without examples might perhaps appear too general may be seen concretely. We add one example, but it will not on that account be superfluous to work out others independently.

Let $p=29$ and $n=7$. The seven roots of $x^7\equiv 1\pmod {29}$ are $1$, $7$, $16$, $20$, $23$, $24$, $25$. Since $n$ is prime, all these roots except $1$ are primitive. If therefore $7=(1)$, the signs have the following meanings:
\[
\begin{array}{ccccccc}
(0)&(1)&(2)&(3)&(4)&(5)&(6)\\
1&7&20&24&23&16&25.
\end{array}
\]
The reader will remember that the signs $(n)$ and $(0)$, $(n+1)$ and $(1)$, and in general $(a)$ and $(b)$, are equivalent whenever $a\equiv b\pmod n$.

\subsection*{241.}
For our purpose, however, we must proceed in another way as well. We retain only those terms $(k)$ for which $k$ is not divisible by $t$; their number is $\frac{t-1}{t}\,n=\lambda$. All these numbers, or numbers congruent to them modulo $n$, can be exhibited by the successive powers of some number. Let this number be $\rho$. Then all the primitive roots of $x^n\equiv 1$ are denoted by
\[
(1),\quad (\rho),\quad (\rho^2),\quad (\rho^3),\quad \ldots,\quad (\rho^{\lambda-1}).
\]
By this device we obtain the complete exclusion of the non-primitive roots; the reasons and advantages of this will become clearer below. In our example we may put $\rho=3$, and the primitive roots of $x^7\equiv 1$ are arranged as
\[
\begin{array}{cccccc}
(1)&(3)&(3^2)&(3^3)&(3^4)&(3^5)\\
\text{that is }(1)&(3)&(2)&(6)&(4)&(5)\\
7&24&20&25&23&16.
\end{array}
\]

\subsection*{242.}
So that the reader may know where the following investigations are headed, it will be useful to state the theorem which we propose to prove and explain.

\textit{If the number $\lambda$, which is $t^{\nu-1}(t-1)$, has the distinct prime factors $a,b,c,d$, etc., and if $\lambda=a^\alpha b^\beta c^\gamma\ldots$, then the resolution of the congruence $x^n-1\equiv 0$ depends on the resolution of $\alpha+\beta+\cdots$ lower congruences, of which $\alpha$ are of degree $a$, $\beta$ of degree $b$, $\gamma$ of degree $c$, etc.}

Thus in our example the resolution of $x^7\equiv 1$ depends on one congruence of the second degree and another of the third degree; and in general it is clear that the degree of these congruences never depends on the modulus $p$. In order to arrive at the proof of this theorem, we must first put forward some propositions concerning the connection between congruences and their roots, although these investigations properly belong to Chapter VIII, where they will be pursued further.

\subsection*{243.}
\textsc{Theorem.} \textit{Suppose that the congruence}
\[
x^m+A x^{m-1}+B x^{m-2}+\cdots+N\equiv 0\pmod{\text{prime}}
\]
\textit{is so constituted that, when the product of the $m$ factors $x-r$, $x-r'$, $x-r''$, $x-r'''\ldots$ is formed, and is $x^m+a x^{m-1}+b x^{m-2}+\cdots+n$, one has $A\equiv a$, $B\equiv b$, $C\equiv c$, etc. modulo $p$. Then the quantities $r$, $r'$, $r''\ldots$ are roots of the proposed congruence, and it has no others.}

\textit{Proof.} I. We always have
\[
x^m+A x^{m-1}+B x^{m-2}+\cdots \equiv x^m+a x^{m-1}+b x^{m-2}+\cdots \pmod p.
\]
The latter expression becomes $0$ on putting $x=r$, $x=r'$, $x=r''$, etc.; hence, for these values of $x$, the former expression becomes $0\pmod p$. This proves the first assertion.

II. If some further value $\varrho$, congruent to none of $r$, $r'$, etc., satisfied the proposed congruence, then
\[
\begin{aligned}
0&\equiv \varrho^m+A\varrho^{m-1}+B\varrho^{m-2}+\cdots \\
&\equiv \varrho^m+a\varrho^{m-1}+b\varrho^{m-2}+\cdots \\
&\equiv (\varrho-r)(\varrho-r')(\varrho-r'')(\varrho-r''')\cdots .
\end{aligned}
\]
But none of the factors $\varrho-r$, $\varrho-r'$, $\varrho-r''$, etc. is $0$ modulo $p$; hence, since $p$ is prime, it is absurd that their product should be $0$. Therefore there are no roots besides $r$, $r'$, etc. This proves the second assertion.

\subsection*{244.}
\textsc{Problem.} \textit{Let $r$, $r'$, $r''\ldots$ be unknown quantities whose number is $m$, whose sum is $\alpha$, whose sum of squares is $\beta$, whose sum of cubes is $\gamma$, and so on, and whose sum of $m$-th powers is $\mu$. Suppose that these numbers themselves, again $m$ in number, are not given, but rather other numbers $\alpha'$, $\beta'$, $\gamma'$, etc. respectively congruent to them modulo a prime $p>m$. Find the congruence of the $m$-th degree whose roots are $r$, $r'$, $r''$, etc.}

\textit{Solution.} Regard $r$, $r'$, $r''$, etc. as the roots of an equation
\[
x^m+A x^{m-1}+B x^{m-2}+C x^{m-3}+\cdots=0
\]
and determine its coefficients $A$, $B$, $C$, etc., using congruence in place of equality, by the known method, namely by setting
\[
\begin{aligned}
-A&\equiv \alpha',\\
-2B&\equiv \beta' + A\alpha',\\
-3C&\equiv \gamma' + A\beta' + B\alpha',\\
-4D&\equiv \delta' + A\gamma' + B\beta' + C\alpha',\\
&\quad\text{etc.}\\
-mN&\equiv \mu' + A\lambda' + \text{etc.}
\end{aligned}
\]
These coefficients cannot be indeterminate, since all the numbers $1,2,3\ldots m$ are less than $p$. I say that the required congruence is
\[
x^m+A x^{m-1}+B x^{m-2}+\cdots+N\equiv 0.
\]

\textit{Proof.} Suppose that the equation whose roots are $r$, $r'$, $r''$, $r'''$, etc. is
\[
x^m+a x^{m-1}+b x^{m-2}+\cdots=0.
\]
Then
\[
\begin{aligned}
-a&=\alpha,\\
-2b&=\beta+a\alpha,\\
-3c&=\gamma+a\beta+b\alpha,\\
-4d&=\delta+a\gamma+b\beta+c\alpha,\\
&\quad\text{etc.}
\end{aligned}
\]
It is therefore evident that
\[
a\equiv A,
\qquad b\equiv B,
\qquad c\equiv C\quad \text{etc.}\pmod p.
\]
Thus, by the preceding section, the numbers $r$, $r'$, $r''$, etc., which are roots of
\[
x^m+a x^{m-1}+b x^{m-2}+\cdots=0,
\]
are also roots of the congruence
\[
x^m+A x^{m-1}+B x^{m-2}+\cdots\equiv 0. \qquad \text{Q. E. D.}
\]
We leave examples for the readers to construct.

\subsection*{245.}
We return to our purpose. Keeping the notation of §§242 and the preceding sections, we set out to show that, if $\lambda$ is the product of any factors $efg$, etc., the primitive roots of $x^n\equiv 1$, whose number is $\lambda$, can be divided into $e$ classes, so that the sums of the roots put into the same class are given by a congruence of degree $e$; and, once these are supposed known, any class can be subdivided into $f$ orders, so that the sums of the terms in each order are given by a congruence of degree $f$. These orders can then again be subdivided, and so on, until one arrives at the individual roots.

\subsection*{246.}
\textit{Definition.} We call the aggregate of all terms contained in a form $(\rho^{ke+\alpha})$ (§241) a \textit{complete period}, or simply a \textit{period}. Here $e$ denotes a divisor of $\lambda$, $\alpha$ any given number, and $k$ all integers from $0$ up to $\frac{\lambda}{e}-1$. For brevity we denote such a period by $(e*\alpha)$. Thus in our example the terms
\[
\begin{array}{rcl}
(1),(2),(4)&\text{constitute the periods}&(2*0),\\
(3),(6),(5)&& (2*1),\\[2pt]
\text{while }(1),(6)&\text{constitute}&(3*0),\\
(3),(4)&& (3*1),\\
(2),(5)&& (3*2).
\end{array}
\]
If all the terms are distributed in any way into periods, and each period is again distributed into smaller periods, and so on, we say that what was promised in the preceding section has been obtained.

Before entering upon this exposition, however, we shall show that although the formation of such a period depends in a certain way on two arbitrary quantities $r$ and $\rho$, it nevertheless contains nothing vague: however these quantities are chosen, the same terms always come together into the same period, provided that the number of terms which the period is to contain has been prescribed.

The criterion for two terms $A$ and $B$ to be in the same period is that both are contained in such a form $(\rho^{ke+\alpha})$, that is,
\[
A\equiv r^{\rho^{ke+\alpha}},\qquad B\equiv r^{\rho^{le+\alpha}}\pmod p.
\]
Here $r$ is a primitive root of $x^n\equiv 1\pmod p$, while $\rho$ is a primitive root of $x^\lambda\equiv 1\pmod n$, as above.

We must prove that if, in place of $r$ and $\rho$, other numbers are chosen, say $s$ and $\sigma$, then $A$ and $B$ are contained in similar forms $s^{\sigma^{ke+\theta}}$ and $s^{\sigma^{le+\theta}}$.

Let $s^m\equiv r\pmod p$, $\sigma^\mu\equiv \rho\pmod n$, and $m\equiv \sigma^\zeta\pmod n$. This can be done because $r$ and $\rho$ are primitive roots. Moreover $m$ is prime to $n$, and $\mu$ to $\lambda$ (Chapter III). By the appropriate substitutions we obtain
\[
A\equiv s^{\sigma^{ke+\mu+\zeta}},\qquad B\equiv s^{\sigma^{le+\mu+\zeta}}. \qquad \text{Q. E. D.}
\]

\subsection*{247.}
\textsc{Theorem.} \textit{The product of two similar periods can be composed, independently of the number $p$, by the addition of similar periods and given numbers.}

We call periods similar when they contain the same number of terms, or equivalently when the number $e$ is the same.

\textit{Example.} Let $n=7$. The product of the periods $(1)+(6)$ and $(2)+(5)$ is, since $(a)\times(b)=(a+b)$, equal to $(3)+(6)+(8)+(11)$, that is, it consists of the periods $(3)+(4)$ and $(1)+(6)$.

\textit{Proof.} Let $\frac{\lambda}{e}=f$, and let the given periods $(e*\alpha)$ and $(e*\delta)$, or sums, be
\[
\begin{aligned}
(\rho^\alpha)+(\rho^{\alpha+e})+(\rho^{\alpha+2e})+\cdots+(\rho^{\alpha+(f-1)e})&=P,\\
(\rho^\delta)+(\rho^{\delta+e})+(\rho^{\delta+2e})+\cdots+(\rho^{\delta+(f-1)e})&=Q.
\end{aligned}
\]
The product $PQ$ consists of $f^2$ terms. These must be arranged as follows. Form $f$ series, each consisting of $f$ terms. The first contains the product of $P$ by $(\rho^\delta)$, the second the product $P(\rho^{\delta+e})$, and so on. In the first series the first place is occupied by the product arising from the part $(\rho^\alpha)$, the second by the product arising from $(\rho^{\alpha+e})$, and so on. In the second series the first place is assigned to the product arising from $(\rho^{\alpha+e})$, the second to the product arising from $(\rho^{\alpha+2e})$, and so on, the last to that arising from $(\rho^\alpha)$. The third begins with the product arising from $(\rho^{\alpha+2e})$, and so forth; after the product from the last part follow the products from the first and second parts, etc. Thus, if the successive parts of the period $P$ are denoted by $1,2,3\ldots z$ and those of the period $Q$ by I, II, III, $\ldots Z$, the parts of $PQ$ are
\[
\begin{array}{l}
1\cdot\mathrm{I}+2\cdot\mathrm{I}+3\cdot\mathrm{I}+4\cdot\mathrm{I}+\cdots+z\cdot\mathrm{I},\\
2\cdot\mathrm{II}+3\cdot\mathrm{II}+4\cdot\mathrm{II}+\cdots+1\cdot\mathrm{II},\\
3\cdot\mathrm{III}+4\cdot\mathrm{III}+\cdots+1\cdot\mathrm{III}+2\cdot\mathrm{III},\\
\text{etc.}
\end{array}
\]
Now collect into $f$ orders all terms occupying the same place in the several series. I say, first, that if one term is $\equiv 1$, then all the others of the same order are also $\equiv 1$; and second, that every order in which no term is $\equiv 1$ forms a period. Once these have been proved, the proposition is plainly obtained.

The general form of such an order is
\[
(\rho^{\alpha+ke}+\rho^\delta),\quad (\rho^{\alpha+(k+1)e}+\rho^{\delta+e}),\quad (\rho^{\alpha+(k+2)e}+\rho^{\delta+2e}),\ldots, (\rho^{\alpha+(k+f-1)e}+\rho^{\delta+(f-1)e}).
\]
For $\rho^{\alpha+(k-1)e}$ one may also write $\rho^{\alpha+(k+f-1)e}$, since $ef=\lambda$ and $\rho^\lambda\equiv 1\pmod n$, and similarly for the preceding terms. Put $\rho^{\alpha+ke}+\rho^\delta\equiv \rho^\chi\pmod n$, which is permitted unless perhaps $\rho^{\alpha+ke}+\rho^\delta$ is divisible by $n$.
\footnote{The proposition must be expressed somewhat differently if $n$ denotes, in general, a power of a prime; when $n$ is prime, nothing need be changed.}
Then the order can be exhibited as
\[
(\rho^\chi),\quad (\rho^{\chi+e}),\quad (\rho^{\chi+2e}),\ldots,(\rho^{\chi+(f-1)e}),
\]
which is plainly the period $(e*\chi)$. If, on the other hand, $\rho^{\alpha+ke}+\rho^\delta$ is divisible by $n$, all terms of the order are $\equiv(0)$, that is, $\equiv 1$. Q. E. D.

\textit{Remark.} This proof also gives a very easy method for expanding the product. We shall give another below which does not have this advantage, but which is not to be despised on account of its simplicity.

\subsection*{248.}
All the smaller periods which constitute a larger period are called a system of periods. Thus the periods
\[
(ef*\alpha),\quad (ef*f+\alpha),\quad (ef*2f+\alpha),\ldots,(ef*(e-1)f+\alpha)
\]
which compose the period $(f*\alpha)$ are designated by this name. A system is \textit{properly ordered} if the numbers placed after the sign $*$, here $\alpha$, $f+\alpha$, $2f+\alpha$, proceed in an arithmetic progression with difference $f$; and systems are \textit{similar} if both the smaller periods and the larger periods are similar.

\textsc{Theorem.} \textit{If the periods of two similar properly ordered systems are multiplied with one another, the first by the first, the second by the second, the third by the third, etc., the sum of all these products can be composed of periods similar to the larger period and of given numbers.}

\textit{Proof.} Let the systems be
\[
\begin{array}{cccc}
(ef*\alpha),&(ef*\alpha+f),&(ef*\alpha+2f),&\ldots,\\
(ef*\delta),&(ef*\delta+f),&(ef*\delta+2f),&\ldots.
\end{array}
\]
The products of each period of the first system with the corresponding periods of the second consist, by the preceding section, of integers and similar periods. A little attention to the genesis of these periods shows that if
\[
(ef*\alpha)\times(ef*\delta)
\]
consists of an integer $N$ and of the periods $(ef*A)$, $(ef*B)$, $(ef*C)$, etc., then the products
\[
\begin{aligned}
(ef*\alpha+f)\times(ef*\delta+f)&\text{ consist of }N\text{ and of }(ef*A+f),(ef*B+f),(ef*C+f),\text{ etc.},\\
(ef*\alpha+2f)\times(ef*\delta+2f)&\text{ consist of }N\text{ and of }(ef*A+2f),(ef*B+2f),(ef*C+2f),\text{ etc.},
\end{aligned}
\]
and in general
\[
(ef*\alpha+\mu f)\times(ef*\delta+\mu f)\text{ consists of }N\text{ and of }(ef*A+\mu f),(ef*B+\mu f),(ef*C+\mu f),\text{ etc.}
\]
It follows at once that the sum of all the products is
\[
eN+(f*A)+(f*B)+(f*C)\text{ etc.}\qquad\text{Q. E. D.}
\]
This proof also supplies a method for finding that sum.

\subsection*{249.}
It is easy to make this theorem still more general. If any number of similar properly ordered systems are given and one forms the products of all first periods, of all second periods, etc., then the sum of all these products consists of numbers and larger periods. If all these systems are taken equal, the sum of any powers whatever of all the periods will consist of numbers and periods similar to the larger one. It is now clear what this is aiming at. Let $\lambda=efgh\ldots$, and let all the first roots be split into $e$ periods $A$, $A'$, $A''$, etc.; let each of these again be split into $f$ periods $B$, $B'$, $B''$, etc.; let each of these be split into $g$ periods $C$, $C'$, $C''$, etc. The sum of all the periods is already given, namely it is $\equiv -1$. But, according to what we have just said, the following are also given:
\[
\begin{gathered}
A^2+(A')^2+(A'')^2+(A''')^2+\text{etc.},\\
A^3+(A')^3+(A'')^3+(A''')^3+\text{etc.},\\
\text{etc.}
\end{gathered}
\]
Hence, by §244, a congruence of degree $e$ can be found whose roots are $A$, $A'$, $A''$, etc. Once these are supposed known, any period is divided into smaller ones,
\[
\begin{array}{llll}
A&\text{into}& B,&B',\ B'',\ldots\\
A'&\text{into}& B^{(n)},&B^{(n+1)},\ B^{(n+2)},\ldots\\
&\text{etc.}
\end{array}
\]
Thus $B+B'+B''+\cdots\equiv A$ is known. But it is known that
\[
\begin{gathered}
B^2+(B')^2+(B'')^2+\cdots,\\
B^3+(B')^3+(B'')^3+\cdots,\\
\text{etc.}
\end{gathered}
\]
can be expressed from units and the periods $A$, $A'$, $A''$, etc. Therefore $B$, $B'$, $B''$, etc. are given by a congruence of degree $f$, from which they can be found. In the same way the periods of which $A'$, $A''$, etc. consist can be determined. From this anyone will see that by a completely similar method any period can be subdivided into smaller ones until the roots themselves are reached.

\subsection*{250.}
In applying these rules, however, a difficulty occurs which must be removed. Since every congruence has several roots, we must see which sign is to be assigned to each so that they can be properly distinguished from one another. Since the designation of periods depends on the numbers $r$ and $\rho$, which may be chosen arbitrarily, something arbitrary must necessarily cling to the designation as well. The number $\rho$ must of course be fixed at the beginning. The essential character of our method consists chiefly in deriving smaller periods from larger ones. But this cannot be done without the proper order of the periods, which we have achieved by the signs. We must therefore strive to distinguish all periods by their signs as soon as they have been found.

Let the period $A$ be designated by $(e*\alpha)$ and split into $f$ periods $B,C,D$, etc., which must be designated. It is clear that each is contained in a form $(ef*ke+\alpha)$; but I say that for one of them, say $B$, the number $k$ may be taken arbitrarily, and the placement of the others can then be derived.

Let $R$ be some primitive root of $x^n\equiv 1$, and suppose $B$ consists of the terms $R^\mu+R^\nu+$ etc. Let $\frac{1}{\mu}\rho^{ke+\alpha}\equiv \frac{1}{\nu}\pmod n$. Since the value of $r$ is arbitrary, provided only that $A$ obtains the sign $(e*\alpha)$, which it is clear can be done, put $r\equiv R^\nu\pmod p$. Then the first term of $B$ will be $r^{\rho^{ke+\alpha}}$, and $B$ may be designated by $(ef*ke+\alpha)$. If, instead of the term $R^\mu$, we had considered the term $R^\nu$, we would have obtained another value of $r$; but it is easily seen that for any root $\rho$, the root $r$ can have $\frac{\lambda}{ef}$ different values.

\subsection*{251.}
We now consider how the remaining periods are distinguished by their signs from the designation of one period. For this purpose another method must be sought for finding the remaining periods, since, insofar as the remaining periods are, like $A$ itself, roots of some congruence, no order is visible among them. Suppose that $A$ is designated by $(ef*0)$. From the preceding results it follows that
\[
\begin{array}{rcl}
A^2&\text{has the form}&M+N(ef*0)+O(ef*1)+P(ef*2)+\cdots,\\
A^3&\text{has the form}&M'+N'(ef*0)+O'(ef*1)+\cdots,\\
&\text{etc.}&\\
A^{ef-1}&\text{has the form}&M^*+N^*(ef*0)+O^*(ef*1)+\cdots.
\end{array}
\]
To these is added the congruence
\[
(ef*0)+(ef*1)+\cdots+(ef*ef-1)\equiv -1.
\]
We therefore have $ef-1$ linear congruences and just as many unknown quantities, which can consequently be determined by elimination.

\textit{Remark.} It may happen that the unknown quantities are given by such expressions as $\frac{V}{Wp}$. How this difficulty may be remedied will be taught below. Here, because this case can occur only very rarely, we shall not dwell on it.

\subsection*{252.}
This is enough in general concerning the solution of pure congruences. Many further points about them will be said below in passing; in particular, much can be transferred here from the solution of pure equations, and we shall not neglect to note this in its proper place. We add one more example, which, when compared with the precepts, will make everything clearer to less experienced readers.

Let $n=31$ and $p=311$; that is, we must investigate the roots of
\[
x^{31}-1\equiv 0\pmod {311}.
\]
First one must seek a primitive root of $y^{30}-1\equiv 0\pmod {31}$; such a root is $y\equiv 3$. Put therefore $\rho\equiv 3$ and first divide all primitive roots of the proposed congruence into five periods, namely
\[
\begin{array}{rcl}
(5*0)&=&(1)+(26)+(25)+(30)+(5)+(6),\\
(5*1)&=&(3)+(16)+(13)+(28)+(15)+(18),\\
(5*2)&=&(9)+(17)+(8)+(22)+(14)+(23),\\
(5*3)&=&(27)+(20)+(24)+(4)+(11)+(7),\\
(5*4)&=&(19)+(29)+(10)+(12)+(2)+(21).
\end{array}
\]
By the required calculations one finds the sum of the periods to be $\equiv -1$, the sum of their squares $\equiv 25$, of their cubes $\equiv 26$, of their fourth powers $\equiv 249$, and of their fifth powers $\equiv 564$. Hence the periods are roots of the congruence
\[
x^5+x^4-12x^3-21x^2+x+5\equiv 0.
\]
Furthermore one finds
\[
\begin{aligned}
(5*0)^2&\equiv 6+2(5*0)+2(5*3)+(5*4),\\
(5*0)^3&\equiv 12+15(5*0)+4(5*1)+3(5*2)+6(5*3)+6(5*4),\\
(5*0)^4&\equiv 90+60(5*0)+28(5*1)+26(5*2)+49(5*3)+38(5*4),
\end{aligned}
\]
and hence, by elimination,
\[
\begin{aligned}
5(5*1)&\equiv 3(5*0)^4-(5*0)^3-33(5*0)^2-24(5*0)+15,\\
5(5*2)&\equiv -2(5*0)^4-(5*0)^3+22(5*0)^2+31(5*0),\\
5(5*3)&\equiv (5*0)^4-2(5*0)^3,\\
5(5*4)&\equiv -2(5*0)^4+4(5*0)^3.
\end{aligned}
\]
One root of the congruence just found is $17$; therefore, if $(5*0)\equiv 17$, then $(5*1)\equiv 183$, $(5*2)\equiv 263$, $(5*3)\equiv 91$, $(5*4)\equiv 67$.

Now each of the periods found is again split into triples, namely
\[
\begin{array}{rcl}
(5*0)&\text{into}&(15*0),(15*5),(15*10),\text{ or into }(1)+(30),(26)+(5),(25)+(6),\\
(5*1)&\text{into}&(15*1),(15*6),(15*11),\text{ or into }(3)+(28),(16)+(15),(13)+(18),\\
&\text{etc.}&\text{etc.}
\end{array}
\]
Suppose that the periods into which the following are divided are roots of the respective congruences
\[
\begin{array}{rcl}
(5*0)&&x^3+A x^2+B x+C\equiv 0,\\
(5*1)&&x^3+A'x^2+B'x+C'\equiv 0,\\
(5*2)&&x^3+A''x^2+B''x+C''\equiv 0,\\
&\text{etc.}&
\end{array}
\]
Then
\[
\begin{array}{lll}
A\equiv -(5*0),&B\equiv (5*0)+(5*3),&C\equiv -2-(5*4),\\
A'\equiv -(5*1),&B'\equiv (5*1)+(5*4),&C'\equiv -2-(5*0),\\
\text{etc.}&\text{etc.}&\text{etc.}
\end{array}
\]
Therefore
\[
\begin{array}{lll}
(15*0),&(15*5),&(15*10)\quad\text{are roots of }x^3-17x^2+108x-60\equiv 0,\\
(15*1),&(15*6),&(15*11),\\
(15*2),&(15*7),&(15*12),\\
(15*3),&(15*8),&(15*13),\\
(15*4),&(15*9),&(15*14).
\end{array}
\]
Here one has
\[
\begin{aligned}
(15*0)^3-3(15*0)&\equiv (15*1),\\
(15*1)^3-3(15*1)&\equiv (15*2),\\
&\quad\text{etc.}
\end{aligned}
\]
Hence, if one root of the first congruence, namely $10$, is set equal to $(15*0)$, one obtains
\[
\begin{array}{lll}
(15*0)\equiv 10,&(15*5)\equiv \phantom{0},&(15*10)\equiv \phantom{0},\\
(15*1)\equiv 37,&(15*6)\equiv \phantom{0},&(15*11)\equiv \phantom{0},\\
(15*2)\equiv -151,&(15*7)\equiv \phantom{0},&(15*12)\equiv \phantom{0},\\
(15*3)\equiv -39,&(15*8)\equiv \phantom{0},&(15*13)\equiv \phantom{0},\\
(15*4)\equiv -112,&(15*9)\equiv \phantom{0},&(15*14)\equiv \phantom{0}.
\end{array}
\]
Finally, when the constituent terms of these individual periods are taken, we have
\[
\begin{array}{rcl}
(1),(30)&\text{roots of}&x^2-(15*0)x+1\equiv 0,\\
(3),(28)&&x^2-(15*1)x+1\equiv 0,\\
&\text{etc.}&
\end{array}
\]
The roots of the first congruence are $126$ and $195$, and these are therefore primitive roots of $x^{31}\equiv 1$; from them the remaining roots can be derived without difficulty.



\clearpage
\begin{center}
{\Large C. F. Gauss}\par
{\large Manuscript Nachlass}\par
{\large \emph{General investigations on congruences}}\par
{\large \emph{Analysis of residues. Chapter eight}}\par
{\large English translation of printed pp. 212--242}
\end{center}

\section*{Disquisitiones generales de congruentiis}
\addcontentsline{toc}{section}{Disquisitiones generales de congruentiis}

\begin{center}
\textbf{Analysis of residues, chapter eight}
\end{center}

\subsection*{330.}
What was set out in the preceding sections about congruences concerns only the simplest cases and was mostly obtained by particular methods. In this section Gauss attempts to reduce the theory of congruences, as far as is still possible, to higher principles, in almost the same manner as the theory of equations is usually treated; the analogy between the two theories has already been repeatedly observed. Since every algebraic congruence in one unknown can be reduced to the form
\[
X\equiv0,
\]
where $X$ is an algebraic function of the unknown $x$ involving no fractions, such functions must first be considered.

\subsection*{331.}
Let $P,Q$ be functions of the indeterminate $x$ of the form
\[
\begin{gathered}
A+Bx+Cxx+D x^3+\cdots,\\
H+Ix+Kxx+Lx^3+\cdots .
\end{gathered}
\]
If in both functions the coefficients of like powers of $x$ are congruent with respect to a given modulus, the functions are called congruent with respect to that modulus. Congruent functions take congruent values when congruent values are substituted for the indeterminate. The elementary rules proved earlier for numbers also hold for functions: from $P\equiv P'$, $Q\equiv Q'$, $R\equiv R'$, etc., one obtains
\[
P+Q+R+\cdots\equiv P'+Q'+R'+\cdots,\qquad
PQ\equiv P'Q',
\]
and similarly for longer products. If $PQ\equiv R$, the function $Q$ will be denoted by $R/P$ after the modulus has been specified; $Q$ is then the quotient obtained by dividing $R$ by $P$ modulo that modulus. All functions congruent to $Q$ are regarded as the same value. It remains to say in which cases such a quotient can have several incongruent values.

\subsection*{332.}
Assume that the modulus is a prime number. If the divisor $Q$ contains a term $Hx^h$ whose coefficient is not divisible by the modulus, a quotient cannot have two distinct values. For if $QA\equiv P$ and $QB\equiv P$, then $Q(A-B)\equiv0$. Writing
\[
Q=\cdots+Hx^h+Ix^{h+1}+\cdots,\qquad
A-B=Lx^l+Mx^{l+1}+\cdots,
\]
with $H$ and $L$ not divisible by $p$, the first non-vanishing term of $Q(A-B)$ is $HLx^{h+l}$, whose coefficient is not congruent to zero.

The rule for division of $P$ by $Q$ is the same as ordinary polynomial division, except that the numerical quotient must always be taken modulo the prime. If
\[
P=\cdots+ax^\alpha+\cdots,\qquad Q=\cdots+kx^\varkappa+\cdots
\]
with $a,k$ prime to the modulus and $\alpha\geq\varkappa$, the leading quotient term is
\[
x^{\alpha-\varkappa}\frac{a}{k},
\]
where $a/k$ is determined modulo the prime. If after all such terms are found a non-zero residue remains, $P$ is not divisible by $Q$ modulo the prime. The operation may likewise be started from the terms of greatest degree.

\subsection*{333.}
If $x\equiv a$ is a root of $X\equiv0$, then $X$ is divisible by $x-a$ modulo the given modulus. Indeed, if not, one would have
\[
X=(x-a)S+b
\]
with $b$ not divisible by the modulus. Substituting $x=a$ gives $X\equiv0$, while $(x-a)S\equiv0$, hence $b\equiv0$, a contradiction.

\subsection*{334.}
Given two functions, find their common divisor of greatest degree modulo a given modulus. Let the functions be $A,B$, with $A$ of at least as high a degree as $B$. Divide $A$ by $B$. If the division is exact, $B$ is the desired common divisor. Otherwise let $C$ be the residue, of lower degree than $B$, and continue:
\[
A=Ba+C,\qquad B=bC+D,\qquad C=cD+E,\quad\ldots
\]
The degrees of $A,B,C,D,\ldots$ continually decrease. If a division is finally exact, the last divisor is the desired greatest common divisor; if the process reaches a constant non-zero residue, the two functions are relatively prime modulo the prime. This is the Euclidean algorithm for functions modulo $p$.

\subsection*{335.}
If $A$ and $B$ are relatively prime functions modulo $p$, with degrees $a$ and $b$, then functions $P,Q$ can be found, of degrees respectively less than $b$ and $a$, such that
\[
PA+QB\equiv1.
\]
The assertion follows by following the Euclidean algorithm backwards. The successive divisions express each residue as a linear combination of $A$ and $B$; since the last non-zero residue is a constant, multiplication by its reciprocal modulo $p$ gives the displayed identity.

\subsection*{336.}
If $M$ is the greatest common divisor of $A$ and $B$, then one can always write
\[
PA+QB=M.
\]
In particular, when $M=1$ the preceding theorem is recovered. The point is that a common divisor of greatest degree may first be divided out; after reducing to the relatively prime case, the preceding argument supplies the required linear combination.

\subsection*{337.}
If $A,B,C,\ldots$ are pairwise relatively prime functions modulo $p$, and a function $M$ is divisible by each of them, then $M$ is divisible by their product. For from $PA+QB\equiv1$ it follows that every function divisible by both $A$ and $B$ is divisible by $AB$; the same argument is then repeated for the remaining factors.

\subsection*{338.}
If the congruence $X\equiv0$ has incongruent roots $a,b,c,\ldots$, then $X$ is divisible by
\[
(x-a)(x-b)(x-c)\cdots .
\]
Thus the number of roots cannot exceed the degree. More importantly, the solution of congruences is only one part of the more general question of resolving functions into factors. A congruence may have no real roots, although its function decomposes into factors of degree two, three, or higher. Such factors behave like imaginary roots in the theory of equations, but Gauss chooses to deduce the subsequent theory from first principles rather than introduce such quantities directly.

\subsection*{339.}
Functions that cannot be divided by any lower-degree function modulo the fixed modulus are called \emph{prime}. Every function of the first degree is prime; a function of the second degree is either prime or a product of two functions of the first degree. Thus a quadratic function is prime precisely when its congruence has no real roots. For example, $x^2+x+1$ is prime modulo $5$, since
\[
x^2+x+1\equiv (x-2)^2+3\pmod5
\]
and $3$ is a quadratic non-residue of $5$.

\subsection*{340.}
Every function is either prime or a product of prime functions, and in the latter case the product is unique. If a function $A$ is not prime it is divisible by a lower-degree function; continuing, the decreasing degrees must lead to a prime divisor. Dividing off one prime divisor after another gives a factorization into primes. If two different prime factorizations existed, delete the common factors in the two products. What remains would give a product of primes on one side relatively prime to a prime factor on the other, which is impossible by the preceding propositions. Hence the factorization is unique.

\subsection*{341.}
The first task is to determine the number of prime functions of a given degree. For a fixed modulus there are only finitely many incongruent functions of each degree; some are composed from lower-degree primes, and the others are prime. Gauss begins with the simple cases, since a completely rigorous general enumeration is delicate.

\subsection*{342.}
Restrict to monic functions, which is sufficient here. There are $p$ functions of the first degree. A second-degree function is either prime or a product of two first-degree functions. The product of two possibly equal first-degree factors gives $p(p+1)/2$ composite quadratics; since there are $p^2$ monic quadratics in all, the number of prime quadratic functions is
\[
\frac12(pp-p).
\]
Similarly, subtracting the products made from lower-degree primes gives the count of prime cubics,
\[
\frac13(p^3-p).
\]

\subsection*{343.}
Let $(1)$ denote the number of prime functions of degree one, $(2)$ the number of degree two, and so on. Counting all monic functions of degree $n$ by their decompositions into prime factors gives relations of the form
\[
p^n=(n)+\text{products built from }(1),(2),\ldots,(n-1).
\]
Expanding these relations for small $n$ reveals the general law governing the numbers $(n)$.

\subsection*{344--346.}
From the beginning of the sequence one sees that the sum of the terms in the expression $(n)$ is $p^n$, with a correction term depending on the divisors of $n$. If the divisors of $n$ are $a,b,c,\ldots$, the general relation is
\[
n(n)=a(a)+b(b)+c(c)+\cdots
\]
in the form appropriate to the divisor lattice. Inverting this relation gives the familiar inclusion--exclusion formula. For example, if $n=a^\alpha b^\beta c^\gamma\cdots$, where $a,b,c,\ldots$ are distinct primes, then
\[
n(n)=p^n-p^{n/a}-p^{n/b}-\cdots+p^{n/ab}+\cdots .
\]
In particular, when $n=a^\alpha$ with $a$ prime,
\[
n(n)=p^n-p^{n/a}.
\]

\subsection*{347.}
The preceding theorem can be expressed by saying that the product of all prime functions of degrees dividing $n$ has degree $p^n$. This product is exactly the function whose roots are all the residues in the finite field generated by such elements; hence it is
\[
x^{p^n}-x
\]
up to the usual normalization. It follows that the prime factors of $x^{p^n}-x$ are precisely the prime functions whose degrees divide $n$.

\subsection*{348.}
Given an equation
\[
x^m+Ax^{m-1}+Bx^{m-2}+\cdots=0
\]
with roots $a,b,c,\ldots$, find an equation whose roots are $a^T,b^T,c^T,\ldots$. One solution is to compute the sums of powers of the roots of the given equation up to the required order, and from them compute the elementary symmetric functions of $a^T,b^T,c^T,\ldots$. A second solution, better adapted to the present purpose, uses a primitive $T$-th root $\theta$ and forms the product over all expressions obtained by multiplying the roots by powers of $\theta$; this also shows that the resulting coefficients are integers when the original coefficients are integers.

\subsection*{349.}
If $T$ is prime, the coefficient of $x^{n-1}$ in $(P,\rho^T)$ is congruent modulo $T$ to the coefficient of $x^{n-1}$ in $P^T$. This gives, in the prime case, a third solution of the preceding problem. The proof uses the fact that the coefficients of the transformed product have the required symmetric form; after putting $\theta^T=1$, the coefficient reduces modulo $T$ to the $T$-th power of the corresponding coefficient of $P$.

\subsection*{350.}
If $r$ is prime, then
\[
(P,\rho^r)\equiv P\pmod r.
\]
Indeed, comparing the coefficient $N'$ of $x^{n-1}$ in $(P,\rho^r)$ with the corresponding coefficient $N$ in $P$, the preceding article gives $N'\equiv N^r$, and Fermat's theorem gives $N^r\equiv N$. The same applies to the remaining coefficients.

\subsection*{351.}
There is a value $v<p^m$ such that $x^v-1$ is divisible modulo $p$ by a given function $P$ of degree $m$, provided the constant term of $P$ is not zero. Divide
\[
1,x,x^2,\ldots,x^{p^m-1}
\]
successively by $P$. The residues all have degree $<m$, so there are only $p^m$ possible residues; two must coincide. Their difference shows that some $x^v-1$ is divisible by $P$.

\subsection*{352.}
If $P$ is prime of degree $m$ and $x^v$ is the least power which leaves residue $1$ after division by $P$, then $v$ is either $p^m-1$ or a divisor of that number, except for the special case $P=x$. The residues
\[
1,x,x^2,\ldots,x^{v-1}
\]
are all distinct and non-zero, and multiplication by powers of $x$ permutes them. Consequently the period length $v$ divides the number of non-zero residues, $p^m-1$.

\subsection*{353.}
Every prime function of degree $m$ divides $x^{p^m-1}-1$ modulo $p$ (apart from the factor $x$ when present). Thus the linear prime functions other than $x$ divide $x^{p-1}-1$, which is Fermat's theorem; the prime quadratic functions divide $x^{p^2-1}-1$, and so forth. More precisely, a prime of degree $m$ divides $x^{p^m}-x$, and not a similar expression of lower exponent unless its degree divides that exponent.

\subsection*{354.}
If $x^v-1$ is divisible by $P$, then the residues recur periodically:
\[
x^{kv+l}\equiv x^l\pmod P.
\]
Using the expansion of the logarithmic derivative of
\[
P=x^m+Ax^{m-1}+Bx^{m-2}+\cdots
\]
one obtains congruences between the sums of equal powers of the roots. These congruences show that the corresponding transformed functions agree modulo $p$ in the stated range.

\subsection*{355.}
Conversely, if in the residue sequence
\[
1,x,x^2,\ldots
\]
the terms after the $v$-th are congruent successively to the first terms, then $x^v-1$ is divisible by $P$, provided $P$ has no repeated factor. This follows by applying the preceding relation to the derivative condition and using the fact that $P$ is relatively prime to its derivative when no factor is repeated.

\subsection*{356.}
Let $P$ be a prime function of degree $m$, and let $X$ be a symmetric function of
\[
x,\;x^p,\;x^{p^2},\ldots,x^{p^{m-1}}.
\]
That is, $X$ remains unchanged when these quantities are permuted. Then on division of $X$ by $P$ the residue is a number. For replacing $x$ by $x^p$ cycles the quantities; the residue therefore must be invariant under this substitution, and the congruence $Y^p\equiv Y$ forces the coefficients to reduce to constants.

\subsection*{357.}
For
\[
P=x^m+Ax^{m-1}+Bx^{m-2}+Cx^{m-3}+\cdots,
\]
the residue of the sum $x+x^p+\cdots+x^{p^{m-1}}$ is $-A$, the residue of the sum of pairwise products is $B$, that of triple products is $-C$, and so on. These are the ordinary elementary symmetric functions of the conjugate residues.

\subsection*{358--359.}
Let $P$ be prime, and let $x^v$ be the least power of $x$ that leaves residue $1$ modulo $P$. If $P'=(P,\rho^{p^\alpha})$, then the exponent occurring is congruent to a power of $p$ modulo $v$. This is obtained by comparing the ordered systems of residues arising from $P$ and from its transforms. The note attached in the source gives the intended replacement for a deleted article.

\subsection*{360.}
When $v$ is divisible by $p$, write $v=p^k v'$. Since
\[
x^v-1\equiv (x^{v'}-1)^{p^k}\pmod p,
\]
one may restrict to the case where $v$ is not divisible by $p$. If $p^m\equiv1\pmod v$ with $m$ minimal, then $x^v-1$ has prime divisors of degree $m$ and also divisors whose degrees divide $m$. The number of such factors is found by subtracting the factors already belonging to lower degrees.

\subsection*{361.}
If a function $X$ is divisible by another function $Y$, and $X'$ is obtained from $X$ by replacing $x$ by $x^p$, then $X'$ is divisible by the corresponding transform $(Y,\rho^{p-1})$. This is immediate from writing $X=EY$ and applying the same substitution to both factors.

\subsection*{362.}
Using these principles one can determine the prime divisors of $x^v-1$. Suppose that all divisors also dividing some $x^{v'}-1$ with $v'<v$ have already been found. The remaining divisors correspond to the periods formed from the roots of $x^v-1$, and their degrees are governed by the least exponent $m$ for which $p^m\equiv1\pmod v$.

\subsection*{363.}
The essential point is to determine the residues of these periods. Let the product of the suitable prime functions be
\[
P=x^m+Ax^{m-1}+Bx^{m-2}+\cdots .
\]
The residue of the greatest period is then expressed in terms of $A$, and the other periods are obtained by the same cyclic process as in the theory of cyclotomy. Hence the auxiliary equations that separate the periods may be formed, and their solvability as congruences follows from the preceding results.

\subsection*{364.}
As an example, resolve $x^{16}-1$ modulo $17$. Here $m=4$. The roots are distributed into periods according to the powers of $17$ modulo the relevant divisor; the resulting products give two prime factors of degree four. The computation shows explicitly how the periods determine factors $P$ and $P'$, and the congruence obtained from them gives the needed decomposition.

\subsection*{365.}
The purpose here is only to show the possibility of the solution; many computational devices that shorten the work are omitted. Still, important consequences must not be passed over. The preceding articles show that the auxiliary equations occurring in the solution of $x^v=1$ have possible roots when the corresponding period has not yet been separated.

\subsection*{366.}
This gives a third proof of the fundamental theorem of Chapter IV, notable because it rests on principles different from the earlier proofs. A fourth proof is obtained in the opposite direction. Let $v$ be a prime of the form $4n+1$. If $p$ is a quadratic residue of $v$, the quadratic auxiliary equation has real roots modulo $p$; conversely, the same reasoning shows that $p$ must be a quadratic residue of $v$. This yields the reciprocity statement.

\subsection*{367.}
Although the preceding discussion was restricted to prime $v$, analogous theorems can be determined for composite $v$. Similar observations apply when the roots are distributed into three or more periods. For example, the cubic auxiliary equation associated with a prime $v$ may be solvable precisely for primes in the required congruence classes.

\subsection*{368.}
A congruence $S\equiv0$ is said to have equal roots, or more generally equal divisors, when its function is divisible by a power of some other function. The criterion is the same as in the theory of equations. If
\[
X=A^aB^bC^c\cdots,
\]
then the common divisor of $X$ and $dX/dx$ contains each prime factor one fewer time than in $X$. Hence $X$ has no repeated factor if and only if $X$ and $dX/dx$ are relatively prime.

\subsection*{369.}
Examples illustrate the criterion. For the first function considered, the greatest common divisor of $X$ and $dX/dx$ is $1$, hence no repeated divisor occurs. In the second example the common divisor is $x^2+x+1$ up to a non-zero numerical factor, so the corresponding squared factor occurs in $X$.

\subsection*{370--371.}
If $X$ has been written in the form
\[
A^aB^bC^c\cdots
\]
with $A,B,C,\ldots$ pairwise relatively prime and $a,b,c,\ldots$ unequal, the resolution can be carried further. The function $x^p-x$ is the product of all linear prime functions; its greatest common divisor with $X$ gives the product of all linear factors of $X$. Similarly, the greatest common divisor of $X/E$ with $x^{p^2}-x$ gives the product of the quadratic prime factors, and so on.

\subsection*{372.}
If a product of several prime functions of the same degree is given, the individual factors must be found by trial. The problem is analogous to factoring composite integers. Here, however, it can be known in advance whether further factorization is possible and what degrees the factors must have. Details of computational shortcuts are left aside.

\subsection*{373.}
Suppose $X$ is resolved modulo $p$ into pairwise relatively prime factors
\[
S,\;S',\;S'',\ldots .
\]
The problem is to resolve $X$ modulo $p^2$ into corresponding factors
\[
S+\varphi p,\quad S'+\varphi' p,\quad S''+\varphi''p,\ldots .
\]
Writing
\[
X=SS'S''\cdots+pR
\]
and comparing the product after adding the corrections, one obtains a linear congruence for $\varphi,\varphi',\ldots$. Because the factors are pairwise relatively prime modulo $p$, the necessary corrections can be determined, and thus each factor lifts modulo $p^2$.

\subsection*{374.}
The same method proves that $X$ can also be resolved modulo $p^3$, $p^4$, and so on. In general, if
\[
X\equiv PQ\pmod{p^\nu},
\]
with $P$ and $Q$ relatively prime modulo $p$, one seeks corrected factors
\[
P'=P+Ap^\nu,\qquad Q'=Q+Bp^\nu .
\]
Comparison of $P'Q'$ with $X$ reduces the next lifting step to a linear congruence of the type already solved.

\subsection*{375.}
If $X$ has no equal divisors modulo $p$, the preceding lifting process decomposes it modulo $p^k$ in the same way as modulo $p$. If equal divisors occur, the matter becomes more complicated. Gauss treats only the frequent case where the equal factors are of the first degree. Suppose
\[
X\equiv X'(x-a)^m\pmod p,\qquad (X',x-a)=1.
\]
After substituting $z+a$ for $x$ one has
\[
Z\equiv Z'z^m\pmod p.
\]
To obtain a divisor of the form $z+\alpha p$ modulo $p^2$ the next term must have the form $zZ''+pB$. For $m=4$ the condition modulo $p^3$ becomes
\[
\alpha Z''\equiv B\pmod{z,p}.
\]
If $Z''\equiv0$ but $B\not\equiv0$, no such divisor exists. If neither $Z''$ nor $B$ is congruent to zero, then $\alpha$ is uniquely determined:
\[
\alpha\equiv \frac{B}{Z''}\pmod{z,p}.
\]
The manuscript then proceeds to the next correction term for a divisor modulo $p^4$ and breaks off with
\[
0\equiv .
\]

\section*{Remarks on the Analysis of Residues}
The two preceding pieces are taken from a large manuscript entitled \emph{Analysis Residuorum}, probably dating from 1797 or 1798. A complete reworking of that manuscript later produced the \emph{Disquisitiones Arithmeticae}. The full title of Chapter Six was:
\begin{quote}
\emph{Solutio congruentiae $x^m-1\equiv 0$ et aequationis $x^m-1=0$; cum dilucidationibus super theoria polygonorum regularium.}
\end{quote}
Its second part, sections 253--278, passed essentially into the seventh section of the \emph{Disquisitiones Arithmeticae}.

Part of Chapter Seven also survives, under the title \emph{Variae quarundam investigationum praecedentium applicationes}, sections 279--302. It comprised subsections on common fractions and their decimal expansion, on the indeterminate equations $xx=a+by$ and $axx+bxy=c$, and on the investigation of divisors of numbers. This material was almost literally incorporated into Section Six of the \emph{Disquisitiones Arithmeticae}.

The two sections printed here deal with the topics that, as the preface and Articles 11, 44, 61, 62, 65, and 84 of the \emph{Disquisitiones Arithmeticae} show, were intended to form the eighth section of that work. The plan was later altered. Among the manuscripts is a fragment headed \emph{Sectio octava: Quarundam disquisitionum ad circuli sectionem pertinentium uberior consideratio}; it begins with Article 367 and would have formed the continuation of the \emph{Disquisitiones Arithmeticae}. Since the few surviving articles later passed in substance into \emph{Summatio quarumdam serierum singularium}, that fragment is excluded from the present edition.

In the preceding printing the text of the original has been essentially preserved, although formally it was not ready for print. The following notes mark the most important alterations and add explanations.

For sections 237 and 230 compare the corresponding articles of the \emph{Disquisitiones Arithmeticae}. In section 241, if $n=2^\nu$ and $\nu>3$, no number $p$ of the indicated kind exists, but the investigation is not essentially changed. In section 251 the difficulty was presumably to be removed by introducing higher powers of $p$ as moduli; compare sections 303, 372, and 373. In section 332 the assumption that the modulus is prime is retained through section 372. In section 338 the incomplete citation may be referred to Article 44 of the \emph{Disquisitiones Arithmeticae}.

For sections 344--346 the manuscript contained two proofs. The first, beginning \emph{iam demonstrare accingimur}, was suppressed according to Gauss's own instruction; he had found an infinitely simpler proof. The printed proof is further shortened by considering the development of a derivative expression, thereby avoiding a reference to the suppressed proof. In section 348 the expression \emph{radix prima} is to be taken in the same sense as \emph{radix propria} in \emph{Summatio quarumdam serierum singularium}, Article 11. The integrality claim about the coefficients of the developed product is supplemented by observing the invariance of the product under replacing a primitive root $\theta$ by $\theta^k$, where $k$ is relatively prime to $T$.

In section 354, multiplication by $x^{v-t}$ gives congruences between the sums of equal powers of the roots of the relevant equations. When $m\ge p$, the Newton formulae no longer determine the coefficient of $x^{m-p}$ modulo $p$ without ambiguity. The general principle is that if $R(x)-S(x)$ is divisible by $P(x)$ modulo $p$, and $a,b,c,\ldots$ are the roots of $P(x)=0$, then
\[
(x-R(a))(x-R(b))(x-R(c))\cdots
\]
and
\[
(x-S(a))(x-S(b))(x-S(c))\cdots
\]
are congruent modulo $p$. Section 355 is clarified by section 368, which shows that $P$ and $dP/dx$ have no common divisor if $P$ has no repeated factor. The note under sections 358--359 comes from a separate sheet, probably intended to replace the deleted section 359. In section 360 an inaccuracy of the manuscript has been corrected. The notes to sections 361, 363, 367, and 371 explain the notation, the reference to moduli that are powers of $p$, the period example connected with a seventh-root equation, and the fact that the intended example in section 371 was not carried out.

\begin{flushright}
R. Dedekind.
\end{flushright}


\clearpage

\section*{Further development of the investigation concerning pure equations}
\addcontentsline{toc}{section}{Further development of the investigation concerning pure equations}

\subsection*{1.}

Since the method by which, in \emph{Disquisitiones Arithmeticae}, art.~360, we taught the solution of the equation
\[
x^n-1=0
\]
constitutes a very fertile and very important theory, and since in that work only its first elements could be touched upon, we hope that geometers will welcome our resuming this subject here. We treat more fully what there had only been briefly sketched, partly with demonstrations suppressed, and we pursue more deeply the advances that have been added since that time.

The exponent $n$ is supposed to be a prime number, and $n-1$ is resolved into the factors $\alpha\delta\gamma$; furthermore, let $g$ denote a primitive root modulo $n$. Let $r$ denote, indeterminately, a root of $x^n-1=0$, and let $R$ denote, indeterminately, a root of $x^\delta-1=0$. If the circumference of the circle whose radius is $1$ is denoted by $P$, and the imaginary quantity $\sqrt{-1}$ by $i$, all roots of $x^\delta-1=0$, that is, all values of $R$, are represented by
\[
\cos \frac{kP}{\delta}+i\sin\frac{kP}{\delta},
\]
where $k$ runs over $0,1,2,3,\ldots,\delta-1$. It is clear that all powers of any root $R$ are again roots. If $R$ corresponds to a value of $k$ prime to $\delta$, then the powers
\[
R^0,\ R,\ R^2,\ R^3,\ldots,\ R^{\delta-1}
\]
are all distinct and exhaust the whole system of roots; in this case we call $R$ a \emph{proper} root of $x^\delta-1=0$. Conversely, a root $R$ corresponding to a value of $k$ not prime to $\delta$ will be called \emph{improper}. If $\delta'$ is the greatest common divisor of $k$ and $\delta$, then $R^{\delta/\delta'}=1$, while
\[
R^0,\ R,\ R^2,\ldots,R^{\delta/\delta'-1}
\]
are distinct; hence $R$ is a proper root of $x^{\delta/\delta'}-1=0$. The same remarks apply to the equation $x^n-1=0$, but all of its roots are necessarily proper except the root $1$.

\subsection*{2.}

With these preliminaries, the investigation will chiefly concern functions of the following form, composed of $\delta\gamma$ terms:
\[
r+Rr^{g^\alpha}+R^2r^{g^{2\alpha}}+R^3r^{g^{3\alpha}}+\cdots+
R^{\delta\gamma-1}r^{g^{\alpha(\delta\gamma-1)}}.
\]
For brevity we denote this expression by $[r,R]$. Each term is the product of a power of $r$ and a power of $R$; the exponents of the former form a geometric progression, while the exponents of the latter form an arithmetic progression. The exponents
\[
1,\ g^\alpha,\ g^{2\alpha},\ g^{3\alpha},\ldots,g^{\alpha(\delta\gamma-1)}
\]
are incongruent modulo $n$, and so the corresponding powers of $r$ are distinct. If the sequence is continued, it begins again, because $g^{\alpha\delta\gamma}\equiv1\pmod n$, whence $r^{g^{\alpha\delta\gamma}}=r$. The other factors
\[
1,\ R,\ R^2,\ R^3,\ldots,R^{\delta\gamma-1}
\]
form $\gamma$ equal periods, since $R^\delta=1$. Thus
\[
\begin{aligned}
{}[r,R]&=
r+R\bigl(r^{g^\alpha}+r^{g^{(\delta+1)\alpha}}+\cdots+r^{g^{(\gamma\delta-\delta+1)\alpha}}\bigr)\\
&\quad+R^2\bigl(r^{g^{2\alpha}}+r^{g^{(\delta+2)\alpha}}+\cdots+r^{g^{(\gamma\delta-\delta+2)\alpha}}\bigr)+\cdots\\
&\quad+R^{\delta-1}\bigl(r^{g^{(\delta-1)\alpha}}+r^{g^{(2\delta-1)\alpha}}+\cdots+r^{g^{(\gamma\delta-1)\alpha}}\bigr),
\end{aligned}
\]
or, using the notation of \emph{Disquisitiones Arithmeticae}, art.~343,
\[
[r,R]=(\gamma,1)+R(\gamma,g^\alpha)+R^2(\gamma,g^{2\alpha})+\cdots+R^{\delta-1}(\gamma,g^{(\delta-1)\alpha}).
\]

\subsection*{3.}

If the root $r$ is taken to be $1$, then
\[
[1,R]=1+R+R^2+R^3+\cdots+R^{\delta\gamma-1}
=\gamma(1+R+R^2+R^3+\cdots+R^{\delta-1}).
\]
This value is $\delta\gamma$ if also $R=1$, and is $0$ for every other value of $R$. Conversely, if $R=1$ while $r$ remains indeterminate, then
\[
[r,1]=r+r^{g^\alpha}+r^{g^{2\alpha}}+\cdots+r^{g^{\alpha\delta\gamma-\alpha}}
=(\delta\gamma,1).
\]
Thus it consists of a period of $\delta\gamma$ roots, one of which is $r$ itself. When $\alpha=1$, these periods contain all roots $r,r^2,r^3,\ldots,r^{n-1}$, only in a changed order.

The following relations are immediate:
\[
[r,R]=R[r^{g^\alpha},R]=R^2[r^{g^{2\alpha}},R]\quad\text{and generally}\quad
[r,R]=R^k[r^{g^{k\alpha}},R].
\]
Consequently $[r^m,R]$ is either $[1,R]$ if $m$ is divisible by $n$, or else reducible to $R^\mu[r,R]$ with $0<\mu<\delta$. If $m$ is not divisible by $n$, let
\[
\operatorname{ind}m=\lambda\alpha+\nu,\qquad \nu<\alpha.
\]
Then
\[
[r^m,R]=[r^{g^{\lambda\alpha+\nu}},R]=R^{-\lambda}[r^{g^\nu},R],
\]
so that $\mu=-\lambda$, or $\mu\equiv-\lambda\pmod\delta$ if a positive exponent is desired.

\subsection*{4.}

\emph{Theorem.} Let $r'$ denote, just as $r$ does, an indeterminate root of $x^n-1=0$, and let $R'$ denote, just as $R$ does, an indeterminate root of $x^\delta-1=0$. Then
\[
\begin{aligned}
[r,R]\,[r',R']={}&[rr',RR']+R[r^{g^\alpha}r',RR']
+R^2[r^{g^{2\alpha}}r',RR']\\
&+R^3[r^{g^{3\alpha}}r',RR']+\cdots+
R^{\delta\gamma-1}[r^{g^{\alpha(\delta\gamma-1)}}r',RR'].
\end{aligned}
\]
Multiplying $[r,R]$ by the single parts of $[r',R']$, and then collecting the first terms of the expanded parts, gives $[rr',RR']$; collecting the second terms gives $R[r^{g^\alpha}r',RR']$, and so on. This yields the formula.

By merely interchanging $r,R$ with $r',R'$, the same product can also be written
\[
\begin{aligned}
[rr',RR']&+R'[rr'^{g^\alpha},RR']+R'^2[rr'^{g^{2\alpha}},RR']\\
&+R'^3[rr'^{g^{3\alpha}},RR']+\cdots+
R'^{\delta\gamma-1}[rr'^{g^{\alpha(\delta\gamma-1)}},RR'].
\end{aligned}
\]
It follows that a product of any number of functions $[r,R]$, $[r',R']$, $[r'',R'']$, and so on is equal to
\[
\sum R^{k'}R'^{k''}R''^{k'''}\cdots
\,[r^{g^{\alpha k'}}r'^{g^{\alpha k''}}r''^{g^{\alpha k'''}}\cdots,\,
RR'R''R'''\cdots],
\]
where $k',k'',k'''$, and so on range independently over $0,1,2,3,\ldots,\delta\gamma-1$.

\subsection*{5.}

The preceding product formula is completely general; it supposes no relation among the roots $r,r',r'',\ldots$ nor among $R,R',R'',\ldots$. If, however, the roots $r',r'',\ldots$ may be regarded as powers of $r$, and $R',R'',\ldots$ as powers of $R$, each part of the product has the form $R^m[r^m,R^\lambda]$, with the same exponent $\lambda$, namely $R^\lambda=RR'R''\cdots$. By art.~3, such a product reduces to
\[
A[1,R^\lambda]+B[r,R^\lambda]+B'[r^g,R^\lambda]+B''[r^{g^2},R^\lambda]
+B'''[r^{g^3},R^\lambda]+\cdots+B^{(\alpha-1)}[r^{g^{\alpha-1}},R^\lambda],
\]
where the coefficients are expressions
\[
h+h'R+h''R^2+h'''R^3+\cdots+h^{(\delta-1)}R^{\delta-1}
\]
with integral coefficients. The simplest case occurs when $r=r'=r''=\cdots$ and $R=R'=R''=\cdots$; then the product becomes a power $[r,R]^\lambda$, which always returns to the form just described.

\subsection*{6.}

Putting $\lambda=\delta$, the power $[r,R]^\delta$ takes the form
\[
\begin{aligned}
&A[1,1]+B[r,1]+B'[r^g,1]+\cdots+B^{(\alpha-1)}[r^{g^{\alpha-1}},1]\\
&=\delta\gamma A+B(\delta\gamma,1)+B'(\delta\gamma,g)+B''(\delta\gamma,g^2)+\cdots+
B^{(\alpha-1)}(\delta\gamma,g^{\alpha-1})=\theta'.
\end{aligned}
\]
Thus, if $R$, the coefficients $A,B,B',\ldots$, and the periods $(\delta\gamma,1)$, $(\delta\gamma,g)$, etc. are known, $\theta'$ is known, and $[r,R]$ may be obtained as $\sqrt[\delta]{\theta'}$. The expression has $\delta$ values. If $R$ is a proper root of $x^\delta-1=0$, this ambiguity is harmless, for the $\delta$ values are precisely
\[
[r,R],\quad [r^{g^\alpha},R],\quad [r^{g^{2\alpha}},R],\ldots,
[r^{g^{(\delta-1)\alpha}},R],
\]
which are proportional to the different $\delta$-th roots. Since the root $r$ itself has only been determined as lying in the period $(\delta\gamma,1)$, any of these values may be chosen.

If $R$ is not a proper root of $x^\delta-1=0$, the preceding conclusion must be modified. Suppose $R$ is a proper root of $x^{\delta'}-1=0$, where $\delta'$ divides $\delta$. Then
\[
[r,R]=[r^{g^{\alpha\delta'}},R],\quad
[r^{g^\alpha},R]=[r^{g^{\alpha\delta'+\alpha}},R],\quad\text{etc.}
\]
Only $\delta'$ of the $\delta$ values are admissible, the rest being spurious. But in that case it is unnecessary to pass all the way to the $\delta$-th power; $[r,R]^{\delta'}$ is already reducible to the same type of expression. Hence $[r,R]$ is obtained from an expression $\sqrt[\delta']{\theta'}$, and the value chosen is immaterial in the same sense.

\subsection*{7.}

The functions $[r,R^2]$, $[r,R^3]$, and generally $[r,R^k]$ may be determined similarly. If, in place of $R$, the powers $R^2,R^3,\ldots,R^k$ are substituted in $\theta'$, giving $\theta'',\theta''',\ldots,\theta^{(k)}$, then
\[
[r,R^2]^\delta=\theta'',\quad [r,R^3]^\delta=\theta''',\quad
[r,R^k]^\delta=\theta^{(k)}.
\]
But radical expressions are inconvenient whenever quantities have to be combined, since they do not indicate which values belong together. Once a definite value of $[r,R]$ has been taken, all the values $[r,R^2]$, $[r,R^3]$, etc. are thereby fixed; the radical expressions do not show this dependence. We therefore seek rational expressions for these functions in terms of $[r,R]$ and known quantities.

The product
\[
[r,R^k][r,R]^{\delta-k}
\]
is reducible to
\[
\delta\gamma A+B(\delta\gamma,1)+B'(\delta\gamma,g)+B''(\delta\gamma,g^2)+\cdots+
B^{(\alpha-1)}(\delta\gamma,g^{\alpha-1}),
\]
where the coefficients are rational functions of $R$. Write
\[
[r,R^2][r,R]^{\delta-2}=\vartheta'',\quad
[r,R^3][r,R]^{\delta-3}=\vartheta''',\quad
[r,R^4][r,R]^{\delta-4}=\vartheta''''.
\]
Then
\[
[r,R^2]=\frac{\vartheta''}{\theta'}[r,R]^2,\quad
[r,R^3]=\frac{\vartheta'''}{\theta'}[r,R]^3,\quad
[r,R^4]=\frac{\vartheta''''}{\theta'}[r,R]^4.
\]
These rationally determine $[r,R^2]$, $[r,R^3]$, etc., provided $[r,R]\neq0$. Gauss notes that $[r,R]$ can never vanish when $r\neq1$, although the proof is reserved for another occasion.

\subsection*{8.}

The preceding results are useful for descending from periods of $\delta\gamma$ terms to periods of $\gamma$ terms. If $R$ is a proper root, then
\[
\begin{aligned}
\delta(\gamma,1)&=(\delta\gamma,1)+[r,R]+[r,R^2]+[r,R^3]+\cdots+[r,R^{\delta-1}],\\
\delta(\gamma,g^\alpha)&=(\delta\gamma,1)+R^{\delta-1}[r,R]+R^{\delta-2}[r,R^2]+R^{\delta-3}[r,R^3]+\cdots+R[r,R^{\delta-1}],\\
\delta(\gamma,g^{2\alpha})&=(\delta\gamma,1)+R^{2\delta-2}[r,R]+R^{2\delta-4}[r,R^2]+R^{2\delta-6}[r,R^3]+\cdots+R^2[r,R^{\delta-1}].
\end{aligned}
\]
If radical values were substituted separately for $[r,R]$, $[r,R^2]$, and so on, the value of each series would remain uncertain. By using rational expressions for $[r,R^k]$, this avoidable ambiguity disappears. Gauss observes that this point seems to have escaped Lagrange, who believed he had simplified the method of \emph{Disquisitiones Arithmeticae}, art.~360 by replacing the rational expressions there by radical expressions. Once the periods $(\gamma,1)$, $(\gamma,g^\alpha)$, etc., or even one of them, have been found, all the remaining periods of $\gamma$ terms can be rationally deduced. Thus the descent requires the solution of $x^\delta=1$ and $x^\delta=\theta'$, while the remaining operations are rational.

\subsection*{9.}

Most of this had already been treated in \emph{Disquisitiones Arithmeticae}; here Gauss supplies a suppressed demonstration. The value of $\sqrt[\delta]{\theta'}$ seems to require both the extraction of a $\delta$-th root of the modulus and a division of the angle into $\delta$ parts. In fact only the latter division is essential; the former always reduces to a square root.

Let $\theta'=P+iQ$ and let one value of $\sqrt[\delta]{\theta'}$ be $p+iq$. If
\[
P=E\cos F,\quad Q=E\sin F,\quad p=e\cos f,\quad q=e\sin f,
\]
then $e=\sqrt[\delta]{E}$ and $f$ is one of
\[
\frac1\delta F,\quad \frac1\delta(F+360^\circ),\quad \frac1\delta(F+720^\circ),\ldots,
\frac1\delta(F+(\delta-1)360^\circ).
\]
But
\[
[r^{-1},R^{-1}]=p-iq=e(\cos f-i\sin f),
\]
and therefore
\[
[r,R][r^{-1},R^{-1}]=e^2.
\]
By the product theorem this equals
\[
\begin{aligned}
&[1,1]+R[r^{g^\alpha-1},1]+R^2[r^{g^{2\alpha}-1},1]+\cdots+
R^{\delta\gamma-1}[r^{g^{\alpha(\delta\gamma-1)}-1},1]\\
&=\delta\gamma+R(\delta\gamma,g^\alpha-1)+R^2(\delta\gamma,g^{2\alpha}-1)+\cdots+
R^{\delta\gamma-1}(\delta\gamma,g^{\alpha(\delta\gamma-1)}-1).
\end{aligned}
\]
Thus $e$ is obtained by a square root. In the important special case $\alpha=1$, this gives simply
\[
ee=\delta\gamma+(R+R^2+\cdots+R^{\delta\gamma-1})(r+r^2+\cdots+r^{n-1})
=\delta\gamma+1=n,
\]
so that $e=\sqrt n$, provided $r$ and $R$ are both different from $1$.

\subsection*{10.}

Gauss now passes to the especially fruitful case $\alpha=1$. Then
\[
[r,R]=r+Rr^g+R^2r^{g^2}+R^3r^{g^3}+\cdots+R^{n-2}r^{g^{n-2}},
\]
where $n$ is prime, $r$ is a root of $x^n-1=0$, $R$ is a root of $x^\delta-1=0$, $\delta$ divides $n-1$, and $g$ is a fixed primitive root modulo $n$. Set
\[
1+r+r^2+\cdots+r^{n-1}=s,\qquad
1+R+R^2+\cdots+R^{n-2}=S.
\]
Then $s=n$ for $r=1$ and $0$ otherwise; similarly $S=n-1$ for $R=1$ and $0$ otherwise. We have
\[
[1,R]=S,\qquad [r,1]=s-1,
\]
and, when $m$ is not divisible by $n$,
\[
[r^m,R^\lambda]=R^{-\lambda\operatorname{ind}m}[r,R^\lambda].
\]
Hence any product of two or more functions of the form $[r^h,R^H]$ can be reduced to
\[
A[1,R^\lambda]+B[r,R^\lambda],
\]
where $A$ and $B$ are rational functions of $R$ with integral coefficients. The next articles develop this reduction into an efficient algorithm.

\subsection*{11.}

For the product $[r,R^\mu][r,R^\nu]$, the theorem gives
\[
\begin{aligned}
{}&[r^2,R^{\mu+\nu}]
+R^\mu[r^{g+1},R^{\mu+\nu}]
+R^{2\mu}[r^{g^2+1},R^{\mu+\nu}]
+\cdots\\
&\quad+R^{(n-2)\mu}[r^{g^{n-2}+1},R^{\mu+\nu}].
\end{aligned}
\]
Among the exponents $2,g+1,g^2+1,\ldots,g^{n-2}+1$, exactly one is divisible by $n$, namely $g^{\frac12(n-1)}+1$. Separating this term and denoting the sum of the rest by $\Omega$, Gauss obtains four equivalent forms:
\[
\text{I.}\quad
\Omega=[r,R^{\mu+\nu}]\sum R^{\mu\,\operatorname{ind}x-(\mu+\nu)\operatorname{ind}(x+1)},
\]
\[
\text{II.}\quad
\Omega=[r,R^{\mu+\nu}]\sum R^{\nu\operatorname{ind}y-(\mu+\nu)\operatorname{ind}(y+1)},
\]
\[
\text{III.}\quad
\Omega=[r,R^{\mu+\nu}]\sum R^{\mu\operatorname{ind}(1-z)+\nu\operatorname{ind}z},
\]
and
\[
\text{IV.}\quad
\Omega=[r,R^{\mu+\nu}]\sum R^{\mu\operatorname{ind}(n+1-z)+\nu\operatorname{ind}z}
=[r,R^{\mu+\nu}]\sum R^{\mu\operatorname{ind}z+\nu\operatorname{ind}(n+1-z)}.
\]
The summation variables range respectively over the residues stated in the source; only their residues modulo $\delta$ are relevant in the exponents of $R$. Combining the exceptional term with $\Omega$ yields
\[
\begin{aligned}
[r,R^\mu][r,R^\nu]&=R^{\frac12(n-1)\mu}\{[1,R^{\mu+\nu}]
+[r,R^{\mu+\nu}]\sum R^{\mu\operatorname{ind}(z-1)+\nu\operatorname{ind}z}\}\\
&=R^{\frac12(n-1)\mu}\{[1,R^{\mu+\nu}]
+[r,R^{\mu+\nu}]\sum R^{\mu\operatorname{ind}z+\nu\operatorname{ind}(z-1)}\}.
\end{aligned}
\]
If $\mu+\nu\equiv0\pmod\delta$, then
\[
[r,R^\mu][r,R^\nu]=(n-1)R^{\frac12(n-1)\mu}
+(r+r^2+\cdots+r^{n-1})(1+R^\mu+R^{2\mu}+\cdots+R^{(n-2)\mu}-R^{\frac12(n-1)\mu}).
\]

\subsection*{12.}

Similarly,
\[
\begin{aligned}
[1,R^\mu][r,R^\nu]
&=[r,R^{\mu+\nu}](1+R^\mu+R^{2\mu}+\cdots+R^{(n-2)\mu})\\
&=[r,R^{\mu+\nu}]\,\frac{n-1}{\delta}(1+R^\mu+R^{2\mu}+\cdots+R^{(\delta-1)\mu}),
\end{aligned}
\]
and
\[
[1,R^\mu][1,R^\nu]
=\frac{n-1}{\delta}[1,R^{\mu+\nu}](1+R^\mu+R^{2\mu}+\cdots+R^{(\delta-1)\mu}).
\]
Thus the general product $[r^m,R^\mu][r^{m'},R^\nu]$ is reduced either to the preceding article or to the cases just described, after applying
\[
[r^m,R^\mu]=R^{-\mu\operatorname{ind}m}[r,R^\mu]
\]
when $m$ is not divisible by $n$.

\subsection*{13.}

Once the product of two factors has been reduced, products of more factors cause no difficulty. If
\[
[r,R^\mu][r,R^\nu]=A[1,R^{\mu+\nu}]+B[r,R^{\mu+\nu}],
\]
and a third factor $[r,R^\pi]$ is added, then
\[
[r,R^{\mu+\nu}][r,R^\pi]=c[1,R^{\mu+\nu+\pi}]+d[r,R^{\mu+\nu+\pi}],
\]
and the product becomes
\[
C[1,R^{\mu+\nu+\pi}]+D[r,R^{\mu+\nu+\pi}],
\]
where
\[
C=Bc,\qquad
D=Bd+A(1+R^{\mu+\nu}+R^{2\mu+2\nu}+\cdots+R^{(n-2)(\mu+\nu)}).
\]

For $n=11$, $\delta=5$, $g=2$, the indices are
\[
\begin{array}{c|cccccccccc}
\text{numbers}&1&2&3&4&5&6&7&8&9&10\\
\text{indices}&0&1&8&2&4&9&7&3&6&5.
\end{array}
\]
Using the preceding formulae,
\[
\Omega=[r,R^2](2+4R+R^3+2R^4)
\]
and hence
\[
[r,R]^2=[1,R^2]+[r,R^2](2+4R+R^3+2R^4).
\]
Likewise
\[
\begin{aligned}
[r,R^2][r,R]&=[1,R^3]+[r,R^3](2+R+4R^2+2R^3),\\
[r,R^3][r,R]&=[1,R^4]+[r,R^4](2+4R+R^3+2R^4),\\
[r,R^4][r,R]&=[1,1]+[r,1](1+2R+2R^2+2R^3+2R^4).
\end{aligned}
\]
Substitution gives
\[
\begin{aligned}
[r,R]^3&=[1,R^3](2+4R+R^3+2R^4)
+[r,R^3](12+22R+18R^2+24R^3+15R^4),\\
[r,R]^4&=[1,R^4](12+22R+18R^2+24R^3+15R^4)\\
&\quad+[r,R^4](164+170R+205R^2+180R^3+190R^4),\\
[r,R]^5&=9188s-98S-(6R+41R^2+16R^3+26R^4)s
+66R+451R^2+176R^3+286R^4.
\end{aligned}
\]

\subsection*{14.}

The preceding calculation simplifies considerably if $R$ is from the start taken to be a proper root of $x^\delta-1=0$. Then $[r,R^\mu][r,R^\nu]$ reduces to a multiple of $[r,R^{\mu+\nu}]$ whenever $\mu+\nu$ is not divisible by $\delta$. If $\mu+\nu$ is divisible by $\delta$, then, for $\mu$ not divisible by $\delta$,
\[
[r,R^\mu][r,R^\nu]=R^{\frac12(n-1)\mu}\{n-1-[r,1]\}.
\]
This is $0$ for $r=1$ and $\mp n$ times the indicated root of unity otherwise. Hence one may set
\[
[r,R]^2=A'[r,R^2],\quad
[r,R^2][r,R]=A''[r,R^3],\quad
[r,R^3][r,R]=A'''[r,R^4],
\]
and so forth, giving
\[
[r,R]^2=A'[r,R^2],\quad
[r,R]^3=A'A''[r,R^3],\quad
[r,R]^4=A'A''A'''[r,R^4],
\]
and finally
\[
[r,R]^\delta=\pm nA'A''A'''\cdots A^{(\delta-2)}.
\]
Therefore, once $[r,R]$ is known,
\[
[r,R^2]=\frac{[r,R]^2}{A'},\qquad
[r,R^3]=\frac{[r,R]^3}{A'A''},\qquad\text{etc.}
\]

\subsection*{15.}

The values may be put in trigonometric form. Write
\[
[r,R]=\sqrt n(\cos f'+i\sin f'),\quad
[r,R^2]=\sqrt n(\cos f''+i\sin f''),\quad
[r,R^3]=\sqrt n(\cos f'''+i\sin f''').
\]
Then
\[
\begin{aligned}
A'&=\sqrt n(\cos(2f'-f'')+i\sin(2f'-f'')),\\
A''&=\sqrt n(\cos(f'+f''-f''')+i\sin(f'+f''-f''')),\\
A'''&=\sqrt n(\cos(f'+f'''-f'''')+i\sin(f'+f'''-f'''')).
\end{aligned}
\]
If
\[
A'=a'(\cos b'+i\sin b'),\quad
A''=a''(\cos b''+i\sin b''),\quad
A'''=a'''(\cos b'''+i\sin b'''),
\]
with the $a$'s positive, then all $a$'s equal $\sqrt n$. According as $(n-1)/\delta$ is even or odd, one has
\[
f'=\frac1\delta(b'+b''+\cdots+b^{(\delta-2)})
\quad\text{or}\quad
f'=\frac1\delta(180^\circ+b'+b''+\cdots+b^{(\delta-2)}).
\]
Then
\[
[r,R^2]=\sqrt n(\cos(2f'-b')+i\sin(2f'-b')),
\]
\[
[r,R^3]=\sqrt n(\cos(3f'-b'-b'')+i\sin(3f'-b'-b'')),
\]
and so on. Substitution in the formulae of art.~8 gives the periods of $(n-1)/\delta$ terms.

\subsection*{16.}

A further simplification follows from the symmetry
\[
\pm [r,R][r,R^{\delta-1}]=[r,R^2][r,R^{\delta-2}]=\pm[r,R^3][r,R^{\delta-3}]=\cdots=n.
\]
In the even case one may put
\[
f^{(\delta-1)}=-f',\quad f^{(\delta-2)}=-f'',\quad f^{(\delta-3)}=-f''',
\]
and in the odd case
\[
f^{(\delta-1)}=180^\circ-f',\quad f^{(\delta-2)}=-f'',\quad
f^{(\delta-3)}=180^\circ-f'''.
\]
Consequently the number of independent functions $A',A'',A''',\ldots$ is reduced by half. In general, the determination of the periods requires the division of the circle into $\delta$ parts, the rational construction of the angles $b',b'',b'''$, the division of an angle composed from these into $\delta$ parts, and finally the square root $\sqrt n$. When $\delta=\frac12(n-1)$, the periods coincide with twice the cosines of the angles
\[
\frac{360^\circ}{n},\quad 2\frac{360^\circ}{n},\quad 3\frac{360^\circ}{n},\ldots,\frac12(n-1)\frac{360^\circ}{n}.
\]
Thus the division of the circle into $n$ parts is reduced to a division into $\frac12(n-1)$ parts, a further angular division, and the square root $\sqrt n$.

\subsection*{17.}

For the example $n=11$, $\delta=5$, Gauss finds
\[
A'=A'''=2+4R+R^3+2R^4=2R-2R^2-R^3,
\]
\[
A''=2+R+4R^2+2R^3=-R+2R^2-2R^4.
\]
Taking $R=\cos72^\circ+i\sin72^\circ$, the angles $b',b''$ satisfy
\[
\tan b'=\frac{2\sin72^\circ-\sin144^\circ}{2\cos72^\circ-3\cos144^\circ},\qquad
\tan b''=\frac{\sin72^\circ+2\sin144^\circ}{-3\cos72^\circ+2\cos144^\circ}.
\]
Numerically,
\[
\begin{array}{rclcrcl}
\tan b'&=&+0.4316226944,&\quad&\log\tan b'&=&9.6351042715,\quad b'=23^\circ20'46''.04603,\\
\tan b''&=&-0.8355819332,&\quad&\log\tan b''&=&9.9219890411n,\quad b''=140^\circ7'6''.52441.
\end{array}
\]
Hence
\[
5f'=186^\circ48'38''.61647,\qquad f'=37^\circ21'43''.723294.
\]
The resulting periods include
\[
\begin{aligned}
(2,1)&=-\frac15+\frac{\sqrt{11}}5\{2\cos37^\circ21'43''.723294+2\cos51^\circ22'41''.400558\},\\
(2,2)&=-\frac15+\frac{\sqrt{11}}5\{2\cos325^\circ21'43''.723294+2\cos267^\circ22'41''.400558\},
\end{aligned}
\]
and yield, for instance,
\[
(2,1)=+1.6825070652=2\cos\frac{360^\circ}{11},\qquad
(2,2)=+0.8308299=2\cos\frac{720^\circ}{11}.
\]

\subsection*{18.}

For $x^{17}-1=0$, already treated in \emph{Disquisitiones Arithmeticae} by another method, put $n=17$, $\delta=8$, $g=3$. The indices are
\[
\begin{array}{c|cccccccccccccccc}
\text{numbers}&1&2&3&4&5&6&7&8&9&10&11&12&13&14&15&16\\
\text{indices}&0&14&1&12&5&15&11&10&2&3&7&13&4&9&6&8.
\end{array}
\]
The calculation gives
\[
A'=A''''=-3-2R-2R^3,\quad
A''=A'''''=1-4R^2,\quad
A'''=A''''''=3+2R+2R^3.
\]
For $R=\cos45^\circ+i\sin45^\circ$ this becomes
\[
A'=A''''=-3-2i\sqrt2,\quad A''=A'''''=1-4i,\quad A'''=A''''''=3+2i\sqrt2.
\]
Thus
\[
b'=223^\circ18'49'',\quad b''=284^\circ2'10,\quad b'''=43^\circ18'49''=b'-180^\circ,
\]
and
\[
4f'=550^\circ39'48'',\qquad f'=137^\circ39'57''.
\]
The octuple periods are then expressed by formulae such as
\[
(2,1)=-\frac18+\frac{\sqrt{17}}8\{2\cos137^\circ39'57''+2\cos52^\circ1'5''+2\cos265^\circ38'52''+1\},
\]
with analogous formulae for the remaining periods. Numerically these identify the half-periods with cosines $\cos(k360^\circ/17)$.

Gauss then turns from these general investigations of the functions $[r,R]$ to special cases, in particular fixed values of $\delta$. Some of the resulting investigations had already appeared in \emph{Disquisitiones Arithmeticae} by a different method; others are new. The deeper connection with higher arithmetic is reserved for another paper.

\subsection*{19.}

The first special case is $\delta=2$, so $R=-1$. Then
\[
[r,R]=r-r^g+r^{g^2}-r^{g^3}+\cdots-r^{g^{n-2}},
\]
and
\[
[r^\lambda,R]=+[r,R]\quad\text{if }\lambda\text{ is a quadratic residue modulo }n,
\]
\[
[r^\lambda,R]=-[r,R]\quad\text{if }\lambda\text{ is a quadratic non-residue modulo }n.
\]
Let the quadratic residues modulo $n$ be denoted by $a$, and the non-residues by $b$. Then
\[
[r,R]=\sum r^a-\sum r^b.
\]
Putting $\omega=360^\circ/n$ and $r=\cos k\omega+i\sin k\omega$, one obtains
\[
[r,R]=\sum\cos ak\omega-\sum\cos bk\omega
+i\sum\sin ak\omega-i\sum\sin bk\omega.
\]
Since by art.~14 the square of $[r,R]$ is $+n$ or $-n$ according as $n$ has the form $4z+1$ or $4z-1$, one obtains:
\[
\text{If }n=4z+1:\quad
\sum\cos ak\omega-\sum\cos bk\omega=\pm\sqrt n,\quad
\sum\sin ak\omega-\sum\sin bk\omega=0.
\]
\[
\text{If }n=4z-1:\quad
\sum\cos ak\omega-\sum\cos bk\omega=0,\quad
\sum\sin ak\omega-\sum\sin bk\omega=\pm\sqrt n.
\]
Together with
\[
\sum\cos ak\omega+\sum\cos bk\omega=-1,\qquad
\sum\sin ak\omega+\sum\sin bk\omega=0,
\]
this gives
\[
\text{for }n=4z+1:\quad
\sum\cos ak\omega=-\frac12\pm\frac12\sqrt n,\quad
\sum\cos bk\omega=-\frac12\mp\frac12\sqrt n,\quad
\sum\sin ak\omega=\sum\sin bk\omega=0,
\]
and
\[
\text{for }n=4z-1:\quad
\sum\cos ak\omega=\sum\cos bk\omega=-\frac12,\quad
\sum\sin ak\omega=\pm\frac12\sqrt n,\quad
\sum\sin bk\omega=\mp\frac12\sqrt n.
\]
These sums had already been derived in \emph{Disquisitiones Arithmeticae}, art.~356, by a not very different method. Neither method removes the ambiguity of sign; Gauss states that he has recently removed this defect in a separate memoir, where it is shown that for $k=1$ the upper signs in all the displayed formulae are to be taken.

\section*{Remarks by Dedekind}

Of the original continuation of this memoir from art.~19 onward, which was devoted to the treatment of special cases, only a few articles survive. They concern the quadratic equation whose roots are the two periods of $\frac{n-1}{2}$ terms; the manuscript breaks off at the beginning of the investigation by which the sign of the square root occurring in the solution was to be determined. From the agreement of this surviving beginning with the memoir \emph{Summatio quarundam serierum singularium}, it follows that the author changed his plan and made the determination of the sign into a separate memoir. Comparing this with the citation in art.~8, the manuscript must belong to the year 1808. A later inserted sheet at art.~19 shows, however, that publication of the preceding material had not been abandoned; on it another continuation begins, already referring to the \emph{Summatio} for the determination of the sign. This second continuation also breaks off soon and is printed here. In preparing the edition the text of the press-ready manuscript has been faithfully retained; only in art.~16 did the formulae for the second case have to be added.

\begin{flushright}
R. Dedekind.
\end{flushright}

\clearpage
\section*{Proof of some theorems concerning the periods of the classes of binary forms of the second degree}

\subsection*{Theorem I}

\emph{The number of classes (pr. pr.) of one and the same determinant which, when raised to the $P$-th dignity, $P$ being either a prime number or a power of a prime number $=p^\pi$, produce the principal class $K$, is equal either to $1$ or to a power of this same prime number $p$.}

\emph{Proof.} Let $(\mathfrak Q)$ be the whole group of all the classes in question, and let $n$ be their number. Since the principal class $K$ is necessarily contained in $(\mathfrak Q)$, the theorem is evident if it is the only one there. But if there are others, the number of classes contained in the period of each will be a power of $p$. Let $A$ be one of them, and suppose that its period $(\mathfrak A)$ contains $p^\alpha$ classes, all of which will be comprised in $(\mathfrak Q)$. If the classes of this period $(\mathfrak A)$ exhaust $(\mathfrak Q)$, then $p^\alpha=n$, and the theorem is proved. Otherwise, let $B$ be any class of $(\mathfrak Q)$ not contained in $(\mathfrak A)$, and suppose that its period is developed until one reaches a class $bB$ that is at the same time among the classes of $(\mathfrak A)$; this must necessarily happen, because at least the principal class is common to this period and to $(\mathfrak A)$. Now suppose that $bB$ is the first class in the period of $B$ common to $(\mathfrak A)$, or that $b$ is as small as possible. I say:

\begin{enumerate}[label=\arabic*\textsuperscript{o}.]
\item That $b$ will be a power of $p$. For it is evident that, on putting $b=p^\beta h$, $bB=iA$, and $hk\equiv 1\pmod {p^\pi}$ (which can be done), one obtains
\[
kbB=p^\beta hkB=p^\beta B=ikA,
\]
that is, $p^\beta B$ will also be among the classes of $(\mathfrak A)$. It follows that $h=1$ and $b=p^\beta$.

\item That, if the classes $K,B,2B,\ldots,(b-1)B$ are denoted by $(\mathfrak B)$, all compositions of a class of $(\mathfrak A)$ with a class of $(\mathfrak B)$ will give $p^{\alpha+\beta}$ distinct classes. For if $mA+nB=m'A+n'B$ and $n=n'$, then necessarily $m=m'$. If $n>n'$, then $(n-n')B=(m'-m)A$, which is impossible unless $n=n'$.

\item That these $p^{\alpha+\beta}$ distinct classes will be comprised in $(\mathfrak Q)$, which is evident.
\end{enumerate}

If these $p^{\alpha+\beta}$ classes exhaust $(\mathfrak Q)$, the theorem is proved. Otherwise, choose another class of $(\mathfrak Q)$ not contained among them, say $C$. Continue its period until one reaches a class already comprised among the classes composed from $(\mathfrak A)$ and $(\mathfrak B)$. By a reasoning similar to the preceding one, one proves that the exponent of this class must be a power of $p$, say $p^\gamma$, and that the composition of the first $p^\gamma$ classes of the period of $C$ with the $p^{\alpha+\beta}$ classes already found will give $p^{\alpha+\beta+\gamma}$ distinct classes, all comprised in $(\mathfrak Q)$. If these classes still do not exhaust $(\mathfrak Q)$, one treats a fourth class $D$ in the same manner, and so on. Since $(\mathfrak Q)$ consists of a finite number of classes, these operations also terminate, and one obtains $n$ equal to a power of $p$. Q.E.D.

\subsection*{Theorem II}

\emph{Suppose that the number of all classes of the principal genus is expressed by $a^\alpha b^\beta c^\gamma$ etc., where $a,b,c$ denote distinct prime numbers. Then in this genus there will be, respectively, $a^\alpha$, $b^\beta$, $c^\gamma$ etc. classes which, when raised to the dignities $a^\alpha$, $b^\beta$, $c^\gamma$ etc., produce the principal class.}

\emph{Proof.} Let $A,A',A''$ etc. be all the classes which, when raised to the dignity $a^\alpha$, produce $K$, and let $(\mathfrak A)$ be their totality. Similarly let $B,B',B''$ etc. be $(\mathfrak B)$, and $C,C',C''$ be $(\mathfrak C)$, and so on. I say that the composition of all classes of $(\mathfrak A)$ with all classes of $(\mathfrak B)$ with all classes of $(\mathfrak C)$, and so on, will produce classes distinct from one another. For if
\[
A+B+C+\cdots=A'+B'+C'+\cdots \quad\text{etc.,}
\]
then, putting $A-A'=A''$, $B-B'=B''$ etc., one has
\[
A''+B''+C''+\cdots=K.
\]
Therefore, raising to the dignity $b^\beta c^\gamma$ etc.,
\[
(b^\beta c^\gamma\cdots)A''=K,
\]
from which it easily follows that $A''=K$ and $A=A'$, and in the same way $B=B'$, $C=C'$, etc. Let the totality of these classes be $(\mathfrak S)$. Moreover it is clear that all these classes will belong to the principal genus. Finally, there can exist no class in the principal genus that is not comprised in $(\mathfrak S)$. Let \ldots

\section*{Remark}

This fragment, written in the year 1801, relates to \emph{Disq. Arithm.} art.~306, IX. The word \emph{dignité} is used here in an otherwise unusual sense.

\begin{flushright}
\textsc{Stern}.
\end{flushright}



\section*{[I.] On the Connection between the Number of Classes into Which Binary Forms of the Second Degree Are Distributed and Their Determinant}

\begin{center}
\textsc{First memoir}\\
\textsc{Presented to the Royal Society, 1834}
\end{center}

\subsection*{1.}

Thirty-three years have already passed since I discovered the principles of the remarkable connection to which this memoir is devoted, as was already announced at the end of the \emph{Disquisitiones Arithmeticae}. Other occupations, however, long drew me away from this investigation, until in more recent times it was possible to return to it and to enlarge it by new efforts. Since this new part of higher arithmetic exceeds the limits of one memoir, the present first memoir will be devoted to forms with negative determinants; forms with positive determinants, which require an entirely special treatment, must remain reserved for another memoir.

\subsection*{2.}

The basis of the whole argument is a special investigation of the number of all combinations of integral values that two indeterminate integers $x,y$ take within a prescribed boundary. Manifestly this problem can also be presented geometrically: one seeks the number of complex integers whose representation falls inside a prescribed figure. The nature of the figure depends on the nature of the line that surrounds it; hence it will depend either on a single equation between the coordinates $x,y$, when the boundary is a curve returning into itself, or on several equations of this kind, when it consists of several curved or straight parts. It also depends on our choice whether points corresponding to complex integers, if any happen to lie on the boundary itself, are to be counted in the total or excluded from it.

In the analytic representation of the problem, the conditions of this limitation can always be displayed so that one or more given functions $P,Q,R$ etc. of the variables $x,y$ must take positive or non-negative values, according as the boundary value is excluded or admitted.

Thus, for example, if the prescribed figure is a circle whose radius is $\sqrt A$, and whose centre falls at a point corresponding to a complex integer, the analytic condition is that $A-x^2-y^2$ should not be negative, since, as we shall always suppose, we agree to retain the points situated on the circumference itself. If the figure is a triangle, the three linear functions
\[
ax+by+c,\qquad a'x+b'y+c',\qquad a''x+b''y+c''
\]
must have non-negative values; and similarly in other cases.

\subsection*{3.}

The exact solution of the problem must, generally speaking, proceed as follows. First, from the nature of the conditions, one variable, say $x$, is confined between limits; then the successive integral values within those limits are run through, and one must find how many integral values of the other variable $y$ correspond to each of them. These numbers must then be collected into a sum. In special cases there will usually be special devices for shortening the labour.

For example, if the figure, as above, is a circle whose radius is $\sqrt A$, let $r$ be the integer immediately less than $\sqrt A$, or equal to $\sqrt A$ itself if $A$ is a square. Likewise let $r',r'',r'''$ etc., $r^{(r)}$ be the integers immediately less than
\[
\sqrt{A-1},\quad \sqrt{A-4},\quad \sqrt{A-9},\quad \ldots,\quad \sqrt{A-r^2}.
\]
Then the required number will be
\[
2r+1+2(2r'+1)+2(2r''+1)+2(2r'''+1)+\cdots,
\]
or
\[
1+4r+4r'+4r''+4r'''+\cdots+4r^{(r)}.
\]

In this example the following method is shorter. Let $q$ be the integer immediately less than $\sqrt{A/4}$, or equal to it whenever it is an integer; and let $r^{(q+1)},r^{(q+2)},r^{(q+3)}$ etc. be the integers immediately less than
\[
\sqrt{A-(q+1)^2},\quad \sqrt{A-(q+2)^2},\quad \sqrt{A-(q+3)^2},\quad\ldots,\quad \sqrt{A-r^2}.
\]
Then the required number is
\[
4q^2+1+4r+8\left(r^{(q+1)}+r^{(q+2)}+r^{(q+3)}+\cdots+r^{(r)}\right).
\]

By this formula the following counts have been obtained:
\[
\begin{array}{rr|rr|rr}
A&M&A&M&A&M\\
100&317&1000&3149&10000&31417\\
200&633&2000&6293&20000&62845\\
300&949&3000&9425&30000&94237\\
400&1257&4000&12581&40000&125629\\
500&1581&5000&15705&50000&157093\\
600&1885&6000&18853&60000&188453\\
700&2209&7000&21993&70000&219901\\
800&2521&8000&25137&80000&251305\\
900&2821&9000&28269&90000&282697\\
1000&3149&10000&31417&100000&314197
\end{array}
\]

\subsection*{4.}

For our purpose the exact determination is not required. Rather, we need to investigate an expression that can come as close as one wishes to the exact number when the limits are enlarged without bound. But since this formulation involves something vague, the matter must first be explained more exactly.

We shall suppose that the functions $P,Q,R$ etc., besides the variables $x,y$, include a constant element $k$, in such a way that each of $P,Q,R$ etc. is a homogeneous function of the three quantities $x,y,k$. In this manner the figure determined by the equations $P=0,Q=0,R=0$ etc. will depend on $k$, so that different values of $k$ correspond to similar figures, similarly placed with respect to the origin of coordinates; the corresponding linear dimensions are proportional to the values of $k$, and the areas are proportional to $k^2$. Let the number of points inside the figure be denoted by $M$, and the area by $V$. It is clear that $M$ and $V$ must increase as $k$ increases. As $k$ grows without bound, $M$ and $V$ approach the ratio of equality as closely as one wishes. More explicitly, given any arbitrarily small quantity $\lambda$, a bound can always be assigned such that, for every value of $k$ exceeding this bound, the quotient $\frac{M}{V}$ must lie between $1-\lambda$ and $1+\lambda$. In the usual manner this may be indicated by saying that $M=V$ for an infinite value of $k$.

In our example the required condition holds when $k=\sqrt A$; the curve is a circle whose area is $\pi A$, where $\pi$ denotes the semicircumference of the circle with radius $1$. The numbers given above show the convergence plainly. If it were worth the effort, we could easily complete the demonstration of this theorem with classical rigour; here, however, we prefer to omit it and hasten toward more difficult matters.

\subsection*{5.}

In this memoir the boundary will be expressed by a single equation of the form
\[
ax^2+2bxy+cy^2=A,
\]
where $a,b,c$ are integers, and where $b^2-ac$ is a negative number, which we shall put equal to $-D$. The curve defining the figure is evidently an ellipse, and it is easy to see that the squares of the semi-axes are the roots of the equation
\[
(ac-b^2)q^2-(a+c)Aq+A^2=0,
\]
or
\[
q=\frac{A\{(a+c)\pm \sqrt{(a-c)^2+4b^2}\}}{2D}.
\]
The product of these roots is $\frac{A^2}{D}$; hence the area of the ellipse is $\frac{\pi A}{\sqrt D}$. It follows that the number of all combinations of integral values of $x,y$ for which $ax^2+2bxy+cy^2$ does not exceed the value $A$ approaches ever more closely, as $A$ increases, to $\frac{\pi A}{\sqrt D}$; for infinite $A$ it may be taken as equal to this value.

It is clear that, for this purpose, it makes no difference whether the combination $x=0,y=0$ is counted with the others or excluded. Thus the number sought, in the latter interpretation, is nothing other than the aggregate of the numbers of representations of the individual numbers $1,2,3,\ldots,A$ by the binary quadratic form $ax^2+2bxy+cy^2$. Since among these numbers some cannot be represented by this form at all, while others admit more representations and others fewer, the quantity $\frac{\pi}{\sqrt D}$ is to be regarded as the mean value of the number of representations of an indeterminate positive integer by any form whose determinant is $-D$.

\subsection*{6.}

Before we examine the general consequences of this, it seems useful to work out a few special cases, so that the mode of argument may be more readily understood. Let us first return to the form $x^2+y^2$. For this form the number of representations of an indeterminate integer receives the mean value $\pi$. The number of actual representations of a given number is not difficult to determine from the general principles established in the \emph{Disquisitiones Arithmeticae}. Denote by $f_A$ the number of representations of the number $A$. It is $4$ if $A=1$ or $2$ or a power of $2$; it is $8$ if $A$ is a prime number of the form $4n+1$, or the product of such a prime number with a power of $2$; it is $0$ if $A$ is a prime number of the form $4n+3$, or if it is divisible by such a prime number but not by its square. In general,
\[
f_A=4(\alpha+1)(\beta+1)(\gamma+1)\cdots
\]
or $0$, according as, after reducing $A$ to the form
\[
2^e S a^\alpha b^\beta c^\gamma\cdots,
\]
where $a,b,c$ etc. denote distinct primes of the form $4n+1$, while $S$ denotes the product of the primes of the form $4n+3$ that occur once or several times among the factors of $A$, the number $S$ is or is not a square. Hence $f_A$ plainly depends only on the way in which the primes $3,5,7,11,13$ etc. occur among the factors of $A$, so that in general one must put
\[
f_A=4(3)(5)(7)(11)(13)\cdots,
\]
provided the values of the characters $(3),(5),(7)$ etc. are taken as follows. If $p$ denotes a prime number, then

\begin{enumerate}[label=\textit{\roman*.},leftmargin=2.2em]
\item $(p)=1$ if $p$ does not divide $A$;
\item $(p)=\alpha+1$ if $p$ is of the form $4n+1$ and $p^\alpha$ is the highest power of $p$ dividing $A$;
\item $(p)=0$ if $p$ is of the form $4n+3$ and the exponent of the highest power of $p$ dividing $A$ is odd;
\item $(p)=1$ if $p$ is of the form $4n+3$ and the exponent of the highest power of $p$ dividing $A$ is even.
\end{enumerate}

The first case is plainly included under the second and fourth cases. In this way the terms of the progression $f_1,f_2,f_3,f_4$ etc. proceed very irregularly, although the larger the number of terms taken, the more accurately the mean value $\pi$ must arise from them. We shall denote the aggregate
\[
f_1+f_2+f_3+\cdots+f_A
\]
by $F_A$.

\subsection*{7.}

Put, in general,
\[
f_m+f_{3m}=f'_m.
\]
It is then easy to see that
\[
f'_A=4(5)(7)(11)(13)\cdots;
\]
that is, $f'_A$ will be independent of the relation of $A$ to the divisor $3$. Thus the irregularity of the series $f'_1,f'_2,f'_3,f'_4,f'_5,f'_6$ etc. begins later and is much smaller. Moreover, if we put
\[
f'_1+f'_2+f'_3+\cdots+f'_m=F'_m,
\]
then
\[
F'_{3A}=F_{3A}+F_A.
\]
It follows easily that, as $A$ increases without bound, one must put
\[
F'_{3A}=4\pi A,
\]
or that the mean value of the terms of the series $f'_1,f'_2,f'_3,f'_4$ etc. is $\frac43\pi$.

In a similar way, putting in general
\[
f'_m+f'_{5m}=f''_m,
\]
one obtains
\[
f''_A=4(7)(11)(13)\cdots;
\]
that is, the fluctuations depending on the relation to the number $5$ disappear from the new series $f''_1,f''_2$ etc. Putting the aggregate
\[
f''_1+f''_2+\cdots+f''_m=F''_m,
\]
one has
\[
F''_{5m}=F'_m+F'_{5m},
\]
from which it follows, as $m$ increases without bound, that one must put
\[
F''_{5m}=\frac43\pi\cdot 6m,
\]
or that the mean value of the terms of this series is $\frac43\pi\cdot\frac65$.

If we proceed in the same way, forming new progressions while successively removing the factors $(7),(11),(13),(17)$ etc., these progressions approach more and more nearly to invariability, and their mean values successively acquire the new factors
\[
\frac87,\quad \frac{10}{11},\quad \frac{12}{13},\quad \ldots.
\]
Here the denominators are the prime numbers in their natural order, while the numerators are either one greater or one less according as the denominators are of the form $4n-1$ or $4n+1$. Since, when this process is continued without bound, the constant value $4$ must become continually closer to the mean value, we have
\[
4=\pi\cdot \frac43\cdot\frac65\cdot\frac87\cdot\frac{10}{11}\cdot\frac{12}{13}\cdots,
\]
or
\[
\pi=4\cdot\frac{3}{3+1}\cdot\frac{5}{5-1}\cdot\frac{7}{7+1}\cdot\frac{11}{11-1}\cdot\frac{13}{13-1}\cdots.
\]

If the individual fractions are expanded into infinite series,
\[
\frac{3}{3+1}=1-\frac13+\frac1{9}-\frac1{27}+\cdots,
\]
\[
\frac{5}{5-1}=1+\frac15+\frac1{25}+\frac1{125}+\cdots,
\]
\[
\frac{7}{7+1}=1-\frac17+\frac1{49}-\frac1{343}+\cdots,
\]
and so forth, the product is easily expanded into
\[
1-\frac13+\frac15-\frac17+\frac19-\frac1{11}+\cdots,
\]
whose sum is commonly known to be $\frac14\pi$. In fact, by the inverse route the equality between $\pi$ and the infinite product
\[
4\cdot\frac34\cdot\frac56\cdot\frac78\cdot\frac{11}{10}\cdot\frac{13}{12}\cdots
\]
had already been obtained from this by the illustrious Euler (\emph{Introductio in analysin infinitorum}, vol. I, chap. XV, art. 285).

\subsection*{8.}

In the second place, consider the form $x^2+2y^2$, for which the number of representations of an indeterminate integer has the mean value $\frac{\pi}{\sqrt2}$. If $f_A$ denotes the number of representations of a given number $A$ by this form, it is $2$ for $A=1$ or $A=2$, and whenever $A$ is a power of $2$; further, $f_A=4$ whenever $A$ is one of the prime numbers for which $2$ is a quadratic residue, that is, primes of the forms $8k+1$ or $8k+3$, for example $A=3,11,17,19,41,43$, etc.; finally, $f_A=0$ whenever $A$ is a prime number for which $2$ is a quadratic non-residue, for example from the series $5,7,13,23,29,31$ etc., that is, of the forms $8k+5$ or $8k+7$. In general one must put
\[
f_A=2(\alpha+1)(\beta+1)(\gamma+1)\cdots
\]
or $f_A=0$, according as, after reducing $A$ to the form
\[
2^e S a^\alpha b^\beta c^\gamma\cdots,
\]
where $a,b,c$ etc. denote distinct primes of the forms $8k+1$ and $8k+3$, while $S$ is the product of the remaining primes, of the forms $8k+5$ and $8k+7$, if such occur among the factors of $A$, the number $S$ is or is not a square.

Hence, by reasonings entirely similar to those in the preceding article, from the series $f_1,f_2,f_3,f_4,f_5$ etc., namely $2,2,4,2,0,2$ etc., we proceed successively to other series increasingly constant, whose mean values are successively
\[
\frac{\pi}{\sqrt2}\cdot\frac23,\quad
\frac{\pi}{\sqrt2}\cdot\frac23\cdot\frac65,\quad
\frac{\pi}{\sqrt2}\cdot\frac23\cdot\frac65\cdot\frac87,\quad \ldots.
\]
Thus we are led to the equation
\[
2=\frac{\pi}{\sqrt2}\cdot\frac23\cdot\frac65\cdot\frac87\cdot\frac{10}{11}\cdot\frac{14}{13}\cdot\frac{16}{17}\cdot\frac{18}{19}\cdots,
\]
where the denominators are the natural sequence of prime numbers, while the numerators are one less whenever the denominators are of the forms $8k+1$ or $8k+3$, and one greater whenever the denominators are of the forms $8k+5$ or $8k+7$.



\section*{[II.] On the Connection between the Number of Classes into Which Binary Forms of the Second Degree Are Distributed and Their Determinant}

\begin{center}
\textsc{First memoir}\\
\textsc{Presented to the Royal Society, 1837}
\end{center}

\subsection*{1.}

Thirty-six years have passed since the principles of the remarkable connection to be treated in this memoir were discovered, as was already announced at the end of the \emph{Disquisitiones Arithmeticae}. Other occupations long drew me away from this investigation, until in more recent times it was possible to return to it and to enlarge it by new efforts. Since the scope of this new part of higher arithmetic crosses the limits of a single memoir, this first memoir will be devoted to forms with negative determinants; forms with positive determinants, which require an entirely special treatment, will remain reserved for another memoir.

\subsection*{2.}

For our purpose we shall need a theorem that is in itself arithmetical, but whose nature can be set before the eyes more conveniently and clearly through considerations presented in geometric form.

Suppose that a figure in an indefinite plane is bounded by a line of any kind whatever. Its area can be assigned approximately by dividing the plane into squares and counting both those squares that lie wholly inside the figure and those that are cut by its boundary. Plainly the resulting area will be too small or too large according as the latter squares are omitted or counted with the former. If, however, it is decided by any rule whatever that the boundary squares are partly excluded and partly counted, the error can be sometimes positive and sometimes negative, but it is necessarily smaller than the aggregate of all the boundary squares. The smaller the squares are taken, the more exactly the area is determined in this way, and this approximation can be carried on indefinitely; that is, the squares may be taken so small that the error becomes smaller than any assigned quantity. Although this may already seem evident in itself, we shall not refrain from strengthening it by a rigorous proof.

Two squares can have one corner point in common, or two points in common, or no point in common. In the first and second cases they will be called contiguous; in the third, disjoint. Manifestly, among squares all of which are mutually contiguous there can be only four; therefore among five distinct squares at least two must be disjoint. Since the distance between two points situated in disjoint squares cannot be less than the side of the squares, which we denote by $a$, it follows that if a point starting from any place in one square successively passes through a second, third, and fourth square, and then reaches a fifth, the length of the path is certainly not less than $a$. By a similar argument, if a line continually passes through other squares, the part between the fifth and ninth square, and also between the ninth and thirteenth, and so on, cannot be less than $a$. Thus we easily infer that any line returning into itself, and touching altogether $n$ distinct squares, certainly cannot have length less than $\frac{(n-4)a}{4}$. Conversely, therefore, a closed line whose length is $l$ certainly cannot have touched more than $4+\frac{4l}{a}$ distinct squares. Since their area is $4a^2+4al$, and since this can be made smaller than any assigned quantity as $a$ decreases without bound, the same conclusion holds a fortiori for the error of quadrature just discussed.

\subsection*{3.}

The rule for admitting or excluding squares placed on the boundary of the figure could be established in many different ways. The simplest, however, seems to be to consider only the position of the centre of each square: squares whose centres are inside the figure are admitted, those whose centres are outside the figure are excluded, and it is left to choice whether centres that happen to lie on the boundary itself are to be counted with the inside or with the outside. Instead of centres, any other similarly placed points in the individual squares could also be adopted.

In this way the matter reduces to imagining, in the plane, equally distant points arranged on equally distant straight lines so as to supply the squares. Then, by the theorem of the preceding article, we may assert that the number of points contained in the figure, multiplied by the square of the distance between two nearest points, becomes as close as one wishes to the area of the figure, provided only that this distance is taken sufficiently small; or, in the usual manner of speaking, that this product gives the area when the distance is infinitely small.

\subsection*{4.}

The curve expressed by the equation in orthogonal coordinates $p,q$,
\[
ap^2+2bpq+cq^2=1,
\]
is a conic section, and indeed an ellipse, if $a,c$ and $ac-b^2$ are positive quantities. The area enclosed by this ellipse is found to be
\[
\frac{\pi}{\sqrt{ac-b^2}}.
\]
The value of the quantity $ap^2+2bpq+cq^2$ is everywhere greater than $1$ outside the ellipse, less than $1$ inside it, and nowhere negative.

Conceive a system of points scattered through the plane in which the ellipse lies, so that they form squares whose sides, equal to $\lambda$, are parallel to the coordinate axes. It is immaterial whether the origin of coordinates, or centre of the ellipse, coincides with one of these points or not. Let the number of points inside the ellipse, including any that lie on the boundary itself, be $m$. Then, by the theorem of the preceding article,
\[
\frac{\pi}{\sqrt{ac-b^2}}
\]
is the limit of the quantity $m\lambda^2$, which it approaches as closely as desired when $\lambda$ decreases without bound.

If we suppose that the origin of coordinates coincides with one point of the system, and put $p=\lambda x,\ q=\lambda y$, then for each point of the system $x$ and $y$ will be integers, and conversely every combination of integral values of the quantities $x,y$ will correspond to some point of the system. Hence the number $m$ is nothing other than the number of all combinations of integral values of $x,y$ for which $F$ is not greater than $M$, if for brevity we denote the function, or form of the second degree, $ax^2+2bxy+cy^2$ by $F$, and denote the quantity $\frac{1}{\lambda^2}$ by $M$. The determinant of this form is $b^2-ac$, for which we shall write $-D$. Our theorem may therefore now be stated as follows.

\textsc{Theorem I.} \emph{The number $m$ of all combinations of integral values of the indeterminates $x,y$ for which the value of a form of negative determinant $-D$ does not exceed the limit $M$ becomes approximately $\frac{\pi M}{\sqrt D}$, with an approximation increasing without bound as $M$ increases without bound. It scarcely needs to be said that the infinite approximation here, and likewise below, is not to be understood as meaning that the difference between $\frac{\pi M}{\sqrt D}$ and $m$ itself decreases without bound, but that the ratio between these quantities approaches equality without bound; that is, $\frac{\pi M}{m\sqrt D}-1$ decreases without bound.}

\subsection*{5.}

To carry out the enumeration in practice, one may proceed as follows. For each integral value of $x$ lying between the limits $-\sqrt{\frac{cM}{D}}$ and $+\sqrt{\frac{cM}{D}}$, compute the two values of $y$ corresponding to the equation $F=M$; the number of integers lying between these then follows immediately. Since this number is the same for opposite values of $x$, almost half the labour is avoided. The same can also be done by enumerating the values of $x$ corresponding to individual values of $y$ lying between the limits $-\sqrt{\frac{aM}{D}}$ and $+\sqrt{\frac{aM}{D}}$. By a suitable combination of the two methods the labour can be reduced further; but we shall not pursue that here in detail. It will suffice to add a few remarks on the simplest case.

Let the form be $F=x^2+y^2$, so that the curve is a circle, and let $r,r',r'',r'''\ldots r^{(r)}$ denote the integers immediately less than
\[
\sqrt M,\quad \sqrt{M-1},\quad \sqrt{M-4},\quad \sqrt{M-9},\quad \ldots,\quad \sqrt{M-r^2},
\]
or those quantities themselves when any of them are integral. Then the required number is
\[
m=2r+1+2(2r'+1)+2(2r''+1)+2(2r'''+1)+\cdots+2(2r^{(r)}+1),
\]
hence
\[
m=1+4r+4r'+4r''+4r'''+\cdots+4r^{(r)}.
\]
The same result is obtained more efficiently by denoting by $q$ the integer immediately less than $\sqrt{\frac12 M}$, or that quantity itself if it is an integer, and using the formula
\[
m=4q^2+1+4r+8\left(r^{(q+1)}+r^{(q+2)}+r^{(q+3)}+\cdots+r^{(r)}\right).
\]

In this way the following numbers were obtained:
\[
\begin{array}{rr|rr|rr}
M&m&M&m&M&m\\
100&317&1000&3149&10000&31417\\
200&633&2000&6293&20000&62845\\
300&949&3000&9425&30000&94237\\
400&1257&4000&12581&40000&125629\\
500&1581&5000&15705&50000&157093\\
600&1885&6000&18853&60000&188453\\
700&2209&7000&21993&70000&219901\\
800&2521&8000&25137&80000&251305\\
900&2821&9000&28269&90000&282697\\
1000&3149&10000&31417&100000&314197
\end{array}
\]

\subsection*{6.}

We give the theorem of article 4 greater generality in the following way.

\textsc{Theorem II.} \emph{Suppose that not all combinations of integral values of $x,y$ for which $F$ does not exceed the value $M$ are to be collected, but only those taken by steps: for instance those for which $x$ is congruent to the given number $G$ modulo the given modulus $g$, and $y$ is congruent to the given number $H$ modulo the given modulus $h$. Then the number $m'$ of these combinations is expressed approximately by $\frac{\pi M}{gh\sqrt D}$, with the approximation increasing without bound as $M$ increases without bound.}

Indeed, putting $x=gx'+G,\ y=hy'+H$, it is clear that $m'$ is the number of all combinations of integral values of $x',y'$ for which
\[
ag^2\left(x'+\frac{G}{g}\right)^2+2bgh\left(x'+\frac{G}{g}\right)\left(y'+\frac{H}{h}\right)+ch^2\left(y'+\frac{H}{h}\right)^2
\]
does not exceed $M$. Therefore, if in the plane we suppose a system of points distributed as in article 4, but so that not the origin of coordinates but the point whose coordinates are $p=\frac{G\lambda}{g},\ q=\frac{H\lambda}{h}$ coincides with one of the points of the system, then $m'$ expresses the number of points lying inside the ellipse whose equation is
\[
ag^2p^2+2bghpq+ch^2q^2=1,
\]
always counting any that lie on the boundary itself. The area of this ellipse is
\[
\frac{\pi}{gh\sqrt{ac-b^2}}=\frac{\pi}{gh\sqrt D};
\]
this is the limit approached without bound by the product $m'\lambda^2=\frac{m'}{M}$, as $\lambda$ decreases or $M$ increases without bound.

Moreover, it is clear that our theorem includes the case in which only one of the two indeterminates $x,y$ is required to progress by steps, while the value of the other is subject to no condition. For this is plainly the same as putting either $h$ or $g$ equal to $1$.

\subsection*{7.}

Everything set out so far is independent of the nature of the coefficients of the form $ax^2+2bxy+cy^2$. From now on, however, we shall suppose that these coefficients are integers. Then every combination of integral values of $x,y$ gives the form an integral value, or corresponds to a representation of some integer by that form. Hence the collection of all combinations of integral values of $x,y$ through which the form $F=ax^2+2bxy+cy^2$ obtains a value not greater than the limit $M$ is the same as the collection of all representations of integers not exceeding the limit $M$, up to and including this limit if it is itself an integer. Thus, for brevity, if we denote by $F(n)$, or simply by $F_n$ when no ambiguity is to be feared, the number of distinct representations of a determinate integer $n$ by the form $F$, then the number previously expressed by $m$ will be
\[
F_0+F_1+F_2+F_3+\cdots+F_M,
\]
and the first theorem takes the following form.

\textsc{Theorem III.} \emph{The aggregate $F_0+F_1+F_2+\cdots+F_M$ is expressed approximately by $\frac{\pi M}{\sqrt D}$, with approximation increasing without bound as $M$ increases without bound.}

\subsection*{8.}

It is useful to add to the third theorem, which concerns representations of all numbers, another concerning only odd numbers. Manifestly, odd numbers cannot be represented by the form $F$ if $a$ and $c$ are both even; therefore the investigation is restricted to the three remaining cases.

I. Whenever $a$ is odd and $c$ is even, an odd number is represented by assigning an odd value to $x$, while the value of $y$ remains arbitrary. Thus Theorem II, by putting $g=2,\ G=1,\ h=1$, shows that the number of all such combinations of values of $x,y$ that give the form an odd value not greater than the limit $M$ is expressed, with infinite approximation as $M$ increases without bound, by $\frac{\pi M}{2\sqrt D}$.

II. Whenever $a$ is even and $c$ is odd, the representation of an odd number requires $y$ to be odd. Hence, putting $g=1,\ h=2,\ H=1$, we are led to the same conclusion.

III. Whenever both $a$ and $c$ are odd, either an odd value of $x$ must be combined with an even value of $y$, or an even value of $x$ with an odd value of $y$, in order that the value of the formula be odd. The number of all combinations of the former kind, and likewise of the latter kind, for which the value of the form does not exceed the limit $M$, is expressed with infinite approximation by $\frac{\pi M}{4\sqrt D}$. Therefore the number of all combinations producing odd values of the form not exceeding $M$ is also, in this case, expressed with infinite approximation by $\frac{\pi M}{2\sqrt D}$.

Since the collection of all such combinations is nothing other than the collection of all representations of all the numbers $1,3,5,7,\ldots,M$ when $M$ is an odd integer, or of $1,3,5,7,\ldots,M-1$ when $M$ is even, we have:

\textsc{Theorem IV.} \emph{The aggregate}
\[
F_1+F_3+F_5+F_7+\cdots+F_M
\]
\emph{or}
\[
F_1+F_3+F_5+F_7+\cdots+F_{M-1},
\]
\emph{according as $M$ is odd or even, is expressed with infinite approximation by $\frac{\pi M}{2\sqrt D}$, provided that $F$ is a form in which one or the other of the coefficients $a,c$, or both, is odd.}




\section*{On the Connection between the Number of Classes: Fragments III--VIII}

\subsection*{[III.]}

Let $C$ be the complex of representatives of all classes of proper primitive forms for determinant $-D$. We denote by $(n)$ the number of all representations of the number $n$ by forms from the complex $C$. Let $p$ be an odd prime number. Then
\[
\begin{array}{rll}
1.&(pn)=(n),&\text{if }p\text{ is a divisor of }D,\\[2mm]
2.&(pn)=(n)+(h),&
\left.
\begin{array}{l}
\text{divisor of }xx+D
\end{array}
\right\}\ \text{if }p\text{ is not a divisor of }D,\\[3mm]
3.&(pn)=-(n)+(h),&
\left.
\begin{array}{l}
\text{non-divisor of }xx+D
\end{array}
\right\}\ \text{if }p\text{ is not a divisor of }D,
\end{array}
\]
where $n=hp^\mu$, $\mu$ is arbitrary, and $h$ is not divisible by $p$.

In case 1 one has
\[
(h)=(ph)=(p^2h)=(p^3h)=\cdots .
\]
In case 2 one has
\[
(ph)=2(h),\qquad (p^2h)=3(h),\qquad (p^3h)=4(h),\quad\text{etc.}
\]
In case 3 one has
\[
(ph)=0,\qquad (p^2h)=(h),\qquad (p^3h)=0,\qquad (p^4h)=(h),\quad\text{etc.}
\]

From each proper primitive class for determinant $-D$, whose number is $\lambda$, let one form be selected, and let $L$ be the complex of these forms. Denote by $f_A$ the number of all representations of the number $A$ by forms from $L$.

Further, let
\[
f(A;p)=f_{\frac{A}{p^\alpha}},
\]
where $p^\alpha$ is the highest power of the prime $p$ that divides $A$; similarly let
\[
f(A;p,q)=f_{\frac{A}{p^\alpha q^\beta}},
\]
where $q$ is another prime and $q^\beta$ is the highest power of $q$ dividing $A$; and likewise
\[
f(A;p,q,r)=f_{\frac{A}{p^\alpha q^\beta r^\gamma}},
\]
where $r^\gamma$ is the highest power of a third prime $r$ dividing $A$, and so on.

\subsection*{[IV.]}

Let $(n)$ denote the number of values $x$ from the complex
\[
0,1,2,3,4,\ldots,p^n-1
\]
for which
\[
xx-D=xx-ap^\mu
\]
is divisible by $p^n$.

\[
\begin{array}{c|c|c}
\mu\ \text{odd, e.g. }=7&\multicolumn{2}{c}{\mu\ \text{even, e.g. }=6}\\
& aNp&aRp\\ \hline
(1)=1&(1)=1&(1)=1\\
(2)=p&(2)=p&(2)=p\\
(3)=p&(3)=p&(3)=p\\
(4)=p^2&(4)=p^2&(4)=p^2\\
(5)=p^2&(5)=p^2&(5)=p^2\\
(6)=p^3&(6)=p^3&(6)=p^3\\
(7)=p^3&(7)=0&(7)=2p^3\\
(8)=0&(8)=0&(8)=2p^3\\
(9)=0&\text{etc.}&\text{etc.}\\
\text{etc.}&&
\end{array}
\]

Put
\[
(1)-\frac{(2)}{p}=(1)',\quad
(2)-\frac{(3)}{p}=(2)',\quad
(3)-\frac{(4)}{p}=(3)',\quad
(4)-\frac{(5)}{p}=(4)',\quad\text{etc.}
\]
Then
\[
fp=(1)',\quad
fp^2=1+(2)',\quad
fp^3=(1)'+(3)',\quad
fp^4=1+(2)'+(4)',\quad\text{etc.}
\]
Consequently, if
\[
\frac{p-1}{p}\left(1+\frac{fp}{p}+\frac{fp^2}{p^2}+\frac{fp^3}{p^3}+\cdots\right)=T,
\]
then
\[
\frac{p+1}{p}T
=
1+\frac{(1)'}{p}+\frac{(2)'}{p^2}+\frac{(3)'}{p^3}+\frac{(4)'}{p^4}+\cdots
=
1+\frac{(1)}{p}
=
1+\frac1p .
\]
Hence $T=1$.

\subsection*{[V.]}

The mean number of classes around the negative determinant $-D$ is approximately
\[
\frac{\pi\sqrt D}{4\left(1+\frac1{2^2}+\frac1{3^2}+\cdots\right)}.
\]
The true number is expressed by the following formulae. For brevity, $m$ is written for the mean number and $M$ for the true number; $p,q$ denote all odd primes not dividing $D$, the former being divisors and the latter non-divisors of $xx+D$; and $r$ denotes odd primes dividing $D$:
\[
\begin{array}{rcl}
\text{I.}\quad
M&=&m\ \mathop{\mathrm{Prod.\,ex}}\,
\dfrac{p^3+p^2}{p^3-1}\,
\dfrac{q^3-q^2}{q^3-1}\,
\dfrac{r^3-r}{r^3-1},
\\[3mm]
\text{II.}\quad
M&=&\dfrac{\pi\sqrt D}{4}\ \mathop{\mathrm{Prod.\,ex}}\,
\dfrac{p+1}{p}\,
\dfrac{q-1}{q}\,
\dfrac{rr-1}{rr},
\\[3mm]
\text{III.}\quad
M&=&\dfrac{2\sqrt D}{\pi}\ \mathop{\mathrm{Prod.\,ex}}\,
\dfrac{p}{p-1}\,
\dfrac{q}{q+1},
\\[3mm]
\text{IV.}\quad
M&=&\sqrt{\dfrac{D}{2}}\ \mathop{\mathrm{Prod.\,ex}}\,
\left\{\dfrac{p+1}{p-1}\,
\dfrac{q-1}{q+1}\,
\dfrac{rr-1}{rr}\right\},
\\[3mm]
\text{V.}\quad
M&=&\dfrac{2\sqrt D}{\pi}
\left\{1\pm\frac13\pm\frac15\pm\frac17\pm\frac19+\cdots\right\}.
\end{array}
\]
Formula III is obtained directly by comparing the two ways of counting the representable numbers up to a given bound.

\subsection*{[VI.]}

\textsc{Theorem.} The number of classes into which all proper primitive binary forms of negative determinant $-D$ are distributed is
\[
\frac{\pi}{4}\sqrt D\,
\mathop{\mathrm{Prod.\,ex}}\,
\frac{p-1}{p}\frac{q+1}{q}\frac{rr-1}{rr},
\]
where $p$ ranges over odd primes for which $-D$ is a non-residue, $q$ over odd primes for which $-D$ is a residue, and $r$ over odd primes dividing $D$. Equivalently,
\[
=
\frac{\frac{\pi}{4}\sqrt D\ \mathop{\mathrm{Prod.\,ex}}\frac{rr-1}{rr}}
{1\pm\frac13\pm\frac15\pm\cdots}.
\]
In the denominator the positive sign is placed before fractions whose denominators are in the non-divisor form, and the negative sign before those whose denominators are in the form of divisors of $xx+D$; terms whose denominators are not prime to $D$ are omitted.

The editorial correction printed with the fragment rewrites the denominator as
\[
\frac{2\sqrt D\left(1\pm\frac13\pm\frac15\pm\cdots\right)}{\pi}
=
\frac{\operatorname{cotg}\theta\pm\operatorname{cotg}3\theta\pm\operatorname{cotg}5\theta\cdots\pm\operatorname{cotg}n\theta}{N:\sqrt D},
\]
with $\theta=\frac{\pi}{N}$, $N=\{\frac14\}D$, and with $n$ ranging over the numbers prime to $D$, signs determined as above.

For a positive determinant the number of classes is
\[
\frac{2\sqrt D\left(1\pm\frac13\pm\frac15\pm\cdots\right)}
{\log(T+U\sqrt D)},
\]
where $T,U$ are the least values of $t,u$ satisfying
\[
tt-Duu=1.
\]
It is also expressed as
\[
\frac{\log\sin\frac12\theta\pm\log\sin\frac32\theta\pm\log\sin\frac52\theta\ \mathrm{etc.}}
{\log(T+U\sqrt D)}.
\]

\subsection*{[VII.]}

For negative determinant $-p$, where $p$ is a positive prime of the form $4n+1$, the number of classes is $(\alpha-\beta)$ in the printed fragment; the note corrects this to $2(\alpha-\beta)$. Here $\alpha$ denotes the number of quadratic residues in the first quadrant
\[
1\cdot2\cdot3\cdots\frac14(p-1),
\]
and $\beta$ denotes the number of non-residues.

The signs in
\[
1\pm\frac13\pm\frac15\cdots
\]
are governed by whether the number $m$ is a product of an even or odd number of prime factors; in the numerator of the following formula the printed note gives
\[
1\pm\frac13\pm\frac15\cdots=\sum\left(\frac{-D}{m}\right)\frac1m .
\]

\subsection*{[VIII.]}

\[
b\equiv 2m+a-1\pmod{8},
\]
where $m$ is the [half] number of classes for determinant $-p$.

{\scriptsize
\[
\begin{array}{r|r|r|r|r|r|r|r}
p&m&a&b&f&\frac{2m+a-1-b}{8}&\alpha&\beta\\ \hline
17&2&+1&-4&-4&+1&3&2\\
41&4&+5&+4&+9&+1&3&4\\
73&2&-3&-8&+27&+1&1&6\\
89&6&+5&-8&+34&+3&9&2\\
97&2&+9&+4&+22&+1&5&6\\
113&4&-7&+8&+15&-1&9&4\\
137&4&-11&+4&+37&-1&3&8\\
193&2&-7&+12&+81&-2&11&6\\
233&6&+13&+8&+144&+2&15&2\\
241&6&+15&+4&+64&-1&13&6\\
257&8&+1&+16&+16&0&15&4\\
281&10&+5&-16&+53&+5&9&10\\
313&4&+13&-12&-25&+1&5&12\\
337&4&+9&+16&-148&0&7&12\\
353&8&+17&+8&+42&+3&15&8\\
5&1&+1&+2&+2&0&&\\
13&1&-3&-2&+5&0&&\\
29&3&+5&+2&+12&+1&&\\
37&1&+1&-6&-6&+1&&\\
53&3&-7&-2&+23&0&&\\
61&3&+5&-6&+11&+2&&\\
101&7&+1&-10&-10&+3&&\\
109&3&-3&+10&+33&-1&&\\
149&7&-7&-10&+44&+2&&\\
157&3&-11&-6&-28&0&&\\
173&7&+13&+2&+80&+3&&\\
181&5&+9&+10&-19&+1&&\\
197&5&+1&-14&-14&+3&&\\
229&5&-15&+2&-107&-1&&\\
269&11&+13&+10&-82&+3&&\\
277&3&+9&+14&-60&0&&\\
293&9&+17&+2&+138&+4&&\\
317&5&-11&+14&+114&-2&&\\
349&7&+5&+18&-136&0&&\\
373&5&-7&+18&+104&-2&&\\
389&11&+17&-10&-115&+6&&\\
397&3&-19&-6&+63&-1&&
\end{array}
\]
}



\section*{[IX.] Distribution of the Quadratic Residues into Octants}

Let $p$ be a prime; $(r)$ denotes the number of quadratic residues of $p$ which lie between
\[
(r-1)\frac{p}{8}\quad \text{and}\quad r\frac{p}{8}.
\]

\subsection*{First case; \(p=8n+1\)}

$2t$ is the number of classes for determinant $-p$;\\
$2u$ is the number of classes for determinant $-2p$.
\[
\begin{aligned}
(1)=(8)&=\frac14(2n+t+u),\\
(2)=(4)=(5)=(7)&=\frac14(2n+t-u),\\
(3)=(6)&=\frac14(2n-3t+u).
\end{aligned}
\]
{\scriptsize
\[
\begin{array}{r|r|r|r|r|r|r}
p&2n&t&u&(1)&(2)&(3)\\\hline
17&4&2&2&2&1&0\\
41&10&4&2&4&3&0\\
73&18&2&8&7&3&5\\
89&22&6&4&8&6&2\\
97&24&2&10&9&4&7\\
113&28&4&4&9&7&5\\
137&34&4&6&11&8&7\\
193&48&2&10&15&10&13
\end{array}
\qquad
\begin{array}{r|r|r|r|r|r|r}
p&2n&t&u&(1)&(2)&(3)\\\hline
233&58&6&4&17&15&11\\
241&60&6&10&19&14&13\\
257&64&8&8&20&16&12\\
281&70&10&4&21&19&11\\
313&78&4&18&25&16&21\\
337&84&4&12&25&19&21\\
353&88&8&12&27&21&19\\
401&100&10&6&29&26&19
\end{array}
\]
}

\subsection*{Second case; \(p=8n+5\)}

$2t$ is the number of classes for determinant $-p$;\\
$2u$ is the number of classes for determinant $-2p$.
\[
\begin{aligned}
(1)=(3)=(6)=(8)&=\frac14(2n-t+u),\\
(2)=(7)&=\frac14(2n+3t-u+2),\\
(4)=(5)&=\frac14(2n-t-u+2).
\end{aligned}
\]
{\scriptsize
\[
\begin{array}{r|r|r|r|r|r|r}
p&2n&t&u&(1)&(2)&(4)\\\hline
5&0&1&1&0&1&0\\
13&2&1&3&1&1&0\\
29&6&3&1&1&4&1\\
37&8&1&5&3&2&1\\
53&12&3&3&3&5&2\\
61&14&3&5&4&5&2\\
101&24&7&3&5&11&4\\
109&26&3&5&7&8&5\\
149&36&7&3&8&14&7\\
157&38&3&13&12&9&6\\
173&42&7&5&10&15&8
\end{array}
\qquad
\begin{array}{r|r|r|r|r|r|r}
p&2n&t&u&(1)&(2)&(4)\\\hline
181&44&5&9&12&13&8\\
197&58&5&5&12&15&10\\
229&56&5&13&16&15&10\\
269&66&11&5&15&24&13\\
277&68&3&11&19&17&14\\
293&72&9&9&18&23&14\\
317&78&5&7&20&22&17\\
349&86&7&13&23&24&17\\
373&92&5&13&25&24&19\\
389&96&11&7&23&31&20\\
397&98&3&21&29&22&19
\end{array}
\]
}

\subsection*{Third case; \(p=8n+3\)}

$t$ is the number of classes for determinant $-p$;\\
$2u$ is the number of classes for determinant $-2p$.
\[
\begin{aligned}
(1)=(4)=(7)&=\frac14(2n+t-u),\\
(2)=(5)=(8)&=\frac14(2n-t+u),\\
(3)&=\frac14(2n+t+u+2),\\
(6)&=\frac14(2n-t-u+2).
\end{aligned}
\]
{\scriptsize
\[
\begin{array}{r|r|r|r|r|r|r|r}
p&2n&t&u&(1)&(2)&(3)&(6)\\\hline
3&0&1&1&0&0&1&0\\
11&2&3&1&1&0&2&0\\
19&4&3&3&1&1&3&0\\
43&10&3&5&2&3&5&1\\
59&14&9&3&5&2&7&1\\
67&16&3&7&3&5&7&2\\
83&20&9&5&6&4&9&2\\
107&26&9&3&8&5&10&4\\
131&32&15&3&11&5&13&4\\
139&34&9&7&9&8&13&5
\end{array}
\qquad
\begin{array}{r|r|r|r|r|r|r|r}
p&2n&t&u&(1)&(2)&(3)&(6)\\\hline
163&40&3&11&8&12&14&7\\
179&44&15&3&14&8&16&7\\
211&52&9&5&14&12&17&10\\
227&56&15&7&16&12&20&9\\
251&62&21&7&19&12&23&9\\
283&70&9&15&16&19&24&12\\
307&76&9&17&17&21&26&13\\
331&82&9&11&20&21&26&16\\
347&86&15&5&24&19&27&17\\
379&94&9&11&23&24&29&19
\end{array}
\]
}

\subsection*{Fourth case; \(p=8n+7\)}

$t$ is the number of classes for determinant $-p$;\\
$2u$ is the number of classes for determinant $-2p$.
\[
\begin{aligned}
(1)&=\frac14(2n+2t-u),\\
(2)=(3)=(5)&=\frac14(2n+u+2),\\
(4)=(6)=(7)&=\frac14(2n-u+2),\\
(8)&=\frac14(2n-2t+u).
\end{aligned}
\]
{\scriptsize
\[
\begin{array}{r|r|r|r|r|r|r|r}
p&2n&t&u&(1)&(2)&(4)&(8)\\\hline
7&0&1&2&0&1&0&0\\
23&4&3&2&2&2&1&0\\
31&6&3&4&2&3&1&1\\
47&10&5&4&4&4&2&1\\
71&16&7&2&7&5&4&1\\
79&18&5&4&6&6&4&3\\
103&24&5&10&6&9&4&6\\
127&30&5&8&8&10&6&7\\
151&36&7&6&11&11&8&7\\
167&40&11&6&14&12&9&6
\end{array}
\qquad
\begin{array}{r|r|r|r|r|r|r|r}
p&2n&t&u&(1)&(2)&(4)&(8)\\\hline
191&46&13&4&17&13&11&6\\
199&48&9&10&14&15&10&10\\
223&54&7&16&13&18&10&14\\
239&58&15&4&21&16&14&8\\
263&64&13&6&21&18&15&11\\
271&66&11&12&19&20&14&14\\
311&76&19&6&27&21&18&11\\
359&88&19&6&30&24&21&14\\
367&90&9&20&22&28&18&23\\
383&94&17&12&29&27&21&18
\end{array}
\]
}

\section*{[X.] Distribution of the Quadratic Residues into Twelfths}

Let $p$ be a prime; $(r)$ denotes the number of quadratic residues of $p$ which lie between
\[
(r-1)\frac{p}{12}\quad \text{and}\quad r\frac{p}{12}.
\]

\subsection*{First case; \(p=24n+1\)}

$2t$ is the number of classes for determinant $-p$;\\
$4u$ is the number of classes for determinant $-3p$.
\[
\begin{aligned}
(1)=(12)&=\frac16(6n+3t+2u),\\
(2)=(4)=(6)=(7)=(9)=(11)&=\frac16(6n-3t+2u),\\
(3)=(5)=(8)=(10)&=\frac16(6n+3t-4u).
\end{aligned}
\]
\[
\begin{array}{r|r|r|r|r|r|r}
p&n&t&u&(1)&(2)&(3)\\\hline
73&3&2&3&5&3&2\\
97&4&2&3&6&4&3\\
193&8&2&6&11&9&5\\
241&10&6&3&14&8&11
\end{array}
\]

\subsection*{Second case; \(p=24n+13\)}

$2t$ is the number of classes for determinant $-p$;\\
$4u$ is the number of classes for determinant $-3p$.
\[
\begin{aligned}
(1)=(3)=(10)=(12)&=\frac12(2n+1+t),\\
(2)=(6)=(7)=(11)&=\frac12(2n+1-t),\\
(4)=(9)&=\frac12(2n+1-t+2u),\\
(5)=(8)&=\frac12(2n+1+t-2u).
\end{aligned}
\]
\[
\begin{array}{r|r|r|r|r|r|r|r}
p&n&t&u&(1)&(2)&(4)&(5)\\\hline
13&0&1&1&1&0&1&0\\
37&1&1&2&2&1&3&0\\
61&2&3&2&4&1&3&2\\
109&4&3&3&6&3&6&3\\
157&6&3&4&8&5&9&4\\
181&7&5&3&10&5&8&7\\
229&9&5&3&12&7&10&9
\end{array}
\]

\subsection*{Third case; \(p=24n+5\)}

$2t$ is the number of classes for determinant $-p$;\\
$2u$ is the number of classes for determinant $-3p$.
\[
\begin{aligned}
(1)=(2)=(6)=(7)=(11)=(12)&=n,\\
(3)=(10)&=\frac12(2n+1+t),\\
(4)=(9)&=\frac12(2n-t+u),\\
(5)=(8)&=\frac12(2n+1-u).
\end{aligned}
\]
\[
\begin{array}{r|r|r|r|r|r|r|r}
p&n&t&u&(1)&(3)&(4)&(5)\\\hline
5&0&1&1&0&1&0&0\\
29&1&3&3&1&3&1&0\\
53&2&3&5&2&4&3&0\\
101&4&7&5&4&8&3&2\\
149&6&7&7&6&10&6&3\\
173&7&7&9&7&11&8&3\\
197&8&5&11&8&11&11&3\\
269&11&11&7&11&17&9&8
\end{array}
\]

\subsection*{Fourth case; \(p=24n+17\)}

$2t$ is the number of classes for determinant $-p$;\\
$2u$ is the number of classes for determinant $-3p$.
\[
\begin{aligned}
(1)=(2)=(6)=(7)=(11)=(12)&=\frac16(6n+3+u),\\
(3)=(10)&=\frac16(6n+6+3t-2u),\\
(4)=(9)&=\frac16(6n+3-3t+u),\\
(5)=(8)&=\frac16(6n+6-2u).
\end{aligned}
\]
\[
\begin{array}{r|r|r|r|r|r|r|r}
p&n&t&u&(1)&(3)&(4)&(5)\\\hline
17&0&2&3&1&1&0&0\\
41&1&4&3&2&3&0&1\\
89&3&6&3&4&6&1&3\\
113&4&4&9&6&4&4&2\\
137&5&4&9&7&5&5&3\\
233&9&6&15&12&8&9&5\\
257&10&8&9&12&12&8&8
\end{array}
\]

\end{document}


